\chapter*{Lieux en litige : vérification sur les scans réservée}
\addcontentsline{toc}{chapter}{Lieux en litige : vérification sur les scans réservée}
643 points sur 443 pages imprimées. Le texte de la source convenue est conservé. Ces notices reproduisent les lectures consignées sans trancher les litiges.
\begingroup\small\raggedright\sloppy\setlength{\parskip}{4pt}
\par\noindent\hypertarget{DP0001}{}\textbf{DP0001 -- I, p. 1 (conflict)}\par
title block, after the authors | Ligne des auteurs | Sous la ligne des auteurs\par
\textbf{Claude Opus 5.5.} Place: title block, after the authors; Reported print: (avec un appendice de N. Bourbaki); Source TeX: A. Grothendieck et J.-L. Verdier \textbackslash{}end\{center\}; Note: The printed title block has the line « (avec un appendice de N. Bourbaki) » under the authors. The TeX's title block omits it.; Proposed edit: A. Grothendieck et J.-L. Verdier \textbackslash{}end\{center\} =\allowbreak > A. Grothendieck et J.-L. Verdier  (avec un appendice de N. Bourbaki) \textbackslash{}end\{center\}\par
\textbf{GPT-6.1 Sol.} Place: Ligne des auteurs; Reported print: par A. Grothendieck et J.L. Verdier; Source TeX: A. Grothendieck et J.-L. Verdier; Note: Le mot « par » manque devant les noms des auteurs.; Proposed edit: \textbackslash{}begin\{center\} A. Grothendieck =\allowbreak > \textbackslash{}begin\{center\} par A. Grothendieck || Place: Sous la ligne des auteurs; Reported print: (avec un appendice de N. Bourbaki); Source TeX: J.-L. Verdier \textbackslash{}end\{center\}; Note: La mention de l'appendice manque dans le bloc de titre.; Proposed edit: J.-L. Verdier \textbackslash{}end\{center\} =\allowbreak > J.-L. Verdier  (avec un appendice de N. Bourbaki) \textbackslash{}end\{center\}\par
the two models differ at this place\par
\par\noindent\hypertarget{DP0002}{}\textbf{DP0002 -- I, p. 11 (doubt)}\par
Preuve de 2.3, the two morphisms defining the couple of arrows\par
\textbf{Claude Opus 5.5.} Place: Preuve de 2.3, the two morphisms defining the couple of arrows; Reported print: \ensuremath{\prod}\_\{i\ensuremath{\in}ob(I)\} G(i) (twice, small 'ob'); Source TeX: \textbackslash{}prod\_\{i\textbackslash{}in \textbackslash{}Ob(I)\} (twice); Note: The print has lower-case 'ob' here, as in the display just above, where the TeX uses \textbackslash{}ob(I). If \textbackslash{}Ob prints a capital 'Ob', both occurrences in the align should be \textbackslash{}ob. There are two occurrences, so no unique 'old' is given.\par
\textbf{GPT-6.1 Sol.} nothing reported here\par
a model flags this place without giving a fix\par
\par\noindent\hypertarget{DP0003}{}\textbf{DP0003 -- I, p. 17 (one\_model\_fix)}\par
Notons maintenant… deux foncteurs C \ensuremath{\rightleftarrows} C', u adjoint à gauche\par
\textbf{Claude Opus 5.5.} Place: Notons maintenant… deux foncteurs C \ensuremath{\rightleftarrows} C', u adjoint à gauche; Reported print: C \ensuremath{\rightleftarrows} C', u adjoint à gauche de v; Source TeX: C' \textbackslash{}ar@<-.6ex>[l]\_v \}\$ \$u\$ \textbackslash{}emph\{adjoint à gauche\}; Note: The print has a comma after the diagram C \ensuremath{\rightleftarrows} C' and before « u adjoint à gauche »; the TeX drops it.; Proposed edit: \textbackslash{}ar@<-.6ex>[l]\_v \}\$ \$u\$ =\allowbreak > \textbackslash{}ar@<-.6ex>[l]\_v \}\$, \$u\$\par
\textbf{GPT-6.1 Sol.} nothing reported here\par
a fix from one model only\par
\par\noindent\hypertarget{DP0004}{}\textbf{DP0004 -- I, p. 26 (cut)}\par
End of the proof of Proposition 5.4\par
\textbf{Claude Opus 5.5.} nothing reported here\par
\textbf{GPT-6.1 Sol.} Place: End of the proof of Proposition 5.4; Reported print: grâce à (2.7).; Source TeX: The page-26 fragment ends with « grâce à »; « (\textbackslash{}hyperref[I.def2.7]\{2.7\}). » opens the page-27 fragment.; Note: The reference (2.7) is printed on page 26. Page 27 begins with Proposition 5.5. This is a cut problem, not missing text in the continuous source.\par
a model reports that the automatic cut gives lines of this page to its neighbour, or the neighbour's lines to this page: check the lines it names against the scan of the page that prints them (that page's checker did not have their TeX)\par
\par\noindent\hypertarget{DP0005}{}\textbf{DP0005 -- I, p. 33 (one\_model\_fix)}\par
Preuve de 5.11, 1), premier isomorphisme canonique\par
\textbf{Claude Opus 5.5.} nothing reported here\par
\textbf{GPT-6.1 Sol.} Place: Preuve de 5.11, 1), premier isomorphisme canonique; Reported print: \$f\textbackslash{}simeq\textbackslash{}varinjlim\_\{Y\textbackslash{}in\textbackslash{}ob C/X\}Y\$; Source TeX: \$f\textbackslash{}xrightarrow\{\textbackslash{}sim\}\textbackslash{}varinjlim\_\{Y\textbackslash{}in\textbackslash{}ob C/X\}Y\$; Note: Le tirage porte une relation d'isomorphisme, sans flèche, comme dans « Or j\_\{X!\}(f)\textbackslash{}simeq X » plus loin.; Proposed edit: f\textbackslash{}xrightarrow\{\textbackslash{}sim\} =\allowbreak > f\textbackslash{}simeq\par
a fix from one model only\par
\par\noindent\hypertarget{DP0006}{}\textbf{DP0006 -- I, p. 34 (one\_model\_fix)}\par
after « On a donc un diagramme commutatif à isomorphisme près: »\par
\textbf{Claude Opus 5.5.} Place: after « On a donc un diagramme commutatif à isomorphisme près: »; Reported print: diagram (5.13.2): (C/H)\textasciicircum{} --(u/H)\_!--> (C'/u\ensuremath{\cdot}H)\textasciicircum{} ; verticals (j\_H)\_!, (j\_\{u!H\})\_! ; C\textasciicircum{} --u\_!--> C'\textasciicircum{} ,; Source TeX: diagram present only as a \% comment; Note: The printed diagram (5.13.2) is commented out in the TeX, so the edition does not print it. The print's top-right corner reads (C'/u\ensuremath{\cdot}H)\textasciicircum{}, a misprint for (C'/u\_!H)\textasciicircum{}.; Proposed edit: \%\textbackslash{}[\textbackslash{}vcenter\{ \%\textbackslash{}xymatrix@C=\allowbreak 2.5cm@R=\allowbreak 1.5cm\{ (C/H)\{\textbackslash{}widehat\{\textbackslash{}phantom\{a\}\}\} \%\textbackslash{}ar[r]\textasciicircum{}\{(u/H)\_\{!\}\} \textbackslash{}ar[d]\_\{(j\_H)\_!\}\& \%  (C'/u\textasciicircum{}\ensuremath{\cdot}H)\{\textbackslash{}widehat\{\textbackslash{}phantom\{a\}\}\} \textbackslash{}ar[d]\_\{(j\_\{u!H\})\_\{!\}\}\textbackslash{}\textbackslash{} \%  C\{\textbackslash{}widehat\{\textbackslash{}phantom\{\textbackslash{};a\}\}\} \textbackslash{}ar[r]\textasciicircum{}\{u\_!\}\& C'\{\textbackslash{}widehat\{\textbackslash{}phantom\{\textbackslash{};a\}\}\}\textbackslash{}quad, \%\}\}\textbackslash{}leqno\{(5.13.2)\} \%\textbackslash{}] =\allowbreak > \textbackslash{}[\textbackslash{}vcenter\{ \textbackslash{}xymatrix@C=\allowbreak 2.5cm@R=\allowbreak 1.5cm\{ (C/H)\textbackslash{}widehat\{\textbackslash{}phantom\{a\}\} \textbackslash{}ar[r]\textasciicircum{}\{(u/H)\_\{!\}\} \textbackslash{}ar[d]\_\{(j\_H)\_!\}\&   (C'/u\_\{!\}H)\textbackslash{}widehat\{\textbackslash{}phantom\{a\}\} \textbackslash{}ar[d]\_\{(j\_\{u\_\{!\}H\})\_\{!\}\}\textbackslash{}\textbackslash{}   C\textbackslash{}widehat\{\textbackslash{}phantom\{a\}\} \textbackslash{}ar[r]\textasciicircum{}\{u\_!\}\& C'\textbackslash{}widehat\{\textbackslash{}phantom\{a\}\}\textbackslash{}quad, \}\}\textbackslash{}leqno\{(5.13.2)\} \textbackslash{}]\par
\textbf{GPT-6.1 Sol.} nothing reported here\par
a fix from one model only\par
\par\noindent\hypertarget{DP0007}{}\textbf{DP0007 -- I, p. 34 (one\_model\_fix)}\par
Après « On a donc un diagramme commutatif à isomorphisme près »\par
\textbf{Claude Opus 5.5.} nothing reported here\par
\textbf{GPT-6.1 Sol.} Place: Après « On a donc un diagramme commutatif à isomorphisme près »; Reported print: Le carré (5.13.2), avec ses quatre objets, ses quatre flèches et leurs étiquettes, est imprimé avant « et comme ».; Source TeX: Le diagramme (5.13.2) est entièrement commenté et ne s'imprime pas.; Note: Rétablir le diagramme et son numéro. La restitution conserve le point exposant imprimé au sommet supérieur droit, signalé séparément comme lapsus du tirage.; Proposed edit: On a donc un diagramme commutatif à isomorphisme près: =\allowbreak > On a donc un diagramme commutatif à isomorphisme près: \textbackslash{}[\textbackslash{}vcenter\{ \textbackslash{}xymatrix@C=\allowbreak 2.5cm@R=\allowbreak 1.5cm\{ (C/H)\textasciicircum{}\{\textbackslash{}widehat\{\textbackslash{}phantom\{,\}\}\} \textbackslash{}ar[r]\textasciicircum{}\{(u/H)\_\{!\}\} \textbackslash{}ar[d]\_\{(j\_H)\_\{!\}\} \& (C'/u\textasciicircum{}\{\textbackslash{}cdot\}H)\textasciicircum{}\{\textbackslash{}widehat\{\textbackslash{}phantom\{,\}\}\} \textbackslash{}ar[d]\_\{(j\_\{u\_\{!\}H\})\_\{!\}\}\textbackslash{}\textbackslash{} \textbackslash{}widehat\{C\} \textbackslash{}ar[r]\textasciicircum{}\{u\_\{!\}\} \& \textbackslash{}widehat\{C'\}\textbackslash{}quad, \}\}\textbackslash{}leqno\{(5.13.2)\} \textbackslash{}]\par
a fix from one model only\par
\par\noindent\hypertarget{DP0008}{}\textbf{DP0008 -- I, p. 40 (one\_model\_fix)}\par
Démonstration (iv), notation for the arrows of E\par
\textbf{Claude Opus 5.5.} nothing reported here\par
\textbf{GPT-6.1 Sol.} Place: Démonstration (iv), notation for the arrows of E; Reported print: u\textbackslash{}in f\textbackslash{}ell\textbackslash{},E; Source TeX: u\textbackslash{}in F\textbackslash{}ell\textbackslash{},E; Note: The printed abbreviation begins with a lowercase f, not a capital F.; Proposed edit: F\textbackslash{}ell\textbackslash{},E =\allowbreak > f\textbackslash{}ell\textbackslash{},E\par
a fix from one model only\par
\par\noindent\hypertarget{DP0009}{}\textbf{DP0009 -- I, p. 55 (agreed\_fix\_held)}\par
proof of 7.6, 'Donc card Hom(T,X)' | Proof of Proposition 7.6, sentence following « les g\_j »\par
\textbf{Claude Opus 5.5.} Place: proof of 7.6, 'Donc card Hom(T,X)'; Reported print: les g\_j . Donc card Hom(T,X); Source TeX: les \$g\_j\$. donc \$\textbackslash{}card\textbackslash{}Hom(T,X)\$; Note: The print begins a new sentence with a capital 'Donc'.; Proposed edit: donc \$\textbackslash{}card\textbackslash{}Hom(T,X)\$ =\allowbreak > Donc \$\textbackslash{}card\textbackslash{}Hom(T,X)\$\par
\textbf{GPT-6.1 Sol.} Place: Proof of Proposition 7.6, sentence following « les g\_j »; Reported print: Donc card Hom(T,X); Source TeX: donc \$\textbackslash{}card\textbackslash{}Hom(T,X)\$; Note: The print begins the sentence with capital D.; Proposed edit: donc \$\textbackslash{}card\textbackslash{}Hom(T,X)\$ =\allowbreak > Donc \$\textbackslash{}card\textbackslash{}Hom(T,X)\$\par
both models give the same fix, held: a model marks it low confidence\par
\par\noindent\hypertarget{DP0010}{}\textbf{DP0010 -- I, p. 56 (conflict)}\par
Comme \ensuremath{\phi}\_i(X') \ensuremath{\simeq} ... (fibre product) | First paragraph, fiber-product formula following « Comme »\par
\textbf{Claude Opus 5.5.} Place: Comme \ensuremath{\phi}\_i(X') \ensuremath{\simeq} ... (fibre product); Reported print: \ensuremath{\phi}\_i(X\_\ensuremath{\alpha}) \ensuremath{\times}\_\{\ensuremath{\phi}\_\{i(Y)\}\} \ensuremath{\phi}\_i(Y'); Source TeX: \textbackslash{}fprod\{\textbackslash{}varphi(X\_\textbackslash{}alpha)\}\{\textbackslash{}varphi\_i(Y')\}\{\textbackslash{}varphi\_\{i(Y)\}\}; Note: The print has the index i on the first factor, \ensuremath{\phi}\_i(X\_\ensuremath{\alpha}); the TeX drops it.; Proposed edit: \textbackslash{}fprod\{\textbackslash{}varphi(X\_\textbackslash{}alpha)\} =\allowbreak > \textbackslash{}fprod\{\textbackslash{}varphi\_i(X\_\textbackslash{}alpha)\}\par
\textbf{GPT-6.1 Sol.} Place: First paragraph, fiber-product formula following « Comme »; Reported print: \$\textbackslash{}varphi\_i(X')\textbackslash{}simeq \textbackslash{}varphi\_i(X\_\textbackslash{}alpha)\textbackslash{}times\_\{\textbackslash{}varphi\_i(Y)\}\textbackslash{}varphi\_i(Y')\$; Source TeX: \textbackslash{}fprod\{\textbackslash{}varphi(X\_\textbackslash{}alpha)\}\{\textbackslash{}varphi\_i(Y')\}\{\textbackslash{}varphi\_\{i(Y)\}\}; Note: The print has subscript i on the first factor, and only i—not i(Y)—as the subscript on the base.; Proposed edit: \textbackslash{}fprod\{\textbackslash{}varphi(X\_\textbackslash{}alpha)\}\{\textbackslash{}varphi\_i(Y')\}\{\textbackslash{}varphi\_\{i(Y)\}\} =\allowbreak > \textbackslash{}fprod\{\textbackslash{}varphi\_i(X\_\textbackslash{}alpha)\}\{\textbackslash{}varphi\_i(Y')\}\{\textbackslash{}varphi\_i(Y)\}\par
the two models differ at this place\par
\par\noindent\hypertarget{DP0011}{}\textbf{DP0011 -- I, p. 61 (one\_model\_fix)}\par
8.1.2, première phrase\par
\textbf{Claude Opus 5.5.} nothing reported here\par
\textbf{GPT-6.1 Sol.} Place: 8.1.2, première phrase; Reported print: \$u\$ donnés; Source TeX: \$u\$ données; Note: Le tirage porte « donnés », sans e supplémentaire.; Proposed edit: \$u\$ données =\allowbreak > \$u\$ donnés\par
a fix from one model only\par
\par\noindent\hypertarget{DP0012}{}\textbf{DP0012 -- I, p. 63 (one\_model\_fix)}\par
first line\par
\textbf{Claude Opus 5.5.} Place: first line; Reported print: foncteur final sur I' (3.4).; Source TeX: sur \$I'\$ \textbackslash{}hyperref[I.prop3.4]\{3.4\}).; Note: The opening parenthesis is missing, so the closing one is unmatched.; Proposed edit: sur \$I'\$ \textbackslash{}hyperref[I.prop3.4]\{3.4\}) =\allowbreak > sur \$I'\$ (\textbackslash{}hyperref[I.prop3.4]\{3.4\})\par
\textbf{GPT-6.1 Sol.} nothing reported here\par
a fix from one model only\par
\par\noindent\hypertarget{DP0013}{}\textbf{DP0013 -- I, p. 88 (agreed\_fix\_held)}\par
end of 8.6.5 proof | Démonstration de 8.6.5, flèche non identique de la catégorie P\par
\textbf{Claude Opus 5.5.} Place: end of 8.6.5 proof; Reported print: \ensuremath{\pi} telle que \ensuremath{\pi}² =\allowbreak  \ensuremath{\pi} (typewriter small pi); Source TeX: \$\textbackslash{}Pi\$ telle que \$\textbackslash{}Pi\textasciicircum{}2=\allowbreak \textbackslash{}Pi\$; Note: The glyph in the scan looks like a lowercase pi.; Proposed edit: identique \$\textbackslash{}Pi\$ telle que \$\textbackslash{}Pi\textasciicircum{}2=\allowbreak \textbackslash{}Pi\$ =\allowbreak > identique \$\textbackslash{}pi\$ telle que \$\textbackslash{}pi\textasciicircum{}2=\allowbreak \textbackslash{}pi\$\par
\textbf{GPT-6.1 Sol.} Place: Démonstration de 8.6.5, flèche non identique de la catégorie P; Reported print: \$\textbackslash{}pi\$ telle que \$\textbackslash{}pi\textasciicircum{}2=\allowbreak \textbackslash{}pi\$; Source TeX: \$\textbackslash{}Pi\$ telle que \$\textbackslash{}Pi\textasciicircum{}2=\allowbreak \textbackslash{}Pi\$; Note: L'imprimé emploie pi minuscule aux trois endroits.; Proposed edit: \$\textbackslash{}Pi\$ telle que \$\textbackslash{}Pi\textasciicircum{}2=\allowbreak \textbackslash{}Pi\$ =\allowbreak > \$\textbackslash{}pi\$ telle que \$\textbackslash{}pi\textasciicircum{}2=\allowbreak \textbackslash{}pi\$\par
both models give the same fix, held: a model marks it low confidence\par
\par\noindent\hypertarget{DP0014}{}\textbf{DP0014 -- I, p. 88 (one\_model\_fix)}\par
Démonstration de 8.6.5, système inductif indexé par les entiers\par
\textbf{Claude Opus 5.5.} nothing reported here\par
\textbf{GPT-6.1 Sol.} Place: Démonstration de 8.6.5, système inductif indexé par les entiers; Reported print: X\textbackslash{}xrightarrow\{p\}X\textbackslash{}xrightarrow\{p\}X\textbackslash{}longrightarrow\textbackslash{}ldots; Source TeX: X\textbackslash{}xrightarrow\{P\}X\textbackslash{}xrightarrow\{P\}X\textbackslash{}longrightarrow\textbackslash{}ldots; Note: Les deux flèches portent le projecteur p minuscule, non P.; Proposed edit: X\textbackslash{}xrightarrow\{P\}X\textbackslash{}xrightarrow\{P\}X =\allowbreak > X\textbackslash{}xrightarrow\{p\}X\textbackslash{}xrightarrow\{p\}X\par
a fix from one model only\par
\par\noindent\hypertarget{DP0015}{}\textbf{DP0015 -- I, p. 98 (one\_model\_fix)}\par
(Comparer aussi IV 7.5.)\par
\textbf{Claude Opus 5.5.} Place: (Comparer aussi IV 7.5.); Reported print: IV 7.5; Source TeX: IV \textbackslash{}hyperref[I.cor7.5]\{7.5\}; Note: The text prints correctly, but the link goes to the Exposé I label I.cor7.5 instead of Exposé IV 7.5. The correct external label for IV 7.5 is not known.; Proposed edit: IV \textbackslash{}hyperref[I.cor7.5]\{7.5\} =\allowbreak > IV 7.5\par
\textbf{GPT-6.1 Sol.} nothing reported here\par
a fix from one model only\par
\par\noindent\hypertarget{DP0016}{}\textbf{DP0016 -- I, p. 102 (one\_model\_fix)}\par
end of Prop. 8.8.3 b)\par
\textbf{Claude Opus 5.5.} Place: end of Prop. 8.8.3 b); Reported print: est cofinal, i.e. satisfaisant aux conditions F 1), F 2) de 8.1.3.; Source TeX: est cofinal. \textbackslash{}ie satisfaisant; Note: The print has a comma; the TeX has a full stop, which breaks the sentence.; Proposed edit: est cofinal. \textbackslash{}ie =\allowbreak > est cofinal, \textbackslash{}ie\par
\textbf{GPT-6.1 Sol.} nothing reported here\par
a fix from one model only\par
\par\noindent\hypertarget{DP0017}{}\textbf{DP0017 -- I, p. 104 (one\_model\_fix)}\par
same sentence, criteria\par
\textbf{Claude Opus 5.5.} Place: same sentence, criteria; Reported print: à l'aide des critères F 1) et F 2) de 8.1.3.; Source TeX: critères (F 1) et (F 2) de; Note: The TeX adds opening parentheses that the print does not have.; Proposed edit: critères (F 1) et (F 2) de =\allowbreak > critères F 1) et F 2) de\par
\textbf{GPT-6.1 Sol.} nothing reported here\par
a fix from one model only\par
\par\noindent\hypertarget{DP0018}{}\textbf{DP0018 -- I, p. 110 (one\_model\_fix)}\par
formula (8.9.1.1)\par
\textbf{Claude Opus 5.5.} Place: formula (8.9.1.1); Reported print: Φ : J\ensuremath{\times}I \ensuremath{\to} C (the typewriter's capital Phi, the letter later called Φ (8.9.1.1)); Source TeX: \textbackslash{}phi:J\textbackslash{}times I\textbackslash{}longrightarrow; Note: The indexed functor is the capital Φ: the corollary that follows refers to it as "Φ (8.9.1.1)". The TeX's lowercase \textbackslash{}phi is another letter.; Proposed edit: \textbackslash{}phi:J\textbackslash{}times I =\allowbreak > \textbackslash{}Phi:J\textbackslash{}times I\par
\textbf{GPT-6.1 Sol.} nothing reported here\par
a fix from one model only\par
\par\noindent\hypertarget{DP0019}{}\textbf{DP0019 -- I, p. 110 (one\_model\_fix)}\par
Démonstration, liste des types de limites projectives finies\par
\textbf{Claude Opus 5.5.} nothing reported here\par
\textbf{GPT-6.1 Sol.} Place: Démonstration, liste des types de limites projectives finies; Reported print: vide, deux points isolés, puis le diagramme à trois objets; Source TeX: vide \$, \textbackslash{}ldots ,; Note: The middle diagram is the discrete category on two objects, not an ellipsis.; Proposed edit: vide \$, \textbackslash{}ldots , =\allowbreak > vide \$, \textbackslash{}bigdot\textbackslash{}quad\textbackslash{}bigdot ,\par
a fix from one model only\par
\par\noindent\hypertarget{DP0020}{}\textbf{DP0020 -- I, p. 116 (one\_model\_fix)}\par
ST), terme initial de la suite\par
\textbf{Claude Opus 5.5.} nothing reported here\par
\textbf{GPT-6.1 Sol.} Place: ST), terme initial de la suite; Reported print: Z\_0=\allowbreak Z; Source TeX: Z\_\textbackslash{}circ=\allowbreak Z; Note: L'indice rond est le zéro initial de la suite indexée par les entiers n\ensuremath{\geq}0, non un symbole de composition.; Proposed edit: Z\_\textbackslash{}circ=\allowbreak Z =\allowbreak > Z\_0=\allowbreak Z\par
a fix from one model only\par
\par\noindent\hypertarget{DP0021}{}\textbf{DP0021 -- I, p. 126 (one\_model\_fix)}\par
(8.11.5.3), both occurrences of X\par
\textbf{Claude Opus 5.5.} nothing reported here\par
\textbf{GPT-6.1 Sol.} Place: (8.11.5.3), both occurrences of X; Reported print: \textbackslash{}Hom\_\{\textbackslash{}Pro(C)\}(g(\textbackslash{}underline\{Y\}'),X)\textbackslash{}simeq\textbackslash{}Hom\_\{\textbackslash{}Pro(C')\}(\textbackslash{}underline\{Y\}'),f(X)); Source TeX: \textbackslash{}Hom\_\{\textbackslash{}Pro(C)\}(g(\ensuremath{𝒴}'),\ensuremath{𝒳})\textbackslash{}simeq \textbackslash{}Hom\_\{\textbackslash{}Pro(C')\}(\ensuremath{𝒴}'\textbackslash{}sisi\{)\}\{\},f(\ensuremath{𝒳})); Note: Both X's are plain in the print, unlike the underlined Y'. The TeX substitutes the distinct script symbol \ensuremath{𝒳}.; Proposed edit: \textbackslash{}Hom\_\{\textbackslash{}Pro(C)\}(g(\ensuremath{𝒴}'),\ensuremath{𝒳})\textbackslash{}simeq \textbackslash{}Hom\_\{\textbackslash{}Pro(C')\}(\ensuremath{𝒴}'\textbackslash{}sisi\{)\}\{\},f(\ensuremath{𝒳})) =\allowbreak > \textbackslash{}Hom\_\{\textbackslash{}Pro(C)\}(g(\ensuremath{𝒴}'),X)\textbackslash{}simeq \textbackslash{}Hom\_\{\textbackslash{}Pro(C')\}(\ensuremath{𝒴}'\textbackslash{}sisi\{)\}\{\},f(X))\par
a fix from one model only\par
\par\noindent\hypertarget{DP0022}{}\textbf{DP0022 -- I, p. 131 (agreed\_fix\_held)}\par
Démonstration c), last sentence | Démonstration, c), dernière phrase\par
\textbf{Claude Opus 5.5.} Place: Démonstration c), last sentence; Reported print: sous-catégories pleines est petit (the s is overtyped); Source TeX: sous-catégories pleine est petit; Note: The print shows 'pleine' followed by an overtyped blot that reads as a final s; the plural is also the grammatical reading.; Proposed edit: sous-catégories pleine est =\allowbreak > sous-catégories pleines est\par
\textbf{GPT-6.1 Sol.} Place: Démonstration, c), dernière phrase; Reported print: sous-catégories pleines; Source TeX: sous-catégories pleine; Note: L'imprimé porte le pluriel « pleines ».; Proposed edit: sous-catégories pleine =\allowbreak > sous-catégories pleines\par
both models give the same fix, held: a model marks it low confidence\par
\par\noindent\hypertarget{DP0023}{}\textbf{DP0023 -- I, p. 134 (one\_model\_fix)}\par
Exemples 8.12.9, référence au site, II 3.0.2\par
\textbf{Claude Opus 5.5.} nothing reported here\par
\textbf{GPT-6.1 Sol.} Place: Exemples 8.12.9, référence au site, II 3.0.2; Reported print: (II 3.0.2)).; Source TeX: (II (3.0.2)).; Note: Le TeX ajoute une parenthèse ouvrante avant 3.0.2 et laisse ainsi ouverte l'incise commencée par « (ou plus généralement ». L'imprimé ferme correctement cette incise.; Proposed edit: (II (3.0.2)). =\allowbreak > (II 3.0.2)).\par
a fix from one model only\par
\par\noindent\hypertarget{DP0024}{}\textbf{DP0024 -- I, p. 136 (one\_model\_fix)}\par
8.13.2, démonstration b), « un cardinal ... »\par
\textbf{Claude Opus 5.5.} nothing reported here\par
\textbf{GPT-6.1 Sol.} Place: 8.13.2, démonstration b), « un cardinal ... »; Reported print: un cardinal \$\textbackslash{}pi\$; Source TeX: un cardinal \$\textbackslash{}Pi\$; Note: Différence récurrente de casse du pi manuscrit. Ce correctif donne un cas représentatif; les autres occurrences et trois exemples sont signalés au coordinateur.; Proposed edit: un cardinal \$\textbackslash{}Pi\$ =\allowbreak > un cardinal \$\textbackslash{}pi\$\par
a fix from one model only\par
\par\noindent\hypertarget{DP0025}{}\textbf{DP0025 -- I, p. 141 (one\_model\_fix)}\par
9.3.1, indice de E ; notation du cardinal dans ces pages\par
\textbf{Claude Opus 5.5.} nothing reported here\par
\textbf{GPT-6.1 Sol.} Place: 9.3.1, indice de E ; notation du cardinal dans ces pages; Reported print: \$E\_\textbackslash{}pi\$; Source TeX: \$E\_\{\textbackslash{}Pi\}\$; Note: Le cardinal est noté \ensuremath{\pi} minuscule dans le tirage, comme dans les passages où le TeX emploie déjà \textbackslash{}pi. Cette différence répétée est regroupée ici ; voir la convention signalée au coordinateur.; Proposed edit: E\_\{\textbackslash{}Pi\} =\allowbreak > E\_\{\textbackslash{}pi\}\par
a fix from one model only\par
\par\noindent\hypertarget{DP0026}{}\textbf{DP0026 -- I, p. 157 (conflict)}\par
proof of Cor. 9.19, 'ce dernier' | Proof of Corollaire 9.19, immediately before « dernier »\par
\textbf{Claude Opus 5.5.} Place: proof of Cor. 9.19, 'ce dernier'; Reported print: ce dernier (two words); Source TeX: ce\%\% dernier; Note: The \% comment swallows the end of line, so the edition prints 'cedernier' as one word.; Proposed edit: montrer que ce\%\% dernier =\allowbreak > montrer que ce dernier\par
\textbf{GPT-6.1 Sol.} Place: Proof of Corollaire 9.19, immediately before « dernier »; Reported print: que de dernier est; Source TeX: que ce\%\% dernier est; Proposed correction: que ce dernier est; Note: The print has « de » for « ce ». The TeX silently corrects it.\par
the two models differ at this place: one model fixes the TeX, the other calls it a misprint of the print\par
\par\noindent\hypertarget{DP0027}{}\textbf{DP0027 -- I, p. 159 (one\_model\_fix)}\par
Exercice 9.20.1 c), 'ensembles d'indices'\par
\textbf{Claude Opus 5.5.} Place: Exercice 9.20.1 c), 'ensembles d'indices'; Reported print: ensemble d'indices (two words); Source TeX: ensembles\%\%  d'indices; Note: The \% comment plus TeX skipping the leading space of the next line make the edition print 'ensemblesd'indices' as one word.; Proposed edit: ensembles\%\% =\allowbreak > ensembles\par
\textbf{GPT-6.1 Sol.} nothing reported here\par
a fix from one model only\par
\par\noindent\hypertarget{DP0028}{}\textbf{DP0028 -- I, p. 160 (agreed\_fix\_held)}\par
(9.21.2), subscript of L | Formula (9.21.2), subscript of L\par
\textbf{Claude Opus 5.5.} Place: (9.21.2), subscript of L; Reported print: L\_\{\ensuremath{\pi}'\_o\} (same letter \ensuremath{\pi} as elsewhere); Source TeX: L\_\{\textbackslash{}prod'\_\{\textbackslash{}circ\}\}; Note: The TeX uses the product operator \textbackslash{}prod here instead of the letter \textbackslash{}Pi used for the cardinal everywhere else, including the same formula.; Proposed edit: L\_\{\textbackslash{}prod'\_\{\textbackslash{}circ\}\} =\allowbreak > L\_\{\textbackslash{}Pi'\_\{\textbackslash{}circ\}\}\par
\textbf{GPT-6.1 Sol.} Place: Formula (9.21.2), subscript of L; Reported print: L\_\{\textbackslash{}Pi'\_\{\textbackslash{}circ\}\}; Source TeX: L\_\{\textbackslash{}prod'\_\{\textbackslash{}circ\}\}; Note: The print has the cardinal \ensuremath{\Pi}′\_\ensuremath{\circ}, as in the preceding sentence, not a product sign.; Proposed edit: L\_\{\textbackslash{}prod'\_\{\textbackslash{}circ\}\} =\allowbreak > L\_\{\textbackslash{}Pi'\_\{\textbackslash{}circ\}\}\par
both models give the same fix, held: a model marks it low confidence\par
\par\noindent\hypertarget{DP0029}{}\textbf{DP0029 -- I, p. 163 (one\_model\_fix)}\par
«Considérons l'ensemble ordonné produit I =\allowbreak  \ensuremath{\prod}\_\ensuremath{\alpha} I\_\ensuremath{\alpha}»\par
\textbf{Claude Opus 5.5.} Place: «Considérons l'ensemble ordonné produit I =\allowbreak  \ensuremath{\prod}\_\ensuremath{\alpha} I\_\ensuremath{\alpha}»; Reported print: I =\allowbreak  \ensuremath{\prod}\_\ensuremath{\alpha} I\_\ensuremath{\alpha} (product sign, \ensuremath{\alpha} set beneath it); Source TeX: \$I=\allowbreak \textbackslash{}Pi\_\{\textbackslash{}alpha\} I\_\{\textbackslash{}alpha\}\$; Note: The print has the product operator. The TeX uses the Greek letter \textbackslash{}Pi, which on these pages is the name of the cardinal \ensuremath{\Pi}.; Proposed edit: \$I=\allowbreak \textbackslash{}Pi\_\{\textbackslash{}alpha\} I\_\{\textbackslash{}alpha\}\$ =\allowbreak > \$I=\allowbreak \textbackslash{}prod\_\{\textbackslash{}alpha\} I\_\{\textbackslash{}alpha\}\$\par
\textbf{GPT-6.1 Sol.} nothing reported here\par
a fix from one model only\par
\par\noindent\hypertarget{DP0030}{}\textbf{DP0030 -- I, p. 165 (cut)}\par
first line of the page's TeX\par
\textbf{Claude Opus 5.5.} Place: first line of the page's TeX; Reported print: «(9.21.12).» is printed at the foot of p. 164, ending a); p. 165 opens with «b) Pour tout i \ensuremath{\in} I»; Source TeX: (9.21.12).; Note: Source line 6861 «(9.21.12).» belongs to p. 164, whose TeX ends with «déduits de». Page 165 really begins at «\textbackslash{}item Pour tout».\par
\textbf{GPT-6.1 Sol.} nothing reported here\par
a model reports that the automatic cut gives lines of this page to its neighbour, or the neighbour's lines to this page: check the lines it names against the scan of the page that prints them (that page's checker did not have their TeX)\par
\par\noindent\hypertarget{DP0031}{}\textbf{DP0031 -- I, p. 165 (cut)}\par
end of c)\par
\textbf{Claude Opus 5.5.} Place: end of c); Reported print: où les flèches verticales sont celles déduites de (9.21.12).; Source TeX: où les flèches verticales sont celles déduites de; Note: «(9.21.12).» is printed on p. 165 but falls outside this page's cut; it is presumably the next source line (6887), now given to p. 166. Page 165 really ends after «(9.21.12).».\par
\textbf{GPT-6.1 Sol.} nothing reported here\par
a model reports that the automatic cut gives lines of this page to its neighbour, or the neighbour's lines to this page: check the lines it names against the scan of the page that prints them (that page's checker did not have their TeX)\par
\par\noindent\hypertarget{DP0032}{}\textbf{DP0032 -- I, p. 165 (cut)}\par
Dernière ligne, après le second diagramme\par
\textbf{Claude Opus 5.5.} nothing reported here\par
\textbf{GPT-6.1 Sol.} Place: Dernière ligne, après le second diagramme; Reported print: où les flèches verticales sont celles déduites de (9.21.12).; Source TeX: où les flèches verticales sont celles déduites de; Note: L'imprimé comprend encore la référence finale « (9.21.12). », tandis que l'extrait fourni s'arrête après « de ». Vérifier le début du lot p.166 et prolonger la coupe de p.165 jusqu'à cette référence, sans l'ajouter en double au source.\par
a model reports that the automatic cut gives lines of this page to its neighbour, or the neighbour's lines to this page: check the lines it names against the scan of the page that prints them (that page's checker did not have their TeX)\par
\par\noindent\hypertarget{DP0033}{}\textbf{DP0033 -- I, p. 166 (cut)}\par
Début de la coupe, avant « On notera d'ailleurs »\par
\textbf{Claude Opus 5.5.} nothing reported here\par
\textbf{GPT-6.1 Sol.} Place: Début de la coupe, avant « On notera d'ailleurs »; Reported print: La page commence par « On notera d'ailleurs ».; Source TeX: (9.21.12). \textbackslash{}end\{enumerate\} \textbackslash{}end\{enonce*\}; Note: Ce fragment appartient à la page précédente, 165. Ce chevauchement initial n'est pas annoncé dans le commentaire de coupe. Corriger le découpage, sans supprimer ni déplacer le texte dans la source.\par
a model reports that the automatic cut gives lines of this page to its neighbour, or the neighbour's lines to this page: check the lines it names against the scan of the page that prints them (that page's checker did not have their TeX)\par
\par\noindent\hypertarget{DP0034}{}\textbf{DP0034 -- I, p. 168 (conflict)}\par
9.21.14 c) | 9.21.14 c), variable quantifiée\par
\textbf{Claude Opus 5.5.} Place: 9.21.14 c); Reported print: pour tout \ensuremath{\alpha} \ensuremath{\in} ob B, X\_J(\ensuremath{\alpha}) =\allowbreak  …; Source TeX: Pour tout \$a\textbackslash{}in \textbackslash{}ob B\$; Note: The print has the Greek \ensuremath{\alpha}, not the Latin a. The print also has a lower-case 'pour'.; Proposed edit: Pour tout \$a\textbackslash{}in \textbackslash{}ob B\$ =\allowbreak > pour tout \$\textbackslash{}alpha\textbackslash{}in \textbackslash{}ob B\$\par
\textbf{GPT-6.1 Sol.} Place: 9.21.14 c), variable quantifiée; Reported print: \textbackslash{}alpha\textbackslash{}in \textbackslash{}ob B; Source TeX: \$a\textbackslash{}in \textbackslash{}ob B\$; Note: Le scan porte un alpha grec, repris dans X\_J(\ensuremath{\alpha}) et E\_\ensuremath{\alpha}, et non un a latin.; Proposed edit: \$a\textbackslash{}in \textbackslash{}ob B\$ =\allowbreak > \$\textbackslash{}alpha\textbackslash{}in \textbackslash{}ob B\$\par
the two models differ at this place\par
\par\noindent\hypertarget{DP0035}{}\textbf{DP0035 -- I, p. 171 (one\_model\_fix)}\par
Cardinal notation throughout pp. 171–175; representative occurrence in the opening paragraph\par
\textbf{Claude Opus 5.5.} nothing reported here\par
\textbf{GPT-6.1 Sol.} Place: Cardinal notation throughout pp. 171–175; representative occurrence in the opening paragraph; Reported print: \$\textbackslash{}card J(S)\textbackslash{}leq \textbackslash{}pi\$; Source TeX: \$\textbackslash{}card J(S)\textbackslash{}leq \textbackslash{}Pi\$; Note: The print uses lowercase pi, whereas the TeX systematically uses uppercase Pi. Reported once as a convention; the coordinator gives three examples and the required global action.; Proposed edit: \textbackslash{}card J(S)\textbackslash{}leq \textbackslash{}Pi =\allowbreak > \textbackslash{}card J(S)\textbackslash{}leq \textbackslash{}pi\par
a fix from one model only\par
\par\noindent\hypertarget{DP0036}{}\textbf{DP0036 -- I, p. 175 (cut)}\par
first line of the page's TeX\par
\textbf{Claude Opus 5.5.} Place: first line of the page's TeX; Reported print: « comme en 9.21.15. » is printed at the foot of p. 174; Source TeX: \textbackslash{}hyperref[I.sec9.21.15]\{9.21.15\}.; Note: This line belongs to p. 174, whose TeX ends at « comme en ». Printed p. 175 begins with Théorème 9.22.\par
\textbf{GPT-6.1 Sol.} nothing reported here\par
a model reports that the automatic cut gives lines of this page to its neighbour, or the neighbour's lines to this page: check the lines it names against the scan of the page that prints them (that page's checker did not have their TeX)\par
\par\noindent\hypertarget{DP0037}{}\textbf{DP0037 -- I, p. 175 (cut)}\par
First line of the supplied TeX cut\par
\textbf{Claude Opus 5.5.} nothing reported here\par
\textbf{GPT-6.1 Sol.} Place: First line of the supplied TeX cut; Reported print: Page 175 begins with « Théorème 9.22 »; « comme en 9.21.15. » ends page 174.; Source TeX: \textbackslash{}hyperref[I.sec9.21.15]\{9.21.15\}.; Note: The initial reference belongs to the end of page 174, not page 175. This is not covered by the cut's comment, which only allows uncertainty at the bottom of page 175. No text should be moved as a content correction.\par
a model reports that the automatic cut gives lines of this page to its neighbour, or the neighbour's lines to this page: check the lines it names against the scan of the page that prints them (that page's checker did not have their TeX)\par
\par\noindent\hypertarget{DP0038}{}\textbf{DP0038 -- I, p. 181 (one\_model\_fix)}\par
diagram of implications, arrow from 'épimorphique universelle'\par
\textbf{Claude Opus 5.5.} Place: diagram of implications, arrow from 'épimorphique universelle'; Reported print: the arrow from 'épimorphique universelle' points to 'épimorphique'; Source TeX: \textbackslash{}txt\{épimorphique \textbackslash{}\textbackslash{}universelle\} \textbackslash{}ar@\{=\allowbreak >\}[llld]; Note: 'épimorphique universelle' sits in column 5 and 'épimorphique' in column 1 of the last row. [llld] ends in the empty column 2, so the arrow misses its target node.; Proposed edit: \textbackslash{}ar@\{=\allowbreak >\}[llld] =\allowbreak > \textbackslash{}ar@\{=\allowbreak >\}[lllld]\par
\textbf{GPT-6.1 Sol.} nothing reported here\par
a fix from one model only\par
\par\noindent\hypertarget{DP0039}{}\textbf{DP0039 -- I, p. 181 (cut)}\par
Dernière ligne, coupure entre les pages 181 et 182\par
\textbf{Claude Opus 5.5.} nothing reported here\par
\textbf{GPT-6.1 Sol.} Place: Dernière ligne, coupure entre les pages 181 et 182; Reported print: si la famille réduite à f est monomorphique; Source TeX: si la famille réduite à \$f\$; Note: Les mots « est monomorphique », fournis au début du fragment 182 (l. 7522), terminent en réalité la page 181. La page 182 commence par « (resp. monomorphique stricte, …) ». Il s’agit d’un problème de coupe, non d’une omission du source.\par
a model reports that the automatic cut gives lines of this page to its neighbour, or the neighbour's lines to this page: check the lines it names against the scan of the page that prints them (that page's checker did not have their TeX)\par
\par\noindent\hypertarget{DP0040}{}\textbf{DP0040 -- I, p. 185 (cut)}\par
Note (*) au bas de la page, coupure 184–185\par
\textbf{Claude Opus 5.5.} nothing reported here\par
\textbf{GPT-6.1 Sol.} Place: Note (*) au bas de la page, coupure 184–185; Reported print: Nous reproduisons ici, avec son accord, des papiers secrets de N. BOURBAKI. Les références de ce texte se rapportent à son savant ouvrage.; Source TeX: \textbackslash{}footnotetext[1]\{Nous reproduisons ici, avec son accord, des papiers secrets de N. \textbackslash{}textsc\{Bourbaki\}. Les références de ce texte se rapportent à son savant ouvrage.\}; Note: Cette note est imprimée page 185, mais son TeX est fourni seulement dans le fragment 184, l. 7634–7636, hors des lignes de chevauchement annoncées. La page 184 finit après la bibliographie; la page 185 commence au titre « II. Appendice : Univers ». Le texte de la note concorde.\par
a model reports that the automatic cut gives lines of this page to its neighbour, or the neighbour's lines to this page: check the lines it names against the scan of the page that prints them (that page's checker did not have their TeX)\par
\par\noindent\hypertarget{DP0041}{}\textbf{DP0041 -- I, p. 194 (one\_model\_fix)}\par
Premier paragraphe, les deux petits univers\par
\textbf{Claude Opus 5.5.} nothing reported here\par
\textbf{GPT-6.1 Sol.} Place: Premier paragraphe, les deux petits univers; Reported print: U\_0 et U\_1; Source TeX: U\_\{\textbackslash{}circ\} et U\_1; Note: Le premier sous-indice est le chiffre zéro, non le symbole cercle.; Proposed edit: \$U\_\{\textbackslash{}circ\}\$ =\allowbreak > \$U\_0\$\par
a fix from one model only\par
\par\noindent\hypertarget{DP0042}{}\textbf{DP0042 -- I, p. 196 (cut)}\par
End of the final sentence\par
\textbf{Claude Opus 5.5.} nothing reported here\par
\textbf{GPT-6.1 Sol.} Place: End of the final sentence; Reported print: tout segment de I distinct de I a un cardinal < c.; Source TeX: \textbackslash{}emph\{tout segment de \$I\$ distinct de \$I\$ a un cardinal; Note: Page 196 includes the final inequality and full stop. Its ending `\$<\textbackslash{}sheaf\{\textbackslash{}sheaf\{c\}\}\$.\}` is instead at the beginning of the supplied p.197 cut, without a neighbouring-page allowance. Page 197 begins « Nous noterons ». Recut only; no content edit.\par
a model reports that the automatic cut gives lines of this page to its neighbour, or the neighbour's lines to this page: check the lines it names against the scan of the page that prints them (that page's checker did not have their TeX)\par
\par\noindent\hypertarget{DP0043}{}\textbf{DP0043 -- I, p. 197 (one\_model\_fix)}\par
Paragraph beginning « On a d'abord »\par
\textbf{Claude Opus 5.5.} nothing reported here\par
\textbf{GPT-6.1 Sol.} Place: Paragraph beginning « On a d'abord »; Reported print: \$A=\allowbreak A\_0\$; Source TeX: \$A=\allowbreak A\_\{\textbackslash{}circ\}\$; Note: The round printed subscript is the zero index of the sequence defined in (2), not a circle symbol.; Proposed edit: \$A=\allowbreak A\_\{\textbackslash{}circ\}\$ =\allowbreak > \$A=\allowbreak A\_0\$\par
a fix from one model only\par
\par\noindent\hypertarget{DP0044}{}\textbf{DP0044 -- I, p. 201 (one\_model\_fix)}\par
Remarque sur le cardinal c(U), end of the first sentence\par
\textbf{Claude Opus 5.5.} Place: Remarque sur le cardinal c(U), end of the first sentence; Reported print: défini par (1), i.e.; Source TeX: par (1). \textbackslash{}ie; Note: The print has a comma before i.e.; the TeX has a full stop.; Proposed edit: par (1). \textbackslash{}ie =\allowbreak > par (1), \textbackslash{}ie\par
\textbf{GPT-6.1 Sol.} nothing reported here\par
a fix from one model only\par
\par\noindent\hypertarget{DP0045}{}\textbf{DP0045 -- I, p. 202 (conflict)}\par
list of symbols in the 'mots significatifs' | Liste des symboles servant à former les mots significatifs\par
\textbf{Claude Opus 5.5.} Place: list of symbols in the 'mots significatifs'; Reported print: avec les symboles \ensuremath{\varnothing} , x\_\ensuremath{\alpha} , \{ , \} et , . Chacun; Source TeX: \$\textbackslash{}\{,\textbackslash{}\}\$ \%\% etc Chacun; Note: The print lists the comma as the fourth symbol ('et ,'), which gives card(I)+4. The TeX has 'etc' and loses the comma symbol.; Proposed edit: \$\textbackslash{}\{,\textbackslash{}\}\$ \%\% etc Chacun =\allowbreak > \$\textbackslash{}\{\$, \$\textbackslash{}\}\$ et \$,\$. Chacun\par
\textbf{GPT-6.1 Sol.} Place: Liste des symboles servant à former les mots significatifs; Reported print: et , . Chacun; Source TeX: etc Chacun; Note: Le dernier symbole énuméré est la virgule, introduite par « et ». Le scan ne porte pas « etc ».; Proposed edit: etc Chacun =\allowbreak > et \$,\$. Chacun\par
the two models differ at this place\par
\par\noindent\hypertarget{DP0046}{}\textbf{DP0046 -- I, p. 202 (one\_model\_fix)}\par
Définition 5, indice du terme initial\par
\textbf{Claude Opus 5.5.} nothing reported here\par
\textbf{GPT-6.1 Sol.} Place: Définition 5, indice du terme initial; Reported print: \$x\_0=\allowbreak x\$; Source TeX: \$x\_\{\textbackslash{}circ\}=\allowbreak x\$; Note: Le petit caractère rond imprimé représente l'indice zéro de la suite indexée par n \ensuremath{\geq} 0. La transcription par \textbackslash{}circ revient également aux pages 204 et 205 ; cette convention est regroupée ici.; Proposed edit: x\_\{\textbackslash{}circ\}=\allowbreak x =\allowbreak > x\_0=\allowbreak x\par
a fix from one model only\par
\par\noindent\hypertarget{DP0047}{}\textbf{DP0047 -- I, p. 206 (conflict)}\par
after Cor. 1, "Il résulte du cor. l qu'on univers artinien" | "cor. l" | Paragraphe suivant la démonstration du corollaire 1 | Même phrase, « Il résulte du cor. 1… »\par
\textbf{Claude Opus 5.5.} Place: after Cor. 1, "Il résulte du cor. l qu'on univers artinien"; Reported print: Il résulte du cor. 1 qu'un univers artinien est déterminé; Source TeX: Il résulte du cor. l qu'on univers artinien; Note: The print has "qu'un"; the TeX has "qu'on".; Proposed edit: qu'on univers artinien =\allowbreak > qu'un univers artinien || Place: "cor. l"; Reported print: cor. 1 (typewriter digit one); Source TeX: cor. l; Note: The typewriter digit 1 has been transcribed as the letter l. The reference is to Corollaire 1. If the 'old' of the previous finding is applied first, this one still matches.; Proposed edit: du cor. l qu' =\allowbreak > du cor. 1 qu'\par
\textbf{GPT-6.1 Sol.} Place: Paragraphe suivant la démonstration du corollaire 1; Reported print: cor. 1; Source TeX: cor. l; Note: Le renvoi imprimé est au corollaire 1, non à la lettre l.; Proposed edit: cor. l =\allowbreak > cor. 1 || Place: Même phrase, « Il résulte du cor. 1… »; Reported print: qu'un univers; Source TeX: qu'on univers; Note: L'imprimé porte bien « qu'un ».; Proposed edit: qu'on univers =\allowbreak > qu'un univers\par
the two models differ at this place\par
\par\noindent\hypertarget{DP0048}{}\textbf{DP0048 -- I, p. 213 (one\_model\_fix)}\par
l'ensemble ... des ensembles artiniens de type fini\par
\textbf{Claude Opus 5.5.} Place: l'ensemble ... des ensembles artiniens de type fini; Reported print: handwritten letter of capital size, resembling a script/Fraktur G or A; Source TeX: \$\textbackslash{}mathfrak\{a\}\$; Note: The handwritten letter appears to be a capital. The scan does not allow a sure identification.; Proposed edit: \$\textbackslash{}mathfrak\{a\}\$ =\allowbreak > \$\textbackslash{}mathfrak\{A\}\$\par
\textbf{GPT-6.1 Sol.} nothing reported here\par
a fix from one model only\par
\par\noindent\hypertarget{DP0049}{}\textbf{DP0049 -- I, p. 215 (one\_model\_fix)}\par
Deuxième formule affichée, entre les deux quantificateurs\par
\textbf{Claude Opus 5.5.} nothing reported here\par
\textbf{GPT-6.1 Sol.} Place: Deuxième formule affichée, entre les deux quantificateurs; Reported print: (\textbackslash{}forall x\textbackslash{}in M)\textbackslash{}quad\textbackslash{}mbox\{et\}\textbackslash{}quad(\textbackslash{}exists x\textbackslash{}in M); Source TeX: (\textbackslash{}forall x \textbackslash{}in M)\textbackslash{}quav (\textbackslash{}exists x \textbackslash{}in M); Note: Le connecteur imprimé entre les deux expressions est « et ».; Proposed edit: \textbackslash{}quav =\allowbreak > \textbackslash{}quad\textbackslash{}mbox\{et\}\textbackslash{}quad\par
a fix from one model only\par
\par\noindent\hypertarget{DP0050}{}\textbf{DP0050 -- II, p. 219 (one\_model\_fix)}\par
Author credit below the title\par
\textbf{Claude Opus 5.5.} nothing reported here\par
\textbf{GPT-6.1 Sol.} Place: Author credit below the title; Reported print: par J.L. Verdier; Source TeX: J.-L. Verdier; Note: The printed word « par » is missing before the author's name.; Proposed edit: \textbackslash{}begin\{center\} J.-L. Verdier =\allowbreak > \textbackslash{}begin\{center\} par J.-L. Verdier\par
a fix from one model only\par
\par\noindent\hypertarget{DP0051}{}\textbf{DP0051 -- II, p. 224 (conflict)}\par
end of Définition 2.1 | Définition 2.1, dernière phrase\par
\textbf{Claude Opus 5.5.} Place: end of Définition 2.1; Reported print: catégorie des faisceaux à valeur dans U-Ens; Source TeX: à valeurs dans; Note: The print reads «à valeur» in the singular here.; Proposed edit: valeurs dans \$\ensuremath{𝒰}-\textbackslash{}Ens\$ ou encore =\allowbreak > valeur dans \$\ensuremath{𝒰}-\textbackslash{}Ens\$ ou encore\par
\textbf{GPT-6.1 Sol.} Place: Définition 2.1, dernière phrase; Reported print: faisceaux à valeur dans U-Ens; Source TeX: faisceaux à valeurs dans \$\ensuremath{𝒰}-\textbackslash{}Ens\$; Proposed correction: faisceaux à valeurs dans \$\ensuremath{𝒰}-\textbackslash{}Ens\$; Note: Le tirage omet le s de « valeurs » ; le TeX corrige silencieusement cette coquille.\par
the two models differ at this place: one model fixes the TeX, the other calls it a misprint of the print\par
\par\noindent\hypertarget{DP0052}{}\textbf{DP0052 -- II, p. 225 (one\_model\_fix)}\par
Preuve de 2.2, affichage des deux applications de source Hom(X,F\_i)\par
\textbf{Claude Opus 5.5.} nothing reported here\par
\textbf{GPT-6.1 Sol.} Place: Preuve de 2.2, affichage des deux applications de source Hom(X,F\_i); Reported print: et; Source TeX: \textbackslash{}quav; Note: Le tirage relie explicitement les deux applications par « et », que le séparateur employé dans le TeX ne restitue pas.; Proposed edit: \textbackslash{}quav =\allowbreak > \textbackslash{}quad\textbackslash{}text\{et\}\textbackslash{}quad\par
a fix from one model only\par
\par\noindent\hypertarget{DP0053}{}\textbf{DP0053 -- II, p. 230 (one\_model\_fix)}\par
Display defining LF(g), immediately before « qui fait »\par
\textbf{Claude Opus 5.5.} nothing reported here\par
\textbf{GPT-6.1 Sol.} Place: Display defining LF(g), immediately before « qui fait »; Reported print: LF(g):LF(X)\textbackslash{}longrightarrow LF(Y); Source TeX: LF(g):LF(X)\textbackslash{}longrightarrow LF(Y)\{\textbackslash{}quoi\}.; Note: The print continues directly with the relative clause « qui fait de … ». The added full stop incorrectly severs that clause from its antecedent.; Proposed edit: LF(g):LF(X)\textbackslash{}longrightarrow LF(Y)\{\textbackslash{}quoi\}. =\allowbreak > LF(g):LF(X)\textbackslash{}longrightarrow LF(Y)\{\textbackslash{}quoi\}\par
a fix from one model only\par
\par\noindent\hypertarget{DP0054}{}\textbf{DP0054 -- II, p. 232 (one\_model\_fix)}\par
Preuve of 3.1, after the diagram\par
\textbf{Claude Opus 5.5.} Place: Preuve of 3.1, after the diagram; Reported print: Par suite (1.3.3) i' est un isomorphisme; Source TeX: Par suite (\textbackslash{}hyperref[II.def1.3]\{1.3\} 3); Note: The print writes the reference as 1.3.3; the TeX prints '1.3 3'.; Proposed edit: (\textbackslash{}hyperref[II.def1.3]\{1.3\} 3) =\allowbreak > (\textbackslash{}hyperref[II.def1.3]\{1.3\}.3)\par
\textbf{GPT-6.1 Sol.} nothing reported here\par
a fix from one model only\par
\par\noindent\hypertarget{DP0055}{}\textbf{DP0055 -- II, p. 234 (one\_model\_fix)}\par
Preuve 3.2, 3), first paragraph\par
\textbf{Claude Opus 5.5.} Place: Preuve 3.2, 3), first paragraph; Reported print: un crible couvrant Y (3.1.2)); Source TeX: (\textbackslash{}hyperref[II.lem3.1]\{3.1\} 2)); Note: The print writes 3.1.2; the TeX prints '3.1 2)'.; Proposed edit: (\textbackslash{}hyperref[II.lem3.1]\{3.1\} 2)) =\allowbreak > (\textbackslash{}hyperref[II.lem3.1]\{3.1\}.2))\par
\textbf{GPT-6.1 Sol.} nothing reported here\par
a fix from one model only\par
\par\noindent\hypertarget{DP0056}{}\textbf{DP0056 -- II, p. 238 (agreed\_fix\_held)}\par
Preuve de 4.3, proof of 3), after 'relation d'équivalence de préfaisceaux' | Preuve de 4.3, renvoi après « relation d'équivalence de préfaisceaux »\par
\textbf{Claude Opus 5.5.} Place: Preuve de 4.3, proof of 3), after 'relation d'équivalence de préfaisceaux'; Reported print: (4.1 1)); Source TeX: (\textbackslash{}hyperref[II.cor4.1.1]\{4.1.1\})); Note: The print refers to 4.1 1) (assertion 1 of 4.1), as at the other two places on these pages. The TeX prints '(4.1.1))', a reference to Corollary 4.1.1.; Proposed edit: (\textbackslash{}hyperref[II.cor4.1.1]\{4.1.1\})) =\allowbreak > (\textbackslash{}hyperref[II.theo4.1]\{4.1\} 1))\par
\textbf{GPT-6.1 Sol.} Place: Preuve de 4.3, renvoi après « relation d'équivalence de préfaisceaux »; Reported print: (4.1 1)); Source TeX: (\textbackslash{}hyperref[II.cor4.1.1]\{4.1.1\})); Note: Le renvoi imprimé vise 4.1, assertion 1), et non le corollaire 4.1.1.; Proposed edit: \textbackslash{}hyperref[II.cor4.1.1]\{4.1.1\})) =\allowbreak > \textbackslash{}hyperref[II.theo4.1]\{4.1\} 1))\par
both models give the same fix, held: it changes the editors' apparatus (the target of a cross-reference): a \textbackslash{}sisi correction, note, label, cross-reference or citation is the editors' record and is not removed or retargeted on the models' word\par
\par\noindent\hypertarget{DP0057}{}\textbf{DP0057 -- II, p. 242 (one\_model\_fix)}\par
Pour tout objet Y de C et tout morphisme \ensuremath{\alpha} … on a, en vertu de (4.6.1)\par
\textbf{Claude Opus 5.5.} Place: Pour tout objet Y de C et tout morphisme \ensuremath{\alpha} … on a, en vertu de (4.6.1); Reported print: on a, en vertu; Source TeX: on \$a\$, en vertu; Note: The verb « a » is set as a math variable.; Proposed edit: on \$a\$, en vertu =\allowbreak > on a, en vertu\par
\textbf{GPT-6.1 Sol.} nothing reported here\par
a fix from one model only\par
\par\noindent\hypertarget{DP0058}{}\textbf{DP0058 -- II, p. 242 (one\_model\_fix)}\par
Après le diagramme (4.6.2), introduction du crible Q\_1\par
\textbf{Claude Opus 5.5.} nothing reported here\par
\textbf{GPT-6.1 Sol.} Place: Après le diagramme (4.6.2), introduction du crible Q\_1; Reported print: J\_\{Q\_1\}:Q\_1\textbackslash{}hookrightarrow Y; Source TeX: J\_\{Q\_i\}:Q\_1\textbackslash{}hookrightarrow Y; Note: Le sous-indice imprimé est le chiffre 1, non la lettre i.; Proposed edit: J\_\{Q\_i\} =\allowbreak > J\_\{Q\_1\}\par
a fix from one model only\par
\par\noindent\hypertarget{DP0059}{}\textbf{DP0059 -- II, p. 243 (one\_model\_fix)}\par
Définition 4.5, the initial object\par
\textbf{Claude Opus 5.5.} Place: Définition 4.5, the initial object; Reported print: \ensuremath{\varnothing}\_A (slashed zero / empty set); Source TeX: \textbackslash{}phi\_A; Note: The print writes the initial object as \ensuremath{\varnothing}; the TeX uses \textbackslash{}phi throughout (here, in \textbackslash{}phi\_\{A\} below, and \textbackslash{}phi\_\{C\textasciicircum{}\{\textbackslash{}widehat\{\textbackslash{}phantom\{,\}\}\}\} on p. 245).; Proposed edit: objet \$\textbackslash{}phi\_A\$ qui =\allowbreak > objet \$\textbackslash{}emptyset\_A\$ qui\par
\textbf{GPT-6.1 Sol.} nothing reported here\par
a fix from one model only\par
\par\noindent\hypertarget{DP0060}{}\textbf{DP0060 -- II, p. 247 (agreed\_fix\_held)}\par
Preuve of 4.6.3, i) \ensuremath{\Rightarrow} ii) | Corollaire 4.6.3, preuve de i) \ensuremath{\Rightarrow} ii), reference after «la condition b)»\par
\textbf{Claude Opus 5.5.} Place: Preuve of 4.6.3, i) \ensuremath{\Rightarrow} ii); Reported print: la condition b) de 4.6 2) entraîne; Source TeX: la condition b) de \textbackslash{}hyperref[II.cor4.6.2]\{4.6.2\}) entraîne; Note: The print refers to 4.6 2), Proposition 4.6 part 2), which carries the conditions b) and c). The TeX refers to 4.6.2).; Proposed edit: \textbackslash{}hyperref[II.cor4.6.2]\{4.6.2\}) entraîne =\allowbreak > \textbackslash{}hyperref[II.prop4.6]\{4.6\} 2) entraîne\par
\textbf{GPT-6.1 Sol.} Place: Corollaire 4.6.3, preuve de i) \ensuremath{\Rightarrow} ii), reference after «la condition b)»; Reported print: la condition b) de 4.6 2); Source TeX: la condition b) de \textbackslash{}hyperref[II.cor4.6.2]\{4.6.2\}); Note: The print refers to proposition 4.6, item 2), not corollary 4.6.2.; Proposed edit: \textbackslash{}hyperref[II.cor4.6.2]\{4.6.2\}) =\allowbreak > \textbackslash{}hyperref[II.prop4.6]\{4.6\} 2)\par
both models give the same fix, held: it changes the editors' apparatus (the target of a cross-reference): a \textbackslash{}sisi correction, note, label, cross-reference or citation is the editors' record and is not removed or retargeted on the models' word\par
\par\noindent\hypertarget{DP0061}{}\textbf{DP0061 -- II, p. 247 (conflict)}\par
Preuve of 4.6.3, ii) \ensuremath{\Rightarrow} i), end | Corollaire 4.6.3, preuve de ii) \ensuremath{\Rightarrow} i), justification of \ensuremath{\gamma})\par
\textbf{Claude Opus 5.5.} Place: Preuve of 4.6.3, ii) \ensuremath{\Rightarrow} i), end; Reported print: résulte de 4.6.1) 2), d'où i); Source TeX: résulte de \textbackslash{}hyperref[II.prop4.6]\{4.6\} 1) 2), d'où i); Proposed correction: \textbackslash{}hyperref[II.cor4.6.1]\{4.6.1\} 1) 2); Note: The printed reference '4.6.1) 2)' is garbled. The TeX silently reads it as '4.6 1) 2)'. The context, where \ensuremath{\gamma}) means covering by the empty sieve for a strict initial object, rather suggests 4.6.1 1) 2). The coordinator should check.\par
\textbf{GPT-6.1 Sol.} Place: Corollaire 4.6.3, preuve de ii) \ensuremath{\Rightarrow} i), justification of \ensuremath{\gamma}); Reported print: résulte de 4.6.1) 2); Source TeX: \textbackslash{}hyperref[II.prop4.6]\{4.6\} 1) 2); Note: The printed reference is to corollary 4.6.1, item 2), not to items 1) and 2) of proposition 4.6. The replacement retains the printed punctuation.; Proposed edit: \textbackslash{}hyperref[II.prop4.6]\{4.6\} 1) 2) =\allowbreak > \textbackslash{}hyperref[II.cor4.6.1]\{4.6.1\}) 2)\par
the two models differ at this place: one model fixes the TeX, the other calls it a misprint of the print\par
\par\noindent\hypertarget{DP0062}{}\textbf{DP0062 -- II, p. 257 (one\_model\_fix)}\par
Définition 6.1, formule définissant le préfaisceau d'ensembles\par
\textbf{Claude Opus 5.5.} nothing reported here\par
\textbf{GPT-6.1 Sol.} Place: Définition 6.1, formule définissant le préfaisceau d'ensembles; Reported print: \textbackslash{}Hom\_D(S,F(X)),\textbackslash{}quad X\textbackslash{}in\textbackslash{}ob(C); Source TeX: \textbackslash{}Hom\_D(S,F(X))\textbackslash{}quav X\textbackslash{}in \textbackslash{}ob(C); Note: L'imprimé porte une simple virgule, sans le mot « avec » introduit par \textbackslash{}quav.; Proposed edit: \textbackslash{}quav X\textbackslash{}in \textbackslash{}ob(C) =\allowbreak > ,\textbackslash{}quad X\textbackslash{}in \textbackslash{}ob(C)\par
a fix from one model only\par
\par\noindent\hypertarget{DP0063}{}\textbf{DP0063 -- II, p. 258 (agreed\_fix\_held)}\par
Remarques 6.2, 2) | Remarques 6.2, 2), end | Remarques 6.2 (2), formulation récurrente « à valeur »\par
\textbf{Claude Opus 5.5.} Place: Remarques 6.2, 2); Reported print: les faisceaux à valeur dans D; Source TeX: les faisceaux à valeurs dans \$D\$ en; Note: The print has the singular «valeur»; the TeX has the plural. Same change at the other places reported below.; Proposed edit: les faisceaux à valeurs dans \$D\$ en =\allowbreak > les faisceaux à valeur dans \$D\$ en || Place: Remarques 6.2, 2), end; Reported print: en faisceaux à valeur dans D'; Source TeX: faisceaux à valeurs dans \$D'\$; Note: The print has the singular «valeur».; Proposed edit: faisceaux à valeurs dans \$D'\$ =\allowbreak > faisceaux à valeur dans \$D'\$\par
\textbf{GPT-6.1 Sol.} Place: Remarques 6.2 (2), formulation récurrente « à valeur »; Reported print: les faisceaux à valeur dans D en faisceaux à valeur dans D'; Source TeX: les faisceaux à valeurs dans D en faisceaux à valeurs dans D'; Note: Le singulier imprimé est systématiquement pluralisé dans le TeX. Cette normalisation est signalée une seule fois ; les trois passages concernés sont indiqués au coordinateur.; Proposed edit: faisceaux à valeurs dans \$D\$ en   faisceaux à valeurs dans \$D'\$. =\allowbreak > faisceaux à valeur dans \$D\$ en   faisceaux à valeur dans \$D'\$.\par
both models give the same fix, held: a model marks it low confidence\par
\par\noindent\hypertarget{DP0064}{}\textbf{DP0064 -- II, p. 258 (one\_model\_fix)}\par
Proposition 6.3.1\par
\textbf{Claude Opus 5.5.} Place: Proposition 6.3.1; Reported print: sur C à valeur dans U-\ensuremath{\gamma}; Source TeX: \$C\$ à valeurs dans \$\ensuremath{𝒰}-\textbackslash{}gamma\$; Note: The print has the singular «valeur».; Proposed edit: \$C\$ à valeurs dans \$\ensuremath{𝒰}-\textbackslash{}gamma\$ =\allowbreak > \$C\$ à valeur dans \$\ensuremath{𝒰}-\textbackslash{}gamma\$\par
\textbf{GPT-6.1 Sol.} nothing reported here\par
a fix from one model only\par
\par\noindent\hypertarget{DP0065}{}\textbf{DP0065 -- II, p. 259 (one\_model\_fix)}\par
6.3.2\par
\textbf{Claude Opus 5.5.} Place: 6.3.2; Reported print: à valeur dans U-\ensuremath{\gamma}; Source TeX: valeurs dans \$\ensuremath{𝒰}-\textbackslash{}gamma\$; Note: The print has the singular «valeur».; Proposed edit: valeurs dans \$\ensuremath{𝒰}-\textbackslash{}gamma\$ =\allowbreak > valeur dans \$\ensuremath{𝒰}-\textbackslash{}gamma\$\par
\textbf{GPT-6.1 Sol.} nothing reported here\par
a fix from one model only\par
\par\noindent\hypertarget{DP0066}{}\textbf{DP0066 -- II, p. 262 (conflict)}\par
Notation 6.8, la famille de faisceaux de A-module A\_X | Notation 6.8, phrase sur la famille génératrice\par
\textbf{Claude Opus 5.5.} Place: Notation 6.8, la famille de faisceaux de A-module A\_X; Reported print: underlined A: \textbackslash{}underline\{A\}\_X; Source TeX: \$\ensuremath{𝒜}\_X\$; Note: The print has an underlined roman A, the same A\_X just defined; the TeX gives a script \ensuremath{𝒜}. This may be the edition's rendering of underlined letters, as with \ensuremath{𝒰}.; Proposed edit: \$\ensuremath{𝒜}\_X\$ =\allowbreak > \$\textbackslash{}underline\{A\}\_X\$\par
\textbf{GPT-6.1 Sol.} Place: Notation 6.8, phrase sur la famille génératrice; Reported print: faisceaux de A-module; Source TeX: faisceaux de \$A\$-module; Proposed correction: faisceaux de \$A\$-modules; Note: Il manque le s de « modules ». Le TeX conserve le singulier fautif de l'imprimé.\par
the two models differ at this place: one model fixes the TeX, the other calls it a misprint of the print\par
\par\noindent\hypertarget{DP0067}{}\textbf{DP0067 -- III, p. 265 (one\_model\_fix)}\par
Ligne du nom de l’auteur\par
\textbf{Claude Opus 5.5.} nothing reported here\par
\textbf{GPT-6.1 Sol.} Place: Ligne du nom de l’auteur; Reported print: par J.L. Verdier; Source TeX: J.-L. Verdier; Note: Le mot « par » manque devant le nom de l’auteur.; Proposed edit: J.-L. Verdier \textbackslash{}end\{center\} =\allowbreak > par J.-L. Verdier \textbackslash{}end\{center\}\par
a fix from one model only\par
\par\noindent\hypertarget{DP0068}{}\textbf{DP0068 -- III, p. 275 (one\_model\_fix)}\par
De plus..., « Enfin si V est un univers »\par
\textbf{Claude Opus 5.5.} Place: De plus..., « Enfin si V est un univers »; Reported print: V underlined (universe); Source TeX: si \$V\$ est; Note: The universe is underlined in the print, as everywhere else, where the TeX uses \ensuremath{𝒱}.; Proposed edit: si \$V\$ est =\allowbreak > si \$\ensuremath{𝒱}\$ est\par
\textbf{GPT-6.1 Sol.} nothing reported here\par
a fix from one model only\par
\par\noindent\hypertarget{DP0069}{}\textbf{DP0069 -- III, p. 279 (one\_model\_fix)}\par
Proposition 2.3: « où a désigne le foncteur »\par
\textbf{Claude Opus 5.5.} Place: Proposition 2.3: « où a désigne le foncteur »; Reported print: où \ensuremath{\underline{a}} désigne (the symbol a, underlined as in u* =\allowbreak  \ensuremath{\underline{a}}û*i'); Source TeX: où a désigne; Note: The TeX sets the sheaf functor a as a plain word here, but as \$\textbackslash{}underline\{a\}\$ in the formula and in 2) (\$\textbackslash{}underline\{a\}'\$). The underline is hard to tell from the statement's own underlining.; Proposed edit: où a désigne =\allowbreak > où \$\textbackslash{}underline\{a\}\$ désigne\par
\textbf{GPT-6.1 Sol.} nothing reported here\par
a fix from one model only\par
\par\noindent\hypertarget{DP0070}{}\textbf{DP0070 -- III, p. 279 (one\_model\_fix)}\par
Proposition 2.3, definition of u*\par
\textbf{Claude Opus 5.5.} nothing reported here\par
\textbf{GPT-6.1 Sol.} Place: Proposition 2.3, definition of u*; Reported print: u\textasciicircum{}\{\textbackslash{}ast\}:\textbackslash{}widetilde\{C\}'\textbackslash{}rightarrow\textbackslash{}widetilde\{C\}; Source TeX: u\textasciicircum{}\{\textbackslash{}ast\}:\textbackslash{}widehat\{C\}'\textbackslash{}rightarrow\textbackslash{}widehat\{C\}; Note: Both categories bear tildes in the print, not hats.; Proposed edit: u\textasciicircum{}\{\textbackslash{}ast\}:\textbackslash{}widehat\{C\}'\textbackslash{}rightarrow \textbackslash{}widehat\{C\} =\allowbreak > u\textasciicircum{}\{\textbackslash{}ast\}:\textbackslash{}widetilde\{C\}'\textbackslash{}rightarrow \textbackslash{}widetilde\{C\}\par
a fix from one model only\par
\par\noindent\hypertarget{DP0071}{}\textbf{DP0071 -- III, p. 282 (one\_model\_fix)}\par
proof of 2.6\par
\textbf{Claude Opus 5.5.} Place: proof of 2.6; Reported print: propriété générale des foncteurs adjoints (I 5.7.1); Source TeX: (\textbackslash{}hyperrefext[I][I.rem5.7]\{5.7\} 1); Note: The print gives the reference as 'I 5.7.1'; the TeX prints '(I 5.7 1)', which is a different reference form and leaves the parenthesis unclosed.; Proposed edit: (\textbackslash{}hyperrefext[I][I.rem5.7]\{5.7\} 1) =\allowbreak > (\textbackslash{}hyperrefext[I][I.rem5.7]\{5.7.1\})\par
\textbf{GPT-6.1 Sol.} nothing reported here\par
a fix from one model only\par
\par\noindent\hypertarget{DP0072}{}\textbf{DP0072 -- III, p. 285 (conflict)}\par
Corollaire 3.4 ii) | Corollaire 3.4, condition ii), référence après « U-topologie »\par
\textbf{Claude Opus 5.5.} Place: Corollaire 3.4 ii); Reported print: est une U-topologie (I 3.0.2); Source TeX: \$U\$-topologie; Note: This is the universe U, set as \ensuremath{𝒰} everywhere else in the file. Its underlining here cannot be told apart from the statement's underlining.; Proposed edit: \$U\$-topologie =\allowbreak > \$\ensuremath{𝒰}\$-topologie\par
\textbf{GPT-6.1 Sol.} Place: Corollaire 3.4, condition ii), référence après « U-topologie »; Reported print: (I 3.0.2); Source TeX: \{\textbackslash{}rm (I 3.0.2)\}; Proposed correction: \{\textbackslash{}rm (II 3.0.2)\}; Note: La définition d'une U-topologie se trouve en II 3.0.2, non en I 3.0.2. Le TeX conserve le mauvais numéro d'exposé imprimé.\par
the two models differ at this place: one model fixes the TeX, the other calls it a misprint of the print\par
\par\noindent\hypertarget{DP0073}{}\textbf{DP0073 -- III, p. 291 (one\_model\_fix)}\par
Diagramme (4.2.2), flèche horizontale inférieure gauche\par
\textbf{Claude Opus 5.5.} nothing reported here\par
\textbf{GPT-6.1 Sol.} Place: Diagramme (4.2.2), flèche horizontale inférieure gauche; Reported print: i\textasciicircum{}s; Source TeX: i\textasciicircum{}\{S\}; Note: L'exposant imprimé est un s minuscule.; Proposed edit: i\textasciicircum{}\{S\} =\allowbreak > i\textasciicircum{}\{s\}\par
a fix from one model only\par
\par\noindent\hypertarget{DP0074}{}\textbf{DP0074 -- III, p. 291 (one\_model\_fix)}\par
Diagramme (4.2.2), flèche horizontale inférieure droite\par
\textbf{Claude Opus 5.5.} nothing reported here\par
\textbf{GPT-6.1 Sol.} Place: Diagramme (4.2.2), flèche horizontale inférieure droite; Reported print: u\textasciicircum{}s; Source TeX: u\textasciicircum{}S; Note: L'exposant imprimé est un s minuscule.; Proposed edit: u\textasciicircum{}S =\allowbreak > u\textasciicircum{}s\par
a fix from one model only\par
\par\noindent\hypertarget{DP0075}{}\textbf{DP0075 -- III, p. 295 (conflict)}\par
item ii) of the Preuve of 5.4: 'on dénote par i et a' | Preuve de 5.4, ii), parenthèse précédant le diagramme\par
\textbf{Claude Opus 5.5.} Place: item ii) of the Preuve of 5.4: 'on dénote par i et a'; Reported print: on dénotè (accent on the final e; the stroke is small and its direction is unclear); Source TeX: on dénoté; Proposed correction: on dénote; Note: The present tense 'dénote' is meant; the print puts an accent on the final e. The TeX keeps an accented form ('dénoté'). The accent is small in the scan, so this reading is not certain.\par
\textbf{GPT-6.1 Sol.} Place: Preuve de 5.4, ii), parenthèse précédant le diagramme; Reported print: on dénote par; Source TeX: on dénoté par; Note: L'imprimé porte « dénote », sans accent sur le e final.; Proposed edit: on dénoté par =\allowbreak > on dénote par\par
the two models differ at this place: one model fixes the TeX, the other calls it a misprint of the print\par
\par\noindent\hypertarget{DP0076}{}\textbf{DP0076 -- IV, p. 299 (one\_model\_fix)}\par
Titre, ligne des auteurs\par
\textbf{Claude Opus 5.5.} nothing reported here\par
\textbf{GPT-6.1 Sol.} Place: Titre, ligne des auteurs; Reported print: par A. Grothendieck et J. L. Verdier; Source TeX: A. Grothendieck et J.-L. Verdier; Note: Le mot « par » manque devant les noms des auteurs.; Proposed edit: \textbackslash{}begin\{center\} A. Grothendieck et J.-L. Verdier =\allowbreak > \textbackslash{}begin\{center\} par A. Grothendieck et J.-L. Verdier\par
a fix from one model only\par
\par\noindent\hypertarget{DP0077}{}\textbf{DP0077 -- IV, p. 302 (agreed\_fix\_held)}\par
1.1.2 c), reference at the end | 1.1.2 c), référence finale\par
\textbf{Claude Opus 5.5.} Place: 1.1.2 c), reference at the end; Reported print: (I 10.6).; Source TeX: \textbackslash{}sisi\{(\textbackslash{}hyperrefext[I][I.sec10.9]\{10.9\}).\}\{(\textbackslash{}hyperrefext[I][I.sec10.10]\{10.10\}).\}; Note: The print's reading, kept in the first argument of \textbackslash{}sisi, is 'I 10.6', not 10.9. The scan is a little blurred before the 1 but '10.6' is legible. The editors' correction (10.10) is not checked.; Proposed edit: \textbackslash{}sisi\{(\textbackslash{}hyperrefext[I][I.sec10.9]\{10.9\}).\} =\allowbreak > \textbackslash{}sisi\{(\textbackslash{}hyperrefext[I][I.sec10.6]\{10.6\}).\}\par
\textbf{GPT-6.1 Sol.} Place: 1.1.2 c), référence finale; Reported print: (I 10.6).; Source TeX: \textbackslash{}hyperrefext[I][I.sec10.9]\{10.9\}; Note: Le premier argument de \textbackslash{}sisi ne reproduit pas l'imprimé : celui-ci porte 10.6, non 10.9. Le second argument reste inchangé.; Proposed edit: \textbackslash{}hyperrefext[I][I.sec10.9]\{10.9\} =\allowbreak > \textbackslash{}hyperrefext[I][I.sec10.6]\{10.6\}\par
both models give the same fix, held: it changes the editors' apparatus (the target of a cross-reference): a \textbackslash{}sisi correction, note, label, cross-reference or citation is the editors' record and is not removed or retargeted on the models' word\par
\par\noindent\hypertarget{DP0078}{}\textbf{DP0078 -- IV, p. 302 (one\_model\_fix)}\par
1.1.2, références introductives\par
\textbf{Claude Opus 5.5.} nothing reported here\par
\textbf{GPT-6.1 Sol.} Place: 1.1.2, références introductives; Reported print: II 4.8, 4.11; Source TeX: \textbackslash{}hyperrefext[II][II.prop4.8]\{4.8\}, \textbackslash{}hyperrefext[II][II.cor4.11]\{4.11\}; Note: L'imprimé ne répète pas le numéro d'exposé II devant 4.11.; Proposed edit: \textbackslash{}hyperrefext[II][II.cor4.11]\{4.11\} =\allowbreak > \textbackslash{}hyperref[II.cor4.11]\{4.11\}\par
a fix from one model only\par
\par\noindent\hypertarget{DP0079}{}\textbf{DP0079 -- IV, p. 305 (one\_model\_fix)}\par
Fin de 1.2.3, première désignation de la topologie induite\par
\textbf{Claude Opus 5.5.} nothing reported here\par
\textbf{GPT-6.1 Sol.} Place: Fin de 1.2.3, première désignation de la topologie induite; Reported print: La topologie \$\textbackslash{}mathcal\{T\}\$ induite sur \$C\$; Source TeX: La topologie \$\textbackslash{}Cc\$ induite sur \$C\$; Note: La lettre manuscrite est un T calligraphique, non un C.; Proposed edit: La topologie \$\textbackslash{}Cc\$ induite =\allowbreak > La topologie \$\textbackslash{}mathcal\{T\}\$ induite\par
a fix from one model only\par
\par\noindent\hypertarget{DP0080}{}\textbf{DP0080 -- IV, p. 305 (one\_model\_fix)}\par
Fin de 1.2.3, dernière phrase\par
\textbf{Claude Opus 5.5.} nothing reported here\par
\textbf{GPT-6.1 Sol.} Place: Fin de 1.2.3, dernière phrase; Reported print: pour la topologie \$\textbackslash{}mathcal\{T\}\$; Source TeX: pour la topologie \$\textbackslash{}Cc\$; Note: Seconde occurrence du même T calligraphique manuscrit.; Proposed edit: pour la topologie \$\textbackslash{}Cc\$ =\allowbreak > pour la topologie \$\textbackslash{}mathcal\{T\}\$\par
a fix from one model only\par
\par\noindent\hypertarget{DP0081}{}\textbf{DP0081 -- IV, p. 305 (cut)}\par
1.2.4.2, fin de la page\par
\textbf{Claude Opus 5.5.} nothing reported here\par
\textbf{GPT-6.1 Sol.} Place: 1.2.4.2, fin de la page; Reported print: La famille; Source TeX: famille \$(u:X\_\{\textbackslash{}alpha\}\textbackslash{}rightarrow H, (u,\textbackslash{}alpha)\textbackslash{}in; Note: La page imprimée s'arrête après « La famille ». La formule commençant par \$(u:X\_\{\textbackslash{}alpha\}\textbackslash{}rightarrow H\$ appartient à la page 306, alors que le commentaire de coupe ne signale aucun débordement final.\par
a model reports that the automatic cut gives lines of this page to its neighbour, or the neighbour's lines to this page: check the lines it names against the scan of the page that prints them (that page's checker did not have their TeX)\par
\par\noindent\hypertarget{DP0082}{}\textbf{DP0082 -- IV, p. 306 (cut)}\par
start of page\par
\textbf{Claude Opus 5.5.} Place: start of page; Reported print: (u : X\_\ensuremath{\alpha} \ensuremath{\longrightarrow} H, (u,\ensuremath{\alpha}) \ensuremath{\in} I(H)) est épimorphique dans Ẽ.; Source TeX: I(H))\$\textbackslash{}pageoriginale est épimorphique; Note: The printed page begins with '(u : X\_\ensuremath{\alpha} \ensuremath{\to} H, (u,\ensuremath{\alpha}) \ensuremath{\in}'. This text is on source line 318, which the cut gives to p. 305. Page 306 really begins on line 318, at '\$(u:X\_\{\textbackslash{}alpha\}\textbackslash{}rightarrow H, (u,\textbackslash{}alpha)\textbackslash{}in'.\par
\textbf{GPT-6.1 Sol.} nothing reported here\par
a model reports that the automatic cut gives lines of this page to its neighbour, or the neighbour's lines to this page: check the lines it names against the scan of the page that prints them (that page's checker did not have their TeX)\par
\par\noindent\hypertarget{DP0083}{}\textbf{DP0083 -- IV, p. 306 (cut)}\par
Début de page\par
\textbf{Claude Opus 5.5.} nothing reported here\par
\textbf{GPT-6.1 Sol.} Place: Début de page; Reported print: \$(u:X\_\{\textbackslash{}alpha\}\textbackslash{}rightarrow H,(u,\textbackslash{}alpha)\textbackslash{}in I(H))\$ est épimorphique; Source TeX: I(H))\$\textbackslash{}pageoriginale est épimorphique; Note: La page imprimée commence par la formule entière, mais la coupe fournie commence à sa fin. Étendre la coupe avant la ligne 319 pour reprendre le début de cette formule, sans modifier le texte source.\par
a model reports that the automatic cut gives lines of this page to its neighbour, or the neighbour's lines to this page: check the lines it names against the scan of the page that prints them (that page's checker did not have their TeX)\par
\par\noindent\hypertarget{DP0084}{}\textbf{DP0084 -- IV, p. 310 (one\_model\_fix)}\par
Note de bas de page, deuxième renvoi\par
\textbf{Claude Opus 5.5.} nothing reported here\par
\textbf{GPT-6.1 Sol.} Place: Note de bas de page, deuxième renvoi; Reported print: I 8.12.8, 8.12.9; Source TeX: \textbackslash{}hyperrefext[I][I.cor8.12.8]\{8.12.8\}, \textbackslash{}hyperrefext[I][I.exer8.12.9]\{8.12.9\}; Note: Le second hyperrefext ajoute un I que l'imprimé ne répète pas devant 8.12.9.; Proposed edit: \textbackslash{}hyperrefext[I][I.exer8.12.9]\{8.12.9\} =\allowbreak > \textbackslash{}hyperref[I.exer8.12.9]\{8.12.9\}\par
a fix from one model only\par
\par\noindent\hypertarget{DP0085}{}\textbf{DP0085 -- IV, p. 313 (one\_model\_fix)}\par
2.2, the unique object of Ouv(X) when X is empty\par
\textbf{Claude Opus 5.5.} nothing reported here\par
\textbf{GPT-6.1 Sol.} Place: 2.2, the unique object of Ouv(X) when X is empty; Reported print: l'unique objet \ensuremath{\varnothing} de Ouv(X); Source TeX: l'unique objet \$\textbackslash{}phi\$ de \$\textbackslash{}Ouv(X)\$; Note: The print has the empty-set sign, as in Top(\ensuremath{\varnothing}) immediately below, not Greek phi.; Proposed edit: l'unique objet \$\textbackslash{}phi\$ =\allowbreak > l'unique objet \$\ensuremath{\varnothing}\$\par
a fix from one model only\par
\par\noindent\hypertarget{DP0086}{}\textbf{DP0086 -- IV, p. 314 (one\_model\_fix)}\par
2.3, empty subobject of the final object of B\_G\par
\textbf{Claude Opus 5.5.} nothing reported here\par
\textbf{GPT-6.1 Sol.} Place: 2.3, empty subobject of the final object of B\_G; Reported print: faisceau vide \ensuremath{\varnothing}; Source TeX: faisceau vide \$\textbackslash{}phi\$; Note: The printed symbol is the empty-set sign, not Greek phi.; Proposed edit: faisceau vide \$\textbackslash{}phi\$ =\allowbreak > faisceau vide \$\ensuremath{\varnothing}\$\par
a fix from one model only\par
\par\noindent\hypertarget{DP0087}{}\textbf{DP0087 -- IV, p. 315 (cut)}\par
End of the supplied TeX, after the final sentence of 2.4\par
\textbf{Claude Opus 5.5.} nothing reported here\par
\textbf{GPT-6.1 Sol.} Place: End of the supplied TeX, after the final sentence of 2.4; Reported print: The page ends with “en un sens convenable qu'il n'y a pas lieu de préciser ici.”; Source TeX: \textbackslash{}subsection\{«~Gros site~»  et «~gros topos~»  d'un espace topologique. Topos classifiant d'un groupe topologique\}\textbackslash{}footnote[1]\{L'introduction de ces sites et topos est; Note: The entire heading 2.5 and its footnote belong to p. 316, not merely the last source line allowed by the opening comment. Correct the extraction boundary; no TeX content should be moved or deleted.\par
a model reports that the automatic cut gives lines of this page to its neighbour, or the neighbour's lines to this page: check the lines it names against the scan of the page that prints them (that page's checker did not have their TeX)\par
\par\noindent\hypertarget{DP0088}{}\textbf{DP0088 -- IV, p. 316 (cut)}\par
top of page: heading 2.5 and footnote (*)\par
\textbf{Claude Opus 5.5.} Place: top of page: heading 2.5 and footnote (*); Reported print: 2.5. "Gros site" et "gros topos" d'un espace topologique. Topos classifiant d'un groupe topologique (*) … footnote (*) L'introduction de ces sites et topos est due à M. GIRAUD, qui a mis également en évidence leurs avantages sur le "petit" site traditionnel.; Source TeX: traditionnel.\}\textbackslash{}etiquette[2.5]\{IV.subsec2.5\}\textbackslash{}pageoriginale; Note: Page 316 prints the whole heading 2.5 and its footnote. In the TeX both stand before source line 704, in the previous page's chunk; only the footnote's last word 'traditionnel.\}' is here. Page 316 really begins with the heading 2.5, and that text should be checked with this page.\par
\textbf{GPT-6.1 Sol.} nothing reported here\par
a model reports that the automatic cut gives lines of this page to its neighbour, or the neighbour's lines to this page: check the lines it names against the scan of the page that prints them (that page's checker did not have their TeX)\par
\par\noindent\hypertarget{DP0089}{}\textbf{DP0089 -- IV, p. 316 (cut)}\par
Beginning of §2.5 and footnote (*)\par
\textbf{Claude Opus 5.5.} nothing reported here\par
\textbf{GPT-6.1 Sol.} Place: Beginning of §2.5 and footnote (*); Reported print: The page opens with “2.5. « Gros site » et « gros topos » d'un espace topologique. Topos classifiant d'un groupe topologique”, with footnote: “L'introduction de ces sites et topos est due à M. GIRAUD, qui a mis également en évidence leurs avantages sur le « petit » site traditionnel.”; Source TeX: traditionnel.\}\textbackslash{}etiquette[2.5]\{IV.subsec2.5\}\textbackslash{}pageoriginale; Note: The cut omits the heading and almost the entire footnote, which precede source line 704. The remaining “traditionnel.\}” belongs to this page's footnote, not to p. 315. Extend the comparison cut upward; do not insert duplicate text.\par
a model reports that the automatic cut gives lines of this page to its neighbour, or the neighbour's lines to this page: check the lines it names against the scan of the page that prints them (that page's checker did not have their TeX)\par
\par\noindent\hypertarget{DP0090}{}\textbf{DP0090 -- IV, p. 319 (one\_model\_fix)}\par
2.7.1, On suppose le système projectif strict\par
\textbf{Claude Opus 5.5.} Place: 2.7.1, On suppose le système projectif strict; Reported print: strict, i.e. les morphismes; Source TeX: \textbackslash{}emph\{strict\}. \textbackslash{}ie les morphismes; Note: The print has a comma, not a full stop.; Proposed edit: \textbackslash{}emph\{strict\}. \textbackslash{}ie =\allowbreak > \textbackslash{}emph\{strict\}, \textbackslash{}ie\par
\textbf{GPT-6.1 Sol.} nothing reported here\par
a fix from one model only\par
\par\noindent\hypertarget{DP0091}{}\textbf{DP0091 -- IV, p. 325 (conflict)}\par
ligne 4 | Opening paragraph, list of algebraic structures\par
\textbf{Claude Opus 5.5.} Place: ligne 4; Reported print: de comodule sur un anneau, de bigèbre sur un anneau; Source TeX: de comodule sur un coanneau; Note: The print reads «sur un anneau». If «coanneau» is meant as an editorial correction, it should go through \textbackslash{}sisi.; Proposed edit: de comodule sur un coanneau =\allowbreak > de comodule sur un anneau\par
\textbf{GPT-6.1 Sol.} Place: Opening paragraph, list of algebraic structures; Reported print: de comodule sur un anneau; Source TeX: de comodule sur un coanneau; Proposed correction: de comodule sur un coanneau; Note: The print says « anneau » where « coanneau » is intended. The TeX silently corrects it.\par
the two models differ at this place: one model fixes the TeX, the other calls it a misprint of the print\par
\par\noindent\hypertarget{DP0092}{}\textbf{DP0092 -- IV, p. 326 (cut)}\par
Last line; boundary with p. 327\par
\textbf{Claude Opus 5.5.} nothing reported here\par
\textbf{GPT-6.1 Sol.} Place: Last line; boundary with p. 327; Reported print: Dans certains cas, il peut arriver que \$f\$; Source TeX: Dans certains cas, il peut arriver que; Note: The final \$f\$ is printed on p. 326 but occurs at the start of the p. 327 TeX cut. P. 327 begins with « satisfasse ». Recut only; do not insert another \$f\$ into the source.\par
a model reports that the automatic cut gives lines of this page to its neighbour, or the neighbour's lines to this page: check the lines it names against the scan of the page that prints them (that page's checker did not have their TeX)\par
\par\noindent\hypertarget{DP0093}{}\textbf{DP0093 -- IV, p. 327 (agreed\_fix\_held)}\par
footnote (*) | Footnote\par
\textbf{Claude Opus 5.5.} Place: footnote (*); Reported print: Cf. plus bas, dans le cas de faisceaux abéliens.; Source TeX: \textbackslash{}Cf plus bas \textbackslash{}hyperref[IV.subsec14.4]\{14.4\}, dans le cas de faisceaux abéliens.; Note: The print gives no reference number in the footnote; the TeX adds « 14.4 » without marking it as an editorial addition (no \textbackslash{}sisi or \textbackslash{}nde).; Proposed edit: plus bas \textbackslash{}hyperref[IV.subsec14.4]\{14.4\}, =\allowbreak > plus bas,\par
\textbf{GPT-6.1 Sol.} Place: Footnote; Reported print: Cf. plus bas, dans le cas de faisceaux abéliens.; Source TeX: \textbackslash{}Cf plus bas \textbackslash{}hyperref[IV.subsec14.4]\{14.4\}, dans le cas de faisceaux abéliens.; Note: The explicit reference 14.4 is absent from the printed footnote.; Proposed edit: plus bas \textbackslash{}hyperref[IV.subsec14.4]\{14.4\}, =\allowbreak > plus bas,\par
both models give the same fix, held: a deletion in the first or last three lines of the page (the page cut may be off); it changes the editors' apparatus (\textbackslash{}hyperref, the target of a cross-reference): a \textbackslash{}sisi correction, note, label, cross-reference or citation is the editors' record and is not removed or retargeted on the models' word\par
\par\noindent\hypertarget{DP0094}{}\textbf{DP0094 -- IV, p. 334 (conflict)}\par
référence à SGA 1 VI 7.4 B) | 4.1.1, parenthesis following the reference to SGA 1 VI 7.4 B)\par
\textbf{Claude Opus 5.5.} Place: référence à SGA 1 VI 7.4 B); Reported print: (pour \ensuremath{𝒞} =\allowbreak  (Esp)\textasciicircum{}o); Source TeX: (pour \$\textbackslash{}epsilon =\allowbreak  (\textbackslash{}Esp)\textasciicircum{}\{\textbackslash{}circ\}\$); Note: The print has a script C (category), not the Greek epsilon.; Proposed edit: \$\textbackslash{}epsilon =\allowbreak  (\textbackslash{}Esp)\textasciicircum{}\{\textbackslash{}circ\}\$ =\allowbreak > \$\textbackslash{}mathcal\{C\} =\allowbreak  (\textbackslash{}Esp)\textasciicircum{}\{\textbackslash{}circ\}\$\par
\textbf{GPT-6.1 Sol.} Place: 4.1.1, parenthesis following the reference to SGA 1 VI 7.4 B); Reported print: \$\textbackslash{}mathcal\{E\} =\allowbreak  (\textbackslash{}Esp)\textasciicircum{}\{\textbackslash{}circ\}\$; Source TeX: \$\textbackslash{}epsilon =\allowbreak  (\textbackslash{}Esp)\textasciicircum{}\{\textbackslash{}circ\}\$; Note: The printed symbol is a calligraphic capital E, not a Greek epsilon.; Proposed edit: \textbackslash{}epsilon =\allowbreak > \textbackslash{}mathcal\{E\}\par
the two models differ at this place\par
\par\noindent\hypertarget{DP0095}{}\textbf{DP0095 -- IV, p. 335 (one\_model\_fix)}\par
4.1.2, display defining the two components of f\par
\textbf{Claude Opus 5.5.} nothing reported here\par
\textbf{GPT-6.1 Sol.} Place: 4.1.2, display defining the two components of f; Reported print: f\textasciicircum{}\{es\}:X \ensuremath{\to} Y et f\textasciicircum{}\{gr\}:G \ensuremath{\to} H; Source TeX: f\textasciicircum{}\{\textbackslash{}es\}:X\textbackslash{}longrightarrow Y\textbackslash{}quav f\textasciicircum{}\{\textbackslash{}gr\}:G\textbackslash{}longrightarrow H; Note: The print joins the two maps with « et »; the TeX instead uses the comma separator \textbackslash{}quav.; Proposed edit: Y\textbackslash{}quav f\textasciicircum{}\{\textbackslash{}gr\} =\allowbreak > Y\textbackslash{}quad\textbackslash{}text\{et\}\textbackslash{}quad f\textasciicircum{}\{\textbackslash{}gr\}\par
a fix from one model only\par
\par\noindent\hypertarget{DP0096}{}\textbf{DP0096 -- IV, p. 344 (conflict)}\par
formula for the action of G on H\ensuremath{\times}X | 4.5.1, formula for the G-action on H \ensuremath{\times} X\par
\textbf{Claude Opus 5.5.} Place: formula for the action of G on H\ensuremath{\times}X; Reported print: g.(h,x) =\allowbreak  (hg\textasciicircum{}\{-1\},gx); Source TeX: g\textbackslash{}circ (h,x); Note: The print has a dot, not a composition circle.; Proposed edit: g\textbackslash{}circ (h,x) =\allowbreak > g.(h,x)\par
\textbf{GPT-6.1 Sol.} Place: 4.5.1, formula for the G-action on H \ensuremath{\times} X; Reported print: \$g.(h,x)=\allowbreak (hg\textasciicircum{}\{-1\},gx)\$; Source TeX: g\textbackslash{}circ (h,x)=\allowbreak (hg\textasciicircum{}\{-1\}, gx); Note: The print denotes the action by a solid point, not an open circle.; Proposed edit: g\textbackslash{}circ (h,x) =\allowbreak > g\textbackslash{}cdot (h,x)\par
the two models differ at this place\par
\par\noindent\hypertarget{DP0097}{}\textbf{DP0097 -- IV, p. 350 (conflict)}\par
4.7, « que le foncteur \ensuremath{\hat{f}}\_* ou f\_* : F ↦ F\ensuremath{\circ}f de Ĉ dans Ĉ' » | 4.7, definition of a cocontinuous functor\par
\textbf{Claude Opus 5.5.} Place: 4.7, « que le foncteur \ensuremath{\hat{f}}\_* ou f\_* : F ↦ F\ensuremath{\circ}f de Ĉ dans Ĉ' »; Reported print: le foncteur \ensuremath{\hat{f}}\_* ou f\_* : F ↦ F\ensuremath{\circ}f de Ĉ dans Ĉ'; Source TeX: le foncteur \$\textbackslash{}widehat\{f\}\_\{\textbackslash{}ast\}\$ de \$\textbackslash{}widehat\{C\}\$; Note: The TeX drops the printed phrase « ou f\_* : F ↦ F\ensuremath{\circ}f ».; Proposed edit: \$\textbackslash{}widehat\{f\}\_\{\textbackslash{}ast\}\$ de \$\textbackslash{}widehat\{C\}\$ =\allowbreak > \$\textbackslash{}widehat\{f\}\_\{\textbackslash{}ast\}\$ ou \$f\_\{\textbackslash{}ast\}:F\textbackslash{}mapsto F\textbackslash{}circ f\$ de \$\textbackslash{}widehat\{C\}\$\par
\textbf{GPT-6.1 Sol.} Place: 4.7, definition of a cocontinuous functor; Reported print: \textbackslash{}widehat\{f\}\_\{\textbackslash{}ast\}\textbackslash{}text\{ ou \}f\_\{\textbackslash{}ast\}; Source TeX: \textbackslash{}widehat\{f\}\_\{\textbackslash{}ast\}; Note: The alternative notation « ou f\_* » is omitted. The erroneous action formula following it in print is treated separately below.; Proposed edit: le foncteur \$\textbackslash{}widehat\{f\}\_\{\textbackslash{}ast\}\$ =\allowbreak > le foncteur \$\textbackslash{}widehat\{f\}\_\{\textbackslash{}ast\}\$ ou \$f\_\{\textbackslash{}ast\}\$\par
the two models differ at this place\par
\par\noindent\hypertarget{DP0098}{}\textbf{DP0098 -- IV, p. 353 (one\_model\_fix)}\par
4.8, notation des préfaisceaux de cohomologie\par
\textbf{Claude Opus 5.5.} nothing reported here\par
\textbf{GPT-6.1 Sol.} Place: 4.8, notation des préfaisceaux de cohomologie; Reported print: \$\textbackslash{}mathcal\{H\}\textasciicircum{}n(F)\$; Source TeX: \$\textbackslash{}Xx\textasciicircum{}n(F)\$; Note: La lettre manuscrite est un H calligraphique, non un X.; Proposed edit: \textbackslash{}Xx\textasciicircum{}n(F) =\allowbreak > \textbackslash{}mathcal\{H\}\textasciicircum{}n(F)\par
a fix from one model only\par
\par\noindent\hypertarget{DP0099}{}\textbf{DP0099 -- IV, p. 354 (one\_model\_fix)}\par
4.9.1, after 'de \$C\$ dans \$C'\$'\par
\textbf{Claude Opus 5.5.} Place: 4.9.1, after 'de \$C\$ dans \$C'\$'; Reported print: C', i.e. tel que; Source TeX: \$C'\$. \textbackslash{}ie tel que; Note: The print has a comma here; the TeX's full stop wrongly ends the sentence.; Proposed edit: dans \$C'\$. \textbackslash{}ie =\allowbreak > dans \$C'\$, \textbackslash{}ie\par
\textbf{GPT-6.1 Sol.} nothing reported here\par
a fix from one model only\par
\par\noindent\hypertarget{DP0100}{}\textbf{DP0100 -- IV, p. 357 (one\_model\_fix)}\par
4.9.6, line 2\par
\textbf{Claude Opus 5.5.} Place: 4.9.6, line 2; Reported print: cf. 4.7.4); Source TeX: \textbackslash{}cf \textbackslash{}hyperref[IV.subsec4.7]\{4.7\} 4); Note: The print reads 4.7.4; the TeX prints '4.7 4'. If the edition deliberately reads it as item 4 of 4.7, this is the edition's style, but the printed number is 4.7.4.; Proposed edit: \textbackslash{}hyperref[IV.subsec4.7]\{4.7\} 4 =\allowbreak > \textbackslash{}hyperref[IV.subsec4.7]\{4.7.4\}\par
\textbf{GPT-6.1 Sol.} nothing reported here\par
a fix from one model only\par
\par\noindent\hypertarget{DP0101}{}\textbf{DP0101 -- IV, p. 359 (one\_model\_fix)}\par
4.10.2, after (4.10.2.1)\par
\textbf{Claude Opus 5.5.} Place: 4.10.2, after (4.10.2.1); Reported print: puis-|que (hyphen at the line break); Source TeX: puis-que; Note: The typescript's line-break hyphen was kept inside the word.; Proposed edit: puis-que =\allowbreak > puisque\par
\textbf{GPT-6.1 Sol.} nothing reported here\par
a fix from one model only\par
\par\noindent\hypertarget{DP0102}{}\textbf{DP0102 -- IV, p. 363 (one\_model\_fix)}\par
Exercice 4.10.6 a), first occurrence of M\par
\textbf{Claude Opus 5.5.} nothing reported here\par
\textbf{GPT-6.1 Sol.} Place: Exercice 4.10.6 a), first occurrence of M; Reported print: \$\textbackslash{}underline\{M\}\$; Source TeX: \$\ensuremath{\mathcal{M}}\$; Note: The print uses the same underlined M as in the statement and the other hypotheses. The TeX changes this one occurrence to a calligraphic M.; Proposed edit: \$\ensuremath{\mathcal{M}}\$ =\allowbreak > \$\textbackslash{}underline\{M\}\$\par
a fix from one model only\par
\par\noindent\hypertarget{DP0103}{}\textbf{DP0103 -- IV, p. 365 (one\_model\_fix)}\par
5.1, description of the objects of \$E\_\{/X\}\$\par
\textbf{Claude Opus 5.5.} nothing reported here\par
\textbf{GPT-6.1 Sol.} Place: 5.1, description of the objects of \$E\_\{/X\}\$; Reported print: au-dessus de \$X\$; Source TeX: au-dessous de \$X\$; Note: The print says « au-dessus », describing objects over X, not « au-dessous ».; Proposed edit: au-dessous de \$X\$ =\allowbreak > au-dessus de \$X\$\par
a fix from one model only\par
\par\noindent\hypertarget{DP0104}{}\textbf{DP0104 -- IV, p. 374 (one\_model\_fix)}\par
5.8.4, les deux occurrences du symbole du pro-groupe\par
\textbf{Claude Opus 5.5.} nothing reported here\par
\textbf{GPT-6.1 Sol.} Place: 5.8.4, les deux occurrences du symbole du pro-groupe; Reported print: Le pro-groupe est noté \$\textbackslash{}underline\{G\}\$ aux deux occurrences.; Source TeX: \$\ensuremath{𝒢}=\allowbreak (G\_i)\$ ; \$\ensuremath{𝒢}\textbackslash{}rightarrow G\_\{i\_\{\textbackslash{}circ\}\}\$; Note: Le soulignement appartient ici au symbole mathématique du pro-groupe, distinct du groupe G ; ce n'est pas un soulignement d'emphase. Le TeX lui substitue une capitale calligraphique.; Proposed edit: \$\ensuremath{𝒢}=\allowbreak (G\_i)\$, dans le cas particulier où \$H\$ est défini comme «~noyau~»  de \$\ensuremath{𝒢}\textbackslash{}rightarrow G\_\{i\_\{\textbackslash{}circ\}\}\$ =\allowbreak > \$\textbackslash{}underline\{G\}=\allowbreak (G\_i)\$, dans le cas particulier où \$H\$ est défini comme «~noyau~»  de \$\textbackslash{}underline\{G\}\textbackslash{}rightarrow G\_\{i\_\{\textbackslash{}circ\}\}\$\par
a fix from one model only\par
\par\noindent\hypertarget{DP0105}{}\textbf{DP0105 -- IV, p. 377 (conflict)}\par
first line, reference for the localisation morphisms | Opening sentence, reference after « morphismes de localisation »\par
\textbf{Claude Opus 5.5.} Place: first line, reference for the localisation morphisms; Reported print: les morphismes de localisation (5.1); Source TeX: localisation (\textbackslash{}hyperref[IV.subsec5.2]\{5.2\}); Note: The print cites 5.1. If the edition deliberately corrected the reference, it should be marked with \textbackslash{}sisi.; Proposed edit: (\textbackslash{}hyperref[IV.subsec5.2]\{5.2\}) =\allowbreak > (\textbackslash{}hyperref[IV.subsec5.1]\{5.1\})\par
\textbf{GPT-6.1 Sol.} Place: Opening sentence, reference after « morphismes de localisation »; Reported print: (5.1); Source TeX: \textbackslash{}hyperref[IV.subsec5.2]\{5.2\}; Proposed correction: \textbackslash{}hyperref[IV.subsec5.2]\{5.2\}; Note: The print gives 5.1; the TeX apparently silently corrects the reference to 5.2, concerning localization morphisms.\par
the two models differ at this place: one model fixes the TeX, the other calls it a misprint of the print\par
\par\noindent\hypertarget{DP0106}{}\textbf{DP0106 -- IV, p. 378 (one\_model\_fix)}\par
Proposition 5.11, product of Homtop categories, first argument of \textbackslash{}sisi\par
\textbf{Claude Opus 5.5.} nothing reported here\par
\textbf{GPT-6.1 Sol.} Place: Proposition 5.11, product of Homtop categories, first argument of \textbackslash{}sisi; Reported print: \textbackslash{}Hhomtop(F,E\_\{/X\})\textbackslash{}times\textasciicircum{}\{2\}\_\{\textbackslash{}Hhomtop(F,E)\}\textbackslash{}Hhomtop(F,E'); Source TeX: \textbackslash{}fprod\{\textbackslash{}Hhomtop(F,E\_\{/X\})\textasciicircum{}2\}\{\textbackslash{}Hhomtop(F,E')\}\{\textbackslash{}Hhomtop(F,E)\}; Note: The printed 2 belongs to the product sign, as the following NB explicitly confirms; it is not an exponent of the first Homtop category. This finding concerns only the original-reading argument A of \textbackslash{}sisi.; Proposed edit: \textbackslash{}fprod\{\textbackslash{}Hhomtop(F,E\_\{/X\})\textasciicircum{}2\}\{\textbackslash{}Hhomtop(F,E')\}\{\textbackslash{}Hhomtop(F,E)\} =\allowbreak > \textbackslash{}Hhomtop(F,E\_\{/X\})\textbackslash{}times\textasciicircum{}\{2\}\_\{\textbackslash{}Hhomtop(F,E)\}\textbackslash{}Hhomtop(F,E')\par
a fix from one model only\par
\par\noindent\hypertarget{DP0107}{}\textbf{DP0107 -- IV, p. 379 (conflict)}\par
(**) and throughout 5.11.1 (also p. 380) | (**) and recurring functor notation on pp. 379–380\par
\textbf{Claude Opus 5.5.} Place: (**) and throughout 5.11.1 (also p. 380); Reported print: \ensuremath{\psi} (small psi); Source TeX: \textbackslash{}Psi; Note: The typescript glyph looks like a small psi; the TeX uses capital \textbackslash{}Psi everywhere on pp. 379-380. Check this against the scan before changing.; Proposed edit: (\textbackslash{}Psi,u),\textbackslash{}mbox\{ avec \} \textbackslash{}Psi=\allowbreak \textbackslash{}varphi\textbackslash{}circ j =\allowbreak > (\textbackslash{}psi,u),\textbackslash{}mbox\{ avec \} \textbackslash{}psi=\allowbreak \textbackslash{}varphi\textbackslash{}circ j\par
\textbf{GPT-6.1 Sol.} Place: (**) and recurring functor notation on pp. 379–380; Reported print: \textbackslash{}psi; Source TeX: \textbackslash{}Psi; Note: The print uses lowercase psi, whereas the transcription consistently uses capital Psi. This is a representative edit for the recurring convention; see the three examples and global action in coordinator.; Proposed edit: \textbackslash{}Psi=\allowbreak \textbackslash{}varphi\textbackslash{}circ j =\allowbreak > \textbackslash{}psi=\allowbreak \textbackslash{}varphi\textbackslash{}circ j\par
the two models differ at this place\par
\par\noindent\hypertarget{DP0108}{}\textbf{DP0108 -- IV, p. 382 (one\_model\_fix)}\par
paragraph after Lemme 5.11.2, 'interpréter u dans le couple'\par
\textbf{Claude Opus 5.5.} Place: paragraph after Lemme 5.11.2, 'interpréter u dans le couple'; Reported print: le couple (\ensuremath{\psi},u), with the same functor letter as everywhere else (rendered \textbackslash{}Psi); Source TeX: \$u\$ dans le couple \$(\textbackslash{}psi,u)\$; Note: The functor is \textbackslash{}Psi everywhere else in the TeX; here the small \textbackslash{}psi gives another letter.; Proposed edit: \$u\$ dans le couple \$(\textbackslash{}psi,u)\$ =\allowbreak > \$u\$ dans le couple \$(\textbackslash{}Psi,u)\$\par
\textbf{GPT-6.1 Sol.} nothing reported here\par
a fix from one model only\par
\par\noindent\hypertarget{DP0109}{}\textbf{DP0109 -- IV, p. 382 (one\_model\_fix)}\par
same paragraph, 'pour \ensuremath{\phi} resp. (\ensuremath{\psi},u) variable'\par
\textbf{Claude Opus 5.5.} Place: same paragraph, 'pour \ensuremath{\phi} resp. (\ensuremath{\psi},u) variable'; Reported print: resp. (\ensuremath{\psi},u) variable, with the same functor letter; Source TeX: \$(\textbackslash{}psi,u)\$ variable; Note: Same capital/small inconsistency as above.; Proposed edit: \$(\textbackslash{}psi,u)\$ variable =\allowbreak > \$(\textbackslash{}Psi,u)\$ variable\par
\textbf{GPT-6.1 Sol.} nothing reported here\par
a fix from one model only\par
\par\noindent\hypertarget{DP0110}{}\textbf{DP0110 -- IV, p. 385 (one\_model\_fix)}\par
Définition 6.1, notation of the category of points\par
\textbf{Claude Opus 5.5.} Place: Définition 6.1, notation of the category of points; Reported print: on note Point(E) ou Pt(E); Source TeX: on note \$\textbackslash{}Ppoints(E)\$ ou \$\textbackslash{}sheaf\{\textbackslash{}Pt\}(E)\$; Note: The print reads 'Point' (singular). The macro name suggests 'Points'; check what \textbackslash{}Ppoints prints (macros.tex not attached).; Proposed edit: \$\textbackslash{}Ppoints(E)\$ =\allowbreak > \$\textbackslash{}underline\{\textbackslash{}mathrm\{Point\}\}(E)\$\par
\textbf{GPT-6.1 Sol.} nothing reported here\par
a fix from one model only\par
\par\noindent\hypertarget{DP0111}{}\textbf{DP0111 -- IV, p. 394 (one\_model\_fix)}\par
Corollaire 6.7.3\par
\textbf{Claude Opus 5.5.} Place: Corollaire 6.7.3; Reported print: Inversement, si; Source TeX: Inversement. si; Note: The print has a comma; the TeX's full stop breaks the sentence.; Proposed edit: Inversement. si =\allowbreak > Inversement, si\par
\textbf{GPT-6.1 Sol.} nothing reported here\par
a fix from one model only\par
\par\noindent\hypertarget{DP0112}{}\textbf{DP0112 -- IV, p. 401 (one\_model\_fix)}\par
Exercice 6.10, qualification de D\par
\textbf{Claude Opus 5.5.} nothing reported here\par
\textbf{GPT-6.1 Sol.} Place: Exercice 6.10, qualification de D; Reported print: D une U-catégorie (univers U souligné, sans tilde); Source TeX: \$\textbackslash{}sheaf\{\textbackslash{}mathcal\{U\}\}\$-catégorie; Note: The print gives the universe U without a tilde; the sheaf accent introduced by the TeX is absent.; Proposed edit: \textbackslash{}sheaf\{\textbackslash{}mathcal\{U\}\} =\allowbreak > \textbackslash{}mathcal\{U\}\par
a fix from one model only\par
\par\noindent\hypertarget{DP0113}{}\textbf{DP0113 -- IV, p. 402 (one\_model\_fix)}\par
Formule (7.1.4), condition sous la limite\par
\textbf{Claude Opus 5.5.} nothing reported here\par
\textbf{GPT-6.1 Sol.} Place: Formule (7.1.4), condition sous la limite; Reported print: U\textbackslash{}cap Z\textbackslash{}neq\textbackslash{}varnothing; Source TeX: U\textbackslash{}cap Z\textbackslash{}neq \textbackslash{}phi; Note: The handwritten sign is the empty-set sign, not the Greek letter phi.; Proposed edit: U\textbackslash{}cap Z\textbackslash{}neq \textbackslash{}phi =\allowbreak > U\textbackslash{}cap Z\textbackslash{}neq \textbackslash{}varnothing\par
a fix from one model only\par
\par\noindent\hypertarget{DP0114}{}\textbf{DP0114 -- IV, p. 404 (one\_model\_fix)}\par
7.1.7, définition de Pt(U)\par
\textbf{Claude Opus 5.5.} nothing reported here\par
\textbf{GPT-6.1 Sol.} Place: 7.1.7, définition de Pt(U); Reported print: U\_x\textbackslash{}neq\textbackslash{}varnothing; Source TeX: U\_x\textbackslash{}neq\textbackslash{}phi; Note: The print has the empty-set sign, not Greek phi.; Proposed edit: U\_x\textbackslash{}neq\textbackslash{}phi =\allowbreak > U\_x\textbackslash{}neq\textbackslash{}varnothing\par
a fix from one model only\par
\par\noindent\hypertarget{DP0115}{}\textbf{DP0115 -- IV, p. 405 (agreed\_fix\_held)}\par
Exercice 7.1.10 b), Bourbaki reference | Exercice 7.1.10 b), référence à Bourbaki\par
\textbf{Claude Opus 5.5.} Place: Exercice 7.1.10 b), Bourbaki reference; Reported print: (Bourbaki, Top. Gen. Chap. III § 4); Source TeX: Top. Gen. Chap. \textbackslash{}hyperrefext[III][III.sec4]\{4\}; Note: The TeX prints 'Chap. III 4' without the '§'. It also links to SGA 4 exposé III, whereas the reference is to Bourbaki, Topologie générale, Chapter III, § 4.; Proposed edit: Chap. \textbackslash{}hyperrefext[III][III.sec4]\{4\} =\allowbreak > Chap. III \textbackslash{}S\textasciitilde{}4\par
\textbf{GPT-6.1 Sol.} Place: Exercice 7.1.10 b), référence à Bourbaki; Reported print: Top. Gen. Chap. III § 4; Source TeX: Top. Gen. Chap. \textbackslash{}hyperrefext[III][III.sec4]\{4\}; Note: The print refers to Bourbaki, Chapter III, §4. The external SGA reference produces « III 4 », omitting the section sign and linking to the wrong work.; Proposed edit: \textbackslash{}hyperrefext[III][III.sec4]\{4\} =\allowbreak > III \textbackslash{}S\textasciitilde{}4\par
both models give the same fix, held: it changes the editors' apparatus (\textbackslash{}hyperrefext, the target of a cross-reference): a \textbackslash{}sisi correction, note, label, cross-reference or citation is the editors' record and is not removed or retargeted on the models' word\par
\par\noindent\hypertarget{DP0116}{}\textbf{DP0116 -- IV, p. 410 (one\_model\_fix)}\par
Exercice 7.2.6 e), line 2\par
\textbf{Claude Opus 5.5.} Place: Exercice 7.2.6 e), line 2; Reported print: de lien G [3]; Source TeX: [\textbackslash{}hyperref[IV.sec3]\{3\}]; Note: The print's '[3]' in square brackets looks like a bibliography reference (presumably Giraud on gerbes), not section 3 of the exposé. The TeX links it to IV.sec3. If so, it should be \textbackslash{}cite with the key of bibliography item [3]; that key cannot be checked here.; Proposed edit: [\textbackslash{}hyperref[IV.sec3]\{3\}] =\allowbreak > [3]\par
\textbf{GPT-6.1 Sol.} nothing reported here\par
a fix from one model only\par
\par\noindent\hypertarget{DP0117}{}\textbf{DP0117 -- IV, p. 416 (cut)}\par
Fin de 7.6 g), coupure entre les pages 416 et 417\par
\textbf{Claude Opus 5.5.} nothing reported here\par
\textbf{GPT-6.1 Sol.} Place: Fin de 7.6 g), coupure entre les pages 416 et 417; Reported print: La page finit par « et la catégorie \$\textbackslash{}Hhom(C,\textbackslash{}Ppointess(E))\$. »; Source TeX: et la catégorie \$\textbackslash{}Hhom(C,; Note: La fin « \textbackslash{}Ppointess(E))\$. » se trouve au début du morceau 417, ligne 4695, sans que ce chevauchement soit annoncé. La page 417 commence par « (Utiliser b) et le fait que… ». Reprendre la coupe, sans modifier le texte source.\par
a model reports that the automatic cut gives lines of this page to its neighbour, or the neighbour's lines to this page: check the lines it names against the scan of the page that prints them (that page's checker did not have their TeX)\par
\par\noindent\hypertarget{DP0118}{}\textbf{DP0118 -- IV, p. 418 (conflict)}\par
7.8 a), first line | 7.8 a), première ligne, condition sur la fibre de U\par
\textbf{Claude Opus 5.5.} Place: 7.8 a), first line; Reported print: U\_p \ensuremath{\ne} \ensuremath{\varnothing}; Source TeX: U\_p\textbackslash{}neq \textbackslash{}phi; Note: The print has the empty set sign, but \textbackslash{}phi prints the Greek letter \ensuremath{\phi}.; Proposed edit: U\_p\textbackslash{}neq \textbackslash{}phi =\allowbreak > U\_p\textbackslash{}neq \ensuremath{\varnothing}\par
\textbf{GPT-6.1 Sol.} Place: 7.8 a), première ligne, condition sur la fibre de U; Reported print: \$U\_p\textbackslash{}neq\textbackslash{}emptyset\$; Source TeX: \$U\_p\textbackslash{}neq \textbackslash{}phi\$; Note: Le cercle barré imprimé est le signe de l'ensemble vide, plutôt qu'une lettre grecque phi.; Proposed edit: \$U\_p\textbackslash{}neq \textbackslash{}phi\$ =\allowbreak > \$U\_p\textbackslash{}neq \textbackslash{}emptyset\$\par
the two models differ at this place\par
\par\noindent\hypertarget{DP0119}{}\textbf{DP0119 -- IV, p. 420 (one\_model\_fix)}\par
7.8 f), « la partie X' de Point(E) »\par
\textbf{Claude Opus 5.5.} nothing reported here\par
\textbf{GPT-6.1 Sol.} Place: 7.8 f), « la partie X' de Point(E) »; Reported print: la partie \$X'\$ de \$\textbackslash{}mathrm\{Point\}(E)\$, avec Point non souligné; Source TeX: la partie \$X'\$ de \$\textbackslash{}Ppoint(E)\$; Note: L'imprimé distingue ici l'ensemble des classes de points de la catégorie des points, soulignée au début de f). Le second emploi doit être \textbackslash{}Point, non \textbackslash{}Ppoint.; Proposed edit: la partie \$X'\$ de \$\textbackslash{}Ppoint(E)\$ =\allowbreak > la partie \$X'\$ de \$\textbackslash{}Point(E)\$\par
a fix from one model only\par
\par\noindent\hypertarget{DP0120}{}\textbf{DP0120 -- IV, p. 421 (conflict)}\par
Remarque 8.2.1, end | Remarque 8.2.1, dernière relation\par
\textbf{Claude Opus 5.5.} Place: Remarque 8.2.1, end; Reported print: la relation \ensuremath{\tilde{T}}(a(F)); Source TeX: \textbackslash{}widehat\{T\}(\textbackslash{}underline\{a\}(F)); Note: The second T carries a tilde in the print, not a hat.; Proposed edit: \textbackslash{}widehat\{T\}(\textbackslash{}underline\{a\}(F)) =\allowbreak > \textbackslash{}widetilde\{T\}(\textbackslash{}underline\{a\}(F))\par
\textbf{GPT-6.1 Sol.} Place: Remarque 8.2.1, dernière relation; Reported print: \$T\textasciicircum{}\{\textbackslash{}sim\}(\textbackslash{}underline\{a\}(F))\$; Source TeX: \$\textbackslash{}widehat\{T\}(\textbackslash{}underline\{a\}(F))\$; Note: Le second T porte un tilde, non un chapeau.; Proposed edit: \textbackslash{}widehat\{T\}(\textbackslash{}underline\{a\}(F)) =\allowbreak > T\textasciicircum{}\{\textbackslash{}sim\}(\textbackslash{}underline\{a\}(F))\par
the two models differ at this place\par
\par\noindent\hypertarget{DP0121}{}\textbf{DP0121 -- IV, p. 422 (one\_model\_fix)}\par
8.4.2, objet vide\par
\textbf{Claude Opus 5.5.} Place: 8.4.2, objet vide; Reported print: \ensuremath{\varnothing}\_E (empty-set sign); Source TeX: \textbackslash{}phi\_E; Note: The print has the empty-set sign, and the TeX uses \ensuremath{\varnothing} elsewhere (p. 421). \textbackslash{}phi prints a Greek phi, unless the edition redefines it.; Proposed edit: \textbackslash{}phi\_E =\allowbreak > \textbackslash{}emptyset\_E\par
\textbf{GPT-6.1 Sol.} nothing reported here\par
a fix from one model only\par
\par\noindent\hypertarget{DP0122}{}\textbf{DP0122 -- IV, p. 423 (one\_model\_fix)}\par
8.4.3, end\par
\textbf{Claude Opus 5.5.} Place: 8.4.3, end; Reported print: savoir \ensuremath{\varnothing}\_\{B\_G\} et e\_G; Source TeX: \textbackslash{}phi\_\{B\_G\}; Note: The print has the empty-set sign, not a Greek phi.; Proposed edit: \textbackslash{}phi\_\{B\_G\} =\allowbreak > \textbackslash{}emptyset\_\{B\_G\}\par
\textbf{GPT-6.1 Sol.} nothing reported here\par
a fix from one model only\par
\par\noindent\hypertarget{DP0123}{}\textbf{DP0123 -- IV, p. 425 (one\_model\_fix)}\par
8.5.5, last line\par
\textbf{Claude Opus 5.5.} Place: 8.5.5, last line; Reported print: (EGA IV 17.7.5, 17.7.7).; Note: The closing reference line of the page is missing from the TeX.; Proposed edit: fini, plat, net, lisse, étale. =\allowbreak > fini, plat, net, lisse, étale.  \textbackslash{}noindent (\textbackslash{}EGA IV 17.7.5, 17.7.7).\par
\textbf{GPT-6.1 Sol.} nothing reported here\par
a fix from one model only\par
\par\noindent\hypertarget{DP0124}{}\textbf{DP0124 -- IV, p. 425 (cut)}\par
Fin de page, après la seconde liste de 8.5.5\par
\textbf{Claude Opus 5.5.} nothing reported here\par
\textbf{GPT-6.1 Sol.} Place: Fin de page, après la seconde liste de 8.5.5; Reported print: (EGA IV 17.7.5, 17.7.7).; Source TeX: L'extrait s'arrête à « fini, plat, net, lisse, étale. »; Note: La page imprimée comprend encore ce renvoi final. Vérifier la coupe avec la p.426 : la p.425 se termine après « (EGA IV 17.7.5, 17.7.7). », non après la liste. Ne pas insérer le renvoi en double s'il figure dans l'extrait suivant.\par
a model reports that the automatic cut gives lines of this page to its neighbour, or the neighbour's lines to this page: check the lines it names against the scan of the page that prints them (that page's checker did not have their TeX)\par
\par\noindent\hypertarget{DP0125}{}\textbf{DP0125 -- IV, p. 426 (cut)}\par
start of page, before Exercice 8.6\par
\textbf{Claude Opus 5.5.} Place: start of page, before Exercice 8.6; Reported print: page begins with « Exercice 8.6. Soit f: E \ensuremath{\to} F … »; Source TeX: \textbackslash{}noindent (\textbackslash{}EGA IV 17.7.5, 17.7.7).; Note: The line « (EGA IV 17.7.5, 17.7.7). » is not on p. 426; it ends p. 425. Page 426 begins at « Exercice 8.6 ». The first comment line does not flag this line as shared.\par
\textbf{GPT-6.1 Sol.} nothing reported here\par
a model reports that the automatic cut gives lines of this page to its neighbour, or the neighbour's lines to this page: check the lines it names against the scan of the page that prints them (that page's checker did not have their TeX)\par
\par\noindent\hypertarget{DP0126}{}\textbf{DP0126 -- IV, p. 426 (cut)}\par
Début de la coupe, avant l'exercice 8.6\par
\textbf{Claude Opus 5.5.} nothing reported here\par
\textbf{GPT-6.1 Sol.} Place: Début de la coupe, avant l'exercice 8.6; Reported print: La page commence par « Exercice 8.6. Soit… ».; Source TeX: \textbackslash{}noindent (\textbackslash{}EGA IV 17.7.5, 17.7.7).; Note: Cette citation ne figure pas sur la page 426. La coupe inclut un fragment de la page précédente sans que son commentaire initial le signale. Corriger la coupe, non le texte source.\par
a model reports that the automatic cut gives lines of this page to its neighbour, or the neighbour's lines to this page: check the lines it names against the scan of the page that prints them (that page's checker did not have their TeX)\par
\par\noindent\hypertarget{DP0127}{}\textbf{DP0127 -- IV, p. 427 (agreed\_fix\_held)}\par
8.7 d), reference after « du type » | Exercice 8.7 d), référence après « équivalences du type »\par
\textbf{Claude Opus 5.5.} Place: 8.7 d), reference after « du type »; Reported print: type (8.7.1)).; Source TeX: type (\textbackslash{}hyperref[IV.exer8.7]\{8.7\} 1)).; Note: The TeX prints « (8.7 1)) », a reference to Exercice 8.7 followed by « 1) ». The print refers to formula (8.7.1).; Proposed edit: type (\textbackslash{}hyperref[IV.exer8.7]\{8.7\} 1)). =\allowbreak > type \textbackslash{}eqhyperref[IV.eq8.7.1]\{8.7.1\}).\par
\textbf{GPT-6.1 Sol.} Place: Exercice 8.7 d), référence après « équivalences du type »; Reported print: (8.7.1); Source TeX: (\textbackslash{}hyperref[IV.exer8.7]\{8.7\} 1); Note: Il s'agit de la formule (8.7.1), non d'une référence à l'exercice 8.7 suivie d'un 1 séparé.; Proposed edit: (\textbackslash{}hyperref[IV.exer8.7]\{8.7\} 1) =\allowbreak > \textbackslash{}eqhyperref[IV.eq8.7.1]\{8.7.1\}\par
both models give the same fix, held: it changes the editors' apparatus (\textbackslash{}eqhyperref, \textbackslash{}hyperref, the target of a cross-reference): a \textbackslash{}sisi correction, note, label, cross-reference or citation is the editors' record and is not removed or retargeted on the models' word\par
\par\noindent\hypertarget{DP0128}{}\textbf{DP0128 -- IV, p. 429 (agreed\_fix\_held)}\par
8.7 l), end | Exercice 8.7 l), référence finale\par
\textbf{Claude Opus 5.5.} Place: 8.7 l), end; Reported print: (comparer 2.7.5).; Source TeX: (comparer \textbackslash{}sisi\{IV \}\{\} \textbackslash{}hyperref[IV.exer2.7.5]\{2.7.5\}); Note: The \textbackslash{}sisi gives « IV » as the print's reading, but the print has only « comparer 2.7.5 », with no « IV ».; Proposed edit: \textbackslash{}sisi\{IV~\}\{\} =\allowbreak > \par
\textbf{GPT-6.1 Sol.} Place: Exercice 8.7 l), référence finale; Reported print: (comparer 2.7.5); Source TeX: (comparer \textbackslash{}sisi\{IV~\}\{\} \textbackslash{}hyperref[IV.exer2.7.5]\{2.7.5\}); Note: La lecture A de \textbackslash{}sisi ajoute « IV », absent de l'imprimé. La correction B n'est pas en cause.; Proposed edit: \textbackslash{}sisi\{IV~\}\{\} =\allowbreak > \par
both models give the same fix, held: it changes the editors' apparatus (\textbackslash{}sisi): a \textbackslash{}sisi correction, note, label, cross-reference or citation is the editors' record and is not removed or retargeted on the models' word\par
\par\noindent\hypertarget{DP0129}{}\textbf{DP0129 -- IV, p. 431 (conflict)}\par
Remarque 9.1.2 b), end | Remarque 9.1.2 b), référence finale\par
\textbf{Claude Opus 5.5.} Place: Remarque 9.1.2 b), end; Reported print: (comparer IV 3.4.1.); Source TeX: (comparer \textbackslash{}hyperref[IV.subsubsec3.4.1]\{3.4.1\}); Note: The print's exposé prefix IV is dropped silently. Elsewhere the edition records it with \textbackslash{}sisi\{IV\}\{\}. The reference itself is the same.; Proposed edit: (comparer \textbackslash{}hyperref[IV.subsubsec3.4.1]\{3.4.1\}) =\allowbreak > (comparer \textbackslash{}sisi\{IV \}\{\}\textbackslash{}hyperref[IV.subsubsec3.4.1]\{3.4.1\})\par
\textbf{GPT-6.1 Sol.} Place: Remarque 9.1.2 b), référence finale; Reported print: « comparer IV 3.4.1. »; Source TeX: \textbackslash{}hyperref[IV.subsubsec3.4.1]\{3.4.1\}; Note: Le numéro d'exposé IV imprimé devant la référence est omis.; Proposed edit: \textbackslash{}hyperref[IV.subsubsec3.4.1]\{3.4.1\} =\allowbreak > IV\textasciitilde{}\textbackslash{}hyperref[IV.subsubsec3.4.1]\{3.4.1\}\par
the two models differ at this place\par
\par\noindent\hypertarget{DP0130}{}\textbf{DP0130 -- IV, p. 433 (conflict)}\par
first lines, reference | Suite de la remarque 9.1.2 d), référence après « d'autre part on a »\par
\textbf{Claude Opus 5.5.} Place: first lines, reference; Reported print: a (II 4.4.1); Source TeX: \textbackslash{}sisi\{\textbackslash{}hyperrefext[II][II.theo4.4]\{4.4\} I\}\{...\}; Note: The print reads 4.4.1, not '4.4 I'. This concerns the A (print) part of \textbackslash{}sisi.; Proposed edit: \{4.4\} I\} =\allowbreak > \{4.4.1\}\}\par
\textbf{GPT-6.1 Sol.} Place: Suite de la remarque 9.1.2 d), référence après « d'autre part on a »; Reported print: (II 4.4.1); Source TeX: \textbackslash{}sisi\{\textbackslash{}hyperrefext[II][II.theo4.4]\{4.4\} I\}\{\textbackslash{}hyperrefext[II][II.cor4.4.4]\{4.4.4\}\}; Note: La référence imprimée se termine par .1, non par un I séparé après 4.4. Cette modification ne concerne que la lecture A ; B n'est pas contrôlé.; Proposed edit: \textbackslash{}hyperrefext[II][II.theo4.4]\{4.4\} I =\allowbreak > II\textasciitilde{}4.4.1\par
the two models differ at this place\par
\par\noindent\hypertarget{DP0131}{}\textbf{DP0131 -- IV, p. 434 (one\_model\_fix)}\par
Exercice 9.1.6 d), formule du 2-produit fibré\par
\textbf{Claude Opus 5.5.} nothing reported here\par
\textbf{GPT-6.1 Sol.} Place: Exercice 9.1.6 d), formule du 2-produit fibré; Reported print: \$F'=\allowbreak F\textbackslash{}mathbin\{\textbackslash{}overset\{2\}\{\textbackslash{}times\}\}\_E E'\$; Source TeX: \textbackslash{}sisi\{\$F'=\allowbreak \textbackslash{}fprod\{F\textasciicircum{}2\}\{E'\}\{E\}\$\}\{\$F'=\allowbreak F\textbackslash{}times\_\{E\}\textasciicircum{}\{2\}E'\$\}; Note: Dans l'imprimé, le 2 surmonte le signe de produit fibré ; A en fait un exposant de F. Seul A est à modifier.; Proposed edit: \textbackslash{}fprod\{F\textasciicircum{}2\}\{E'\}\{E\} =\allowbreak > F\textbackslash{}times\_\{E\}\textasciicircum{}\{2\}E'\par
a fix from one model only\par
\par\noindent\hypertarget{DP0132}{}\textbf{DP0132 -- IV, p. 435 (one\_model\_fix)}\par
Exercice 9.1.7 a)\par
\textbf{Claude Opus 5.5.} Place: Exercice 9.1.7 a); Reported print: (Utiliser [2, Chap I 1.3.]); Source TeX: [2, Chap 1 1.3]; Note: The print has the roman numeral I, as in the earlier '[2, Chap I]'.; Proposed edit: Chap 1 1.3 =\allowbreak > Chap I 1.3\par
\textbf{GPT-6.1 Sol.} nothing reported here\par
a fix from one model only\par
\par\noindent\hypertarget{DP0133}{}\textbf{DP0133 -- IV, p. 441 (cut)}\par
9.1.9 c), last line of the TeX\par
\textbf{Claude Opus 5.5.} Place: 9.1.9 c), last line of the TeX; Reported print: Page 441 ends with '... si et seulement si'; page 442 begins 'g est surjectif, et g est surjectif sur les graphes ...'; Source TeX: si et seulement si \$g\$ est surjectif, \textbackslash{}emph\{et\}; Note: '\$g\$ est surjectif, \textbackslash{}emph\{et\}' belongs to p. 442. Page 442 really begins at the first '\$g\$', and the \textbackslash{}pageoriginale on p. 442's first line belongs before 'g est surjectif, et'. The text itself is complete.\par
\textbf{GPT-6.1 Sol.} nothing reported here\par
a model reports that the automatic cut gives lines of this page to its neighbour, or the neighbour's lines to this page: check the lines it names against the scan of the page that prints them (that page's checker did not have their TeX)\par
\par\noindent\hypertarget{DP0134}{}\textbf{DP0134 -- IV, p. 441 (cut)}\par
Fin de 9.1.9 c), coupure 441/442\par
\textbf{Claude Opus 5.5.} nothing reported here\par
\textbf{GPT-6.1 Sol.} Place: Fin de 9.1.9 c), coupure 441/442; Reported print: La page 441 finit par « plongement si et seulement si ». La page 442 commence par « g est surjectif, et g est surjectif sur les graphes… ».; Source TeX: plongement si et seulement si \$g\$ est surjectif, \textbackslash{}emph\{et\}; Note: Le fragment « \$g\$ est surjectif, \textbackslash{}emph\{et\} » appartient à la page 442. Ce débordement n'est pas annoncé dans les commentaires des extraits ; le début de l'extrait 442 est donc incomplet.\par
a model reports that the automatic cut gives lines of this page to its neighbour, or the neighbour's lines to this page: check the lines it names against the scan of the page that prints them (that page's checker did not have their TeX)\par
\par\noindent\hypertarget{DP0135}{}\textbf{DP0135 -- IV, p. 442 (one\_model\_fix)}\par
9.1.9 d), première condition\par
\textbf{Claude Opus 5.5.} nothing reported here\par
\textbf{GPT-6.1 Sol.} Place: 9.1.9 d), première condition; Reported print: QO 1); Source TeX: \textbackslash{}item[Q0 1)]; Note: Le sigle imprimé comporte la lettre O, non le chiffre zéro.; Proposed edit: \textbackslash{}item[Q0 1)] =\allowbreak > \textbackslash{}item[QO 1)]\par
a fix from one model only\par
\par\noindent\hypertarget{DP0136}{}\textbf{DP0136 -- IV, p. 442 (one\_model\_fix)}\par
9.1.9 d), deuxième condition\par
\textbf{Claude Opus 5.5.} nothing reported here\par
\textbf{GPT-6.1 Sol.} Place: 9.1.9 d), deuxième condition; Reported print: QO 2); Source TeX: \textbackslash{}item[Q0 2)]; Note: Même confusion entre la lettre O et le chiffre zéro.; Proposed edit: \textbackslash{}item[Q0 2)] =\allowbreak > \textbackslash{}item[QO 2)]\par
a fix from one model only\par
\par\noindent\hypertarget{DP0137}{}\textbf{DP0137 -- IV, p. 442 (one\_model\_fix)}\par
9.1.10 a), renvoi aux conditions de 9.1.9 d)\par
\textbf{Claude Opus 5.5.} nothing reported here\par
\textbf{GPT-6.1 Sol.} Place: 9.1.10 a), renvoi aux conditions de 9.1.9 d); Reported print: conditions QO 1) et QO 2); Source TeX: conditions Q0 1) et Q0 2); Note: Les deux renvois reprennent le sigle QO, avec la lettre O.; Proposed edit: conditions Q0 1) et Q0 2) =\allowbreak > conditions QO 1) et QO 2)\par
a fix from one model only\par
\par\noindent\hypertarget{DP0138}{}\textbf{DP0138 -- IV, p. 445 (cut)}\par
Première ligne de l'extrait, coupure 444/445\par
\textbf{Claude Opus 5.5.} nothing reported here\par
\textbf{GPT-6.1 Sol.} Place: Première ligne de l'extrait, coupure 444/445; Reported print: La référence « (9.1.13) » clôt la page 444 ; la page 445 commence directement par « d) Soit X un espace sobre noethérien ».; Source TeX: (\textbackslash{}hyperref[IV.exer9.1.13]\{9.1.13\}).; Note: Cette première ligne appartient entièrement à la page 444, contrairement au découpage annoncé. Elle manque à l'extrait 444 et ne doit pas être contrôlée comme texte de la page 445.\par
a model reports that the automatic cut gives lines of this page to its neighbour, or the neighbour's lines to this page: check the lines it names against the scan of the page that prints them (that page's checker did not have their TeX)\par
\par\noindent\hypertarget{DP0139}{}\textbf{DP0139 -- IV, p. 447 (one\_model\_fix)}\par
d), last sentence\par
\textbf{Claude Opus 5.5.} Place: d), last sentence; Reported print: et le Inf y est distributif; Source TeX: le Inf \$y\$ est distributif; Note: Here « y » is the French adverb, not a variable. Setting it in math makes it read as the point y.; Proposed edit: le Inf \$y\$ est =\allowbreak > le Inf y est\par
\textbf{GPT-6.1 Sol.} nothing reported here\par
a fix from one model only\par
\par\noindent\hypertarget{DP0140}{}\textbf{DP0140 -- IV, p. 449 (conflict)}\par
9.1.13 d) | 9.1.13 d), début\par
\textbf{Claude Opus 5.5.} Place: 9.1.13 d); Reported print: Enoncer les duals des observations; Source TeX: Énoncer les duales des observations; Proposed correction: Énoncer les duales des observations; Note: The print has « duals »; the agreement with « observations » requires « duales ». The TeX silently corrects it.\par
\textbf{GPT-6.1 Sol.} Place: 9.1.13 d), début; Reported print: Énoncer les duals des observations; Source TeX: Énoncer les duales des observations; Note: Le scan porte «duals», non «duales».; Proposed edit: les duales des observations =\allowbreak > les duals des observations\par
the two models differ at this place: one model fixes the TeX, the other calls it a misprint of the print\par
\par\noindent\hypertarget{DP0141}{}\textbf{DP0141 -- IV, p. 449 (conflict)}\par
9.1.14, « Cependant la réponse à a) et b) est affirmative si E est de la forme Top(X) » | 9.1.14, cas où E est de la forme Top(X)\par
\textbf{Claude Opus 5.5.} Place: 9.1.14, « Cependant la réponse à a) et b) est affirmative si E est de la forme Top(X) »; Reported print: X est espace localement noethérien; Source TeX: \$X\$ est un espace localement noethérien; Proposed correction: \$X\$ est un espace localement noethérien; Note: The print lacks « un ». The TeX silently adds it.\par
\textbf{GPT-6.1 Sol.} Place: 9.1.14, cas où E est de la forme Top(X); Reported print: \$X\$ est espace localement noethérien; Source TeX: \$X\$ est un espace localement noethérien; Note: L'article «un» n'est pas imprimé.; Proposed edit: \$X\$ est un espace =\allowbreak > \$X\$ est espace\par
the two models differ at this place: one model fixes the TeX, the other calls it a misprint of the print\par
\par\noindent\hypertarget{DP0142}{}\textbf{DP0142 -- IV, p. 450 (conflict)}\par
Paragraph after c), reference to Deligne | Paragraphe suivant 9.1.14 c), référence au résultat de Deligne\par
\textbf{Claude Opus 5.5.} Place: Paragraph after c), reference to Deligne; Reported print: grâce au résultat de DELIGNE (     ) affirmant; Source TeX: grâce au résultat de \textbackslash{}textsc\{Deligne\} (voir [16] proposition 4.47); Note: The print leaves the parentheses empty. The reference « voir [16] proposition 4.47 » is an editorial addition, and it is not marked as \textbackslash{}nde.; Proposed edit: (voir [16] proposition 4.47) =\allowbreak > (\textbackslash{}quad)\par
\textbf{GPT-6.1 Sol.} Place: Paragraphe suivant 9.1.14 c), référence au résultat de Deligne; Reported print: DELIGNE ( ); Source TeX: \textbackslash{}textsc\{Deligne\} (voir [16] proposition 4.47); Note: Les parenthèses imprimées sont vides; le renvoi bibliographique est un ajout du TeX.; Proposed edit: (voir [16] proposition 4.47) =\allowbreak > ( )\par
the two models differ at this place\par
\par\noindent\hypertarget{DP0143}{}\textbf{DP0143 -- IV, p. 454 (one\_model\_fix)}\par
9.3.2, first line\par
\textbf{Claude Opus 5.5.} Place: 9.3.2, first line; Reported print: Soient donc, comme 9.2.1 ,; Source TeX: Soient donc, comme dans \textbackslash{}hyperref[IV.subsubsec9.2.1]\{9.2.1\},; Note: The TeX adds the word «dans», which the print does not have.; Proposed edit: comme dans \textbackslash{}hyperref[IV.subsubsec9.2.1] =\allowbreak > comme \textbackslash{}hyperref[IV.subsubsec9.2.1]\par
\textbf{GPT-6.1 Sol.} nothing reported here\par
a fix from one model only\par
\par\noindent\hypertarget{DP0144}{}\textbf{DP0144 -- IV, p. 455 (one\_model\_fix)}\par
Proposition 9.3.3 c), premiers mots\par
\textbf{Claude Opus 5.5.} nothing reported here\par
\textbf{GPT-6.1 Sol.} Place: Proposition 9.3.3 c), premiers mots; Reported print: Le foncteur; Source TeX: \textbackslash{}item le foncteur; Note: L'imprimé commence le texte du c) par une capitale.; Proposed edit: \textbackslash{}item le foncteur =\allowbreak > \textbackslash{}item Le foncteur\par
a fix from one model only\par
\par\noindent\hypertarget{DP0145}{}\textbf{DP0145 -- IV, p. 456 (one\_model\_fix)}\par
9.3.5, notation de l'objet initial\par
\textbf{Claude Opus 5.5.} nothing reported here\par
\textbf{GPT-6.1 Sol.} Place: 9.3.5, notation de l'objet initial; Reported print: \$i\_\{\textbackslash{}ast\}(\textbackslash{}varnothing\_\{\textbackslash{}mathcal\{F\}\})\$; Source TeX: \$i\_\{\textbackslash{}ast\}(\textbackslash{}phi\_\{\textbackslash{}mathcal\{F\}\})\$; Note: Convention récurrente : le rond barré imprimé pour l'objet initial est transcrit par \textbackslash{}phi. Trois exemples : ici, dans la proposition 9.4.2 p.458 et à la fin de sa démonstration p.459. La même substitution concerne les deux autres occurrences de \textbackslash{}phi dans ce paragraphe de p.456.; Proposed edit: \$i\_\{\textbackslash{}ast\}(\textbackslash{}phi\_\{\textbackslash{}mathcal\{F\}\})\$ =\allowbreak > \$i\_\{\textbackslash{}ast\}(\textbackslash{}varnothing\_\{\textbackslash{}mathcal\{F\}\})\$\par
a fix from one model only\par
\par\noindent\hypertarget{DP0146}{}\textbf{DP0146 -- IV, p. 457 (conflict)}\par
9.3.6, 'ou ce qui revient au même (...)' | 9.3.6, référence après « ou ce qui revient au même »\par
\textbf{Claude Opus 5.5.} Place: 9.3.6, 'ou ce qui revient au même (...)'; Reported print: (I.8    ); Source TeX: \textbackslash{}sisi\{\textbackslash{}hyperref[IV.rem1.8]\{1.8\}\textbackslash{}quad\}; Note: The print reads 'I.8' (Exposé I, followed by a blank). The \textbackslash{}sisi print-reading argument gives '1.8' (a reference to IV 1.8) instead.; Proposed edit: \textbackslash{}sisi\{\textbackslash{}hyperref[IV.rem1.8]\{1.8\}\textbackslash{}quad\} =\allowbreak > \textbackslash{}sisi\{I.8\textbackslash{}quad\}\par
\textbf{GPT-6.1 Sol.} Place: 9.3.6, référence après « ou ce qui revient au même »; Reported print: (I.8); Source TeX: \textbackslash{}sisi\{\textbackslash{}hyperref[IV.rem1.8]\{1.8\}\textbackslash{}quad\}\{\textbackslash{}hyperrefext[I][I.sec8]\{8\}\}; Note: Le premier argument de \textbackslash{}sisi donne 1.8, tandis que le print porte I.8. Seul ce premier argument est à rectifier.; Proposed edit: \textbackslash{}hyperref[IV.rem1.8]\{1.8\}\textbackslash{}quad =\allowbreak > \textbackslash{}hyperrefext[I][I.sec8]\{8\}\textbackslash{}quad\par
the two models differ at this place\par
\par\noindent\hypertarget{DP0147}{}\textbf{DP0147 -- IV, p. 458 (one\_model\_fix)}\par
proof of 9.4.1, accessibility of j\_*\par
\textbf{Claude Opus 5.5.} Place: proof of 9.4.1, accessibility of j\_*; Reported print: (I 9.5;on note qu'un topos est une catégorie accessible (I 9.11.3)).; Source TeX: (\textbackslash{}hyperrefext[I][I.prop9.5]\{9.5\}); on note; Note: The TeX adds a closing parenthesis after 9.5 that the print does not have. The parenthesis opened before 'I 9.5' closes only after '(I 9.11.3)'.; Proposed edit: \{9.5\}); on note =\allowbreak > \{9.5\}; on note\par
\textbf{GPT-6.1 Sol.} nothing reported here\par
a fix from one model only\par
\par\noindent\hypertarget{DP0148}{}\textbf{DP0148 -- IV, p. 461 (one\_model\_fix)}\par
9.4.8 a), symbole de l'objet initial ; convention répétée\par
\textbf{Claude Opus 5.5.} nothing reported here\par
\textbf{GPT-6.1 Sol.} Place: 9.4.8 a), symbole de l'objet initial ; convention répétée; Reported print: \$i\_\{\textbackslash{}ast\}(\textbackslash{}varnothing\_\{E'\})\$; Source TeX: \$i\_\{\textbackslash{}ast\}(\textbackslash{}phi\_\{E'\})\$; Note: Le scan emploie le signe \ensuremath{\varnothing} pour l'objet initial, tandis que le TeX emploie \textbackslash{}phi. Trois exemples : \ensuremath{\varnothing}\_\{E'\} ici, \ensuremath{\varnothing}\_E dans U\ensuremath{\cap}V=\allowbreak \ensuremath{\varnothing}\_E en 9.4.8 b), et \ensuremath{\varnothing}\_\{\textbackslash{}mathcal\{F\}\} en 9.5.4 b), p. 464. Un seul signalement pour cette convention répétée.; Proposed edit: \textbackslash{}phi\_\{E'\} =\allowbreak > \textbackslash{}varnothing\_\{E'\}\par
a fix from one model only\par
\par\noindent\hypertarget{DP0149}{}\textbf{DP0149 -- IV, p. 462 (conflict)}\par
9.4.8 b), end: «s'appelle la frontière» | Première ligne, définition de la frontière\par
\textbf{Claude Opus 5.5.} Place: 9.4.8 b), end: «s'appelle la frontière»; Reported print: s'appele; Source TeX: s'appelle la; Proposed correction: s'appelle; Note: Typing error in the print. The TeX corrects it silently.\par
\textbf{GPT-6.1 Sol.} Place: Première ligne, définition de la frontière; Reported print: \$\textbackslash{}Sup(U,V)=\allowbreak U\textbackslash{}amalg V\$; Source TeX: \$\textbackslash{}Sup(U,V)=\allowbreak \textbackslash{}mathcal\{U\}\textbackslash{}amalg \textbackslash{}mathcal\{V\}\$; Note: Le membre de droite imprimé porte les lettres ordinaires U et V, désignant les objets ouverts, non les lettres scriptes désignant les sous-topos.; Proposed edit: \$\textbackslash{}Sup(U,V)=\allowbreak \textbackslash{}mathcal\{U\}\textbackslash{}amalg \textbackslash{}mathcal\{V\}\$ =\allowbreak > \$\textbackslash{}Sup(U,V)=\allowbreak U\textbackslash{}amalg V\$\par
the two models differ at this place: one model fixes the TeX, the other calls it a misprint of the print\par
\par\noindent\hypertarget{DP0150}{}\textbf{DP0150 -- IV, p. 463 (one\_model\_fix)}\par
9.4.10: «(EGA 0\_III 9.1, EGA IV 10.1)»\par
\textbf{Claude Opus 5.5.} Place: 9.4.10: «(EGA 0\_III 9.1, EGA IV 10.1)»; Reported print: EGA 0\_III 9.1; Source TeX: \textbackslash{}hyperref[IV.subsec9.1]\{9.1\}; Note: The number is a reference to EGA 0\_III 9.1. The TeX links it to section 9.1 of Exposé IV, which is a wrong reference.; Proposed edit: \textbackslash{}hyperref[IV.subsec9.1]\{9.1\} =\allowbreak > 9.1\par
\textbf{GPT-6.1 Sol.} nothing reported here\par
a fix from one model only\par
\par\noindent\hypertarget{DP0151}{}\textbf{DP0151 -- IV, p. 465 (one\_model\_fix)}\par
9.5.6, second diagram, bottom right corner\par
\textbf{Claude Opus 5.5.} Place: 9.5.6, second diagram, bottom right corner; Reported print: i*j\_*j*Y with no punctuation after it; Source TeX: i\textasciicircum{}\{\textbackslash{}ast\}j\_\{\textbackslash{}ast\}j\textasciicircum{}\{\textbackslash{}ast\}Y ; \}\}; Note: The TeX adds a semicolon that the print does not have.; Proposed edit: i\textasciicircum{}\{\textbackslash{}ast\}j\_\{\textbackslash{}ast\}j\textasciicircum{}\{\textbackslash{}ast\}Y ; \}\} =\allowbreak > i\textasciicircum{}\{\textbackslash{}ast\}j\_\{\textbackslash{}ast\}j\textasciicircum{}\{\textbackslash{}ast\}Y \}\}\par
\textbf{GPT-6.1 Sol.} nothing reported here\par
a fix from one model only\par
\par\noindent\hypertarget{DP0152}{}\textbf{DP0152 -- IV, p. 467 (conflict)}\par
On peut aussi éviter le recours à ... | d), the two references to I 9.25\par
\textbf{Claude Opus 5.5.} Place: On peut aussi éviter le recours à ...; Reported print: recours à 1.9.25; Source TeX: recours à \textbackslash{}sisi\{(1.9.2 5)\}; Note: The second occurrence in the print reads 1.9.25, without parentheses or a space; only the first occurrence reads (1.9.2 5). The \textbackslash{}sisi A reading is copied from the first one.; Proposed edit: recours à \textbackslash{}sisi\{(1.9.2 5)\} =\allowbreak > recours à \textbackslash{}sisi\{1.9.25\}\par
\textbf{GPT-6.1 Sol.} Place: d), the two references to I 9.25; Reported print: « (1.9.2 5) » and « 1.9.25 »; Source TeX: \textbackslash{}sisi\{(1.9.2 5)\}\{\textbackslash{}hyperrefext[I][I.cor9.25]\{9.25\}\}; Proposed correction: I\textasciitilde{}9.25; Note: Both printed references use 1 in place of the exposé numeral I. Both occurrences are corrected through \textbackslash{}sisi; the spacing between 2 and 5 is not a separate finding.\par
the two models differ at this place: one model fixes the TeX, the other calls it a misprint of the print\par
\par\noindent\hypertarget{DP0153}{}\textbf{DP0153 -- IV, p. 470 (conflict)}\par
foncteur restriction (...) | Exercise 9.5.10 a), reference after « foncteur restriction »\par
\textbf{Claude Opus 5.5.} Place: foncteur restriction (...); Reported print: (1.5.11), the first character bold/overstruck (perhaps I); Source TeX: \textbackslash{}sisi\{(4.5.11)\}; Note: The print does not read 4.5.11. The \textbackslash{}sisi A reading must give the print's 1.5.11 (or I.5.11).; Proposed edit: \textbackslash{}sisi\{(4.5.11)\} =\allowbreak > \textbackslash{}sisi\{(1.5.11)\}\par
\textbf{GPT-6.1 Sol.} Place: Exercise 9.5.10 a), reference after « foncteur restriction »; Reported print: (…5.11); the first numeral has a faint or broken stroke.; Source TeX: \textbackslash{}sisi\{(4.5.11)\}\{\textbackslash{}ref\{VI.lem5.4\}\}; Proposed correction: VI\textasciitilde{}5.4; Note: The printed restriction-functor reference ending in 5.11 is corrected by the TeX to VI 5.4 through \textbackslash{}sisi. The damaged first numeral does not support a separate asserted 1/4 transcription difference.\par
the two models differ at this place: one model fixes the TeX, the other calls it a misprint of the print\par
\par\noindent\hypertarget{DP0154}{}\textbf{DP0154 -- IV, p. 470 (one\_model\_fix)}\par
End of the first paragraph, references to exposé III\par
\textbf{Claude Opus 5.5.} nothing reported here\par
\textbf{GPT-6.1 Sol.} Place: End of the first paragraph, references to exposé III; Reported print: (III 1.2 et 1.7); Source TeX: (\textbackslash{}hyperrefext[III][III.prop1.2]\{1.2\} et \textbackslash{}hyperrefext[III][III.prop1.7]\{1.7\}); Note: The second external-reference command repeats III, which the print gives only before 1.2.; Proposed edit: \textbackslash{}hyperrefext[III][III.prop1.7]\{1.7\} =\allowbreak > \textbackslash{}hyperref[III.prop1.7]\{1.7\}\par
a fix from one model only\par
\par\noindent\hypertarget{DP0155}{}\textbf{DP0155 -- IV, p. 471 (one\_model\_fix)}\par
b), références à III 1.2 et 1.7\par
\textbf{Claude Opus 5.5.} nothing reported here\par
\textbf{GPT-6.1 Sol.} Place: b), références à III 1.2 et 1.7; Reported print: (III 1.2 et 1.7); Source TeX: (\textbackslash{}hyperrefext[III][III.prop1.2]\{1.2\} et \textbackslash{}hyperrefext[III][III.prop1.7]\{1.7\}); Note: Le second renvoi ne répète pas « III » dans l'imprimé ; le second \textbackslash{}hyperrefext l'ajoute.; Proposed edit: \textbackslash{}hyperrefext[III][III.prop1.7]\{1.7\} =\allowbreak > \textbackslash{}hyperref[III.prop1.7]\{1.7\}\par
a fix from one model only\par
\par\noindent\hypertarget{DP0156}{}\textbf{DP0156 -- IV, p. 474 (one\_model\_fix)}\par
9.7.1, first display\par
\textbf{Claude Opus 5.5.} Place: 9.7.1, first display; Reported print: Ob C'\ensuremath{\cap}Ob C" =\allowbreak  \ensuremath{\varnothing}; Source TeX: \textbackslash{}Ob C'\textbackslash{}cap \textbackslash{}Ob C''=\allowbreak \textbackslash{}phi; Note: This is the empty set. The TeX's \textbackslash{}phi (the Greek letter, also used for the map \ensuremath{\phi} of 9.5.12 b) and 9.6.5) is not right here, unless macros.tex redefines \textbackslash{}phi.; Proposed edit: \textbackslash{}Ob C'\textbackslash{}cap \textbackslash{}Ob C''=\allowbreak \textbackslash{}phi =\allowbreak > \textbackslash{}Ob C'\textbackslash{}cap \textbackslash{}Ob C''=\allowbreak \textbackslash{}emptyset\par
\textbf{GPT-6.1 Sol.} nothing reported here\par
a fix from one model only\par
\par\noindent\hypertarget{DP0157}{}\textbf{DP0157 -- IV, p. 474 (one\_model\_fix)}\par
9.7.1, second display\par
\textbf{Claude Opus 5.5.} Place: 9.7.1, second display; Reported print: Hom(X",X') =\allowbreak  \ensuremath{\varnothing}; Source TeX: \textbackslash{}Hom(X'',X')=\allowbreak \textbackslash{}phi; Note: The empty set is set as \textbackslash{}phi; see the previous finding.; Proposed edit: \textbackslash{}Hom(X'',X')=\allowbreak \textbackslash{}phi =\allowbreak > \textbackslash{}Hom(X'',X')=\allowbreak \textbackslash{}emptyset\par
\textbf{GPT-6.1 Sol.} nothing reported here\par
a fix from one model only\par
\par\noindent\hypertarget{DP0158}{}\textbf{DP0158 -- IV, p. 474 (one\_model\_fix)}\par
9.7.1, third display\par
\textbf{Claude Opus 5.5.} Place: 9.7.1, third display; Reported print: Hom(v(A"),u(A')) =\allowbreak  \ensuremath{\varnothing}; Source TeX: u(A'))=\allowbreak \textbackslash{}phi; Note: The empty set is set as \textbackslash{}phi.; Proposed edit: u(A'))=\allowbreak \textbackslash{}phi =\allowbreak > u(A'))=\allowbreak \textbackslash{}emptyset\par
\textbf{GPT-6.1 Sol.} nothing reported here\par
a fix from one model only\par
\par\noindent\hypertarget{DP0159}{}\textbf{DP0159 -- IV, p. 475 (one\_model\_fix)}\par
(9.7.2.2)\par
\textbf{Claude Opus 5.5.} Place: (9.7.2.2); Reported print: Hom(q,p) =\allowbreak  \ensuremath{\varnothing}; Source TeX: \textbackslash{}Hom(q,p)=\allowbreak \textbackslash{}phi; Note: The empty set is set as \textbackslash{}phi.; Proposed edit: \textbackslash{}Hom(q,p)=\allowbreak \textbackslash{}phi =\allowbreak > \textbackslash{}Hom(q,p)=\allowbreak \textbackslash{}emptyset\par
\textbf{GPT-6.1 Sol.} nothing reported here\par
a fix from one model only\par
\par\noindent\hypertarget{DP0160}{}\textbf{DP0160 -- IV, p. 476 (conflict)}\par
proof c), top row of the diagram | proof c), bottom row of the diagram | Démonstration, c), premier diagramme\par
\textbf{Claude Opus 5.5.} Place: proof c), top row of the diagram; Reported print: Point(\ensuremath{𝔉}) (the same handwritten F used throughout the page); Source TeX: \textbackslash{}Point(\textbackslash{}Gg); Note: The diagram uses \textbackslash{}Gg, a script G, where the rest of the page has \textbackslash{}mathcal\{F\} for the same closed subtopos. The print's letter is the same everywhere.; Proposed edit: \textbackslash{}Point(\textbackslash{}Gg) =\allowbreak > \textbackslash{}Point(\textbackslash{}mathcal\{F\}) || Place: proof c), bottom row of the diagram; Reported print: Pro(\ensuremath{𝔉}); Source TeX: \textbackslash{}Pro(\textbackslash{}Gg); Note: Same as the top row: the print has the letter F here, not G.; Proposed edit: \textbackslash{}Pro(\textbackslash{}Gg) =\allowbreak > \textbackslash{}Pro(\textbackslash{}mathcal\{F\})\par
\textbf{GPT-6.1 Sol.} Place: Démonstration, c), premier diagramme; Reported print: \textbackslash{}Ppoint(\textbackslash{}mathcal\{U\})\textbackslash{}to\textbackslash{}Ppoint(E)\textbackslash{}leftarrow\textbackslash{}Ppoint(\textbackslash{}mathcal\{F\}), avec \textbackslash{}Pro(\textbackslash{}mathcal\{F\}) en bas à droite; Source TeX: \textbackslash{}Point(\textbackslash{}Uu), \textbackslash{}Point(E), \textbackslash{}Point(\textbackslash{}Gg), \textbackslash{}Pro(\textbackslash{}Gg); Note: Both right-hand objects contain script F, not G. The three underlined Point symbols denote categories, distinguished from the un-underlined sets of isomorphism classes used in 9.7.5; this is mathematical notation, not emphasis.; Proposed edit: \textbackslash{}Point(\textbackslash{}Uu) \textbackslash{}ar@\{\textasciicircum{}\{(\}->\}[d] \textbackslash{}ar[r]\textasciicircum{}\{u\}\& \textbackslash{}Point(E) \textbackslash{}ar@\{\textasciicircum{}\{(\}->\}[d]\& \textbackslash{}Point(\textbackslash{}Gg) \textbackslash{}ar@\{\textasciicircum{}\{(\}->\}[d] \textbackslash{}ar[l]\_\{v\}\textbackslash{}\textbackslash{} \textbackslash{}Pro(\textbackslash{}Uu) \textbackslash{}ar[r]\textasciicircum{}\{\textbackslash{}Pro(j\_\{\textbackslash{}ast\})\}  \& \textbackslash{}Pro(E)   \& \textbackslash{}Pro(\textbackslash{}Gg) =\allowbreak > \textbackslash{}Ppoint(\textbackslash{}Uu) \textbackslash{}ar@\{\textasciicircum{}\{(\}->\}[d] \textbackslash{}ar[r]\textasciicircum{}\{u\}\& \textbackslash{}Ppoint(E) \textbackslash{}ar@\{\textasciicircum{}\{(\}->\}[d]\& \textbackslash{}Ppoint(\textbackslash{}mathcal\{F\}) \textbackslash{}ar@\{\textasciicircum{}\{(\}->\}[d] \textbackslash{}ar[l]\_\{v\}\textbackslash{}\textbackslash{} \textbackslash{}Pro(\textbackslash{}Uu) \textbackslash{}ar[r]\textasciicircum{}\{\textbackslash{}Pro(j\_\{\textbackslash{}ast\})\}  \& \textbackslash{}Pro(E)   \& \textbackslash{}Pro(\textbackslash{}mathcal\{F\})\par
the two models differ at this place\par
\par\noindent\hypertarget{DP0161}{}\textbf{DP0161 -- IV, p. 476 (conflict)}\par
proof c), reference for the adjoint Pro(i\textasciicircum{}*) | Démonstration, c), référence après Pro(i\textasciicircum{}\{\textbackslash{}ast\})\par
\textbf{Claude Opus 5.5.} Place: proof c), reference for the adjoint Pro(i\textasciicircum{}*); Reported print: (I 8    ); Source TeX: \textbackslash{}sisi\{\textbackslash{}hyperref[IV.rem1.8]\{1.8\}\}\{\textbackslash{}ref\{I.sec8.11.5\}\}; Note: The print refers to exposé I, section 8 (the rest of the number is left blank). The \textbackslash{}sisi's A reading gives an internal reference IV 1.8 instead.; Proposed edit: \textbackslash{}hyperref[IV.rem1.8]\{1.8\} =\allowbreak > \textbackslash{}hyperrefext[I][I.sec8]\{8\}\par
\textbf{GPT-6.1 Sol.} Place: Démonstration, c), référence après Pro(i\textasciicircum{}\{\textbackslash{}ast\}); Reported print: (I 8); Source TeX: \textbackslash{}sisi\{\textbackslash{}hyperref[IV.rem1.8]\{1.8\}\}\{\textbackslash{}ref\{I.sec8.11.5\}\}; Note: The print cites Exposé I, section 8, not paragraph 1.8 of Exposé IV. Only the first argument of \textbackslash{}sisi is changed.; Proposed edit: \textbackslash{}hyperref[IV.rem1.8]\{1.8\} =\allowbreak > I 8\par
the two models differ at this place\par
\par\noindent\hypertarget{DP0162}{}\textbf{DP0162 -- IV, p. 476 (one\_model\_fix)}\par
Démonstration, b), condition sur Hom(p',p)\par
\textbf{Claude Opus 5.5.} nothing reported here\par
\textbf{GPT-6.1 Sol.} Place: Démonstration, b), condition sur Hom(p',p); Reported print: \textbackslash{}Hom(p',p)\textbackslash{}neq\textbackslash{}varnothing; Source TeX: \textbackslash{}Hom(p',p)\textbackslash{}neq \textbackslash{}phi; Note: The handwritten sign is a slashed circle denoting the empty set, rather than the letter phi.; Proposed edit: \textbackslash{}Hom(p',p)\textbackslash{}neq \textbackslash{}phi =\allowbreak > \textbackslash{}Hom(p',p)\textbackslash{}neq \textbackslash{}varnothing\par
a fix from one model only\par
\par\noindent\hypertarget{DP0163}{}\textbf{DP0163 -- IV, p. 476 (one\_model\_fix)}\par
Lemme 9.7.2.4, ligne supérieure du diagramme (9.7.2.4.1)\par
\textbf{Claude Opus 5.5.} nothing reported here\par
\textbf{GPT-6.1 Sol.} Place: Lemme 9.7.2.4, ligne supérieure du diagramme (9.7.2.4.1); Reported print: \textbackslash{}Ppoint(F)\textbackslash{}longrightarrow\textbackslash{}Ppoint(E); Source TeX: \textbackslash{}Point(F)\textbackslash{}longrightarrow\textbackslash{}Point(E); Note: The printed symbols are the underlined categories of points, not the sets denoted by \textbackslash{}Point.; Proposed edit: \textbackslash{}Point(F)\textbackslash{}ar[r]\textasciicircum{}\{p \textbackslash{}mapsto f \textbackslash{}circ p\} \textbackslash{}ar[d] \& \textbackslash{}Point(E)\textbackslash{}ar[d] =\allowbreak > \textbackslash{}Ppoint(F)\textbackslash{}ar[r]\textasciicircum{}\{p \textbackslash{}mapsto f \textbackslash{}circ p\} \textbackslash{}ar[d] \& \textbackslash{}Ppoint(E)\textbackslash{}ar[d]\par
a fix from one model only\par
\par\noindent\hypertarget{DP0164}{}\textbf{DP0164 -- IV, p. 480 (agreed\_fix\_held)}\par
9.8.2 b), formula for X\_2 | Exercice 9.8.2 b), formule définissant X\_2\par
\textbf{Claude Opus 5.5.} Place: 9.8.2 b), formula for X\_2; Reported print: X\_2=\allowbreak (R,p\_2)x\_X(R,p\_1); Source TeX: \textbackslash{}fprod\{(R,p\_2)\}\{R,p\_1\}\{X\}; Note: The print puts parentheses around both factors; the TeX gives them only for the first. This finding is not certain: \textbackslash{}fprod's definition is not attached, and it may supply them.; Proposed edit: \textbackslash{}fprod\{(R,p\_2)\}\{R,p\_1\}\{X\} =\allowbreak > \textbackslash{}fprod\{(R,p\_2)\}\{(R,p\_1)\}\{X\}\par
\textbf{GPT-6.1 Sol.} Place: Exercice 9.8.2 b), formule définissant X\_2; Reported print: X\_2=\allowbreak (R,p\_2)\textbackslash{}times\_X(R,p\_1); Source TeX: X\_2=\allowbreak \textbackslash{}fprod\{(R,p\_2)\}\{R,p\_1\}\{X\}; Note: The print parenthesizes the second factor (R,p\_1), as well as the first.; Proposed edit: \textbackslash{}fprod\{(R,p\_2)\}\{R,p\_1\}\{X\} =\allowbreak > \textbackslash{}fprod\{(R,p\_2)\}\{(R,p\_1)\}\{X\}\par
both models give the same fix, held: a model marks it low confidence\par
\par\noindent\hypertarget{DP0165}{}\textbf{DP0165 -- IV, p. 480 (one\_model\_fix)}\par
Exercice 9.8.2 a), premier terme du diagramme de catégories Homtop\par
\textbf{Claude Opus 5.5.} nothing reported here\par
\textbf{GPT-6.1 Sol.} Place: Exercice 9.8.2 a), premier terme du diagramme de catégories Homtop; Reported print: \textbackslash{}Hhomtop(\textbackslash{}Top(X\_0),E); Source TeX: \textbackslash{}Hhomtop(\textbackslash{}Top(X\_o),E); Note: The subscript is the index zero of the semi-simplicial diagram, not the letter o.; Proposed edit: \textbackslash{}Top(X\_o) =\allowbreak > \textbackslash{}Top(X\_0)\par
a fix from one model only\par
\par\noindent\hypertarget{DP0166}{}\textbf{DP0166 -- IV, p. 492 (conflict)}\par
Lemme 10.4.2, last formula | Lemma 10.4.2, final formula\par
\textbf{Claude Opus 5.5.} Place: Lemme 10.4.2, last formula; Reported print: \ensuremath{𝓗}om\_\{E\textasciicircum{}\textasciitilde{}\_V\}(X,Y) with underlined V (universe); Source TeX: \textbackslash{}scrHom\_\{E\textasciicircum{}\{\textbackslash{}sim\}\_\{V\}\}(X,Y); Note: The universe is \ensuremath{𝒱} everywhere else. Here the TeX has a plain V.; Proposed edit: \textbackslash{}scrHom\_\{E\textasciicircum{}\{\textbackslash{}sim\}\_\{V\}\}(X,Y) =\allowbreak > \textbackslash{}scrHom\_\{E\textasciicircum{}\{\textbackslash{}sim\}\_\{\ensuremath{𝒱}\}\}(X,Y)\par
\textbf{GPT-6.1 Sol.} Place: Lemma 10.4.2, final formula; Reported print: \textbackslash{}scrHom\_\{E\textasciicircum{}\{\textbackslash{}sim\}\_\{V\}\}(X,Y), with V not underlined; Source TeX: \textbackslash{}scrHom\_\{E\textasciicircum{}\{\textbackslash{}sim\}\_\{V\}\}(X,Y); Proposed correction: \textbackslash{}scrHom\_\{E\textasciicircum{}\{\textbackslash{}sim\}\_\{\ensuremath{𝒱}\}\}(X,Y); Note: The subscript should be the universe \ensuremath{𝒱} defined immediately before. Its distinguishing underline is absent here in the print; the TeX keeps the plain V.\par
the two models differ at this place: one model fixes the TeX, the other calls it a misprint of the print\par
\par\noindent\hypertarget{DP0167}{}\textbf{DP0167 -- IV, p. 492 (one\_model\_fix)}\par
10.4.1, reference following the equivalence of topoi\par
\textbf{Claude Opus 5.5.} nothing reported here\par
\textbf{GPT-6.1 Sol.} Place: 10.4.1, reference following the equivalence of topoi; Reported print: (III 4 et IV 1); Source TeX: (\textbackslash{}hyperrefext[III][III.sec4]\{4\} et \textbackslash{}hyperref[IV.sec1]\{1\}); Note: The exposé numeral IV is omitted. This recurring reference convention is reported once; three examples are listed in the coordinator.; Proposed edit: \textbackslash{}hyperref[IV.sec1]\{1\} =\allowbreak > IV\textasciitilde{}\textbackslash{}hyperref[IV.sec1]\{1\}\par
a fix from one model only\par
\par\noindent\hypertarget{DP0168}{}\textbf{DP0168 -- IV, p. 492 (doubt)}\par
10.4.3 (four occurrences)\par
\textbf{Claude Opus 5.5.} Place: 10.4.3 (four occurrences); Reported print: E\textasciicircum{}\textasciitilde{}\_V (tilde as superscript, V as subscript); Source TeX: E\_\{\textbackslash{}widetilde\{\ensuremath{𝒱}\}\}; Note: Throughout 10.4.3 the print writes E\textasciicircum{}\textasciitilde{}\_V, which is the notation of 10.4.2. The TeX writes E\_\{\textbackslash{}widetilde\{\ensuremath{𝒱}\}\}, the tilde over the subscript. Notational only. If it is changed, replace all four occurrences with E\textasciicircum{}\{\textbackslash{}sim\}\_\{\ensuremath{𝒱}\}. The fix is a placeholder that changes nothing.; Proposed edit: est pleinement fidèle et exact à gauche =\allowbreak > est pleinement fidèle et exact à gauche\par
\textbf{GPT-6.1 Sol.} nothing reported here\par
a model flags this place without giving a fix\par
\par\noindent\hypertarget{DP0169}{}\textbf{DP0169 -- IV, p. 493 (conflict)}\par
10.5.3, second Hom | 10.5.3, end of the same inline right-hand Hom\par
\textbf{Claude Opus 5.5.} Place: 10.5.3, second Hom; Reported print: Hom\_\{E'\}(v\_!(Xxv*Z,Y) ;; Source TeX: \textbackslash{}Hom\_\{E'\}(v\_\{!\}(X\textbackslash{}times v\textasciicircum{}\{\textbackslash{}ast\}Z),Y))\$; Note: The print lacks a closing parenthesis after Z, a misprint. The TeX corrects it but leaves an extra final parenthesis, so the formula is unbalanced.; Proposed edit: v\textasciicircum{}\{\textbackslash{}ast\}Z),Y))\$ =\allowbreak > v\textasciicircum{}\{\textbackslash{}ast\}Z),Y)\$\par
\textbf{GPT-6.1 Sol.} Place: 10.5.3, end of the same inline right-hand Hom; Reported print: Only one closing parenthesis follows Y.; Source TeX: \textbackslash{}Hom\_\{E'\}(v\_\{!\}(X\textbackslash{}times v\textasciicircum{}\{\textbackslash{}ast\}Z),Y)); Note: There is an extra final parenthesis in the TeX. The replacement records the printed reading in the first argument of \textbackslash{}sisi and the balanced correction in the second.; Proposed edit: \textbackslash{}Hom\_\{E'\}(v\_\{!\}(X\textbackslash{}times v\textasciicircum{}\{\textbackslash{}ast\}Z),Y)) =\allowbreak > \textbackslash{}sisi\{\textbackslash{}Hom\_\{E'\}(v\_\{!\}(X\textbackslash{}times v\textasciicircum{}\{\textbackslash{}ast\}Z,Y)\}\{\textbackslash{}Hom\_\{E'\}(v\_\{!\}(X\textbackslash{}times v\textasciicircum{}\{\textbackslash{}ast\}Z),Y)\}\par
the two models differ at this place\par
\par\noindent\hypertarget{DP0170}{}\textbf{DP0170 -- IV, p. 495 (cut)}\par
Opening formula; boundary between pages 494 and 495\par
\textbf{Claude Opus 5.5.} nothing reported here\par
\textbf{GPT-6.1 Sol.} Place: Opening formula; boundary between pages 494 and 495; Reported print: Page 495 begins with the complete chain \textbackslash{}scrHom(\textbackslash{}amalg\_I v\textasciicircum{}\{\textbackslash{}ast\}(e'),v\textasciicircum{}\{\textbackslash{}ast\}Y)\textbackslash{}simeq\textbackslash{}Pi\_i\textbackslash{}scrHom(v\textasciicircum{}\{\textbackslash{}ast\}(e'),v\textasciicircum{}\{\textbackslash{}ast\}Y)\textbackslash{}simeq\textbackslash{}Pi\_I v\textasciicircum{}\{\textbackslash{}ast\}Y.; Source TeX: v\textasciicircum{}\{\textbackslash{}ast\}Y\$.\textbackslash{}pageoriginale Le morphisme fonctoriel \$\textbackslash{}Phi\_v\$ induit; Note: Page 494 actually ends after « \textbackslash{}Pi\_I Y\$ et ». The opening Hom chain of page 495 has its first two TeX lines assigned to page 494, beyond the overlap allowed by the comments. The cut must be corrected; no text should be inserted or duplicated.\par
a model reports that the automatic cut gives lines of this page to its neighbour, or the neighbour's lines to this page: check the lines it names against the scan of the page that prints them (that page's checker did not have their TeX)\par
\par\noindent\hypertarget{DP0171}{}\textbf{DP0171 -- IV, p. 497 (one\_model\_fix)}\par
11.1.5, égalité initiale\par
\textbf{Claude Opus 5.5.} nothing reported here\par
\textbf{GPT-6.1 Sol.} Place: 11.1.5, égalité initiale; Reported print: A=\allowbreak \textbackslash{}underline\{\textbackslash{}underline\{\textbackslash{}ZZ\}\}; Source TeX: \$A=\allowbreak \textbackslash{}ZZ\$; Note: Le Z porte les mêmes deux soulignements que le faisceau constant dans la parenthèse suivante. Il s'agit d'une notation mathématique, non d'un soulignement d'emphase.; Proposed edit: \$A=\allowbreak \textbackslash{}ZZ\$ =\allowbreak > \$A=\allowbreak \textbackslash{}underline\{\textbackslash{}underline\{\textbackslash{}ZZ\}\}\$\par
a fix from one model only\par
\par\noindent\hypertarget{DP0172}{}\textbf{DP0172 -- IV, p. 497 (one\_model\_fix)}\par
11.2.1, renvoi à IV 3.1.2\par
\textbf{Claude Opus 5.5.} nothing reported here\par
\textbf{GPT-6.1 Sol.} Place: 11.2.1, renvoi à IV 3.1.2; Reported print: III 1.7 ou IV 3.1.2; Source TeX: \textbackslash{}hyperref[IV.subsubsec3.1.2]\{3.1.2\}; Note: Cas représentatif de l'omission récurrente du préfixe IV dans les renvois internes. Cette convention est signalée une seule fois, avec trois exemples dans coordinator.; Proposed edit: \textbackslash{}hyperref[IV.subsubsec3.1.2]\{3.1.2\} =\allowbreak > IV\textasciitilde{}\textbackslash{}hyperref[IV.subsubsec3.1.2]\{3.1.2\}\par
a fix from one model only\par
\par\noindent\hypertarget{DP0173}{}\textbf{DP0173 -- IV, p. 499 (conflict)}\par
proof of 11.3.1, cas général | Même phrase, numéro du renvoi\par
\textbf{Claude Opus 5.5.} Place: proof of 11.3.1, cas général; Reported print: provient d'un objet de C ((IV 1). Le topos; Source TeX: \$C\$ (\textbackslash{}hyperref[IV.cor1.2.1]\{1.2.1\}); Note: The print refers to IV 1, not to 1.2.1. The print also has a doubled opening parenthesis '((', which is a misprint and is not kept in the fix.; Proposed edit: \$C\$ (\textbackslash{}hyperref[IV.cor1.2.1]\{1.2.1\}) =\allowbreak > \$C\$ (1)\par
\textbf{GPT-6.1 Sol.} Place: Même phrase, numéro du renvoi; Reported print: IV 1; Source TeX: \textbackslash{}hyperref[IV.cor1.2.1]\{1.2.1\}; Note: L'imprimé renvoie au §1, sans la précision 1.2.1 ajoutée par le TeX.; Proposed edit: \textbackslash{}hyperref[IV.cor1.2.1]\{1.2.1\} =\allowbreak > IV\textasciitilde{}1\par
the two models differ at this place\par
\par\noindent\hypertarget{DP0174}{}\textbf{DP0174 -- IV, p. 499 (conflict)}\par
proof of 11.3.1 | Démonstration de 11.3.1, construction du prolongement par zéro pour les faisceaux\par
\textbf{Claude Opus 5.5.} Place: proof of 11.3.1; Reported print: Il résulte alors de III 1.7 ii) que le prolongement; Source TeX: \textbackslash{}hyperrefext[III][III.prop1.7]\{1.7\}) que; Note: The print reads 'III 1.7ii)' (shown as '1.711)' in typewriter type). The TeX drops the 'ii' and keeps only the closing parenthesis.; Proposed edit: \{1.7\}) que le prolongement =\allowbreak > \{1.7\}\textasciitilde{}ii) que le prolongement\par
\textbf{GPT-6.1 Sol.} Place: Démonstration de 11.3.1, construction du prolongement par zéro pour les faisceaux; Reported print: Il résulte alors de III 1.7 11) que; Source TeX: \textbackslash{}hyperrefext[III][III.prop1.7]\{1.7\}) que; Note: Le TeX omet le numéro du sous-alinéa 11).; Proposed edit: \textbackslash{}hyperrefext[III][III.prop1.7]\{1.7\}) que =\allowbreak > \textbackslash{}hyperrefext[III][III.prop1.7]\{1.7\} 11) que\par
the two models differ at this place\par
\par\noindent\hypertarget{DP0175}{}\textbf{DP0175 -- IV, p. 500 (one\_model\_fix)}\par
Remarque 11.3.2, premier renvoi parenthétique\par
\textbf{Claude Opus 5.5.} nothing reported here\par
\textbf{GPT-6.1 Sol.} Place: Remarque 11.3.2, premier renvoi parenthétique; Reported print: (III 5.3 et IV 5.2); Source TeX: (\textbackslash{}ref\{III.subsec5.3\} et \textbackslash{}hyperref[IV.subsec5.2]\{5.2\}); Note: Le préfixe III manque devant 5.3. Ce renvoi est extérieur à l'exposé IV, contrairement à la convention des renvois internes signalée séparément.; Proposed edit: \textbackslash{}ref\{III.subsec5.3\} =\allowbreak > \textbackslash{}hyperrefext[III][III.subsec5.3]\{5.3\}\par
a fix from one model only\par
\par\noindent\hypertarget{DP0176}{}\textbf{DP0176 -- IV, p. 501 (one\_model\_fix)}\par
Exercice 11.3.5, renvoi à 6.1\par
\textbf{Claude Opus 5.5.} nothing reported here\par
\textbf{GPT-6.1 Sol.} Place: Exercice 11.3.5, renvoi à 6.1; Reported print: un point de E (IV 6.1); Source TeX: point de \$E\$ (\textbackslash{}hyperref[IV.def6.1]\{6.1\}); Note: Le préfixe IV est imprimé. Cette omission récurrente est signalée une seule fois comme convention, avec trois exemples dans le coordinator.; Proposed edit: point de \$E\$ (\textbackslash{}hyperref[IV.def6.1]\{6.1\}) =\allowbreak > point de \$E\$ (IV \textbackslash{}hyperref[IV.def6.1]\{6.1\})\par
a fix from one model only\par
\par\noindent\hypertarget{DP0177}{}\textbf{DP0177 -- IV, p. 501 (one\_model\_fix)}\par
Exercice 11.3.5, fin de la parenthèse sur la fibre\par
\textbf{Claude Opus 5.5.} nothing reported here\par
\textbf{GPT-6.1 Sol.} Place: Exercice 11.3.5, fin de la parenthèse sur la fibre; Reported print: avec la fibre X\_x IV 6.7.2).; Source TeX: avec la fibre \$X\_x\$ (\textbackslash{}hyperref[IV.eq6.7.2]\{6.7.2\}).; Note: Le print n'ouvre pas une seconde parenthèse avant le renvoi : la parenthèse ouverte devant « en correspondance » se ferme après 6.7.2. Le fix rétablit aussi le préfixe IV imprimé.; Proposed edit: avec la fibre \$X\_x\$ (\textbackslash{}hyperref[IV.eq6.7.2]\{6.7.2\}). =\allowbreak > avec la fibre \$X\_x\$ IV \textbackslash{}hyperref[IV.eq6.7.2]\{6.7.2\}).\par
a fix from one model only\par
\par\noindent\hypertarget{DP0178}{}\textbf{DP0178 -- IV, p. 504 (conflict)}\par
proof of 12.7, 'injection canonique de Hom\_A(M,\ensuremath{𝓗}om\_Z(N,P))' | Preuve de 12.7, première occurrence après l'énoncé\par
\textbf{Claude Opus 5.5.} Place: proof of 12.7, 'injection canonique de Hom\_A(M,\ensuremath{𝓗}om\_Z(N,P))'; Reported print: \ensuremath{𝓗}om\_Z (plain Z, not underlined); Source TeX: \textbackslash{}scrHom\_\{\textbackslash{}sisi\{\textbackslash{}underline\{\textbackslash{}underline\{\textbackslash{}ZZ\}\}\}\{\textbackslash{}underline\{\textbackslash{}ZZ\}\}\}(N,P))\$ dans; Note: The \textbackslash{}sisi reading A (double-underlined Z) does not match the print, which has a plain Z here. Only the statement of 12.7 has the double underline.; Proposed edit: \textbackslash{}scrHom\_\{\textbackslash{}sisi\{\textbackslash{}underline\{\textbackslash{}underline\{\textbackslash{}ZZ\}\}\}\{\textbackslash{}underline\{\textbackslash{}ZZ\}\}\}(N,P))\$ dans =\allowbreak > \textbackslash{}scrHom\_Z(N,P))\$ dans\par
\textbf{GPT-6.1 Sol.} Place: Preuve de 12.7, première occurrence après l'énoncé; Reported print: \textbackslash{}Hom\_A(M,\textbackslash{}scrHom\_Z(N,P)); Source TeX: \textbackslash{}Hom\_A(M,\textbackslash{}scrHom\_\{\textbackslash{}sisi\{\textbackslash{}underline\{\textbackslash{}underline\{\textbackslash{}ZZ\}\}\}\{\textbackslash{}underline\{\textbackslash{}ZZ\}\}\}(N,P)); Note: L'indice mathématique imprimé est ici Z sans soulignement. La branche A de \textbackslash{}sisi lui attribue deux soulignements absents du print ; B n'est pas contrôlé et reste inchangé. Même convention répétée en 12.8 : trois exemples sont donnés dans le coordinator. En revanche, les deux soulignements sont effectivement imprimés dans l'énoncé de 12.7.; Proposed edit: injection canonique de \$\textbackslash{}Hom\_A(M,\textbackslash{}scrHom\_\{\textbackslash{}sisi\{\textbackslash{}underline\{\textbackslash{}underline\{\textbackslash{}ZZ\}\}\}\{\textbackslash{}underline\{\textbackslash{}ZZ\}\}\}(N,P))\$ =\allowbreak > injection canonique de \$\textbackslash{}Hom\_A(M,\textbackslash{}scrHom\_\{\textbackslash{}sisi\{Z\}\{\textbackslash{}underline\{\textbackslash{}ZZ\}\}\}(N,P))\$\par
the two models differ at this place\par
\par\noindent\hypertarget{DP0179}{}\textbf{DP0179 -- IV, p. 505 (one\_model\_fix)}\par
12.8, first line\par
\textbf{Claude Opus 5.5.} Place: 12.8, first line; Reported print: P ↦ Hom\_A(N,\ensuremath{𝓗}om\_Z(M,P)) (plain Z); Source TeX: \textbackslash{}scrHom\_\{\textbackslash{}sisi\{\textbackslash{}underline\{\textbackslash{}underline\{\textbackslash{}ZZ\}\}\}\{\textbackslash{}underline\{\textbackslash{}ZZ\}\}\}(M,P))\$ est; Note: The \textbackslash{}sisi reading A (double-underlined Z) differs from the print, which has a plain Z.; Proposed edit: \textbackslash{}scrHom\_\{\textbackslash{}sisi\{\textbackslash{}underline\{\textbackslash{}underline\{\textbackslash{}ZZ\}\}\}\{\textbackslash{}underline\{\textbackslash{}ZZ\}\}\}(M,P))\$ est =\allowbreak > \textbackslash{}scrHom\_Z(M,P))\$ est\par
\textbf{GPT-6.1 Sol.} nothing reported here\par
a fix from one model only\par
\par\noindent\hypertarget{DP0180}{}\textbf{DP0180 -- IV, p. 505 (one\_model\_fix)}\par
12.8, 'représente aussi le foncteur Hom\_A(N,\ensuremath{𝓗}om\_Z(M,P))'\par
\textbf{Claude Opus 5.5.} Place: 12.8, 'représente aussi le foncteur Hom\_A(N,\ensuremath{𝓗}om\_Z(M,P))'; Reported print: \ensuremath{𝓗}om\_Z (plain Z); Source TeX: \textbackslash{}scrHom\_\{\textbackslash{}sisi\{\textbackslash{}underline\{\textbackslash{}underline\{\textbackslash{}ZZ\}\}\}\{\textbackslash{}underline\{\textbackslash{}ZZ\}\}\}(M,P))\$. On a; Note: The \textbackslash{}sisi reading A differs from the print, which has a plain Z.; Proposed edit: \textbackslash{}scrHom\_\{\textbackslash{}sisi\{\textbackslash{}underline\{\textbackslash{}underline\{\textbackslash{}ZZ\}\}\}\{\textbackslash{}underline\{\textbackslash{}ZZ\}\}\}(M,P))\$. On a =\allowbreak > \textbackslash{}scrHom\_Z(M,P))\$. On a\par
\textbf{GPT-6.1 Sol.} nothing reported here\par
a fix from one model only\par
\par\noindent\hypertarget{DP0181}{}\textbf{DP0181 -- IV, p. 505 (one\_model\_fix)}\par
(12.8.1)\par
\textbf{Claude Opus 5.5.} Place: (12.8.1); Reported print: Hom\_A(M,\ensuremath{𝓗}om\_Z(N,P)) (plain Z); Source TeX: \textbackslash{}scrHom\_\{\textbackslash{}sisi\{\textbackslash{}underline\{\textbackslash{}underline\{\textbackslash{}ZZ\}\}\}\{\textbackslash{}underline\{\textbackslash{}ZZ\}\}\}(N,P))\{\textbackslash{}quoi\},\textbackslash{}tag\{12.8.1\}; Note: The \textbackslash{}sisi reading A differs from the print, which has a plain Z.; Proposed edit: \textbackslash{}scrHom\_\{\textbackslash{}sisi\{\textbackslash{}underline\{\textbackslash{}underline\{\textbackslash{}ZZ\}\}\}\{\textbackslash{}underline\{\textbackslash{}ZZ\}\}\}(N,P))\{\textbackslash{}quoi\},\textbackslash{}tag\{12.8.1\} =\allowbreak > \textbackslash{}scrHom\_Z(N,P))\{\textbackslash{}quoi\},\textbackslash{}tag\{12.8.1\}\par
\textbf{GPT-6.1 Sol.} nothing reported here\par
a fix from one model only\par
\par\noindent\hypertarget{DP0182}{}\textbf{DP0182 -- IV, p. 505 (one\_model\_fix)}\par
(12.8.2)\par
\textbf{Claude Opus 5.5.} Place: (12.8.2); Reported print: Hom\_A(N,\ensuremath{𝓗}om\_Z(M,P)) (plain Z); Source TeX: \textbackslash{}scrHom\_\{\textbackslash{}sisi\{\textbackslash{}underline\{\textbackslash{}underline\{\textbackslash{}ZZ\}\}\}\{\textbackslash{}underline\{\textbackslash{}ZZ\}\}\}(M,P))\{\textbackslash{}quoi\}.\textbackslash{}tag\{12.8.2\}; Note: The \textbackslash{}sisi reading A differs from the print, which has a plain Z.; Proposed edit: \textbackslash{}scrHom\_\{\textbackslash{}sisi\{\textbackslash{}underline\{\textbackslash{}underline\{\textbackslash{}ZZ\}\}\}\{\textbackslash{}underline\{\textbackslash{}ZZ\}\}\}(M,P))\{\textbackslash{}quoi\}.\textbackslash{}tag\{12.8.2\} =\allowbreak > \textbackslash{}scrHom\_Z(M,P))\{\textbackslash{}quoi\}.\textbackslash{}tag\{12.8.2\}\par
\textbf{GPT-6.1 Sol.} nothing reported here\par
a fix from one model only\par
\par\noindent\hypertarget{DP0183}{}\textbf{DP0183 -- IV, p. 508 (one\_model\_fix)}\par
13.1.1, last line\par
\textbf{Claude Opus 5.5.} Place: 13.1.1, last line; Reported print: \ensuremath{\theta} : A' \ensuremath{\to} m\_*A (same typewriter theta as in (m,\ensuremath{\theta})); Source TeX: \textbackslash{}Theta:A'\textbackslash{}rightarrow m\_\{\textbackslash{}ast\}A; Note: The print uses the same glyph throughout; the TeX switches to a capital Theta here only.; Proposed edit: \textbackslash{}Theta:A' =\allowbreak > \textbackslash{}theta:A'\par
\textbf{GPT-6.1 Sol.} nothing reported here\par
a fix from one model only\par
\par\noindent\hypertarget{DP0184}{}\textbf{DP0184 -- IV, p. 508 (one\_model\_fix)}\par
Définition 13.1, renvoi à 3.1\par
\textbf{Claude Opus 5.5.} nothing reported here\par
\textbf{GPT-6.1 Sol.} Place: Définition 13.1, renvoi à 3.1; Reported print: (IV.3.1); Source TeX: (\textbackslash{}hyperref[IV.def3.1]\{3.1\}); Note: Omission du numéro d'exposé dans les renvois internes : convention signalée une seule fois, avec trois exemples dans le message au coordinateur.; Proposed edit: (\textbackslash{}hyperref[IV.def3.1]\{3.1\}) =\allowbreak > (\textbackslash{}hyperref[IV.def3.1]\{IV.3.1\})\par
a fix from one model only\par
\par\noindent\hypertarget{DP0185}{}\textbf{DP0185 -- IV, p. 509 (one\_model\_fix)}\par
13.2.1\par
\textbf{Claude Opus 5.5.} Place: 13.2.1; Reported print: \ensuremath{\theta}' : A' \ensuremath{\to} m\_*A; Source TeX: \textbackslash{}Theta':A'\textbackslash{}rightarrow m\_\{\textbackslash{}ast\}A; Note: The print uses the same theta glyph as in (m,\ensuremath{\theta}); the TeX has a capital Theta.; Proposed edit: \textbackslash{}Theta':A' =\allowbreak > \textbackslash{}theta':A'\par
\textbf{GPT-6.1 Sol.} nothing reported here\par
a fix from one model only\par
\par\noindent\hypertarget{DP0186}{}\textbf{DP0186 -- IV, p. 510 (agreed\_fix\_held)}\par
proof of 13.4.2, second line of the sequence | Suite d'isomorphismes finale, deuxième ligne, renvoi à 12.1\par
\textbf{Claude Opus 5.5.} Place: proof of 13.4.2, second line of the sequence; Reported print: Hom\_A(u\textasciicircum{}*N, \ensuremath{\mathcal{H}}om\_E(u\textasciicircum{}\{-1\}X,M))   "  "  "  "   (12.1),; Source TeX: \textbackslash{}Hom\_A(u\textasciicircum{}\{\textbackslash{}ast\}N,\textbackslash{}scrHom\_E (u\textasciicircum{}\{-1\}X, M)) \textbackslash{}simeq \&\&; Note: The print has only ditto marks after the left-hand side; the TeX adds an explicit \ensuremath{\simeq}.; Proposed edit: (u\textasciicircum{}\{-1\}X, M)) \textbackslash{}simeq \&\& =\allowbreak > (u\textasciicircum{}\{-1\}X, M)) \&\&\par
\textbf{GPT-6.1 Sol.} Place: Suite d'isomorphismes finale, deuxième ligne, renvoi à 12.1; Reported print: \textbackslash{}Hom\_A(u\textasciicircum{}\{\textbackslash{}ast\}N,\textbackslash{}scrHom\_E(u\textasciicircum{}\{-1\}X,M)), puis les guillemets de répétition sans signe d'isomorphisme; Source TeX: \textbackslash{}Hom\_A(u\textasciicircum{}\{\textbackslash{}ast\}N,\textbackslash{}scrHom\_E (u\textasciicircum{}\{-1\}X, M)) \textbackslash{}simeq; Note: Le TeX ajoute un signe d'isomorphisme explicite avant les guillemets de répétition ; aucun n'est imprimé sur cette ligne.; Proposed edit: \textbackslash{}scrHom\_E (u\textasciicircum{}\{-1\}X, M)) \textbackslash{}simeq =\allowbreak > \textbackslash{}scrHom\_E (u\textasciicircum{}\{-1\}X, M))\par
both models give the same fix, held: a model marks it low confidence\par
\par\noindent\hypertarget{DP0187}{}\textbf{DP0187 -- IV, p. 510 (cut)}\par
end of page\par
\textbf{Claude Opus 5.5.} Place: end of page; Reported print: the page ends with the line Hom\_A(u\textasciicircum{}*N,\ensuremath{\mathcal{H}}om\_E(u\textasciicircum{}\{-1\}X,M)) " " " " (12.1) ,; Source TeX: three further alignat lines: (13.4.1), (10.3.1), (12.1); Note: The last three lines of the alignat are not on page 510; they belong to page 511 (IV-212). Page 510 ends after the line ending in (12.1).\par
\textbf{GPT-6.1 Sol.} nothing reported here\par
a model reports that the automatic cut gives lines of this page to its neighbour, or the neighbour's lines to this page: check the lines it names against the scan of the page that prints them (that page's checker did not have their TeX)\par
\par\noindent\hypertarget{DP0188}{}\textbf{DP0188 -- IV, p. 510 (cut)}\par
Fin de la page, suite d'isomorphismes\par
\textbf{Claude Opus 5.5.} nothing reported here\par
\textbf{GPT-6.1 Sol.} Place: Fin de la page, suite d'isomorphismes; Reported print: La page s'arrête après la deuxième ligne de la suite, celle portant le renvoi (12.1).; Source TeX: L'alignat contient encore trois lignes mathématiques après cette deuxième ligne.; Note: Ces trois dernières lignes appartiennent à la page suivante, 511, et non à 510. Le commentaire de découpage ne les signale pas comme débordement ; elles ne sont pas contrôlées ici.\par
a model reports that the automatic cut gives lines of this page to its neighbour, or the neighbour's lines to this page: check the lines it names against the scan of the page that prints them (that page's checker did not have their TeX)\par
\par\noindent\hypertarget{DP0189}{}\textbf{DP0189 -- IV, p. 511 (conflict)}\par
Corollaire 13.5 | Corollaire 13.5, référence du foncteur fibre\par
\textbf{Claude Opus 5.5.} Place: Corollaire 13.5; Reported print: le foncteur fibre associé (IV 6.1); Source TeX: (\textbackslash{}hyperref[IV.def6.1]\{6.1\}); Note: The print gives the exposé numeral 'IV' in this internal reference; the edition drops it.; Proposed edit: (\textbackslash{}hyperref[IV.def6.1]\{6.1\}) =\allowbreak > (IV\textasciitilde{}\textbackslash{}hyperref[IV.def6.1]\{6.1\})\par
\textbf{GPT-6.1 Sol.} Place: Corollaire 13.5, référence du foncteur fibre; Reported print: (IV 6.1); Source TeX: (\textbackslash{}hyperref[IV.def6.1]\{6.1\}); Note: Le numéro d'exposé IV, explicitement imprimé, manque dans le renvoi.; Proposed edit: (\textbackslash{}hyperref[IV.def6.1]\{6.1\}) =\allowbreak > (\textbackslash{}hyperrefext[IV][IV.def6.1]\{6.1\})\par
the two models differ at this place\par
\par\noindent\hypertarget{DP0190}{}\textbf{DP0190 -- IV, p. 511 (cut)}\par
top of page, before 'La formule 13.4.5'\par
\textbf{Claude Opus 5.5.} Place: top of page, before 'La formule 13.4.5'; Reported print: The page opens with the three-line display: Hom\_A(u\textasciicircum{}*N, \ensuremath{\mathcal{H}}om\_E(u\textasciicircum{}\{-1\}X,M)) \ensuremath{\simeq} Hom\_\{A'\}(N, u\_*\ensuremath{\mathcal{H}}om\_E(u\textasciicircum{}\{-1\}X,M)) (13.4.1), Hom\_\{A'\}(N, \ensuremath{\mathcal{H}}om\_\{E'\}(X,u\_*M)) " " " " (10.3.1), " " " " Hom\_\{A'\}(X, \ensuremath{\mathcal{H}}om\_\{A'\}(N,u\_*M)) (12.1).; Source TeX: \textbackslash{}end\{alignat*\}; Note: Page 511 begins with this display. In the TeX only its \textbackslash{}end\{alignat*\} falls in this page's cut; the display's lines are in the source lines before 9382, i.e. in p. 510's cut. Check the display against this scan and move the page start back to the display's \textbackslash{}begin\{alignat*\}.\par
\textbf{GPT-6.1 Sol.} nothing reported here\par
a model reports that the automatic cut gives lines of this page to its neighbour, or the neighbour's lines to this page: check the lines it names against the scan of the page that prints them (that page's checker did not have their TeX)\par
\par\noindent\hypertarget{DP0191}{}\textbf{DP0191 -- IV, p. 511 (cut)}\par
Début de page, avant « La formule 13.4.5 »\par
\textbf{Claude Opus 5.5.} nothing reported here\par
\textbf{GPT-6.1 Sol.} Place: Début de page, avant « La formule 13.4.5 »; Reported print: Trois lignes d'isomorphismes portant les renvois (13.4.1), (10.3.1) et (12.1).; Source TeX: \textbackslash{}end\{alignat*\} La\textbackslash{}pageoriginale formule (13.4.5); Note: La coupe omet les trois lignes du display imprimées en haut de cette page et n'en fournit que la fermeture. Le début réel de p. 511 est ce display, avant « La formule 13.4.5 ». Recouper l'extrait avec la page précédente, sans ajouter ces lignes au texte source.\par
a model reports that the automatic cut gives lines of this page to its neighbour, or the neighbour's lines to this page: check the lines it names against the scan of the page that prints them (that page's checker did not have their TeX)\par
\par\noindent\hypertarget{DP0192}{}\textbf{DP0192 -- IV, p. 513 (conflict)}\par
d) | Exercice 13.9 d), référence du topos associé à un espace annelé\par
\textbf{Claude Opus 5.5.} Place: d); Reported print: espace topologique annelé (X,O\_X) (IV 2.1); Source TeX: (\textbackslash{}hyperref[IV.subsec2.1]\{2.1\}); Note: The print gives 'IV 2.1'; the edition drops the exposé numeral.; Proposed edit: (\textbackslash{}hyperref[IV.subsec2.1]\{2.1\}) =\allowbreak > (IV\textasciitilde{}\textbackslash{}hyperref[IV.subsec2.1]\{2.1\})\par
\textbf{GPT-6.1 Sol.} Place: Exercice 13.9 d), référence du topos associé à un espace annelé; Reported print: (IV 2.1); Source TeX: (\textbackslash{}hyperref[IV.subsec2.1]\{2.1\}); Note: Le numéro d'exposé IV, présent dans l'imprimé, est omis.; Proposed edit: (\textbackslash{}hyperref[IV.subsec2.1]\{2.1\}) =\allowbreak > (\textbackslash{}hyperrefext[IV][IV.subsec2.1]\{2.1\})\par
the two models differ at this place\par
\par\noindent\hypertarget{DP0193}{}\textbf{DP0193 -- IV, p. 517 (one\_model\_fix)}\par
14.7, référence après le monomorphisme X \ensuremath{\times} U \ensuremath{\to} X\par
\textbf{Claude Opus 5.5.} nothing reported here\par
\textbf{GPT-6.1 Sol.} Place: 14.7, référence après le monomorphisme X \ensuremath{\times} U \ensuremath{\to} X; Reported print: (IV 5); Source TeX: (\textbackslash{}hyperref[IV.sec5]\{5\}); Note: Le numéro d'exposé IV manque dans la référence transcrite.; Proposed edit: \textbackslash{}hyperref[IV.sec5]\{5\} =\allowbreak > IV\textasciitilde{}\textbackslash{}hyperref[IV.sec5]\{5\}\par
a fix from one model only\par
\par\noindent\hypertarget{DP0194}{}\textbf{DP0194 -- IV, p. 519 (one\_model\_fix)}\par
[6]\par
\textbf{Claude Opus 5.5.} Place: [6]; Reported print: North Holland Pub. Cie, 1969; Source TeX: North Holland Pub, Cie, 1969; Note: The print has a period after 'Pub', not a comma.; Proposed edit: North Holland Pub, Cie =\allowbreak > North Holland Pub. Cie\par
\textbf{GPT-6.1 Sol.} nothing reported here\par
a fix from one model only\par
\par\noindent\hypertarget{DP0195}{}\textbf{DP0195 -- IV, p. 519 (agreed\_fix\_held)}\par
bibliography, entries after [14] | Fin de la bibliographie, après [14]\par
\textbf{Claude Opus 5.5.} Place: bibliography, entries after [14]; Reported print: (the bibliography ends with [14] Mitchell); Source TeX: \textbackslash{}bibitem\{SGA4-04:15\} J. Lurie ... \textbackslash{}bibitem\{SGA4-04:16\} P. Johnstone ... \textbackslash{}bibitem\{SGA4-04:17\} F. Borceux, G. Van den Bossche ...; Note: Entries 15-17 are editors' additions, not in the print, and are not marked as N.D.E. They may be cited in editors' notes elsewhere.; Proposed edit: \textbackslash{}bibitem\{SGA4-04:15\} J. Lurie, Higher Topos Theory, Annals of Mathematics Studies, 170. Princeton University Press,  Princeton, NJ, 2009. xviii+925 pp.  \textbackslash{}bibitem\{SGA4-04:16\}  P. Johnstone, Topos theory,  London Mathematical Society Monographs, Vol. 10.  Academic Press [Harcourt Brace Jovanovich, Publishers], London-New York, 1977. xxiii+367 pp.  \textbackslash{}bibitem\{SGA4-04:17\} F. Borceux, G. Van den Bossche, Structure des topologies d'un topos, Cahiers de topologie et géométrie différentielle catégoriques, t. 25, n°1 (1984), p 37-39.  =\allowbreak > \par
\textbf{GPT-6.1 Sol.} Place: Fin de la bibliographie, après [14]; Reported print: La bibliographie se termine à [14], B. Mitchell.; Source TeX: Entrées supplémentaires [15] J. Lurie, [16] P. Johnstone et [17] F. Borceux–G. Van den Bossche.; Note: Ces trois références ajoutées par l'édition ne figurent pas dans l'imprimé.; Proposed edit: \textbackslash{}bibitem\{SGA4-04:15\} J. Lurie, Higher Topos Theory, Annals of Mathematics Studies, 170. Princeton University Press,  Princeton, NJ, 2009. xviii+925 pp.  \textbackslash{}bibitem\{SGA4-04:16\}  P. Johnstone, Topos theory,  London Mathematical Society Monographs, Vol. 10.  Academic Press [Harcourt Brace Jovanovich, Publishers], London-New York, 1977. xxiii+367 pp.  \textbackslash{}bibitem\{SGA4-04:17\} F. Borceux, G. Van den Bossche, Structure des topologies d'un topos, Cahiers de topologie et géométrie différentielle catégoriques, t. 25, n°1 (1984), p 37-39. =\allowbreak > \par
both models give the same fix, held: a deletion in the first or last three lines of the page (the page cut may be off)\par
\par\noindent\hypertarget{DP0196}{}\textbf{DP0196 -- V, p. 2 (cut)}\par
Opening block before \textbackslash{}section*\{Introduction\}\par
\textbf{Claude Opus 5.5.} nothing reported here\par
\textbf{GPT-6.1 Sol.} Place: Opening block before \textbackslash{}section*\{Introduction\}; Reported print: The page begins with « Introduction », followed by « On présente… ».; Source TeX: \textbackslash{}chapter\{Cohomologie dans les topos\}\textbackslash{}etiquette[V]\{V\} … J.-L. Verdier; Note: The initial chapter-title and author block is not printed on this page, and the extraction comment does not allow preceding-page material here. Page 2 begins at \textbackslash{}section*\{Introduction\}. Check the extraction boundary rather than deleting this text.\par
a model reports that the automatic cut gives lines of this page to its neighbour, or the neighbour's lines to this page: check the lines it names against the scan of the page that prints them (that page's checker did not have their TeX)\par
\par\noindent\hypertarget{DP0197}{}\textbf{DP0197 -- V, p. 7 (one\_model\_fix)}\par
Prop. 1.2, 3)\par
\textbf{Claude Opus 5.5.} Place: Prop. 1.2, 3); Reported print: de modules plats; Source TeX: de Modules plats; Note: The print has a lower-case 'modules' here; elsewhere the text tells Module from module by the capital.; Proposed edit: acyclique de Modules plats =\allowbreak > acyclique de modules plats\par
\textbf{GPT-6.1 Sol.} nothing reported here\par
a fix from one model only\par
\par\noindent\hypertarget{DP0198}{}\textbf{DP0198 -- V, p. 7 (one\_model\_fix)}\par
Beginning of the proof of c) implies a)\par
\textbf{Claude Opus 5.5.} nothing reported here\par
\textbf{GPT-6.1 Sol.} Place: Beginning of the proof of c) implies a); Reported print: c) \$\textbackslash{}to\$ a); Source TeX: c) \$\textbackslash{}Rightarrow\$ a); Note: The printed implication has a single shaft, not a double shaft.; Proposed edit: c) \$\textbackslash{}Rightarrow\$ a) =\allowbreak > c) \$\textbackslash{}to\$ a)\par
a fix from one model only\par
\par\noindent\hypertarget{DP0199}{}\textbf{DP0199 -- V, p. 8 (one\_model\_fix)}\par
Proof of Proposition 1.2, second occurrence of the equivalence (i), (ii)\par
\textbf{Claude Opus 5.5.} nothing reported here\par
\textbf{GPT-6.1 Sol.} Place: Proof of Proposition 1.2, second occurrence of the equivalence (i), (ii); Reported print: d'après l'équivalence (i) \$\textbackslash{}leftrightarrow\$ (ii); Source TeX: d'après l'équivalence (i) \$\textbackslash{}Leftrightarrow\$ (ii); Note: This handwritten arrow has a single shaft, unlike the double-shaft equivalence arrow earlier in the proof.; Proposed edit: d'après l'équivalence (i) \$\textbackslash{}Leftrightarrow\$ (ii) =\allowbreak > d'après l'équivalence (i) \$\textbackslash{}leftrightarrow\$ (ii)\par
a fix from one model only\par
\par\noindent\hypertarget{DP0200}{}\textbf{DP0200 -- V, p. 10 (one\_model\_fix)}\par
after (1.5.2)\par
\textbf{Claude Opus 5.5.} Place: after (1.5.2); Reported print: F(U,V,m) .  On obtient un isomorphisme; Source TeX: F(U, V, m)\{\textbackslash{}quoi\},\textbackslash{}tag\{1.5.2\} ... on obtient; Note: The print ends (1.5.2) with a full stop and starts a new sentence with 'On obtient'. If this is applied, 'on obtient' must also take the capital 'On'.; Proposed edit: F(U, V, m)\{\textbackslash{}quoi\},\textbackslash{}tag\{1.5.2\} =\allowbreak > F(U, V, m)\{\textbackslash{}quoi\}.\textbackslash{}tag\{1.5.2\}\par
\textbf{GPT-6.1 Sol.} nothing reported here\par
a fix from one model only\par
\par\noindent\hypertarget{DP0201}{}\textbf{DP0201 -- V, p. 10 (agreed\_fix\_held)}\par
(1.5.4), right-hand side | Formula (1.5.4), index of the last direct limit\par
\textbf{Claude Opus 5.5.} Place: (1.5.4), right-hand side; Reported print: (vois(x)/V)° (the x looks overstruck); Source TeX: (\textbackslash{}Vois(-x)/V)\textasciicircum{}\{\textbackslash{}circ\}; Note: The print shows a typist's overstrike on the x, not a minus sign. The TeX's '-x' is spurious.; Proposed edit: (\textbackslash{}Vois(-x)/V) =\allowbreak > (\textbackslash{}Vois(x)/V)\par
\textbf{GPT-6.1 Sol.} Place: Formula (1.5.4), index of the last direct limit; Reported print: \textbackslash{}varinjlim\_\{(\textbackslash{}Vois(x)/V)\textasciicircum{}\{\textbackslash{}circ\}\}; Source TeX: \textbackslash{}varinjlim\_\{(\textbackslash{}Vois(-x)/V)\textasciicircum{}\{\textbackslash{}circ\}\}; Note: The printed category index is vois(x)/V, not vois(-x)/V; the minus in the TeX is extra.; Proposed edit: \textbackslash{}Vois(-x) =\allowbreak > \textbackslash{}Vois(x)\par
both models give the same fix, held: a model marks it low confidence\par
\par\noindent\hypertarget{DP0202}{}\textbf{DP0202 -- V, p. 12 (conflict)}\par
1.6.2, third term | 1.6.2, produit tensoriel à l'intérieur de x\_V\textasciicircum{}*\par
\textbf{Claude Opus 5.5.} Place: 1.6.2, third term; Reported print: x\textasciicircum{}*\_V(x\_\{V*\}N \ensuremath{\otimes}\_\{A\_x\} j\textasciicircum{}*\_V M); Source TeX: x\textasciicircum{}\textbackslash{}ast\_V (x\_\{V\_\textbackslash{}ast\} N; Note: The print has x\_\{V*\}, as everywhere else; the TeX nests the star as a subscript of V.; Proposed edit: x\_\{V\_\textbackslash{}ast\} N =\allowbreak > x\_\{V\textbackslash{}ast\} N\par
\textbf{GPT-6.1 Sol.} Place: 1.6.2, produit tensoriel à l'intérieur de x\_V\textasciicircum{}*; Reported print: x\_V\textasciicircum{}*(x\_\{V*\}N \textbackslash{}otimes\_\{A\_x\} j\_V\textasciicircum{}*M); Source TeX: x\textasciicircum{}\textbackslash{}ast\_V (x\_\{V\_\textbackslash{}ast\} N \textbackslash{}otimes\_\{A\_x\}   j\textasciicircum{}\textbackslash{}ast\_V M); Proposed correction: x\textasciicircum{}\textbackslash{}ast\_V (x\_\{V\textbackslash{}ast\} N \textbackslash{}otimes\_\{A/V\} j\textasciicircum{}\textbackslash{}ast\_V M); Note: Ce produit tensoriel est formé dans E/V : son anneau de base doit être A/V, non A\_x. Le TeX conserve l'indice erroné de l'imprimé.\par
the two models differ at this place: one model fixes the TeX, the other calls it a misprint of the print\par
\par\noindent\hypertarget{DP0203}{}\textbf{DP0203 -- V, p. 20 (conflict)}\par
Corollaire 2.3.5, Les morphismes Φ et Φ' induisent les morphismes Φ\textasciicircum{}q et Φ'\textasciicircum{}q : Φ'\textasciicircum{}q : | Corollaire 2.3.5, introduction des morphismes cohomologiques | Corollaire 2.3.5, sens des morphismes cohomologiques\par
\textbf{Claude Opus 5.5.} Place: Corollaire 2.3.5, Les morphismes Φ et Φ' induisent les morphismes Φ\textasciicircum{}q et Φ'\textasciicircum{}q : Φ'\textasciicircum{}q :; Reported print: Φ et Φ' ... Φ\textasciicircum{}q et Φ'\textasciicircum{}q : Φ'\textasciicircum{}q; Source TeX: \$\textbackslash{}phi\$ et \$\textbackslash{}phi'\$ ... \$\textbackslash{}phi\textasciicircum{}\{q\}\$ et \$\{\textbackslash{}phi'\}\textasciicircum{}\{q\}\$; Note: Capital Φ in the print.; Proposed edit: Les morphismes \$\textbackslash{}phi\$ et \$\textbackslash{}phi'\$ induisent les   morphismes \$\textbackslash{}phi\textasciicircum{}\{q\}\$ et \$\{\textbackslash{}phi'\}\textasciicircum{}\{q\} : \textbackslash{}sisi\{\{\textbackslash{}phi'\}\textasciicircum{}\{q\} :\}\{\} =\allowbreak > Les morphismes \$\textbackslash{}Phi\$ et \$\textbackslash{}Phi'\$ induisent les   morphismes \$\textbackslash{}Phi\textasciicircum{}\{q\}\$ et \$\{\textbackslash{}Phi'\}\textasciicircum{}\{q\} : \textbackslash{}sisi\{\{\textbackslash{}Phi'\}\textasciicircum{}\{q\} :\}\{\}\par
\textbf{GPT-6.1 Sol.} Place: Corollaire 2.3.5, introduction des morphismes cohomologiques; Reported print: \textbackslash{}Phi\textasciicircum{}q et \{\textbackslash{}Phi'\}\textasciicircum{}q : \{\textbackslash{}Phi'\}\textasciicircum{}q : H\textasciicircum{}q; Source TeX: \textbackslash{}phi\textasciicircum{}\{q\}\$ et \$\{\textbackslash{}phi'\}\textasciicircum{}\{q\} : \textbackslash{}sisi\{\{\textbackslash{}phi'\}\textasciicircum{}\{q\} :\}\{\} H\textasciicircum{}q; Proposed correction: \textbackslash{}Phi\textasciicircum{}\{q\}\textbackslash{}text\{ et \}\{\textbackslash{}Phi'\}\textasciicircum{}\{q\} : H\textasciicircum{}q; Note: L'étiquette Φ'\textasciicircum{}q suivie de deux-points est répétée dans l'imprimé. L'édition supprime la répétition dans la seconde branche de \textbackslash{}sisi. || Place: Corollaire 2.3.5, sens des morphismes cohomologiques; Reported print: H\textasciicircum{}q(\textbackslash{}mathcal\{U\},M) \textbackslash{}to H\textasciicircum{}q(\textbackslash{}mathcal\{U\}',M); Source TeX: H\textasciicircum{}q(\textbackslash{}mathcal\{U\},   M) \textbackslash{}to H\textasciicircum{}q(\textbackslash{}mathcal\{U\}', M); Proposed correction: H\textasciicircum{}q(\textbackslash{}mathcal\{U\}',M) \textbackslash{}to H\textasciicircum{}q(\textbackslash{}mathcal\{U\},M); Note: Les morphismes de familles U vers U' induisent des morphismes de cohomologie dans le sens inverse, comme le confirme la preuve par Hom appliqué aux résolutions. Le TeX conserve le sens fautif de l'imprimé.\par
the two models differ at this place: one model fixes the TeX, the other calls it a misprint of the print\par
\par\noindent\hypertarget{DP0204}{}\textbf{DP0204 -- V, p. 20 (one\_model\_fix)}\par
Corollaire 2.3.5, Les morphismes Φ\textasciicircum{}q et Φ'\textasciicircum{}q sont égaux\par
\textbf{Claude Opus 5.5.} Place: Corollaire 2.3.5, Les morphismes Φ\textasciicircum{}q et Φ'\textasciicircum{}q sont égaux; Reported print: Φ\textasciicircum{}q et Φ'\textasciicircum{}q; Source TeX: \$\textbackslash{}phi\textasciicircum{}q\$ et \$\{\textbackslash{}phi'\}\textasciicircum{}\{q\}\$; Note: Capital Φ in the print.; Proposed edit: morphismes \$\textbackslash{}phi\textasciicircum{}q\$ et \$\{\textbackslash{}phi'\}\textasciicircum{}\{q\}\$ sont égaux =\allowbreak > morphismes \$\textbackslash{}Phi\textasciicircum{}q\$ et \$\{\textbackslash{}Phi'\}\textasciicircum{}\{q\}\$ sont égaux\par
\textbf{GPT-6.1 Sol.} nothing reported here\par
a fix from one model only\par
\par\noindent\hypertarget{DP0205}{}\textbf{DP0205 -- V, p. 20 (one\_model\_fix)}\par
Preuve, Les morphismes Φ et Φ' définissent\par
\textbf{Claude Opus 5.5.} Place: Preuve, Les morphismes Φ et Φ' définissent; Reported print: Φ et Φ'; Source TeX: \$\textbackslash{}phi\$ et \$\textbackslash{}phi'\$; Note: Capital Φ in the print.; Proposed edit: Les morphismes \$\textbackslash{}phi\$ et \$\textbackslash{}phi'\$ définissent =\allowbreak > Les morphismes \$\textbackslash{}Phi\$ et \$\textbackslash{}Phi'\$ définissent\par
\textbf{GPT-6.1 Sol.} nothing reported here\par
a fix from one model only\par
\par\noindent\hypertarget{DP0206}{}\textbf{DP0206 -- V, p. 21 (agreed\_fix\_held)}\par
item c) | item c), presheaf ES(C) | c), notation du préfaisceau semi-simplicial | c), valeur du préfaisceau\par
\textbf{Claude Opus 5.5.} Place: item c); Reported print: dont la valeur en tout objet X; Source TeX: dont la valeurs en tout objet; Proposed edit: la valeurs en tout =\allowbreak > la valeur en tout || Place: item c), presheaf ES(C); Reported print: préfaisceau semi-simplicial ES(C) with ES underlined (also in Z\_\{ES(C)\} throughout); Source TeX: semi-simplicial \$ES(C)\$ dont; Note: The print tells the presheaf (underlined ES) from the simplicial set ES(C) of a); the TeX does not.; Proposed edit: semi-simplicial \$ES(C)\$ dont =\allowbreak > semi-simplicial \$\textbackslash{}underline\{ES\}(C)\$ dont\par
\textbf{GPT-6.1 Sol.} Place: c), notation du préfaisceau semi-simplicial; Reported print: \textbackslash{}underline\{ES\}(C); Source TeX: semi-simplicial \$ES(C)\$; Note: Ce soulignement est une notation mathématique : il distingue le préfaisceau de l'ensemble semi-simplicial ES(C), notamment de sa valeur ES(\{\}\_\{X\textbackslash{}backslash\}C). Son omission est systématique pour le préfaisceau et ses occurrences en indice ; trois exemples sont indiqués au coordinateur.; Proposed edit: semi-simplicial \$ES(C)\$ =\allowbreak > semi-simplicial \$\textbackslash{}underline\{ES\}(C)\$ || Place: c), valeur du préfaisceau; Reported print: dont la valeur en tout objet X; Source TeX: dont la valeurs en tout objet \$X\$; Note: Le s final de « valeurs » n'est pas imprimé.; Proposed edit: dont la valeurs =\allowbreak > dont la valeur\par
both models give the same fix, held: a model marks it low confidence\par
\par\noindent\hypertarget{DP0207}{}\textbf{DP0207 -- V, p. 21 (one\_model\_fix)}\par
g), toute fin de l'isomorphisme\par
\textbf{Claude Opus 5.5.} nothing reported here\par
\textbf{GPT-6.1 Sol.} Place: g), toute fin de l'isomorphisme; Reported print: H\_q(ST\_\{∙\}(M) \textbackslash{}sim H\_q(C,M); Source TeX: H\_q(ST\_\{∙\}(M) \textbackslash{}sim H\_q(C, M)); Note: Le TeX ajoute une parenthèse après le second H\_q, absente de l'imprimé. Cette position ne corrige pas la parenthèse manquante après le premier H\_q, signalée comme coquille de l'imprimé.; Proposed edit: H\_q(ST\_\{∙\}(M) \textbackslash{}sim H\_q(C, M)) =\allowbreak > H\_q(ST\_\{∙\}(M) \textbackslash{}sim H\_q(C, M)\par
a fix from one model only\par
\par\noindent\hypertarget{DP0208}{}\textbf{DP0208 -- V, p. 22 (one\_model\_fix)}\par
heading 2.4\par
\textbf{Claude Opus 5.5.} Place: heading 2.4; Reported print: Cas des petits sites. Cohomologie de Čech.; Source TeX: Cas des petits sites, Cohomologie de Čech.; Proposed edit: Cas des petits sites, Cohomologie =\allowbreak > Cas des petits sites. Cohomologie\par
\textbf{GPT-6.1 Sol.} nothing reported here\par
a fix from one model only\par
\par\noindent\hypertarget{DP0209}{}\textbf{DP0209 -- V, p. 22 (one\_model\_fix)}\par
i), dernier complexe de la parenthèse « resp. »\par
\textbf{Claude Opus 5.5.} nothing reported here\par
\textbf{GPT-6.1 Sol.} Place: i), dernier complexe de la parenthèse « resp. »; Reported print: ST\_\{∙\}(C\textasciicircum{}\{\textbackslash{}circ\},\textbackslash{}check\{M\}); Source TeX: \$ST.(C\textasciicircum{}\{\textbackslash{}circ\}, \textbackslash{}check\{M\})\$; Note: Le point est l'indice homologique de ST, comme dans le complexe immédiatement précédent.; Proposed edit: \$ST.(C\textasciicircum{}\{\textbackslash{}circ\}, \textbackslash{}check\{M\})\$ =\allowbreak > \$ST\_\{∙\}(C\textasciicircum{}\{\textbackslash{}circ\}, \textbackslash{}check\{M\})\$\par
a fix from one model only\par
\par\noindent\hypertarget{DP0210}{}\textbf{DP0210 -- V, p. 23 (one\_model\_fix)}\par
2.4.3, family \ensuremath{𝒰} =\allowbreak  (\ensuremath{𝒰}\_i \ensuremath{\to} X)\par
\textbf{Claude Opus 5.5.} Place: 2.4.3, family \ensuremath{𝒰} =\allowbreak  (\ensuremath{𝒰}\_i \ensuremath{\to} X); Reported print: script \ensuremath{𝒰}\_i; Source TeX: (U\_i \textbackslash{}to X); Note: The handwritten U\_i appears to be script, like the family \ensuremath{𝒰}.; Proposed edit: (U\_i \textbackslash{}to X) =\allowbreak > (\textbackslash{}Uu\_i \textbackslash{}to X)\par
\textbf{GPT-6.1 Sol.} nothing reported here\par
a fix from one model only\par
\par\noindent\hypertarget{DP0211}{}\textbf{DP0211 -- V, p. 23 (conflict)}\par
2.4.3, 'complexe de Čech de G relatif à la famille couvrante \ensuremath{𝒰} (ou du recouvrement \ensuremath{𝒰})' | 2.4.3, appellation du complexe de Čech\par
\textbf{Claude Opus 5.5.} Place: 2.4.3, 'complexe de Čech de G relatif à la famille couvrante \ensuremath{𝒰} (ou du recouvrement \ensuremath{𝒰})'; Reported print: relatif à la famille couvrante \ensuremath{𝒰} (ou; Source TeX: relatifs à la famille couvrante \$\textbackslash{}mathcal\{U\}\$.\} (ou; Note: The print has the singular 'relatif' and no period before the parenthesis.; Proposed edit: relatifs à la famille couvrante \$\textbackslash{}mathcal\{U\}\$.\} (ou =\allowbreak > relatif à la famille couvrante \$\textbackslash{}mathcal\{U\}\$\} (ou\par
\textbf{GPT-6.1 Sol.} Place: 2.4.3, appellation du complexe de Čech; Reported print: complexe de Čech de G relatif à la famille couvrante; Source TeX: relatifs à la famille couvrante \$\textbackslash{}mathcal\{U\}\$; Note: L'imprimé porte le singulier « relatif », accordé à « complexe ».; Proposed edit: relatifs à la famille couvrante \$\textbackslash{}mathcal\{U\}\$ =\allowbreak > relatif à la famille couvrante \$\textbackslash{}mathcal\{U\}\$\par
the two models differ at this place\par
\par\noindent\hypertarget{DP0212}{}\textbf{DP0212 -- V, p. 27 (one\_model\_fix)}\par
Proof of Corollary 3.4, equality preceding the reference II 3.4\par
\textbf{Claude Opus 5.5.} nothing reported here\par
\textbf{GPT-6.1 Sol.} Place: Proof of Corollary 3.4, equality preceding the reference II 3.4; Reported print: \textbackslash{}check\{\textbackslash{}Hh\}\textasciicircum{}\{0\}\textbackslash{}check\{\textbackslash{}Hh\}\textasciicircum{}\{0\}\textbackslash{}Hh\textasciicircum{}\{q\}(M)=\allowbreak 0; Source TeX: \textbackslash{}check\{\textbackslash{}Hh\}\textasciicircum{}\{0\}\textbackslash{}check\{\textbackslash{}Hh\}\textasciicircum{}\{0\}\textbackslash{}check\{\textbackslash{}Hh\}\textasciicircum{}\{q\}(M)=\allowbreak 0; Note: The scan gives Čech checks to the two degree-zero functors, not to the final degree-q factor.; Proposed edit: \textbackslash{}check\{\textbackslash{}Hh\}\textasciicircum{}\{q\} (M) =\allowbreak  0\$ pour tout =\allowbreak > \textbackslash{}Hh\textasciicircum{}\{q\} (M) =\allowbreak  0\$ pour tout\par
a fix from one model only\par
\par\noindent\hypertarget{DP0213}{}\textbf{DP0213 -- V, p. 27 (one\_model\_fix)}\par
Corollary 3.5, notation for the underlying abelian Groupe; continuing on page 28\par
\textbf{Claude Opus 5.5.} nothing reported here\par
\textbf{GPT-6.1 Sol.} Place: Corollary 3.5, notation for the underlying abelian Groupe; continuing on page 28; Reported print: Notons \textbackslash{}underline\{M\}; Source TeX: Notons \$\ensuremath{\mathcal{M}}\$; Note: The underline is mathematical notation distinguishing the underlying Groupe from M, not emphasis. The substitution of script \ensuremath{\mathcal{M}} recurs throughout the corollary and its proof; the coordinator entry gives the scope and three examples.; Proposed edit: Notons \$\ensuremath{\mathcal{M}}\$ =\allowbreak > Notons \$\textbackslash{}underline\{M\}\$\par
a fix from one model only\par
\par\noindent\hypertarget{DP0214}{}\textbf{DP0214 -- V, p. 31 (one\_model\_fix)}\par
end of page, reference to the spectral sequence\par
\textbf{Claude Opus 5.5.} Place: end of page, reference to the spectral sequence; Reported print: du morphisme u (5.3.2). L'assertion; Source TeX: \textbackslash{}eqhyperref[V.5.3.2]\{5.3.2\}).; Note: \textbackslash{}eqhyperref already prints the parentheses, so the TeX prints '(5.3.2))' with an extra closing parenthesis.; Proposed edit: \textbackslash{}eqhyperref[V.5.3.2]\{5.3.2\}). =\allowbreak > \textbackslash{}eqhyperref[V.5.3.2]\{5.3.2\}.\par
\textbf{GPT-6.1 Sol.} nothing reported here\par
a fix from one model only\par
\par\noindent\hypertarget{DP0215}{}\textbf{DP0215 -- V, p. 31 (one\_model\_fix)}\par
last line\par
\textbf{Claude Opus 5.5.} Place: last line; Reported print: L'assertion 2) ne sera pas utili-; Source TeX: L'assertion 2 ne sera pas utilisée; Note: The ')' after 2 is missing in the TeX.; Proposed edit: L'assertion 2 ne =\allowbreak > L'assertion 2) ne\par
\textbf{GPT-6.1 Sol.} nothing reported here\par
a fix from one model only\par
\par\noindent\hypertarget{DP0216}{}\textbf{DP0216 -- V, p. 31 (one\_model\_fix)}\par
Preuve de 4.9, dernier terme de la chaîne d'égalités\par
\textbf{Claude Opus 5.5.} nothing reported here\par
\textbf{GPT-6.1 Sol.} Place: Preuve de 4.9, dernier terme de la chaîne d'égalités; Reported print: H\textasciicircum{}q(u\textasciicircum{}\textbackslash{}ast(\ensuremath{𝒳}),F) =\allowbreak  0; Source TeX: H\textasciicircum{}q (u\textasciicircum{}\textbackslash{}ast(\ensuremath{𝒳} ), F)) =\allowbreak  0; Note: L'imprimé n'a qu'une parenthèse fermante après F dans ce dernier terme.; Proposed edit: (u\textasciicircum{}\textbackslash{}ast(\ensuremath{𝒳} ), F)) =\allowbreak  0 =\allowbreak > (u\textasciicircum{}\textbackslash{}ast(\ensuremath{𝒳} ), F) =\allowbreak  0\par
a fix from one model only\par
\par\noindent\hypertarget{DP0217}{}\textbf{DP0217 -- V, p. 32 (one\_model\_fix)}\par
proof of 4.10 1)\par
\textbf{Claude Opus 5.5.} Place: proof of 4.10 1); Reported print: une suite exacte. Il nous faut; Source TeX: une suite exacte, Il nous faut; Note: The print has a full stop, not a comma.; Proposed edit: une suite exacte, Il =\allowbreak > une suite exacte. Il\par
\textbf{GPT-6.1 Sol.} nothing reported here\par
a fix from one model only\par
\par\noindent\hypertarget{DP0218}{}\textbf{DP0218 -- V, p. 32 (one\_model\_fix)}\par
Diagramme définissant C\_∙(\ensuremath{𝒳}), terme indexé par i,j\par
\textbf{Claude Opus 5.5.} nothing reported here\par
\textbf{GPT-6.1 Sol.} Place: Diagramme définissant C\_∙(\ensuremath{𝒳}), terme indexé par i,j; Reported print: \textbackslash{}coprod\_\{i,j\}\textbackslash{}ZZ\_\{X\_i\textbackslash{}times\_X X\_j\}; Source TeX: \textbackslash{}displaystyle\{\textbackslash{}mathop\{\textbackslash{}coprod\}\_\{i,j\}\}\textbackslash{}ZZ\_\{X\_\{i\}\textbackslash{}times X\_j\}; Note: La base X du produit fibré est imprimée en indice du signe de produit.; Proposed edit: \textbackslash{}ZZ\_\{X\_\{i\}\textbackslash{}times X\_j\} =\allowbreak > \textbackslash{}ZZ\_\{X\_\{i\}\textbackslash{}times\_X X\_j\}\par
a fix from one model only\par
\par\noindent\hypertarget{DP0219}{}\textbf{DP0219 -- V, p. 34 (agreed\_fix\_held)}\par
4.12, first line | 4.12, référence après « se localise »\par
\textbf{Claude Opus 5.5.} Place: 4.12, first line; Reported print: se localise (4.11 1)).; Source TeX: se localise \textbackslash{}eqhyperref[V.4.11.1]\{4.11.1\}.; Note: The print refers to assertion 1) of Proposition 4.11, not to formula (4.11.1).; Proposed edit: se localise \textbackslash{}eqhyperref[V.4.11.1]\{4.11.1\}. =\allowbreak > se localise (\textbackslash{}hyperref[V.prop4.11]\{4.11\} 1)).\par
\textbf{GPT-6.1 Sol.} Place: 4.12, référence après « se localise »; Reported print: (4.11 1)); Source TeX: \textbackslash{}eqhyperref[V.4.11.1]\{4.11.1\}; Note: L'imprimé renvoie à l'assertion 1) de la proposition 4.11, non à la formule (4.11.1).; Proposed edit: \textbackslash{}eqhyperref[V.4.11.1]\{4.11.1\} =\allowbreak > (\textbackslash{}hyperref[V.prop4.11]\{4.11\} 1))\par
both models give the same fix, held: it changes the editors' apparatus (\textbackslash{}eqhyperref, \textbackslash{}hyperref, the target of a cross-reference): a \textbackslash{}sisi correction, note, label, cross-reference or citation is the editors' record and is not removed or retargeted on the models' word\par
\par\noindent\hypertarget{DP0220}{}\textbf{DP0220 -- V, p. 38 (one\_model\_fix)}\par
title of section 6\par
\textbf{Claude Opus 5.5.} Place: title of section 6; Reported print: \ensuremath{\mathcal{E}}xt locaux et cohomologie à supports (script E); Source TeX: \textbackslash{}section\{Ext locaux; Note: The print's title uses the script \ensuremath{\mathcal{E}}xt of the sheaf Ext, not the group Ext.; Proposed edit: \textbackslash{}section\{Ext locaux =\allowbreak > \textbackslash{}section\{\$\textbackslash{}Eext\$ locaux\par
\textbf{GPT-6.1 Sol.} nothing reported here\par
a fix from one model only\par
\par\noindent\hypertarget{DP0221}{}\textbf{DP0221 -- V, p. 40 (one\_model\_fix)}\par
footnote\par
\textbf{Claude Opus 5.5.} Place: footnote; Reported print: cité p.80 plus bas; Source TeX: cité p-80 plus bas; Note: The print abbreviates 'p.' with a period, not a hyphen.; Proposed edit: cité p-80 =\allowbreak > cité p.80\par
\textbf{GPT-6.1 Sol.} nothing reported here\par
a fix from one model only\par
\par\noindent\hypertarget{DP0222}{}\textbf{DP0222 -- V, p. 42 (conflict)}\par
Proposition 6.7, first term of the display | Proposition 6.7, first term of the displayed formula\par
\textbf{Claude Opus 5.5.} Place: Proposition 6.7, first term of the display; Reported print: H\textasciicircum{}o\_Z(\ensuremath{\mathcal{H}}om\_A(M,N)), H not underlined; Source TeX: H\textasciicircum{}\{0\}\_Z(\textbackslash{}Hhom\_A(M, N)); Proposed correction: \ensuremath{\mathcal{H}}\textasciicircum{}\{0\}\_Z(\textbackslash{}Hhom\_A(M, N)); Note: The isomorphisms are between sheaves, so this should be the sheaf \ensuremath{\mathcal{H}}\textasciicircum{}0\_Z. The print lacks the underline and the TeX keeps the group H.\par
\textbf{GPT-6.1 Sol.} Place: Proposition 6.7, first term of the displayed formula; Reported print: \textbackslash{}underline\{H\}\textasciicircum{}\{0\}\_Z(\textbackslash{}Hhom\_A(M,N)); Source TeX: H\textasciicircum{}\{0\}\_Z(\textbackslash{}Hhom\_A(M, N)); Note: The printed H is underlined: this is the sheaf-valued functor, not the cohomology group.; Proposed edit: H\textasciicircum{}\{0\}\_Z(\textbackslash{}Hhom\_A(M, N)) =\allowbreak > \ensuremath{\mathcal{H}}\textasciicircum{}\{0\}\_Z(\textbackslash{}Hhom\_A(M, N))\par
the two models differ at this place: one model fixes the TeX, the other calls it a misprint of the print\par
\par\noindent\hypertarget{DP0223}{}\textbf{DP0223 -- V, p. 42 (one\_model\_fix)}\par
Proposition 6.7, the two isomorphism arrows\par
\textbf{Claude Opus 5.5.} nothing reported here\par
\textbf{GPT-6.1 Sol.} Place: Proposition 6.7, the two isomorphism arrows; Reported print: \textbackslash{}xrightarrow\{\textbackslash{}sim\}\textbackslash{}Hhom\_A(M,\ensuremath{\mathcal{H}}\textasciicircum{}0\_Z(N))\textbackslash{}xrightarrow\{\textbackslash{}sim\}; Source TeX: \textbackslash{}simeq \textbackslash{}Hhom\_A(M, \ensuremath{\mathcal{H}}\textasciicircum{}0\_Z(N)) \textbackslash{}simeq; Note: The print has two right-pointing arrows labelled with a tilde, rather than two isomorphism relation signs.; Proposed edit: \textbackslash{}simeq \textbackslash{}Hhom\_A(M, \ensuremath{\mathcal{H}}\textasciicircum{}0\_Z(N)) \textbackslash{}simeq =\allowbreak > \textbackslash{}xrightarrow\{\textbackslash{}sim\} \textbackslash{}Hhom\_A(M, \ensuremath{\mathcal{H}}\textasciicircum{}0\_Z(N)) \textbackslash{}xrightarrow\{\textbackslash{}sim\}\par
a fix from one model only\par
\par\noindent\hypertarget{DP0224}{}\textbf{DP0224 -- V, p. 44 (agreed\_fix\_held)}\par
(6.9.2), third line | (6.9.2), label of the third spectral sequence\par
\textbf{Claude Opus 5.5.} Place: (6.9.2), third line; Reported print: "E\textasciicircum{}\{p,q\}\_2 (double prime); Source TeX: ``E\textasciicircum{}\{p, q\}\_2; Note: The print has a double prime. In math mode `` gives left quotes.; Proposed edit: ``E\textasciicircum{}\{p, q\}\_2 =\allowbreak > ''E\textasciicircum{}\{p, q\}\_2\par
\textbf{GPT-6.1 Sol.} Place: (6.9.2), label of the third spectral sequence; Reported print: ''E\textasciicircum{}\{p,q\}\_2; Source TeX: ``E\textasciicircum{}\{p, q\}\_2; Note: The printed label has two primes before E; the TeX uses two grave accents.; Proposed edit: ``E\textasciicircum{}\{p, q\}\_2 =\allowbreak > ''E\textasciicircum{}\{p, q\}\_2\par
both models give the same fix, held: a model marks it low confidence\par
\par\noindent\hypertarget{DP0225}{}\textbf{DP0225 -- V, p. 47 (one\_model\_fix)}\par
Sentence preceding (6.13.8)\par
\textbf{Claude Opus 5.5.} nothing reported here\par
\textbf{GPT-6.1 Sol.} Place: Sentence preceding (6.13.8); Reported print: foncteurs dérivés droits; Source TeX: foncteurs dérivés droites; Note: The print reads « droits », not « droites ».; Proposed edit: foncteurs dérivés droites =\allowbreak > foncteurs dérivés droits\par
a fix from one model only\par
\par\noindent\hypertarget{DP0226}{}\textbf{DP0226 -- V, p. 55 (conflict)}\par
Théorème 7.3.2 3), «faisceau constant Z» | Théorème 7.3.2, assertion 3, dernière lettre\par
\textbf{Claude Opus 5.5.} Place: Théorème 7.3.2 3), «faisceau constant Z»; Reported print: Z underlined (the integers); Source TeX: \$\ensuremath{𝒵}\$; Note: The print has the underlined Z of the integers (bold in the edition); the TeX has a script Z.; Proposed edit: \$\ensuremath{𝒵}\$ =\allowbreak > \$\textbackslash{}mathbf\{Z\}\$\par
\textbf{GPT-6.1 Sol.} Place: Théorème 7.3.2, assertion 3, dernière lettre; Reported print: \$\textbackslash{}underline\{Z\}\$; Source TeX: \$\ensuremath{𝒵}\$; Note: Le scan porte un Z ordinaire souligné, non une capitale calligraphique.; Proposed edit: \$\ensuremath{𝒵}\$ =\allowbreak > \$\textbackslash{}underline\{Z\}\$\par
the two models differ at this place\par
\par\noindent\hypertarget{DP0227}{}\textbf{DP0227 -- V, p. 58 (conflict)}\par
Preuve of 7.3.6, 'le morphisme (x) est couvrant' | Preuve of 7.3.6, 'Or le morphisme (x) est isomorphe' | Preuve du lemme 7.3.6, les deux renvois suivant le carré\par
\textbf{Claude Opus 5.5.} Place: Preuve of 7.3.6, 'le morphisme (x) est couvrant'; Reported print: (x); Source TeX: \$(\textbackslash{}ast)\$; Note: The print writes the label as (x), as on the diagram arrow (which the TeX also gives as (x)). Apply together with the next finding.; Proposed edit: le morphisme \$(\textbackslash{}ast)\$ est couvrant =\allowbreak > le morphisme \$(\textbackslash{}mathrm\{x\})\$ est couvrant || Place: Preuve of 7.3.6, 'Or le morphisme (x) est isomorphe'; Reported print: (x); Source TeX: \$(\textbackslash{}ast)\$; Note: Same label (x) as in the diagram.; Proposed edit: morphisme \$(\textbackslash{}ast)\$ est isomorphe =\allowbreak > morphisme \$(\textbackslash{}mathrm\{x\})\$ est isomorphe\par
\textbf{GPT-6.1 Sol.} Place: Preuve du lemme 7.3.6, les deux renvois suivant le carré; Reported print: le morphisme (x) est couvrant. Or le morphisme (x) est isomorphe; Source TeX: le morphisme \$(\textbackslash{}ast)\$ est couvrant. Or le morphisme \$(\textbackslash{}ast)\$ est isomorphe; Note: Les deux renvois imprimés portent (x), comme l'étiquette de la flèche dans le carré, et non une astérisque.; Proposed edit: \$(\textbackslash{}ast)\$ est couvrant. Or le morphisme \$(\textbackslash{}ast)\$ =\allowbreak > (x) est couvrant. Or le morphisme (x)\par
the two models differ at this place\par
\par\noindent\hypertarget{DP0228}{}\textbf{DP0228 -- V, p. 62 (one\_model\_fix)}\par
'Supposons maintenant que le préfaisceau F soit un faisceau'\par
\textbf{Claude Opus 5.5.} Place: 'Supposons maintenant que le préfaisceau F soit un faisceau'; Reported print: soit un faisceau. Comme le faisceau; Source TeX: soit un faisceau, Comme; Note: The print ends the sentence with a full stop; the TeX has a comma before the capitalised 'Comme'.; Proposed edit: soit un faisceau, Comme =\allowbreak > soit un faisceau. Comme\par
\textbf{GPT-6.1 Sol.} nothing reported here\par
a fix from one model only\par
\par\noindent\hypertarget{DP0229}{}\textbf{DP0229 -- V, p. 63 (conflict)}\par
Remarque 7.4.2 1), 'les foncteurs ... ne sont autres que' | Remarque 7.4.2, 1), foncteurs associés aux recouvrements de type 0\par
\textbf{Claude Opus 5.5.} Place: Remarque 7.4.2 1), 'les foncteurs ... ne sont autres que'; Reported print: H with check and a small circle (0) above it; Source TeX: \textbackslash{}Hvi\{q\}; Note: The print puts a small 0 over the check, not \ensuremath{\infty} (the \ensuremath{\infty} elsewhere on the page is clearly printed). This fits the sense: hyper-coverings of type 0 give Čech cohomology.; Proposed edit: \$\textbackslash{}Hvi\{q\}\$ ne sont =\allowbreak > \$\textbackslash{}overset\{0\}\{\textbackslash{}check\{H\}\}\{\}\textasciicircum{}\{q\}\$ ne sont\par
\textbf{GPT-6.1 Sol.} Place: Remarque 7.4.2, 1), foncteurs associés aux recouvrements de type 0; Reported print: \textbackslash{}overset\{0\}\{\textbackslash{}check\{H\}\}\textasciicircum{}\{q\}; Source TeX: \textbackslash{}Hvi\{q\}; Note: Le scan porte 0 au-dessus du H accentué, non l'infini.; Proposed edit: \$\textbackslash{}Hvi\{q\}\$ =\allowbreak > \$\textbackslash{}overset\{0\}\{\textbackslash{}check\{H\}\}\textasciicircum{}\{q\}\$\par
the two models differ at this place\par
\par\noindent\hypertarget{DP0230}{}\textbf{DP0230 -- V, p. 64 (conflict)}\par
Définition 8.1.1 (L1) | Définition 8.1.1, axiome (L1), famille finie\par
\textbf{Claude Opus 5.5.} Place: Définition 8.1.1 (L1); Reported print: la famille finie (\ensuremath{𝒮}\_i) d'objets de L\_U; Source TeX: \$(\textbackslash{}Ss\_i)\$; Proposed correction: \$(\textbackslash{}lambda\_i)\$; Note: The family should be (\ensuremath{\lambda}\_i), as used later in Hom\_\{f\_j\}(\ensuremath{\mu}\_j,\ensuremath{\lambda}\_i). The TeX keeps the print's script S.\par
\textbf{GPT-6.1 Sol.} Place: Définition 8.1.1, axiome (L1), famille finie; Reported print: (\textbackslash{}lambda\_i); Source TeX: (\textbackslash{}Ss\_i); Note: Le caractère manuscrit est lambda, non le S cursif désignant le site.; Proposed edit: finie \$(\textbackslash{}Ss\_i)\$ =\allowbreak > finie \$(\textbackslash{}lambda\_i)\$\par
the two models differ at this place: one model fixes the TeX, the other calls it a misprint of the print\par
\par\noindent\hypertarget{DP0231}{}\textbf{DP0231 -- V, p. 65 (conflict)}\par
(L2), 'telles que p(\ensuremath{\psi}\_j) ... f\_j' | Définition 8.1.1, axiome (L2), projection des psi\_j\par
\textbf{Claude Opus 5.5.} Place: (L2), 'telles que p(\ensuremath{\psi}\_j) ... f\_j'; Reported print: the sign between p(\ensuremath{\psi}\_j) and f\_j is blurred and looks like =\allowbreak  (unlike the clear \ensuremath{\ne} of (L0) and (L1)); Source TeX: \textbackslash{}neq f\_j; Note: The stroke is unclear on the scan. The sense requires equality (the \ensuremath{\psi}\_j lie over f\_j).; Proposed edit: \textbackslash{}neq f\_j =\allowbreak > =\allowbreak  f\_j\par
\textbf{GPT-6.1 Sol.} Place: Définition 8.1.1, axiome (L2), projection des psi\_j; Reported print: p(\textbackslash{}psi\_j) \textbackslash{}neq f\_j; Source TeX: p(\textbackslash{}psi\_j) \textbackslash{}neq f\_j; Proposed correction: p(\textbackslash{}psi\_j) =\allowbreak  f\_j; Note: Le TeX conserve le signe d'inégalité imprimé. La condition requise est que psi\_j soit au-dessus de f\_j.\par
the two models differ at this place: one model fixes the TeX, the other calls it a misprint of the print\par
\par\noindent\hypertarget{DP0232}{}\textbf{DP0232 -- V, p. 65 (one\_model\_fix)}\par
8.1.1.1\par
\textbf{Claude Opus 5.5.} Place: 8.1.1.1; Reported print: et (L1), (L2) peuvent s'énoncer; Source TeX: \$(\textbackslash{}Ll\_1)\$,   \$(\textbackslash{}Ll\_2)\$; Note: These are the axiom names (L1), (L2) on the line; the TeX writes subscripts.; Proposed edit: \$(\textbackslash{}Ll\_1)\$,   \$(\textbackslash{}Ll\_2)\$ =\allowbreak > \$(\textbackslash{}Ll1)\$,   \$(\textbackslash{}Ll2)\$\par
\textbf{GPT-6.1 Sol.} nothing reported here\par
a fix from one model only\par
\par\noindent\hypertarget{DP0233}{}\textbf{DP0233 -- V, p. 65 (conflict)}\par
8.1.2 and following (the champ) | 8.1.2 et occurrences suivantes du nom du champ\par
\textbf{Claude Opus 5.5.} Place: 8.1.2 and following (the champ); Reported print: handwritten script \ensuremath{𝒞} (champ \ensuremath{𝒞}, \ensuremath{\Gamma}\ensuremath{𝒞}, \ensuremath{𝒞}\_S, Hom\_\ensuremath{𝒮}(L,\ensuremath{𝒞})); Source TeX: \textbackslash{}mathsf\{e\}; Note: The handwritten script capital C was read as the letter e. The edition presumably writes the script C as \textbackslash{}Cc (by analogy with \textbackslash{}Aa, \textbackslash{}Ff; not checked). Every \textbackslash{}mathsf\{e\} on the page should become \textbackslash{}Cc.; Proposed edit: \$\textbackslash{}mathsf\{e\}\$ un champ sur \$\textbackslash{}Ss\$ =\allowbreak > \$\textbackslash{}Cc\$ un champ sur \$\textbackslash{}Ss\$\par
\textbf{GPT-6.1 Sol.} Place: 8.1.2 et occurrences suivantes du nom du champ; Reported print: \$\textbackslash{}mathcal\{C\}\$ un champ; Source TeX: \$\textbackslash{}mathsf\{e\}\$ un champ; Note: Le caractère cursif du scan est un C, non un e. Cette confusion se répète dans les sept occurrences de \textbackslash{}mathsf\{e\} sur la page ; elle est signalée une seule fois.; Proposed edit: \$\textbackslash{}mathsf\{e\}\$ un champ =\allowbreak > \$\textbackslash{}mathcal\{C\}\$ un champ\par
the two models differ at this place\par
\par\noindent\hypertarget{DP0234}{}\textbf{DP0234 -- V, p. 65 (conflict)}\par
(L2), double arrow | Définition 8.1.1, axiome (L2), source de la double flèche\par
\textbf{Claude Opus 5.5.} Place: (L2), double arrow; Reported print: (\ensuremath{\phi}\_o,\ensuremath{\phi}\_1): \ensuremath{𝒮} \ensuremath{\rightrightarrows} \ensuremath{\mu}; Source TeX: (\textbackslash{}phi\_0, \textbackslash{}phi\_1): \textbackslash{}Ss \textbackslash{}rightrightarrows \textbackslash{}mu; Proposed correction: (\textbackslash{}phi\_0, \textbackslash{}phi\_1): \textbackslash{}lambda \textbackslash{}rightrightarrows \textbackslash{}mu; Note: The source of the double arrow should be \ensuremath{\lambda} (compare (L'2)). The TeX keeps the print's script S.\par
\textbf{GPT-6.1 Sol.} Place: Définition 8.1.1, axiome (L2), source de la double flèche; Reported print: (\textbackslash{}phi\_0, \textbackslash{}phi\_1): \textbackslash{}lambda \textbackslash{}rightrightarrows \textbackslash{}mu; Source TeX: (\textbackslash{}phi\_0, \textbackslash{}phi\_1): \textbackslash{}Ss \textbackslash{}rightrightarrows \textbackslash{}mu; Note: La source manuscrite est lambda, également nommé comme but des psi\_j dans la suite.; Proposed edit: (\textbackslash{}phi\_0, \textbackslash{}phi\_1): \textbackslash{}Ss \textbackslash{}rightrightarrows \textbackslash{}mu =\allowbreak > (\textbackslash{}phi\_0, \textbackslash{}phi\_1): \textbackslash{}lambda \textbackslash{}rightrightarrows \textbackslash{}mu\par
the two models differ at this place: one model fixes the TeX, the other calls it a misprint of the print\par
\par\noindent\hypertarget{DP0235}{}\textbf{DP0235 -- V, p. 66 (one\_model\_fix)}\par
diagram of Lemme 8.1.6, arrows from F(\ensuremath{\lambda}'\_j)\par
\textbf{Claude Opus 5.5.} Place: diagram of Lemme 8.1.6, arrows from F(\ensuremath{\lambda}'\_j); Reported print: two dashed arrows with arrowheads, from F(\ensuremath{\lambda}'\_j) to F(\ensuremath{\lambda}\_1) and to F(\ensuremath{\lambda}\_2); Source TeX: \textbackslash{}ar@\{--\}[ru] \textbackslash{}ar@\{--\}[rd]; Note: @\{--\} draws dashed lines without heads; the print's dashed arrows have arrowheads.; Proposed edit: \textbackslash{}ar@\{--\}[ru] \textbackslash{}ar@\{--\}[rd] =\allowbreak > \textbackslash{}ar@\{-->\}[ru] \textbackslash{}ar@\{-->\}[rd]\par
\textbf{GPT-6.1 Sol.} nothing reported here\par
a fix from one model only\par
\par\noindent\hypertarget{DP0236}{}\textbf{DP0236 -- V, p. 67 (one\_model\_fix)}\par
display defining \ensuremath{\phi}\par
\textbf{Claude Opus 5.5.} Place: display defining \ensuremath{\phi}; Reported print: Hom(c\_\ensuremath{\mathcal{M}} X,(X\_\ensuremath{\mu})\_\{\ensuremath{\mu}\ensuremath{\in}\ensuremath{\mathcal{L}}\})   \ensuremath{\phi}\_\ensuremath{\mu} : X | p(\ensuremath{\mu}); Source TeX: (X\_\{\textbackslash{}mu\})\_\{\textbackslash{}mu \textbackslash{}in \textbackslash{}Ll\}) : \textbackslash{}phi\_\{\textbackslash{}mu\}; Note: The print has no colon before \ensuremath{\phi}\_\ensuremath{\mu} in this display; the parallel \ensuremath{\psi} display does have one.; Proposed edit: (X\_\{\textbackslash{}mu\})\_\{\textbackslash{}mu \textbackslash{}in \textbackslash{}Ll\}) : \textbackslash{}phi\_\{\textbackslash{}mu\} =\allowbreak > (X\_\{\textbackslash{}mu\})\_\{\textbackslash{}mu \textbackslash{}in \textbackslash{}Ll\}) \textbackslash{}quad \textbackslash{}phi\_\{\textbackslash{}mu\}\par
\textbf{GPT-6.1 Sol.} nothing reported here\par
a fix from one model only\par
\par\noindent\hypertarget{DP0237}{}\textbf{DP0237 -- V, p. 72 (one\_model\_fix)}\par
proof of 8.1.12, first line\par
\textbf{Claude Opus 5.5.} Place: proof of 8.1.12, first line; Reported print: t\_x (typewriter subscript x =\allowbreak  \ensuremath{\ast}); Source TeX: t\_x; Note: The typescript's subscript x is the star, which the edition writes \textbackslash{}ast everywhere else (e.g. t\_\{\textbackslash{}ast\} on p. 75). Here the TeX has a literal x.; Proposed edit: t\_x : =\allowbreak > t\_\{\textbackslash{}ast\} :\par
\textbf{GPT-6.1 Sol.} nothing reported here\par
a fix from one model only\par
\par\noindent\hypertarget{DP0238}{}\textbf{DP0238 -- V, p. 72 (one\_model\_fix)}\par
proof of 8.1.12\par
\textbf{Claude Opus 5.5.} Place: proof of 8.1.12; Reported print: crible couvrant dans \ensuremath{𝒮}), donc par le; Source TeX: crible couvrant\} dans \$\textbackslash{}Ss\$). donc par le; Note: The print has a comma; the TeX has a period, which breaks the sentence.; Proposed edit: \$\textbackslash{}Ss\$). donc =\allowbreak > \$\textbackslash{}Ss\$), donc\par
\textbf{GPT-6.1 Sol.} nothing reported here\par
a fix from one model only\par
\par\noindent\hypertarget{DP0239}{}\textbf{DP0239 -- V, p. 72 (one\_model\_fix)}\par
Proposition 8.1.12, proof (i), object names; recurring on p. 73\par
\textbf{Claude Opus 5.5.} nothing reported here\par
\textbf{GPT-6.1 Sol.} Place: Proposition 8.1.12, proof (i), object names; recurring on p. 73; Reported print: \textbackslash{}lambda\textasciicircum{}1 et \textbackslash{}lambda\textasciicircum{}2; Source TeX: \$\textbackslash{}Ss\textasciicircum{}1\$ et   \$\textbackslash{}Ss\textasciicircum{}2\$; Note: The long handwritten object symbols are lambdas, not the script S naming the site. This recurring substitution is reported once; the coordinator gives the scope for the remaining occurrences.; Proposed edit: \$\textbackslash{}Ss\textasciicircum{}1\$ et   \$\textbackslash{}Ss\textasciicircum{}2\$ =\allowbreak > \$\textbackslash{}lambda\textasciicircum{}1\$ et   \$\textbackslash{}lambda\textasciicircum{}2\$\par
a fix from one model only\par
\par\noindent\hypertarget{DP0240}{}\textbf{DP0240 -- V, p. 74 (conflict)}\par
last diagram | Final diagram, vertical arrow below F(\ensuremath{\mu}⁰)\par
\textbf{Claude Opus 5.5.} Place: last diagram; Reported print: vertical arrow \ensuremath{\mathcal{F}}(\ensuremath{\mu}°) ↓ \ensuremath{\prod}\ensuremath{\mathcal{F}}(\ensuremath{\mu}\_j); Source TeX: \textbackslash{}Ff(\textbackslash{}mu\textasciicircum{}0) \textbackslash{}ar[rd]\textbackslash{}ar[rr]\textasciicircum{}\{\textbackslash{}bar\{h\}\} (no downward arrow); Note: The vertical arrow from \ensuremath{\mathcal{F}}(\ensuremath{\mu}°) down to \ensuremath{\prod}\ensuremath{\mathcal{F}}(\ensuremath{\mu}\_j) is missing from the TeX.; Proposed edit: \textbackslash{}Ff(\textbackslash{}mu\textasciicircum{}0) \textbackslash{}ar[rd] =\allowbreak > \textbackslash{}Ff(\textbackslash{}mu\textasciicircum{}0) \textbackslash{}ar[dd] \textbackslash{}ar[rd]\par
\textbf{GPT-6.1 Sol.} Place: Final diagram, vertical arrow below F(\ensuremath{\mu}⁰); Reported print: \textbackslash{}Ff(\textbackslash{}mu\textasciicircum{}0)\textbackslash{}longrightarrow\textbackslash{}prod\textbackslash{}Ff(\textbackslash{}mu\_j), directed downwards; Source TeX: \textbackslash{}Ff(\textbackslash{}mu\textasciicircum{}0) \textbackslash{}ar[rd]\textbackslash{}ar[rr]\textasciicircum{}\{\textbackslash{}bar\{h\}\}; Note: The printed downward arrow to the product of F(\ensuremath{\mu}\_j) is missing.; Proposed edit: \textbackslash{}Ff(\textbackslash{}mu\textasciicircum{}0) \textbackslash{}ar[rd]\textbackslash{}ar[rr]\textasciicircum{}\{\textbackslash{}bar\{h\}\} =\allowbreak > \textbackslash{}Ff(\textbackslash{}mu\textasciicircum{}0) \textbackslash{}ar[dd]\textbackslash{}ar[rd]\textbackslash{}ar[rr]\textasciicircum{}\{\textbackslash{}bar\{h\}\}\par
the two models differ at this place\par
\par\noindent\hypertarget{DP0241}{}\textbf{DP0241 -- V, p. 75 (one\_model\_fix)}\par
(iv)\par
\textbf{Claude Opus 5.5.} Place: (iv); Reported print: f : \ensuremath{\mathcal{F}} \ensuremath{\longrightarrow} \ensuremath{𝒢} ; f\_o : \ensuremath{\mathcal{F}}\_o \ensuremath{\longrightarrow} \ensuremath{𝒢}\_o; Source TeX: \textbackslash{}mathfrak\{g\} ; \textbackslash{}mathfrak\{g\}\_\{0\}; Note: The sheaf is a capital script G, in the same style as \ensuremath{\mathcal{F}}; the TeX has a small Fraktur g. The same applies to \textbackslash{}mathfrak\{g\}\_\{0\}. The \textbackslash{}Gg macro is presumed to exist.; Proposed edit: \textbackslash{}Ff \textbackslash{}to \textbackslash{}mathfrak\{g\}\$ =\allowbreak > \textbackslash{}Ff \textbackslash{}to \textbackslash{}Gg\$\par
\textbf{GPT-6.1 Sol.} nothing reported here\par
a fix from one model only\par
\par\noindent\hypertarget{DP0242}{}\textbf{DP0242 -- V, p. 75 (conflict)}\par
8.1.13.0 | 8.1.13.0, definition of the fibres of q\_*C\par
\textbf{Claude Opus 5.5.} Place: 8.1.13.0; Reported print: (q\_x\ensuremath{\mathfrak{C}})\_U =\allowbreak  \ensuremath{\mathfrak{C}}\_\{q\textasciicircum{}xU\}; Source TeX: \textbackslash{}mathsf\{e\}\_q x\_\{U\}; Note: The print's subscript q\textasciicircum{}xU is broken up in the TeX.; Proposed edit: \textbackslash{}mathsf\{e\}\_q x\_\{U\} =\allowbreak > \textbackslash{}mathsf\{e\}\_\{q\textasciicircum{}\{\textbackslash{}ast\}U\}\par
\textbf{GPT-6.1 Sol.} Place: 8.1.13.0, definition of the fibres of q\_*C; Reported print: (q\_\{\textbackslash{}ast\}\textbackslash{}mathcal\{C\})\_U=\allowbreak \textbackslash{}mathcal\{C\}\_\{q\textasciicircum{}\{\textbackslash{}ast\}U\}; Source TeX: \$(q\_\{\textbackslash{}ast\} \textbackslash{}mathsf\{e\})\_U =\allowbreak  \textbackslash{}mathsf\{e\}\_q x\_\{U\}\$; Note: The entire q*U is the subscript of C. The TeX instead has a subscript q followed by a separate x\_U.; Proposed edit: \$(q\_\{\textbackslash{}ast\} \textbackslash{}mathsf\{e\})\_U =\allowbreak  \textbackslash{}mathsf\{e\}\_q x\_\{U\}\$ =\allowbreak > \$(q\_\{\textbackslash{}ast\} \textbackslash{}mathcal\{C\})\_U =\allowbreak  \textbackslash{}mathcal\{C\}\_\{q\textasciicircum{}\{\textbackslash{}ast\} U\}\$\par
the two models differ at this place\par
\par\noindent\hypertarget{DP0243}{}\textbf{DP0243 -- V, p. 75 (conflict)}\par
8.1.13.0, second paragraph | 8.1.13.0, index of (X\_L) | 8.1.13.0, construction of the projective system indexed by L*\par
\textbf{Claude Opus 5.5.} Place: 8.1.13.0, second paragraph; Reported print: indexé par \ensuremath{\mathcal{L}} à valeur dans q\_x\ensuremath{\mathfrak{C}} , on définit un système projectif local indexé par \ensuremath{\mathcal{L}}\textasciicircum{}x et à valeur dans \ensuremath{\mathfrak{C}} : (X\_L)\_\{L\ensuremath{\in}\ensuremath{\mathcal{L}}\textasciicircum{}x\}; Source TeX: indexé par \$\textbackslash{}Ll\$ et à valeur dans \$\textbackslash{}mathsf\{e\} : (X\_L)\_\{L \textbackslash{}in \textbackslash{}mathsf\{L\textasciicircum{}\{x\}\}\}\$; Note: A whole line of the print is dropped: 'à valeur dans q\_*\ensuremath{\mathfrak{C}}, on définit un système projectif local indexé par \ensuremath{\mathcal{L}}\textasciicircum{}x'.; Proposed edit: indexé par \$\textbackslash{}Ll\$ et à valeur dans \$\textbackslash{}mathsf\{e\} : =\allowbreak > indexé par \$\textbackslash{}Ll\$ à valeur dans \$q\_\{\textbackslash{}ast\}\textbackslash{}mathsf\{e\}\$, on définit un système projectif local indexé par \$\textbackslash{}Ll\textasciicircum{}\{\textbackslash{}ast\}\$ et à valeur dans \$\textbackslash{}mathsf\{e\} : || Place: 8.1.13.0, index of (X\_L); Reported print: L \ensuremath{\in} \ensuremath{\mathcal{L}}\textasciicircum{}x; Source TeX: L \textbackslash{}in \textbackslash{}mathsf\{L\textasciicircum{}\{x\}\}; Note: The index set is script \ensuremath{\mathcal{L}} with the star; the TeX has a sans-serif L with a literal x.; Proposed edit: \textbackslash{}mathsf\{L\textasciicircum{}\{x\}\} =\allowbreak > \textbackslash{}Ll\textasciicircum{}\{\textbackslash{}ast\}\par
\textbf{GPT-6.1 Sol.} Place: 8.1.13.0, construction of the projective system indexed by L*; Reported print: à valeur dans q\_*C, on définit un système projectif local indexé par L* et à valeur dans C : (X\_L)\_\{L\ensuremath{\in}L*\}; Source TeX: \$\textbackslash{}mathsf\{e\} : (X\_L)\_\{L \textbackslash{}in \textbackslash{}mathsf\{L\textasciicircum{}\{x\}\}\}\$; Note: The initial system takes values in q\_*C. The clause defining the new local projective system is omitted. The printed indexing category is the same script L* defined earlier, not sans-serif L\textasciicircum{}x.; Proposed edit: \$\textbackslash{}mathsf\{e\} : (X\_L)\_\{L \textbackslash{}in \textbackslash{}mathsf\{L\textasciicircum{}\{x\}\}\}\$ =\allowbreak > \$q\_\{\textbackslash{}ast\}\textbackslash{}mathcal\{C\}\$, on définit un système projectif local indexé par \$\textbackslash{}Ll\textasciicircum{}\{\textbackslash{}ast\}\$ et à valeur dans \$\textbackslash{}mathcal\{C\} : (X\_L)\_\{L \textbackslash{}in \textbackslash{}Ll\textasciicircum{}\{\textbackslash{}ast\}\}\$\par
the two models differ at this place\par
\par\noindent\hypertarget{DP0244}{}\textbf{DP0244 -- V, p. 75 (conflict)}\par
8.1.13.0 and following paragraph, the champ | 8.1.13.0, letter naming the champ; recurring throughout the page\par
\textbf{Claude Opus 5.5.} Place: 8.1.13.0 and following paragraph, the champ; Reported print: \ensuremath{\mathfrak{C}} (capital gothic C) : \ensuremath{\mathfrak{C}} un champ, q\_x\ensuremath{\mathfrak{C}}, \ensuremath{\mathfrak{C}}\_\{q\textasciicircum{}xU\}, à valeur dans \ensuremath{\mathfrak{C}}, Si \ensuremath{\mathfrak{C}} est un champ, p\_x\ensuremath{\mathfrak{C}}; Source TeX: \textbackslash{}mathsf\{e\} throughout; Note: The champ is a capital C in the print; the TeX writes a small e throughout the page. Change every \textbackslash{}mathsf\{e\} to the C letter; this fix gives only the first occurrence.; Proposed edit: et \$\textbackslash{}mathsf\{e\}\$ un champ sur =\allowbreak > et \$\textbackslash{}mathsf\{C\}\$ un champ sur\par
\textbf{GPT-6.1 Sol.} Place: 8.1.13.0, letter naming the champ; recurring throughout the page; Reported print: \textbackslash{}mathcal\{C\}; Source TeX: \textbackslash{}mathsf\{e\}; Note: The printed letter is script capital C, not sans-serif e. This repeated substitution is reported once; apply the coordinator's convention to the remaining occurrences.; Proposed edit: et \$\textbackslash{}mathsf\{e\}\$ un champ =\allowbreak > et \$\textbackslash{}mathcal\{C\}\$ un champ\par
the two models differ at this place\par
\par\noindent\hypertarget{DP0245}{}\textbf{DP0245 -- V, p. 76 (one\_model\_fix)}\par
8.2.0, the champ\par
\textbf{Claude Opus 5.5.} Place: 8.2.0, the champ; Reported print: Soit \ensuremath{𝒞} le champ sur \ensuremath{𝒮} (script C, also in (i), (iii), \ensuremath{𝒞}\_U, Hom\_\ensuremath{𝒮}(\ensuremath{\mathcal{L}},\ensuremath{𝒞})°, (\ensuremath{\Gamma}\ensuremath{𝒞})°); Source TeX: Soit \$\textbackslash{}mathsf\{e\}\$; Note: The print's letter is a script C. The TeX writes \textbackslash{}mathsf\{e\} at every occurrence on this page; see coordinator.; Proposed edit: Soit \$\textbackslash{}mathsf\{e\}\$ =\allowbreak > Soit \$\textbackslash{}mathcal\{C\}\$\par
\textbf{GPT-6.1 Sol.} nothing reported here\par
a fix from one model only\par
\par\noindent\hypertarget{DP0246}{}\textbf{DP0246 -- V, p. 79 (one\_model\_fix)}\par
Lemme 8.2.8\par
\textbf{Claude Opus 5.5.} Place: Lemme 8.2.8; Reported print: \ensuremath{𝒪}\_S|U-modules (small m); Source TeX: \$\textbackslash{}Oo\_S | U\$-Modules; Note: The print has a small 'modules' here.; Proposed edit: \$\textbackslash{}Oo\_S | U\$-Modules =\allowbreak > \$\textbackslash{}Oo\_S | U\$-modules\par
\textbf{GPT-6.1 Sol.} nothing reported here\par
a fix from one model only\par
\par\noindent\hypertarget{DP0247}{}\textbf{DP0247 -- V, p. 79 (one\_model\_fix)}\par
Lemme 8.2.8\par
\textbf{Claude Opus 5.5.} Place: Lemme 8.2.8; Reported print: \ensuremath{𝒪}\_T|V-modules (small m); Source TeX: \$\textbackslash{}Oo\_T | V\$-Modules; Note: The print has a small 'modules' here.; Proposed edit: \$\textbackslash{}Oo\_T | V\$-Modules =\allowbreak > \$\textbackslash{}Oo\_T | V\$-modules\par
\textbf{GPT-6.1 Sol.} nothing reported here\par
a fix from one model only\par
\par\noindent\hypertarget{DP0248}{}\textbf{DP0248 -- V, p. 82 (fix\_unplaced)}\par
Bibliography [14], author's initial\par
\textbf{Claude Opus 5.5.} nothing reported here\par
\textbf{GPT-6.1 Sol.} Place: Bibliography [14], author's initial; Reported print: S. Lubkin; Source TeX: G. Lubkin; Note: The print has S., not G.; Proposed edit: G. Lubkin =\allowbreak > S. Lubkin\par
a fix whose old text is not found once on the page\par
\par\noindent\hypertarget{DP0249}{}\textbf{DP0249 -- Vbis, p. 84 (one\_model\_fix)}\par
'(ibid (5.2.1))'\par
\textbf{Claude Opus 5.5.} Place: '(ibid (5.2.1))'; Reported print: ibid (5.2.1), meaning [5] (5.2.1); Source TeX: \textbackslash{}hyperref[Vbis.subsubsec5.2.1]\{5.2.1\}; Note: The reference is to the cited work [5] ('ibid'), but the link points to this exposé's own 5.2.1. The printed number is right; only the link target is wrong.; Proposed edit: \textbackslash{}hyperref[Vbis.subsubsec5.2.1]\{5.2.1\} =\allowbreak > 5.2.1\par
\textbf{GPT-6.1 Sol.} nothing reported here\par
a fix from one model only\par
\par\noindent\hypertarget{DP0250}{}\textbf{DP0250 -- Vbis, p. 84 (one\_model\_fix)}\par
'(ibid (4.3))'\par
\textbf{Claude Opus 5.5.} Place: '(ibid (4.3))'; Reported print: ibid (4.3), meaning [5] (4.3); Source TeX: \textbackslash{}hyperref[Vbis.subsec4.3]\{4.3\}; Note: Same problem: the reference is to [5], but the link points to Vbis 4.3.; Proposed edit: \textbackslash{}hyperref[Vbis.subsec4.3]\{4.3\} =\allowbreak > 4.3\par
\textbf{GPT-6.1 Sol.} nothing reported here\par
a fix from one model only\par
\par\noindent\hypertarget{DP0251}{}\textbf{DP0251 -- Vbis, p. 84 (cut)}\par
before 'Introduction' (chapter title and author)\par
\textbf{Claude Opus 5.5.} Place: before 'Introduction' (chapter title and author); Reported print: page 84 begins with the heading INTRODUCTION; no title or author is printed on it; Source TeX: \textbackslash{}chapter\{Techniques de descente cohomologique\} ... B. Saint-Donat; Note: The exposé title and the author's name are not printed on page 84; they presumably belong to the preceding title page (83). Page 84 really begins at 'INTRODUCTION'.\par
\textbf{GPT-6.1 Sol.} nothing reported here\par
a model reports that the automatic cut gives lines of this page to its neighbour, or the neighbour's lines to this page: check the lines it names against the scan of the page that prints them (that page's checker did not have their TeX)\par
\par\noindent\hypertarget{DP0252}{}\textbf{DP0252 -- Vbis, p. 84 (cut)}\par
TeX prefix before Introduction\par
\textbf{Claude Opus 5.5.} nothing reported here\par
\textbf{GPT-6.1 Sol.} Place: TeX prefix before Introduction; Reported print: The page begins with INTRODUCTION; the chapter title and author are not printed here.; Source TeX: \textbackslash{}chapter\{Techniques de descente cohomologique\} and the centered B. Saint-Donat precede \textbackslash{}section*\{Introduction\}.; Note: The cut includes material from the preceding title page without declaring an overlap. Page 84 begins at \textbackslash{}section*\{Introduction\}. No textual edit is needed.\par
a model reports that the automatic cut gives lines of this page to its neighbour, or the neighbour's lines to this page: check the lines it names against the scan of the page that prints them (that page's checker did not have their TeX)\par
\par\noindent\hypertarget{DP0253}{}\textbf{DP0253 -- Vbis, p. 85 (one\_model\_fix)}\par
Formula (3.1), superscript on Delta\par
\textbf{Claude Opus 5.5.} nothing reported here\par
\textbf{GPT-6.1 Sol.} Place: Formula (3.1), superscript on Delta; Reported print: \textbackslash{}Delta\textasciicircum{}\{o\}; Source TeX: \textbackslash{}Delta\textasciicircum{}\{0\}; Note: The print has the superscript o denoting the opposite category, not the numeral zero.; Proposed edit: \textbackslash{}Delta\textasciicircum{}\{0\} =\allowbreak > \textbackslash{}Delta\textasciicircum{}\{o\}\par
a fix from one model only\par
\par\noindent\hypertarget{DP0254}{}\textbf{DP0254 -- Vbis, p. 90 (conflict)}\par
1.2.8 | 1.2.8, définition de la catégorie des sections\par
\textbf{Claude Opus 5.5.} Place: 1.2.8; Reported print: la catégorie Hom\_D(D,E) with Hom underlined; Source TeX: \$\textbackslash{}Hom\_D (D, E)\$; Note: The print underlines this Hom, as it underlines \ensuremath{\Gamma} just before; the TeX renders that underlined \ensuremath{\Gamma} as \textbackslash{}Ysheaf\{\textbackslash{}Gamma\}.; Proposed edit: catégorie \$\textbackslash{}Hom\_D (D, E)\$ =\allowbreak > catégorie \$\textbackslash{}Hhom\_D (D, E)\$\par
\textbf{GPT-6.1 Sol.} Place: 1.2.8, définition de la catégorie des sections; Reported print: \textbackslash{}underline\{\textbackslash{}Hom\}\_D(D,E); Source TeX: \textbackslash{}Hom\_D (D, E); Note: Le soulignement de Hom désigne ici une catégorie de foncteurs, non un ensemble de morphismes ; ce n'est pas une marque d'emphase.; Proposed edit: \textbackslash{}Hom\_D (D, E) =\allowbreak > \textbackslash{}underline\{\textbackslash{}Hom\}\_D (D, E)\par
the two models differ at this place\par
\par\noindent\hypertarget{DP0255}{}\textbf{DP0255 -- Vbis, p. 91 (agreed\_fix\_held)}\par
Lemme 1.2.10, display | Lemme 1.2.10, display defining f\textasciicircum{}\{\textbackslash{}ast\}\par
\textbf{Claude Opus 5.5.} Place: Lemme 1.2.10, display; Reported print: \textbackslash{}underline\{Hom\}\_B(J,A) \ensuremath{\longrightarrow} \textbackslash{}underline\{Hom\}\_B(I,A) (underlined Hom: category of B-functors); Source TeX: \textbackslash{}Hom\_B (J, A) \textbackslash{}longrightarrow \textbackslash{}Hom\_B (I, A); Note: The print underlines Hom for the category of functors; the TeX uses the plain \textbackslash{}Hom. The same convention recurs on pp. 92 and 95. The print is not consistent: on p. 92, F \ensuremath{\in} Hom\_B(I,A) is not underlined.; Proposed edit: \textbackslash{}Hom\_B (J, A) \textbackslash{}longrightarrow \textbackslash{}Hom\_B (I, A) =\allowbreak > \textbackslash{}underline\{\textbackslash{}Hom\}\_B (J, A) \textbackslash{}longrightarrow \textbackslash{}underline\{\textbackslash{}Hom\}\_B (I, A)\par
\textbf{GPT-6.1 Sol.} Place: Lemme 1.2.10, display defining f\textasciicircum{}\{\textbackslash{}ast\}; Reported print: \textbackslash{}underline\{\textbackslash{}Hom\}\_B(J,A) \textbackslash{}longrightarrow \textbackslash{}underline\{\textbackslash{}Hom\}\_B(I,A); Source TeX: \textbackslash{}Hom\_B (J, A) \textbackslash{}longrightarrow \textbackslash{}Hom\_B (I, A); Note: Soulignement de notation : il désigne les catégories de foncteurs, contrairement aux ensembles Hom\_u du lemme suivant. Cette convention récurrente est signalée une seule fois.; Proposed edit: \textbackslash{}Hom\_B (J, A) \textbackslash{}longrightarrow \textbackslash{}Hom\_B (I, A) =\allowbreak > \textbackslash{}underline\{\textbackslash{}Hom\}\_B (J, A) \textbackslash{}longrightarrow \textbackslash{}underline\{\textbackslash{}Hom\}\_B (I, A)\par
both models give the same fix, held: a model marks it low confidence\par
\par\noindent\hypertarget{DP0256}{}\textbf{DP0256 -- Vbis, p. 92 (one\_model\_fix)}\par
Hom(x,y) cases, first line\par
\textbf{Claude Opus 5.5.} Place: Hom(x,y) cases, first line; Reported print: si x , y \ensuremath{\in} ob(I); Source TeX: x, y, \textbackslash{}in \textbackslash{}ob (I); Note: The TeX has an extra comma after y.; Proposed edit: x, y, \textbackslash{}in \textbackslash{}ob (I) =\allowbreak > x, y \textbackslash{}in \textbackslash{}ob (I)\par
\textbf{GPT-6.1 Sol.} nothing reported here\par
a fix from one model only\par
\par\noindent\hypertarget{DP0257}{}\textbf{DP0257 -- Vbis, p. 92 (one\_model\_fix)}\par
Hom(x,y) cases, second line\par
\textbf{Claude Opus 5.5.} Place: Hom(x,y) cases, second line; Reported print: si x , y \ensuremath{\in} ob(J); Source TeX: x, y, \textbackslash{}in \textbackslash{}ob (J); Note: The TeX has an extra comma after y.; Proposed edit: x, y, \textbackslash{}in \textbackslash{}ob (J) =\allowbreak > x, y \textbackslash{}in \textbackslash{}ob (J)\par
\textbf{GPT-6.1 Sol.} nothing reported here\par
a fix from one model only\par
\par\noindent\hypertarget{DP0258}{}\textbf{DP0258 -- Vbis, p. 93 (one\_model\_fix)}\par
1.2.13\par
\textbf{Claude Opus 5.5.} Place: 1.2.13; Reported print: 1 \ensuremath{\longrightarrow} T\ensuremath{\circ}S; Source TeX: 1 \textbackslash{}to T\textasciicircum{}\{\textbackslash{}circ\}S; Note: This is composition, as in S\ensuremath{\circ}T just after. The TeX sets the circle as an exponent.; Proposed edit: T\textasciicircum{}\{\textbackslash{}circ\}S =\allowbreak > T \textbackslash{}circ S\par
\textbf{GPT-6.1 Sol.} nothing reported here\par
a fix from one model only\par
\par\noindent\hypertarget{DP0259}{}\textbf{DP0259 -- Vbis, p. 94 (one\_model\_fix)}\par
Scholie 1.2.17\par
\textbf{Claude Opus 5.5.} Place: Scholie 1.2.17; Reported print: le couple (Φ,\ensuremath{\Psi}) : E \ensuremath{\longrightarrow} E'; Source TeX: (\textbackslash{}Phi, \textbackslash{}psi); Note: The print has a capital \ensuremath{\Psi} in the couple, although it writes \ensuremath{\psi} just before. The print is inconsistent; the TeX does not follow it.; Proposed edit: (\textbackslash{}Phi, \textbackslash{}psi) =\allowbreak > (\textbackslash{}Phi, \textbackslash{}Psi)\par
\textbf{GPT-6.1 Sol.} nothing reported here\par
a fix from one model only\par
\par\noindent\hypertarget{DP0260}{}\textbf{DP0260 -- Vbis, p. 96 (conflict)}\par
1.3.5, end: s'identifie à la catégorie Hom\_D(D,Mod(E,A)) | Opening paragraph, identification with the category of sections\par
\textbf{Claude Opus 5.5.} Place: 1.3.5, end: s'identifie à la catégorie Hom\_D(D,Mod(E,A)); Reported print: Hom\_D underlined (the category of functors), like Mod; Source TeX: \$\textbackslash{}Hom\_D; Note: The print underlines Hom\_D as it does Mod; the TeX uses the plain \textbackslash{}Hom.; Proposed edit: \$\textbackslash{}Hom\_D (D, =\allowbreak > \$\textbackslash{}Ysheaf\{\textbackslash{}Hom\}\_D (D,\par
\textbf{GPT-6.1 Sol.} Place: Opening paragraph, identification with the category of sections; Reported print: \textbackslash{}underline\{\textbackslash{}Hom\}\_D(D,\textbackslash{}underline\{\textbackslash{}Mod\}(E,A)); Source TeX: \textbackslash{}Hom\_D (D,\textbackslash{}Ysheaf\{\textbackslash{}Mod\}(E,A)); Note: The printed underline marks the category-valued Hom, rather than emphasis.; Proposed edit: \textbackslash{}Hom\_D (D, =\allowbreak > \textbackslash{}underline\{\textbackslash{}Hom\}\_D (D,\par
the two models differ at this place\par
\par\noindent\hypertarget{DP0261}{}\textbf{DP0261 -- Vbis, p. 96 (conflict)}\par
proof of 1.3.7: identification Hom\_D(D,Mod(E,A)) \ensuremath{\simeq} Mod(\ensuremath{\Gamma}(E),A) | Final identification following Proposition 1.3.7\par
\textbf{Claude Opus 5.5.} Place: proof of 1.3.7: identification Hom\_D(D,Mod(E,A)) \ensuremath{\simeq} Mod(\ensuremath{\Gamma}(E),A); Reported print: Hom\_D underlined; Source TeX: \$\textbackslash{}Hom\_D(D, \textbackslash{}Ysheaf\{\textbackslash{}Mod\}(E, A)); Note: Same as above.; Proposed edit: \$\textbackslash{}Hom\_D(D, \textbackslash{}Ysheaf =\allowbreak > \$\textbackslash{}Ysheaf\{\textbackslash{}Hom\}\_D(D, \textbackslash{}Ysheaf\par
\textbf{GPT-6.1 Sol.} Place: Final identification following Proposition 1.3.7; Reported print: \textbackslash{}underline\{\textbackslash{}Hom\}\_D(D,\textbackslash{}underline\{\textbackslash{}Mod\}(E,A)); Source TeX: \textbackslash{}Hom\_D(D, \textbackslash{}Ysheaf\{\textbackslash{}Mod\}(E, A)); Note: The print again uses the underlined, category-valued Hom.; Proposed edit: \textbackslash{}Hom\_D(D, =\allowbreak > \textbackslash{}underline\{\textbackslash{}Hom\}\_D(D,\par
the two models differ at this place\par
\par\noindent\hypertarget{DP0262}{}\textbf{DP0262 -- Vbis, p. 97 (agreed\_fix\_held)}\par
proof of 1.3.10 (i): expression explicite de e\_\{i*\} | Proof of Proposition 1.3.10(i), functor whose explicit expression is cited\par
\textbf{Claude Opus 5.5.} Place: proof of 1.3.10 (i): expression explicite de e\_\{i*\}; Reported print: e\_\{i*\}; Source TeX: \$e\_\{i\textasciicircum{}\{\textbackslash{}ast\}\}\$; Note: The star is on the subscript line in the print, not an exponent of i.; Proposed edit: \$e\_\{i\textasciicircum{}\{\textbackslash{}ast\}\}\$ =\allowbreak > \$e\_\{i\textbackslash{}ast\}\$\par
\textbf{GPT-6.1 Sol.} Place: Proof of Proposition 1.3.10(i), functor whose explicit expression is cited; Reported print: e\_\{i\textbackslash{}ast\}; Source TeX: \$e\_\{i\textasciicircum{}\{\textbackslash{}ast\}\}\$; Note: In the print the star is part of the subscript i*, not a superscript on i.; Proposed edit: \$e\_\{i\textasciicircum{}\{\textbackslash{}ast\}\}\$ =\allowbreak > \$e\_\{i\textbackslash{}ast\}\$\par
both models give the same fix, held: a model marks it low confidence\par
\par\noindent\hypertarget{DP0263}{}\textbf{DP0263 -- Vbis, p. 98 (one\_model\_fix)}\par
(1.3.10.2)\par
\textbf{Claude Opus 5.5.} Place: (1.3.10.2); Reported print: heavy arrow (abutment of the spectral sequence) \ensuremath{\Longrightarrow}; Source TeX: \textbackslash{}longrightarrow; Note: The print uses a thick arrow for the spectral sequence, not an ordinary arrow.; Proposed edit: F\_t) \textbackslash{}longrightarrow =\allowbreak > F\_t) \textbackslash{}Longrightarrow\par
\textbf{GPT-6.1 Sol.} nothing reported here\par
a fix from one model only\par
\par\noindent\hypertarget{DP0264}{}\textbf{DP0264 -- Vbis, p. 99 (conflict)}\par
Lemme 1.3.10.5, the sequence, second term | Lemme 1.3.10.5, F and G terms of the displayed exact sequence\par
\textbf{Claude Opus 5.5.} Place: Lemme 1.3.10.5, the sequence, second term; Reported print: \ensuremath{\Gamma}(\ensuremath{\phi}\_*)(G)); Source TeX: \textbackslash{}Gamma((\textbackslash{}varphi\_\{\textbackslash{}ast\})(G)); Note: The TeX moves the parentheses. The print's own parentheses are unbalanced (misprint for \ensuremath{\Gamma}(\ensuremath{\phi}\_*)(G)).; Proposed edit: \textbackslash{}Gamma((\textbackslash{}varphi\_\{\textbackslash{}ast\})(G)) =\allowbreak > \textbackslash{}Gamma(\textbackslash{}varphi\_\{\textbackslash{}ast\})(G))\par
\textbf{GPT-6.1 Sol.} Place: Lemme 1.3.10.5, F and G terms of the displayed exact sequence; Reported print: \textbackslash{}Gamma(\textbackslash{}varphi\_\{\textbackslash{}ast\})(F)) \textbackslash{}to \textbackslash{}Gamma(\textbackslash{}varphi\_\{\textbackslash{}ast\})(G)); Source TeX: \textbackslash{}Gamma(\textbackslash{}varphi\_\{\textbackslash{}ast\})(F)) \textbackslash{}to \textbackslash{}Gamma((\textbackslash{}varphi\_\{\textbackslash{}ast\})(G)); Proposed correction: \textbackslash{}Gamma(\textbackslash{}varphi\_\{\textbackslash{}ast\})(F) \textbackslash{}to \textbackslash{}Gamma(\textbackslash{}varphi\_\{\textbackslash{}ast\})(G); Note: Both printed terms have an unmatched final parenthesis. The TeX keeps it for F; for G it restores balance by adding an opening parenthesis instead.\par
the two models differ at this place: one model fixes the TeX, the other calls it a misprint of the print\par
\par\noindent\hypertarget{DP0265}{}\textbf{DP0265 -- Vbis, p. 99 (one\_model\_fix)}\par
proof of 1.3.11: théorème (5.1) de [[7] - Chap.1]\par
\textbf{Claude Opus 5.5.} Place: proof of 1.3.11: théorème (5.1) de [[7] - Chap.1]; Reported print: théorème (5.1) de [[7] - Chap.1]; Source TeX: \textbackslash{}hyperref[Vbis.subsec5.1]\{5.1\}; Note: It prints 5.1 correctly, but the link points to this exposé's 5.1, whereas the theorem is in reference [7].; Proposed edit: \textbackslash{}hyperref[Vbis.subsec5.1]\{5.1\} =\allowbreak > 5.1\par
\textbf{GPT-6.1 Sol.} nothing reported here\par
a fix from one model only\par
\par\noindent\hypertarget{DP0266}{}\textbf{DP0266 -- Vbis, p. 101 (one\_model\_fix)}\par
2.1.3, category of sections\par
\textbf{Claude Opus 5.5.} nothing reported here\par
\textbf{GPT-6.1 Sol.} Place: 2.1.3, category of sections; Reported print: \textbackslash{}underline\{\textbackslash{}Hom\}\_\{B\textasciicircum{}\{\textbackslash{}circ\}\}(B\textasciicircum{}\{\textbackslash{}circ\},\textbackslash{}mathcal\{E\}\textasciicircum{}\{\textbackslash{}circ\}); Source TeX: \textbackslash{}Hom\_\{B\textasciicircum{}\{\textbackslash{}circ\}\}(B\textasciicircum{}\{\textbackslash{}circ\}, \textbackslash{}mathcal\{E\}\textasciicircum{}\{\textbackslash{}circ\}); Note: Systematic notation difference, reported once: the print underlines Hom denoting functor categories, distinguishing it from Hom denoting morphism groups. This is mathematical notation, not emphatic underlining. The edit is representative; three examples are identified in coordinator.; Proposed edit: \textbackslash{}Hom\_\{B\textasciicircum{}\{\textbackslash{}circ\}\} =\allowbreak > \textbackslash{}underline\{\textbackslash{}Hom\}\_\{B\textasciicircum{}\{\textbackslash{}circ\}\}\par
a fix from one model only\par
\par\noindent\hypertarget{DP0267}{}\textbf{DP0267 -- Vbis, p. 104 (conflict)}\par
2.3.1.1, target of \ensuremath{\bar{F}} | 2.3.1.1, target of the functor \textbackslash{}bar\{F\}\par
\textbf{Claude Opus 5.5.} Place: 2.3.1.1, target of \ensuremath{\bar{F}}; Reported print: (Hom(D,Ab))° \ensuremath{\to} A°Ens; Source TeX: \textbackslash{}to \textbackslash{}Hom(A\textasciicircum{}\{\textbackslash{}circ\}, \textbackslash{}Ens); Proposed correction: \textbackslash{}Hom(A\textasciicircum{}\{\textbackslash{}circ\},\textbackslash{}Ens); Note: The print omits Hom( , ); the TeX silently corrects to Hom(A°,Ens), as on p. 105.\par
\textbf{GPT-6.1 Sol.} Place: 2.3.1.1, target of the functor \textbackslash{}bar\{F\}; Reported print: \{\}\textasciicircum{}\{A\textasciicircum{}\{\textbackslash{}circ\}\}\textbackslash{}Ens; Source TeX: \textbackslash{}Hom(A\textasciicircum{}\{\textbackslash{}circ\}, \textbackslash{}Ens); Note: The print uses the left-exponent notation for the functor category. The TeX replaces it with an equivalent Hom expression, a rewriting of the printed formula.; Proposed edit: \textbackslash{}Hom(A\textasciicircum{}\{\textbackslash{}circ\}, \textbackslash{}Ens) =\allowbreak > \{\}\textasciicircum{}\{A\textasciicircum{}\{\textbackslash{}circ\}\}\textbackslash{}Ens\par
the two models differ at this place: one model fixes the TeX, the other calls it a misprint of the print\par
\par\noindent\hypertarget{DP0268}{}\textbf{DP0268 -- Vbis, p. 107 (one\_model\_fix)}\par
2.3.4, 'complexe de chaînes dans Hom(D,Ab)'\par
\textbf{Claude Opus 5.5.} Place: 2.3.4, 'complexe de chaînes dans Hom(D,Ab)'; Reported print: Ho\ensuremath{\underline{m}}(D,Ab) (Hom underlined); Source TeX: \$\textbackslash{}Hom(D, \textbackslash{}Ab)\$; Note: The print underlines this Hom (the functor category), which may be meant to distinguish it from the plain Hom groups used below.; Proposed edit: complexe de chaînes dans \$\textbackslash{}Hom(D, =\allowbreak > complexe de chaînes dans \$\textbackslash{}Hhom(D,\par
\textbf{GPT-6.1 Sol.} nothing reported here\par
a fix from one model only\par
\par\noindent\hypertarget{DP0269}{}\textbf{DP0269 -- Vbis, p. 107 (agreed\_fix\_held)}\par
Proposition 2.3.5, augmented complex\par
\textbf{Claude Opus 5.5.} Place: Proposition 2.3.5, augmented complex; Reported print: Z\textasciicircum{}n \ensuremath{\to} Z\textasciicircum{}\{n-1\} ..... Z\textasciicircum{}1 \ensuremath{\to} Z\textasciicircum{}o \ensuremath{\to}t Z; Source TeX: \textbackslash{}ldots Z\textasciicircum{}n \textbackslash{}to Z\textasciicircum{}\{n - 1\}\textbackslash{}to \textbackslash{}ldots Z\textasciicircum{}\{1\}; Note: The print has no leading dots and no arrow after Z\textasciicircum{}\{n-1\}.; Proposed edit: \textbackslash{}ldots Z\textasciicircum{}n \textbackslash{}to Z\textasciicircum{}\{n - 1\}\textbackslash{}to \textbackslash{}ldots Z\textasciicircum{}\{1\} =\allowbreak > Z\textasciicircum{}n \textbackslash{}to Z\textasciicircum{}\{n - 1\} \textbackslash{}ldots Z\textasciicircum{}\{1\}\par
\textbf{GPT-6.1 Sol.} Place: Proposition 2.3.5, augmented complex; Reported print: Z\textasciicircum{}n \textbackslash{}to Z\textasciicircum{}\{n-1\} \textbackslash{}ldots Z\textasciicircum{}1 \textbackslash{}to Z\textasciicircum{}0 \textbackslash{}xrightarrow\{t\} \textbackslash{}ZZ; Source TeX: \textbackslash{}ldots Z\textasciicircum{}n \textbackslash{}to Z\textasciicircum{}\{n - 1\}\textbackslash{}to \textbackslash{}ldots Z\textasciicircum{}\{1\} \textbackslash{}to Z\textasciicircum{}0 \textbackslash{}xrightarrow\{t\} \textbackslash{}ZZ; Note: The print has neither an initial ellipsis nor an arrow between Z\textasciicircum{}\{n-1\} and the following ellipsis.; Proposed edit: \textbackslash{}ldots Z\textasciicircum{}n \textbackslash{}to Z\textasciicircum{}\{n - 1\}\textbackslash{}to \textbackslash{}ldots Z\textasciicircum{}\{1\} =\allowbreak > Z\textasciicircum{}n \textbackslash{}to Z\textasciicircum{}\{n - 1\} \textbackslash{}ldots Z\textasciicircum{}\{1\}\par
both models give the same fix, held: a model marks it low confidence\par
\par\noindent\hypertarget{DP0270}{}\textbf{DP0270 -- Vbis, p. 110 (one\_model\_fix)}\par
(2.3.10.2), first defining equation\par
\textbf{Claude Opus 5.5.} nothing reported here\par
\textbf{GPT-6.1 Sol.} Place: (2.3.10.2), first defining equation; Reported print: s\textasciicircum{}\{n,m\}\_1; Source TeX: S\textasciicircum{}\{n, m\}\_1; Note: The print uses lowercase s, as in the diagram.; Proposed edit: S\textasciicircum{}\{n, m\}\_1 =\allowbreak > s\textasciicircum{}\{n, m\}\_1\par
a fix from one model only\par
\par\noindent\hypertarget{DP0271}{}\textbf{DP0271 -- Vbis, p. 110 (one\_model\_fix)}\par
(2.3.10.2), second defining equation\par
\textbf{Claude Opus 5.5.} nothing reported here\par
\textbf{GPT-6.1 Sol.} Place: (2.3.10.2), second defining equation; Reported print: s\textasciicircum{}\{n,m\}\_2; Source TeX: S\textasciicircum{}\{n, m\}\_2; Note: The print again uses lowercase s.; Proposed edit: S\textasciicircum{}\{n, m\}\_2 =\allowbreak > s\textasciicircum{}\{n, m\}\_2\par
a fix from one model only\par
\par\noindent\hypertarget{DP0272}{}\textbf{DP0272 -- Vbis, p. 113 (one\_model\_fix)}\par
Lemme 2.4.4.2, De plus ...\par
\textbf{Claude Opus 5.5.} Place: Lemme 2.4.4.2, De plus ...; Reported print: \ensuremath{\subset} D\textasciicircum{}+\_G(\ensuremath{\Gamma}(\ensuremath{\bar{X}}'),O).; Source TeX: \textbackslash{}subset (D\textasciicircum{}+\_G(\textbackslash{}Ysheaf\{\textbackslash{}Gamma\}(\textbackslash{}bar\{X\}'), \textbackslash{}Oo)); Note: The TeX puts parentheses around the right-hand side; the print has none.; Proposed edit: \textbackslash{}subset (D\textasciicircum{}+\_G(\textbackslash{}Ysheaf\{\textbackslash{}Gamma\}(\textbackslash{}bar\{X\}'), \textbackslash{}Oo)) =\allowbreak > \textbackslash{}subset D\textasciicircum{}+\_G(\textbackslash{}Ysheaf\{\textbackslash{}Gamma\}(\textbackslash{}bar\{X\}'), \textbackslash{}Oo)\par
\textbf{GPT-6.1 Sol.} nothing reported here\par
a fix from one model only\par
\par\noindent\hypertarget{DP0273}{}\textbf{DP0273 -- Vbis, p. 114 (conflict)}\par
2.4.7, du foncteur L\textasciicircum{}+\ensuremath{\bar{\theta}}* | 2.4.7, nous noterons L\textasciicircum{}+\_G \ensuremath{\bar{\theta}}* | 2.4.7 and the repeated left-derived-functor notation on pp. 114–115\par
\textbf{Claude Opus 5.5.} Place: 2.4.7, du foncteur L\textasciicircum{}+\ensuremath{\bar{\theta}}*; Reported print: du foncteur L\textasciicircum{}+ \ensuremath{\bar{\theta}}*; Source TeX: du foncteur \$\textbackslash{}bar\{\textbackslash{}theta\}\textasciicircum{}\{\textbackslash{}ast\}\$; Note: The print has the derived functor L\textasciicircum{}+ (double underlined L); the TeX drops it.; Proposed edit: du foncteur \$\textbackslash{}bar\{\textbackslash{}theta\}\textasciicircum{}\{\textbackslash{}ast\}\$ =\allowbreak > du foncteur \$L\textasciicircum{}\{+\}\textbackslash{}bar\{\textbackslash{}theta\}\textasciicircum{}\{\textbackslash{}ast\}\$ || Place: 2.4.7, nous noterons L\textasciicircum{}+\_G \ensuremath{\bar{\theta}}*; Reported print: nous noterons L\textasciicircum{}+\_G \ensuremath{\bar{\theta}}*; Source TeX: nous noterons \$\textbackslash{}bar\textbackslash{}theta\_\{G\}\textasciicircum{}\{\textbackslash{}ast\}\$; Note: The print reads L\textasciicircum{}+\_G \ensuremath{\bar{\theta}}*: the TeX drops L\textasciicircum{}+ and moves G onto \ensuremath{\theta}.; Proposed edit: nous noterons \$\textbackslash{}bar\textbackslash{}theta\_\{G\}\textasciicircum{}\{\textbackslash{}ast\}\$ =\allowbreak > nous noterons \$L\textasciicircum{}\{+\}\_\{G\}\textbackslash{}bar\textbackslash{}theta\textasciicircum{}\{\textbackslash{}ast\}\$\par
\textbf{GPT-6.1 Sol.} Place: 2.4.7 and the repeated left-derived-functor notation on pp. 114–115; Reported print: \textbackslash{}mathbf\{L\}\textasciicircum{}\{+\}\textbackslash{}bar\{\textbackslash{}theta\}\textasciicircum{}\{\textbackslash{}ast\}; \textbackslash{}mathbf\{L\}\textasciicircum{}\{+\}\_\{G\}\textbackslash{}bar\textbackslash{}theta\textasciicircum{}\{\textbackslash{}ast\}; Source TeX: \textbackslash{}bar\{\textbackslash{}theta\}\textasciicircum{}\{\textbackslash{}ast\}; \textbackslash{}bar\textbackslash{}theta\_\{G\}\textasciicircum{}\{\textbackslash{}ast\}; Note: Systematic notation change, reported once: the TeX suppresses L\textasciicircum{}+ and transfers G from L to \ensuremath{\theta}. This representative edit restores both introductory forms. The same restoration is required in both adjunction maps of 2.4.7.1, condition 2° of 2.4.8, and Theorem 2.4.12 on p. 115.; Proposed edit: \$\textbackslash{}bar\{\textbackslash{}theta\}\textasciicircum{}\{\textbackslash{}ast\}\$ prend ses valeurs dans \$D\textasciicircum{}+\_G(\textbackslash{}Ysheaf\{\textbackslash{}Gamma\}(\textbackslash{}bar\{X\}), \textbackslash{}mathcal\{O\})\$: nous noterons \$\textbackslash{}bar\textbackslash{}theta\_\{G\}\textasciicircum{}\{\textbackslash{}ast\}\$ =\allowbreak > \$\textbackslash{}mathbf\{L\}\textasciicircum{}\{+\}\textbackslash{}bar\{\textbackslash{}theta\}\textasciicircum{}\{\textbackslash{}ast\}\$ prend ses valeurs dans \$D\textasciicircum{}+\_G(\textbackslash{}Ysheaf\{\textbackslash{}Gamma\}(\textbackslash{}bar\{X\}), \textbackslash{}mathcal\{O\})\$: nous noterons \$\textbackslash{}mathbf\{L\}\textasciicircum{}\{+\}\_\{G\}\textbackslash{}bar\textbackslash{}theta\textasciicircum{}\{\textbackslash{}ast\}\$\par
the two models differ at this place\par
\par\noindent\hypertarget{DP0274}{}\textbf{DP0274 -- Vbis, p. 114 (one\_model\_fix)}\par
Lemme 2.4.7.1, \ensuremath{\alpha}\par
\textbf{Claude Opus 5.5.} Place: Lemme 2.4.7.1, \ensuremath{\alpha}; Reported print: R\textasciicircum{}+\_G \ensuremath{\bar{\theta}}\_* \ensuremath{\circ} L\textasciicircum{}+\_G \ensuremath{\bar{\theta}}*; Source TeX: \textbackslash{}circ \textbackslash{}bar\textbackslash{}theta\_\{G\}\textasciicircum{}\{\textbackslash{}ast\}\$ et; Note: The print has L\textasciicircum{}+\_G \ensuremath{\bar{\theta}}*.; Proposed edit: \textbackslash{}circ \textbackslash{}bar\textbackslash{}theta\_\{G\}\textasciicircum{}\{\textbackslash{}ast\}\$ et =\allowbreak > \textbackslash{}circ L\textasciicircum{}\{+\}\_\{G\}\textbackslash{}bar\textbackslash{}theta\textasciicircum{}\{\textbackslash{}ast\}\$ et\par
\textbf{GPT-6.1 Sol.} nothing reported here\par
a fix from one model only\par
\par\noindent\hypertarget{DP0275}{}\textbf{DP0275 -- Vbis, p. 114 (one\_model\_fix)}\par
Lemme 2.4.7.1, \ensuremath{\beta}\par
\textbf{Claude Opus 5.5.} Place: Lemme 2.4.7.1, \ensuremath{\beta}; Reported print: \ensuremath{\beta} : L\textasciicircum{}+\_G \ensuremath{\bar{\theta}}* \ensuremath{\circ} R\textasciicircum{}+\_G \ensuremath{\bar{\theta}}\_*; Source TeX: \$\textbackslash{}beta : \textbackslash{}bar\textbackslash{}theta\_\{G\}\textasciicircum{}\{\textbackslash{}ast\} \textbackslash{}circ; Note: The print has L\textasciicircum{}+\_G \ensuremath{\bar{\theta}}*.; Proposed edit: \$\textbackslash{}beta : \textbackslash{}bar\textbackslash{}theta\_\{G\}\textasciicircum{}\{\textbackslash{}ast\} \textbackslash{}circ =\allowbreak > \$\textbackslash{}beta : L\textasciicircum{}\{+\}\_\{G\}\textbackslash{}bar\textbackslash{}theta\textasciicircum{}\{\textbackslash{}ast\} \textbackslash{}circ\par
\textbf{GPT-6.1 Sol.} nothing reported here\par
a fix from one model only\par
\par\noindent\hypertarget{DP0276}{}\textbf{DP0276 -- Vbis, p. 114 (conflict)}\par
Déf. 2.4.8, 2°) | Definition 2.4.8, condition 2°\par
\textbf{Claude Opus 5.5.} Place: Déf. 2.4.8, 2°); Reported print: Le foncteur L\textasciicircum{}+\_G \ensuremath{\bar{\theta}}* est pleinement fidèle.; Source TeX: Le foncteur \$\textbackslash{}bar\textbackslash{}theta\_G\textasciicircum{}\{\textbackslash{}star\}\$ pleinement fidèle.; Note: The word «est» is missing from the TeX, and the TeX drops L\textasciicircum{}+ as elsewhere.; Proposed edit: \$\textbackslash{}bar\textbackslash{}theta\_G\textasciicircum{}\{\textbackslash{}star\}\$ pleinement =\allowbreak > \$L\textasciicircum{}\{+\}\_\{G\}\textbackslash{}bar\textbackslash{}theta\textasciicircum{}\{\textbackslash{}ast\}\$ est pleinement\par
\textbf{GPT-6.1 Sol.} Place: Definition 2.4.8, condition 2°; Reported print: est pleinement fidèle.; Source TeX: pleinement fidèle.; Note: The TeX omits the printed verb « est ».; Proposed edit: pleinement fidèle. =\allowbreak > est pleinement fidèle.\par
the two models differ at this place\par
\par\noindent\hypertarget{DP0277}{}\textbf{DP0277 -- Vbis, p. 115 (one\_model\_fix)}\par
Théorème 2.4.12\par
\textbf{Claude Opus 5.5.} Place: Théorème 2.4.12; Reported print: l'image essentielle de L\textasciicircum{}+\_G \ensuremath{\bar{\theta}}*; Source TeX: l'image essentielle de \$\textbackslash{}bar\{\textbackslash{}theta\}\_G\textasciicircum{}\{\textbackslash{}ast\}\$; Note: The print has L\textasciicircum{}+\_G \ensuremath{\bar{\theta}}*.; Proposed edit: l'image essentielle de \$\textbackslash{}bar\{\textbackslash{}theta\}\_G\textasciicircum{}\{\textbackslash{}ast\}\$ =\allowbreak > l'image essentielle de \$L\textasciicircum{}\{+\}\_\{G\}\textbackslash{}bar\{\textbackslash{}theta\}\textasciicircum{}\{\textbackslash{}ast\}\$\par
\textbf{GPT-6.1 Sol.} nothing reported here\par
a fix from one model only\par
\par\noindent\hypertarget{DP0278}{}\textbf{DP0278 -- Vbis, p. 117 (conflict)}\par
Prop. 2.5.3, label of the first arrow of the factorisation | Proposition 2.5.3, first arrow of the factorization\par
\textbf{Claude Opus 5.5.} Place: Prop. 2.5.3, label of the first arrow of the factorisation; Reported print: \_F R\textasciicircum{}+\ensuremath{\varepsilon}\_\{Z\ensuremath{\cdot}*\}; Source TeX: F\textasciicircum{}\{\textbackslash{}sisi\{\textbackslash{}RR\textasciicircum{}+\}\{R\} \textbackslash{}varepsilon\} Z\textasciicircum{}\{\textbackslash{}bigdot\}\textbackslash{}ast; Note: In the TeX, Z\textasciicircum{}\ensuremath{\bullet} * is not the subscript of \ensuremath{\varepsilon} and F carries an exponent. The print reads \_F R\textasciicircum{}+\ensuremath{\varepsilon}\_\{Z\ensuremath{\cdot}*\}, as in the TeX of p. 119.; Proposed edit: \textbackslash{}xrightarrow\{F\textasciicircum{}\{\textbackslash{}sisi\{\textbackslash{}RR\textasciicircum{}+\}\{R\} \textbackslash{}varepsilon\} Z\textasciicircum{}\{\textbackslash{}bigdot\}\textbackslash{}ast\} =\allowbreak > \textbackslash{}xrightarrow\{\{\}\_F\textbackslash{}sisi\{\textbackslash{}RR\textasciicircum{}+\}\{R\} \textbackslash{}varepsilon\_\{Z\textasciicircum{}\{\textbackslash{}bigdot\}\textbackslash{}ast\}\}\par
\textbf{GPT-6.1 Sol.} Place: Proposition 2.5.3, first arrow of the factorization; Reported print: \{\}\_F\textbackslash{}RR\textasciicircum{}+\textbackslash{}varepsilon\_\{Z\textasciicircum{}\{\textbackslash{}bigdot\}\textbackslash{}ast\}; Source TeX: F\textasciicircum{}\{\textbackslash{}sisi\{\textbackslash{}RR\textasciicircum{}+\}\{R\} \textbackslash{}varepsilon\} Z\textasciicircum{}\{\textbackslash{}bigdot\}\textbackslash{}ast; Note: In print F is the left subscript of the derived functor, and Z\textasciicircum{}\{\textbackslash{}bigdot\}\textbackslash{}ast is the subscript of \ensuremath{\varepsilon}. The TeX instead makes the derived functor an exponent of F and detaches its index.; Proposed edit: F\textasciicircum{}\{\textbackslash{}sisi\{\textbackslash{}RR\textasciicircum{}+\}\{R\} \textbackslash{}varepsilon\} Z\textasciicircum{}\{\textbackslash{}bigdot\}\textbackslash{}ast =\allowbreak > \{\}\_F\textbackslash{}sisi\{\textbackslash{}RR\textasciicircum{}+\}\{R\}\textbackslash{}varepsilon\_\{Z\textasciicircum{}\{\textbackslash{}bigdot\}\textbackslash{}ast\}\par
the two models differ at this place\par
\par\noindent\hypertarget{DP0279}{}\textbf{DP0279 -- Vbis, p. 117 (one\_model\_fix)}\par
Proposition 2.5.3, factorization display, source category\par
\textbf{Claude Opus 5.5.} nothing reported here\par
\textbf{GPT-6.1 Sol.} Place: Proposition 2.5.3, factorization display, source category; Reported print: D\textasciicircum{}\{+\}(\textbackslash{}underline\{\textbackslash{}Gamma\}(S\textbackslash{}times D),\textbackslash{}mathcal\{O\}\_S); Source TeX: D\textasciicircum{}\{+\} (\textbackslash{}Gamma (S \textbackslash{}times D), \textbackslash{}mathcal\{O\}\_S); Note: The mathematical underline on \ensuremath{\Gamma} identifies the category of sections; it is not textual emphasis.; Proposed edit: D\textasciicircum{}\{+\} (\textbackslash{}Gamma (S \textbackslash{}times D), \textbackslash{}mathcal\{O\}\_S) =\allowbreak > D\textasciicircum{}\{+\} (\textbackslash{}Ysheaf\{\textbackslash{}Gamma\} (S \textbackslash{}times D), \textbackslash{}mathcal\{O\}\_S)\par
a fix from one model only\par
\par\noindent\hypertarget{DP0280}{}\textbf{DP0280 -- Vbis, p. 118 (one\_model\_fix)}\par
diagram of Prop. 2.5.3, top arrow\par
\textbf{Claude Opus 5.5.} Place: diagram of Prop. 2.5.3, top arrow; Reported print: R\textasciicircum{}+(e\_i\textasciicircum{}*); Source TeX: \textbackslash{}sisi\{\textbackslash{}RR\textasciicircum{}\{+\}\}\{R\}e\textasciicircum{}\{\textbackslash{}ast\}\_i; Note: The print has parentheses around e\_i\textasciicircum{}*.; Proposed edit: \textbackslash{}sisi\{\textbackslash{}RR\textasciicircum{}\{+\}\}\{R\}e\textasciicircum{}\{\textbackslash{}ast\}\_i\} =\allowbreak > \textbackslash{}sisi\{\textbackslash{}RR\textasciicircum{}\{+\}\}\{R\}(e\textasciicircum{}\{\textbackslash{}ast\}\_i)\}\par
\textbf{GPT-6.1 Sol.} nothing reported here\par
a fix from one model only\par
\par\noindent\hypertarget{DP0281}{}\textbf{DP0281 -- Vbis, p. 118 (one\_model\_fix)}\par
Opening triangular diagram, downward arrow\par
\textbf{Claude Opus 5.5.} nothing reported here\par
\textbf{GPT-6.1 Sol.} Place: Opening triangular diagram, downward arrow; Reported print: \{\}\_F\textbackslash{}RR\textasciicircum{}\{+\}\textbackslash{}epsilon\_\{Z\textasciicircum{}\{\textbackslash{}bigdot\}\textbackslash{}ast\}; Source TeX: F\textasciicircum{}\{\textbackslash{}sisi\{\textbackslash{}RR\textasciicircum{}\{+\}\}\{R\}\textbackslash{}epsilon\_\{Z\textasciicircum{}\{\textbackslash{}bigdot\}\textbackslash{}ast\}\}; Note: The print uses a left-subscript F on the derived functor, not F with that functor as its exponent.; Proposed edit: F\textasciicircum{}\{\textbackslash{}sisi\{\textbackslash{}RR\textasciicircum{}\{+\}\}\{R\}\textbackslash{}epsilon\_\{Z\textasciicircum{}\{\textbackslash{}bigdot\}\textbackslash{}ast\}\} =\allowbreak > \{\}\_F\textbackslash{}sisi\{\textbackslash{}RR\textasciicircum{}\{+\}\}\{R\}\textbackslash{}epsilon\_\{Z\textasciicircum{}\{\textbackslash{}bigdot\}\textbackslash{}ast\}\par
a fix from one model only\par
\par\noindent\hypertarget{DP0282}{}\textbf{DP0282 -- Vbis, p. 118 (one\_model\_fix)}\par
Diagram (2.5.3.3), arrow from D\textasciicircum{}+F(S,O\_S) to D\textasciicircum{}+(S,O\_S)\par
\textbf{Claude Opus 5.5.} nothing reported here\par
\textbf{GPT-6.1 Sol.} Place: Diagram (2.5.3.3), arrow from D\textasciicircum{}+F(S,O\_S) to D\textasciicircum{}+(S,O\_S); Reported print: \textbackslash{}iota; Source TeX: \textbackslash{}imath; Note: The printed label is the Greek iota denoting the forgetful functor, not dotless Latin i.; Proposed edit: \textbackslash{}ar[ru]\_\{\textbackslash{}imath\} =\allowbreak > \textbackslash{}ar[ru]\_\{\textbackslash{}iota\}\par
a fix from one model only\par
\par\noindent\hypertarget{DP0283}{}\textbf{DP0283 -- Vbis, p. 118 (one\_model\_fix)}\par
Diagram (2.5.3.3), arrow from K\textasciicircum{}+(F(S,O\_S)) to K\textasciicircum{}+(S,O\_S)\par
\textbf{Claude Opus 5.5.} nothing reported here\par
\textbf{GPT-6.1 Sol.} Place: Diagram (2.5.3.3), arrow from K\textasciicircum{}+(F(S,O\_S)) to K\textasciicircum{}+(S,O\_S); Reported print: \textbackslash{}iota; Source TeX: \textbackslash{}imath; Note: This printed forgetful-functor label is also Greek iota.; Proposed edit: \textbackslash{}ar[rd]\textasciicircum{}\{\textbackslash{}imath\} =\allowbreak > \textbackslash{}ar[rd]\textasciicircum{}\{\textbackslash{}iota\}\par
a fix from one model only\par
\par\noindent\hypertarget{DP0284}{}\textbf{DP0284 -- Vbis, p. 118 (one\_model\_fix)}\par
Diagram (2.5.3.3), bottom horizontal arrow\par
\textbf{Claude Opus 5.5.} nothing reported here\par
\textbf{GPT-6.1 Sol.} Place: Diagram (2.5.3.3), bottom horizontal arrow; Reported print: K\textasciicircum{}\{+\}(\textbackslash{}epsilon\_\{Z\textasciicircum{}\{\textbackslash{}bigdot\}\textbackslash{}ast\}); Source TeX: K\textasciicircum{}\{+\}(\textbackslash{}epsilon\_\{Z\textasciicircum{}\{\textbackslash{}bigdot\}\}\textbackslash{}ast); Note: The printed star is part of \ensuremath{\varepsilon}'s subscript.; Proposed edit: K\textasciicircum{}\{+\}(\textbackslash{}epsilon\_\{Z\textasciicircum{}\{\textbackslash{}bigdot\}\}\textbackslash{}ast) =\allowbreak > K\textasciicircum{}\{+\}(\textbackslash{}epsilon\_\{Z\textasciicircum{}\{\textbackslash{}bigdot\}\textbackslash{}ast\})\par
a fix from one model only\par
\par\noindent\hypertarget{DP0285}{}\textbf{DP0285 -- Vbis, p. 120 (conflict)}\par
after (2.5.5.2), second isomorphism | Explanatory display following (2.5.5.2), final factor\par
\textbf{Claude Opus 5.5.} Place: after (2.5.5.2), second isomorphism; Reported print: id\_\{D\textasciicircum{}+(S,O\_S)\} \textasciitilde{}> R\textasciicircum{}+\ensuremath{\bar{\theta}}\_* \ensuremath{\circ} L\textasciicircum{}+ \ensuremath{\bar{\theta}}\textasciicircum{}*; Source TeX: \textbackslash{}bar\{\textbackslash{}theta\}\_\textbackslash{}ast \textbackslash{}circ \textbackslash{}bar\{\textbackslash{}theta\}\_\textbackslash{}ast\{\textbackslash{}quoi\}; Note: The last factor is L\textasciicircum{}+ \ensuremath{\bar{\theta}}\textasciicircum{}* in the print, not \ensuremath{\bar{\theta}}\_*.; Proposed edit: \textbackslash{}circ \textbackslash{}bar\{\textbackslash{}theta\}\_\textbackslash{}ast\{\textbackslash{}quoi\} =\allowbreak > \textbackslash{}circ \textbackslash{}LL\textasciicircum{}\{+\}\textbackslash{}bar\{\textbackslash{}theta\}\textasciicircum{}\textbackslash{}ast\{\textbackslash{}quoi\}\par
\textbf{GPT-6.1 Sol.} Place: Explanatory display following (2.5.5.2), final factor; Reported print: \textbackslash{}RR\textasciicircum{}+\textbackslash{}bar\{\textbackslash{}theta\}\_\{\textbackslash{}ast\}\textbackslash{}circ\textbackslash{}LL\textasciicircum{}\{+\}\textbackslash{}bar\{\textbackslash{}theta\}\textasciicircum{}\{\textbackslash{}ast\}); Source TeX: \textbackslash{}sisi\{\textbackslash{}RR\textasciicircum{}+\}\{R\} \textbackslash{}bar\{\textbackslash{}theta\}\_\textbackslash{}ast \textbackslash{}circ \textbackslash{}bar\{\textbackslash{}theta\}\_\textbackslash{}ast\{\textbackslash{}quoi\}.; Note: The final factor in print is the derived inverse image L\textasciicircum{}\{+\}\textbackslash{}bar\{\textbackslash{}theta\}\textasciicircum{}\{\textbackslash{}ast\}, not another direct image. This edit also closes the printed parenthetical explanation.; Proposed edit: \textbackslash{}bar\{\textbackslash{}theta\}\_\textbackslash{}ast\{\textbackslash{}quoi\}. =\allowbreak > \textbackslash{}LL\textasciicircum{}\{+\}\textbackslash{}bar\{\textbackslash{}theta\}\textasciicircum{}\{\textbackslash{}ast\})\{\textbackslash{}quoi\}.\par
the two models differ at this place\par
\par\noindent\hypertarget{DP0286}{}\textbf{DP0286 -- Vbis, p. 120 (conflict)}\par
«est essentiellement commutatif car» | Sentence immediately following (2.5.5.2)\par
\textbf{Claude Opus 5.5.} Place: «est essentiellement commutatif car»; Reported print: est essentiellement commutatif.  ( R\textasciicircum{}+… \ensuremath{\circ} L\textasciicircum{}+\ensuremath{\bar{\theta}}\textasciicircum{}*) .; Source TeX: est essentiellement commutatif car; Note: The print has no «car»: the sentence ends with a full stop and the following display is put in parentheses.; Proposed edit: est essentiellement commutatif car =\allowbreak > est essentiellement commutatif.\par
\textbf{GPT-6.1 Sol.} Place: Sentence immediately following (2.5.5.2); Reported print: est essentiellement commutatif. (…); Source TeX: est essentiellement commutatif car; Note: The print has a complete sentence followed by a parenthetical explanation, without «car». Apply together with the closing-parenthesis edit in the following finding.; Proposed edit: est essentiellement commutatif car \$\$ =\allowbreak > est essentiellement commutatif. \$\$ (\par
the two models differ at this place\par
\par\noindent\hypertarget{DP0287}{}\textbf{DP0287 -- Vbis, p. 121 (conflict)}\par
diagram (2.5.5.3), label of the top-left arrow | Diagram (2.5.5.3), upper-left horizontal arrow\par
\textbf{Claude Opus 5.5.} Place: diagram (2.5.5.3), label of the top-left arrow; Reported print: L\textasciicircum{}+ \ensuremath{\bar{\theta}}\textasciicircum{}* (L double underlined); Source TeX: \textbackslash{}bar\{\textbackslash{}theta\}\textasciicircum{}\{\textbackslash{}ast\}; Note: L\textasciicircum{}+ is dropped. \textbackslash{}LL is assumed by analogy with \textbackslash{}RR.; Proposed edit: \textbackslash{}ar[r]\textasciicircum{}\{\textbackslash{}bar\{\textbackslash{}theta\}\textasciicircum{}\{\textbackslash{}ast\}\} =\allowbreak > \textbackslash{}ar[r]\textasciicircum{}\{\textbackslash{}LL\textasciicircum{}\{+\}\textbackslash{}bar\{\textbackslash{}theta\}\textasciicircum{}\{\textbackslash{}ast\}\}\par
\textbf{GPT-6.1 Sol.} Place: Diagram (2.5.5.3), upper-left horizontal arrow; Reported print: \textbackslash{}mathbf\{L\}\textasciicircum{}\{+\}\textbackslash{}bar\{\textbackslash{}theta\}\textasciicircum{}\{\textbackslash{}ast\}; Source TeX: \textbackslash{}bar\{\textbackslash{}theta\}\textasciicircum{}\{\textbackslash{}ast\}; Note: The edition systematically suppresses the left-derived-functor symbol, transferring G to the pullback notation where applicable. This is a representative edit; three examples are given in coordinator.; Proposed edit: \textbackslash{}ar[r]\textasciicircum{}\{\textbackslash{}bar\{\textbackslash{}theta\}\textasciicircum{}\{\textbackslash{}ast\}\} =\allowbreak > \textbackslash{}ar[r]\textasciicircum{}\{\textbackslash{}mathbf\{L\}\textasciicircum{}\{+\}\textbackslash{}bar\{\textbackslash{}theta\}\textasciicircum{}\{\textbackslash{}ast\}\}\par
the two models differ at this place\par
\par\noindent\hypertarget{DP0288}{}\textbf{DP0288 -- Vbis, p. 121 (conflict)}\par
Le foncteur spectral annoncé ... compose avec | Composition defining the spectral functor, immediately below the diagram\par
\textbf{Claude Opus 5.5.} Place: Le foncteur spectral annoncé ... compose avec; Reported print: \_F R\textasciicircum{}+ \textbackslash{}overline\{H\ensuremath{\circ}\ensuremath{\theta}\}\_* \ensuremath{\circ} L\textasciicircum{}+(\ensuremath{\bar{\theta}}\textasciicircum{}*); Source TeX: \{\}\_F R(\textbackslash{}bar\{H \textbackslash{}circ \textbackslash{}theta\})\_\{\textbackslash{}ast\} \textbackslash{}circ\textbackslash{}bar\{\textbackslash{}theta\}\textasciicircum{}\{\textbackslash{}ast\}; Note: The print has R\textasciicircum{}+, and here the TeX has no \textbackslash{}sisi. The print also has L\textasciicircum{}+(\ensuremath{\bar{\theta}}\textasciicircum{}*).; Proposed edit: \$\{\}\_F R(\textbackslash{}bar\{H \textbackslash{}circ \textbackslash{}theta\})\_\{\textbackslash{}ast\} \textbackslash{}circ\textbackslash{}bar\{\textbackslash{}theta\}\textasciicircum{}\{\textbackslash{}ast\}\$ =\allowbreak > \$\{\}\_F \textbackslash{}sisi\{\textbackslash{}RR\textasciicircum{}\{+\}\}\{R\}(\textbackslash{}bar\{H \textbackslash{}circ \textbackslash{}theta\})\_\{\textbackslash{}ast\} \textbackslash{}circ \textbackslash{}LL\textasciicircum{}\{+\}(\textbackslash{}bar\{\textbackslash{}theta\}\textasciicircum{}\{\textbackslash{}ast\})\$\par
\textbf{GPT-6.1 Sol.} Place: Composition defining the spectral functor, immediately below the diagram; Reported print: \{\}\_F\textbackslash{}RR\textasciicircum{}\{+\}(\textbackslash{}overline\{H\textbackslash{}circ\textbackslash{}theta\})\_\{\textbackslash{}ast\}; Source TeX: \{\}\_F R(\textbackslash{}bar\{H \textbackslash{}circ \textbackslash{}theta\})\_\{\textbackslash{}ast\}; Note: The printed derived functor has the superscript +. Unlike the diagram, this occurrence has no sisi recording that reading.; Proposed edit: \{\}\_F R(\textbackslash{}bar\{H \textbackslash{}circ \textbackslash{}theta\})\_\{\textbackslash{}ast\} =\allowbreak > \{\}\_F \textbackslash{}RR\textasciicircum{}\{+\}(\textbackslash{}bar\{H \textbackslash{}circ \textbackslash{}theta\})\_\{\textbackslash{}ast\}\par
the two models differ at this place\par
\par\noindent\hypertarget{DP0289}{}\textbf{DP0289 -- Vbis, p. 121 (one\_model\_fix)}\par
Proposition 2.5.6, displayed E\_2 term\par
\textbf{Claude Opus 5.5.} Place: Proposition 2.5.6, displayed E\_2 term; Reported print: Ȟ\textasciicircum{}p \ensuremath{\circ} R\textasciicircum{}q\ensuremath{\Gamma}((H\ensuremath{\circ}\ensuremath{\theta})\_*) \ensuremath{\circ} L\textasciicircum{}+ \ensuremath{\bar{\theta}}\textasciicircum{}*; Source TeX: \textbackslash{}bar\{\textbackslash{}theta\}\textasciicircum{}\textbackslash{}ast; Note: L\textasciicircum{}+ is missing.; Proposed edit: \textbackslash{}bar\{\textbackslash{}theta\}\textasciicircum{}\textbackslash{}ast =\allowbreak > \textbackslash{}LL\textasciicircum{}\{+\}\textbackslash{}bar\{\textbackslash{}theta\}\textasciicircum{}\textbackslash{}ast\par
\textbf{GPT-6.1 Sol.} nothing reported here\par
a fix from one model only\par
\par\noindent\hypertarget{DP0290}{}\textbf{DP0290 -- Vbis, p. 122 (one\_model\_fix)}\par
Proposition 2.5.7, E\_1 term\par
\textbf{Claude Opus 5.5.} Place: Proposition 2.5.7, E\_1 term; Reported print: R\textasciicircum{}q\_G(h\ensuremath{\circ}\ensuremath{\theta})\_* \ensuremath{\circ} L\textasciicircum{}+\_G \ensuremath{\theta}\_p\textasciicircum{}*; Source TeX: \textbackslash{}theta\_\{p,G\}\textasciicircum{}\textbackslash{}ast; Note: L\textasciicircum{}+\_G has been replaced by a subscript G on \ensuremath{\theta}.; Proposed edit: \textbackslash{}theta\_\{p,G\}\textasciicircum{}\textbackslash{}ast =\allowbreak > \textbackslash{}LL\textasciicircum{}\{+\}\_G \textbackslash{}theta\_\{p\}\textasciicircum{}\textbackslash{}ast\par
\textbf{GPT-6.1 Sol.} nothing reported here\par
a fix from one model only\par
\par\noindent\hypertarget{DP0291}{}\textbf{DP0291 -- Vbis, p. 122 (conflict)}\par
Proposition 2.5.7, arrow of the spectral functor | Proposition 2.5.7, spectral-sequence abutment arrow\par
\textbf{Claude Opus 5.5.} Place: Proposition 2.5.7, arrow of the spectral functor; Reported print: \ensuremath{\Longrightarrow} R\textasciicircum{}\{p+q\}\_G h\_*; Source TeX: \textbackslash{}to R\textasciicircum{}\{p + q\}\_G; Note: The print has a double arrow (abutment).; Proposed edit: \textbackslash{}to R\textasciicircum{}\{p + q\}\_G =\allowbreak > \textbackslash{}Longrightarrow R\textasciicircum{}\{p + q\}\_G\par
\textbf{GPT-6.1 Sol.} Place: Proposition 2.5.7, spectral-sequence abutment arrow; Reported print: \textbackslash{}Rightarrow R\textasciicircum{}\{p+q\}\_G h\_\{\textbackslash{}ast\}; Source TeX: \textbackslash{}to R\textasciicircum{}\{p + q\}\_G h\_\{\textbackslash{}ast\}; Note: The print has a double convergence arrow, not a single arrow.; Proposed edit: \textbackslash{}to R\textasciicircum{}\{p + q\}\_G h\_\{\textbackslash{}ast\} =\allowbreak > \textbackslash{}Rightarrow R\textasciicircum{}\{p + q\}\_G h\_\{\textbackslash{}ast\}\par
the two models differ at this place\par
\par\noindent\hypertarget{DP0292}{}\textbf{DP0292 -- Vbis, p. 122 (one\_model\_fix)}\par
Proposition 2.5.7, E\_2 term\par
\textbf{Claude Opus 5.5.} Place: Proposition 2.5.7, E\_2 term; Reported print: \ensuremath{\circ} L\textasciicircum{}+\_G \ensuremath{\bar{\theta}}\textasciicircum{}*; Source TeX: \textbackslash{}bar\textbackslash{}theta\_G\textasciicircum{}\{\textbackslash{}star\}; Note: L\textasciicircum{}+\_G has been replaced by a subscript G; the print has * (\textbackslash{}ast).; Proposed edit: \textbackslash{}bar\textbackslash{}theta\_G\textasciicircum{}\{\textbackslash{}star\} =\allowbreak > \textbackslash{}LL\textasciicircum{}\{+\}\_G \textbackslash{}bar\textbackslash{}theta\textasciicircum{}\{\textbackslash{}ast\}\par
\textbf{GPT-6.1 Sol.} nothing reported here\par
a fix from one model only\par
\par\noindent\hypertarget{DP0293}{}\textbf{DP0293 -- Vbis, p. 123 (agreed\_fix\_held)}\par
first paragraph, 'pour le foncteur i\_\{n*\}' | First paragraph, final occurrence of the right adjoint\par
\textbf{Claude Opus 5.5.} Place: first paragraph, 'pour le foncteur i\_\{n*\}'; Reported print: i\_\{n*\}; Source TeX: \$i\_\{n\_\textbackslash{}ast\}\$; Note: The * stands beside n, not as a subscript of n.; Proposed edit: \$i\_\{n\_\textbackslash{}ast\}\$ =\allowbreak > \$i\_\{n\textbackslash{}ast\}\$\par
\textbf{GPT-6.1 Sol.} Place: First paragraph, final occurrence of the right adjoint; Reported print: i\_\{n\textbackslash{}ast\}; Source TeX: i\_\{n\_\textbackslash{}ast\}; Note: The printed star accompanies n in the index of i; the TeX instead nests it as a subscript of n.; Proposed edit: i\_\{n\_\textbackslash{}ast\} =\allowbreak > i\_\{n\textbackslash{}ast\}\par
both models give the same fix, held: a model marks it low confidence\par
\par\noindent\hypertarget{DP0294}{}\textbf{DP0294 -- Vbis, p. 124 (conflict)}\par
Lemme 3.0.2.4 | Lemma 3.0.2.4, category containing X and X'\par
\textbf{Claude Opus 5.5.} Place: Lemme 3.0.2.4; Reported print: deux objets de \ensuremath{\Delta}\_n\textasciicircum{}oE; Source TeX: \$\textbackslash{}Delta\textasciicircum{}\{\textbackslash{}circ\}E\$. Soient; Note: The subscript n is missing.; Proposed edit: \$\textbackslash{}Delta\textasciicircum{}\{\textbackslash{}circ\}E\$. Soient =\allowbreak > \$\{\textbackslash{}mathop\{\textbackslash{}Delta\}\_n\}\textasciicircum{}\{\textbackslash{}circ\}E\$. Soient\par
\textbf{GPT-6.1 Sol.} Place: Lemma 3.0.2.4, category containing X and X'; Reported print: \textbackslash{}Delta\_n\textasciicircum{}\{\textbackslash{}circ\}E; Source TeX: \textbackslash{}Delta\textasciicircum{}\{\textbackslash{}circ\}E; Note: The printed category is n-truncated; its subscript n is missing from the TeX.; Proposed edit: \$\textbackslash{}Delta\textasciicircum{}\{\textbackslash{}circ\}E\$ =\allowbreak > \$\textbackslash{}Delta\_n\textasciicircum{}\{\textbackslash{}circ\}E\$\par
the two models differ at this place\par
\par\noindent\hypertarget{DP0295}{}\textbf{DP0295 -- Vbis, p. 124 (one\_model\_fix)}\par
Proposition 3.1.1, \ensuremath{\eta}\textasciicircum{}f\_G, source\par
\textbf{Claude Opus 5.5.} Place: Proposition 3.1.1, \ensuremath{\eta}\textasciicircum{}f\_G, source; Reported print: R\textasciicircum{}+\_G ū'\_* \ensuremath{\circ} L\textasciicircum{}+\_G ū'\textasciicircum{}*; Source TeX: \{\textbackslash{}bar\{u\}'\}\_G\textasciicircum{}\textbackslash{}ast; Note: L\textasciicircum{}+\_G has been replaced by a subscript G.; Proposed edit: \{\textbackslash{}bar\{u\}'\}\_G\textasciicircum{}\textbackslash{}ast \textbackslash{}longrightarrow =\allowbreak > \textbackslash{}LL\textasciicircum{}\{+\}\_G \textbackslash{}bar\{u\}'\textasciicircum{}\textbackslash{}ast \textbackslash{}longrightarrow\par
\textbf{GPT-6.1 Sol.} nothing reported here\par
a fix from one model only\par
\par\noindent\hypertarget{DP0296}{}\textbf{DP0296 -- Vbis, p. 124 (one\_model\_fix)}\par
Proposition 3.1.1, \ensuremath{\eta}\textasciicircum{}f\_G, target\par
\textbf{Claude Opus 5.5.} Place: Proposition 3.1.1, \ensuremath{\eta}\textasciicircum{}f\_G, target; Reported print: R\textasciicircum{}+\_G ū\_* \ensuremath{\circ} L\textasciicircum{}+ ū\textasciicircum{}*; Source TeX: \textbackslash{}circ  \textbackslash{}bar\{u\}\_G\textasciicircum{}\textbackslash{}ast; Note: L\textasciicircum{}+ is missing. The print lacks the G there; see the misprint finding.; Proposed edit: \textbackslash{}circ  \textbackslash{}bar\{u\}\_G\textasciicircum{}\textbackslash{}ast =\allowbreak > \textbackslash{}circ  \textbackslash{}LL\textasciicircum{}\{+\}\_G \textbackslash{}bar\{u\}\textasciicircum{}\textbackslash{}ast\par
\textbf{GPT-6.1 Sol.} nothing reported here\par
a fix from one model only\par
\par\noindent\hypertarget{DP0297}{}\textbf{DP0297 -- Vbis, p. 125 (one\_model\_fix)}\par
diagram after Proposition 3.1.1, upper object\par
\textbf{Claude Opus 5.5.} Place: diagram after Proposition 3.1.1, upper object; Reported print: R\textasciicircum{}+\_G ū'\_* \ensuremath{\circ} L\textasciicircum{}+\_G ū'\textasciicircum{}*; Source TeX: \{\textbackslash{}overline\{u\}'\}\_G\textasciicircum{}\{\textbackslash{}ast\}; Note: L\textasciicircum{}+\_G has been replaced by a subscript G.; Proposed edit: \{\textbackslash{}overline\{u\}'\}\_G\textasciicircum{}\{\textbackslash{}ast\} =\allowbreak > \textbackslash{}LL\textasciicircum{}\{+\}\_G\textbackslash{}overline\{u\}'\textasciicircum{}\{\textbackslash{}ast\}\par
\textbf{GPT-6.1 Sol.} nothing reported here\par
a fix from one model only\par
\par\noindent\hypertarget{DP0298}{}\textbf{DP0298 -- Vbis, p. 125 (one\_model\_fix)}\par
diagram after Proposition 3.1.1, lower object\par
\textbf{Claude Opus 5.5.} Place: diagram after Proposition 3.1.1, lower object; Reported print: R\textasciicircum{}+\_G ū\_* \ensuremath{\circ} L\textasciicircum{}+\_G ū\textasciicircum{}*; Source TeX: \textbackslash{}overline\{u\}\_G\textasciicircum{}\{\textbackslash{}ast\}; Note: L\textasciicircum{}+\_G has been replaced by a subscript G.; Proposed edit: \textbackslash{}overline\{u\}\_G\textasciicircum{}\{\textbackslash{}ast\} =\allowbreak > \textbackslash{}LL\textasciicircum{}\{+\}\_G\textbackslash{}overline\{u\}\textasciicircum{}\{\textbackslash{}ast\}\par
\textbf{GPT-6.1 Sol.} nothing reported here\par
a fix from one model only\par
\par\noindent\hypertarget{DP0299}{}\textbf{DP0299 -- Vbis, p. 125 (one\_model\_fix)}\par
Functor-category Hom symbols, especially (3.1.1.2)\par
\textbf{Claude Opus 5.5.} nothing reported here\par
\textbf{GPT-6.1 Sol.} Place: Functor-category Hom symbols, especially (3.1.1.2); Reported print: \textbackslash{}underline\{\textbackslash{}Hom\}((I\textbackslash{}times D)\textasciicircum{}\{\textbackslash{}circ\},B/S); Source TeX: \textbackslash{}Hom((I \textbackslash{}times D)\textasciicircum{}\{\textbackslash{}circ\}, B/S); Note: The print underlines categorical Hom, distinguishing it from the un-underlined morphism sets on page 123. This is a representative edit for the convention; three examples are given in coordinator.; Proposed edit: \textbackslash{}Hom((I \textbackslash{}times D)\textasciicircum{}\{\textbackslash{}circ\}, B/S) =\allowbreak > \textbackslash{}underline\{\textbackslash{}Hom\}((I \textbackslash{}times D)\textasciicircum{}\{\textbackslash{}circ\}, B/S)\par
a fix from one model only\par
\par\noindent\hypertarget{DP0300}{}\textbf{DP0300 -- Vbis, p. 125 (one\_model\_fix)}\par
Isomorphism (3.1.1.2), right-hand side\par
\textbf{Claude Opus 5.5.} nothing reported here\par
\textbf{GPT-6.1 Sol.} Place: Isomorphism (3.1.1.2), right-hand side; Reported print: \textbackslash{}Hom(I\textasciicircum{}\{\textbackslash{}circ\},\textbackslash{}Hom(D\textasciicircum{}\{\textbackslash{}circ\},B/S)); Source TeX: \textbackslash{}Hom(I\textasciicircum{}\{0\},\textbackslash{}Hom(D\textasciicircum{}0,B/S)); Note: Both printed superscripts are opposite-category circles, as on the left-hand side, rather than numerical zeros.; Proposed edit: \textbackslash{}Hom(I\textasciicircum{}\{0\}, \textbackslash{}Hom (D\textasciicircum{}0, B/S)) =\allowbreak > \textbackslash{}Hom(I\textasciicircum{}\{\textbackslash{}circ\}, \textbackslash{}Hom (D\textasciicircum{}\{\textbackslash{}circ\}, B/S))\par
a fix from one model only\par
\par\noindent\hypertarget{DP0301}{}\textbf{DP0301 -- Vbis, p. 127 (conflict)}\par
display after « On obtient par suite un morphisme », source | Morphisme précédant (3.1.1.8), membre de gauche\par
\textbf{Claude Opus 5.5.} Place: display after « On obtient par suite un morphisme », source; Reported print: R\textasciicircum{}+ ū'\_* \ensuremath{\circ} L\textasciicircum{}+ ū'*(F\ensuremath{\cdot}); Source TeX: \textbackslash{}sisi\{\textbackslash{}RR\textasciicircum{}+\}\{R\} \textbackslash{}bar\{u\}'\_\textbackslash{}ast \textbackslash{}circ \{\textbackslash{}bar\{u\}'\}\textasciicircum{}\textbackslash{}ast(F\textasciicircum{}\{\textbackslash{}bigdot\}); Note: The print has L\textasciicircum{}+ (double-underlined L) before ū'*; the TeX drops it. If \textbackslash{}LL is not defined in macros.tex, use the edition's macro for the derived L.; Proposed edit: \textbackslash{}bar\{u\}'\_\textbackslash{}ast \textbackslash{}circ \{\textbackslash{}bar\{u\}'\}\textasciicircum{}\textbackslash{}ast(F\textasciicircum{}\{\textbackslash{}bigdot\}) =\allowbreak > \textbackslash{}bar\{u\}'\_\textbackslash{}ast \textbackslash{}circ \textbackslash{}LL\textasciicircum{}+\{\textbackslash{}bar\{u\}'\}\textasciicircum{}\textbackslash{}ast(F\textasciicircum{}\{\textbackslash{}bigdot\})\par
\textbf{GPT-6.1 Sol.} Place: Morphisme précédant (3.1.1.8), membre de gauche; Reported print: \textbackslash{}RR\textasciicircum{}+\textbackslash{}bar\{u\}'\_\textbackslash{}ast\textbackslash{}circ\textbackslash{}mathbf\{L\}\textasciicircum{}+\{\textbackslash{}bar\{u\}'\}\textasciicircum{}\textbackslash{}ast(F\textasciicircum{}\{\textbackslash{}bigdot\}); Source TeX: \textbackslash{}sisi\{\textbackslash{}RR\textasciicircum{}+\}\{R\} \textbackslash{}bar\{u\}'\_\textbackslash{}ast \textbackslash{}circ \{\textbackslash{}bar\{u\}'\}\textasciicircum{}\textbackslash{}ast(F\textasciicircum{}\{\textbackslash{}bigdot\}); Note: Le foncteur dérivé L\textasciicircum{}+ imprimé devant l'image inverse est absent du TeX.; Proposed edit: \textbackslash{}circ \{\textbackslash{}bar\{u\}'\}\textasciicircum{}\textbackslash{}ast(F\textasciicircum{}\{\textbackslash{}bigdot\}) =\allowbreak > \textbackslash{}circ \textbackslash{}mathbf\{L\}\textasciicircum{}\{+\} \{\textbackslash{}bar\{u\}'\}\textasciicircum{}\textbackslash{}ast(F\textasciicircum{}\{\textbackslash{}bigdot\})\par
the two models differ at this place\par
\par\noindent\hypertarget{DP0302}{}\textbf{DP0302 -- Vbis, p. 127 (conflict)}\par
display after « On obtient par suite un morphisme », target | Morphisme précédant (3.1.1.8), membre de droite\par
\textbf{Claude Opus 5.5.} Place: display after « On obtient par suite un morphisme », target; Reported print: R\textasciicircum{}+ū\_* \ensuremath{\circ} L\textasciicircum{}+ ū*(F\ensuremath{\cdot}); Source TeX: \textbackslash{}sisi\{\textbackslash{}RR\textasciicircum{}+\}\{R\} \textbackslash{}bar\{u\}\_\textbackslash{}ast \textbackslash{}circ \textbackslash{}bar\{u\}\textasciicircum{}\textbackslash{}ast(F\textasciicircum{}\{\textbackslash{}bigdot\}); Note: L\textasciicircum{}+ is missing before ū*.; Proposed edit: \textbackslash{}bar\{u\}\_\textbackslash{}ast \textbackslash{}circ \textbackslash{}bar\{u\}\textasciicircum{}\textbackslash{}ast(F\textasciicircum{}\{\textbackslash{}bigdot\}) =\allowbreak > \textbackslash{}bar\{u\}\_\textbackslash{}ast \textbackslash{}circ \textbackslash{}LL\textasciicircum{}+\textbackslash{}bar\{u\}\textasciicircum{}\textbackslash{}ast(F\textasciicircum{}\{\textbackslash{}bigdot\})\par
\textbf{GPT-6.1 Sol.} Place: Morphisme précédant (3.1.1.8), membre de droite; Reported print: \textbackslash{}RR\textasciicircum{}+\textbackslash{}bar\{u\}\_\textbackslash{}ast\textbackslash{}circ\textbackslash{}mathbf\{L\}\textasciicircum{}+\textbackslash{}bar\{u\}\textasciicircum{}\textbackslash{}ast(F\textasciicircum{}\{\textbackslash{}bigdot\}); Source TeX: \textbackslash{}sisi\{\textbackslash{}RR\textasciicircum{}+\}\{R\} \textbackslash{}bar\{u\}\_\textbackslash{}ast \textbackslash{}circ \textbackslash{}bar\{u\}\textasciicircum{}\textbackslash{}ast(F\textasciicircum{}\{\textbackslash{}bigdot\}); Note: Le tirage porte également L\textasciicircum{}+ devant cette image inverse.; Proposed edit: \textbackslash{}circ \textbackslash{}bar\{u\}\textasciicircum{}\textbackslash{}ast(F\textasciicircum{}\{\textbackslash{}bigdot\}) =\allowbreak > \textbackslash{}circ \textbackslash{}mathbf\{L\}\textasciicircum{}\{+\} \textbackslash{}bar\{u\}\textasciicircum{}\textbackslash{}ast(F\textasciicircum{}\{\textbackslash{}bigdot\})\par
the two models differ at this place\par
\par\noindent\hypertarget{DP0303}{}\textbf{DP0303 -- Vbis, p. 127 (conflict)}\par
diagram (3.1.1.8), middle-right node | Diagramme (3.1.1.8), sommet médian supérieur à droite\par
\textbf{Claude Opus 5.5.} Place: diagram (3.1.1.8), middle-right node; Reported print: R\textasciicircum{}+ū'\_*\ensuremath{\circ}L\textasciicircum{}+ū'*(F\ensuremath{\cdot}); Source TeX: \textbackslash{}bar\{u'\}\_\{\textbackslash{}ast\}\textbackslash{}circ \textbackslash{}bar\{u'\}\textasciicircum{}\{\textbackslash{}ast\}(F\textasciicircum{}\{\textbackslash{}bigdot\}); Note: L\textasciicircum{}+ is missing in the diagram node.; Proposed edit: \textbackslash{}circ      \textbackslash{}bar\{u'\}\textasciicircum{}\{\textbackslash{}ast\}(F\textasciicircum{}\{\textbackslash{}bigdot\}) =\allowbreak > \textbackslash{}circ      \textbackslash{}LL\textasciicircum{}\{+\}\textbackslash{}bar\{u'\}\textasciicircum{}\{\textbackslash{}ast\}(F\textasciicircum{}\{\textbackslash{}bigdot\})\par
\textbf{GPT-6.1 Sol.} Place: Diagramme (3.1.1.8), sommet médian supérieur à droite; Reported print: \textbackslash{}RR\textasciicircum{}+\textbackslash{}bar\{u\}'\_\textbackslash{}ast\textbackslash{}circ\textbackslash{}mathbf\{L\}\textasciicircum{}+\textbackslash{}bar\{u'\}\textasciicircum{}\{\textbackslash{}ast\}(F\textasciicircum{}\{\textbackslash{}bigdot\}); Source TeX: \textbackslash{}sisi\{\textbackslash{}RR\textasciicircum{}\{+\}\}\{R\}\textbackslash{}bar\{u'\}\_\{\textbackslash{}ast\}\textbackslash{}circ      \textbackslash{}bar\{u'\}\textasciicircum{}\{\textbackslash{}ast\}(F\textasciicircum{}\{\textbackslash{}bigdot\}); Note: L\textasciicircum{}+ est imprimé dans ce sommet du diagramme.; Proposed edit: \textbackslash{}bar\{u'\}\textasciicircum{}\{\textbackslash{}ast\}(F\textasciicircum{}\{\textbackslash{}bigdot\}) =\allowbreak > \textbackslash{}mathbf\{L\}\textasciicircum{}\{+\} \textbackslash{}bar\{u'\}\textasciicircum{}\{\textbackslash{}ast\}(F\textasciicircum{}\{\textbackslash{}bigdot\})\par
the two models differ at this place\par
\par\noindent\hypertarget{DP0304}{}\textbf{DP0304 -- Vbis, p. 127 (conflict)}\par
diagram (3.1.1.8), bottom-right node | Diagramme (3.1.1.8), sommet inférieur à droite\par
\textbf{Claude Opus 5.5.} Place: diagram (3.1.1.8), bottom-right node; Reported print: R\textasciicircum{}+ū\_*\ensuremath{\circ}L\textasciicircum{}+ū *(F\ensuremath{\cdot}); Source TeX: \textbackslash{}bar\{u\}\_\{\textbackslash{}ast\}\textbackslash{}circ\textbackslash{}bar\{u\} \{\}\textasciicircum{}\{\textbackslash{}ast\}(F\textasciicircum{}\{\textbackslash{}bigdot\}); Note: L\textasciicircum{}+ is missing in the diagram node.; Proposed edit: \textbackslash{}circ\textbackslash{}bar\{u\} \{\}\textasciicircum{}\{\textbackslash{}ast\}(F\textasciicircum{}\{\textbackslash{}bigdot\}) =\allowbreak > \textbackslash{}circ\textbackslash{}LL\textasciicircum{}\{+\}\textbackslash{}bar\{u\} \{\}\textasciicircum{}\{\textbackslash{}ast\}(F\textasciicircum{}\{\textbackslash{}bigdot\})\par
\textbf{GPT-6.1 Sol.} Place: Diagramme (3.1.1.8), sommet inférieur à droite; Reported print: \textbackslash{}RR\textasciicircum{}+\textbackslash{}bar\{u\}\_\textbackslash{}ast\textbackslash{}circ\textbackslash{}mathbf\{L\}\textasciicircum{}+\textbackslash{}bar\{u\}\textasciicircum{}\{\textbackslash{}ast\}(F\textasciicircum{}\{\textbackslash{}bigdot\}); Source TeX: \textbackslash{}sisi\{\textbackslash{}RR\textasciicircum{}\{+\}\}\{R\}\textbackslash{}bar\{u\}\_\{\textbackslash{}ast\}\textbackslash{}circ\textbackslash{}bar\{u\} \{\}\textasciicircum{}\{\textbackslash{}ast\}(F\textasciicircum{}\{\textbackslash{}bigdot\}); Note: L\textasciicircum{}+ manque aussi dans ce sommet.; Proposed edit: \textbackslash{}circ\textbackslash{}bar\{u\} \{\}\textasciicircum{}\{\textbackslash{}ast\}(F\textasciicircum{}\{\textbackslash{}bigdot\}) =\allowbreak > \textbackslash{}circ\textbackslash{}mathbf\{L\}\textasciicircum{}\{+\} \textbackslash{}bar\{u\} \{\}\textasciicircum{}\{\textbackslash{}ast\}(F\textasciicircum{}\{\textbackslash{}bigdot\})\par
the two models differ at this place\par
\par\noindent\hypertarget{DP0305}{}\textbf{DP0305 -- Vbis, p. 128 (one\_model\_fix)}\par
Proposition 3.1.4, « on suppose que f*id\_\{L\_i\} »\par
\textbf{Claude Opus 5.5.} Place: Proposition 3.1.4, « on suppose que f*id\_\{L\_i\} »; Reported print: f*id\_\{L\_i\}; Source TeX: f\_\{\textbackslash{}ast\} \textbackslash{}id\_\{L\_\{i\}\}; Note: Horizontal composition f*id, as u*id in 3.1.6 (where the TeX writes u \textbackslash{}ast \textbackslash{}id); f\_\{\textbackslash{}ast\} prints a subscript star.; Proposed edit: que \$f\_\{\textbackslash{}ast\} =\allowbreak > que \$f \textbackslash{}ast\par
\textbf{GPT-6.1 Sol.} nothing reported here\par
a fix from one model only\par
\par\noindent\hypertarget{DP0306}{}\textbf{DP0306 -- Vbis, p. 128 (one\_model\_fix)}\par
Proposition 3.1.4, « (resp. f*id\_\{c\_i\} »\par
\textbf{Claude Opus 5.5.} Place: Proposition 3.1.4, « (resp. f*id\_\{c\_i\} »; Reported print: f*id\_\{c\_i\}; Source TeX: f\_\{\textbackslash{}ast\} \textbackslash{}id\_\{C\_i\}; Note: Same as the previous item.; Proposed edit: (\textbackslash{}resp \$f\_\{\textbackslash{}ast\} =\allowbreak > (\textbackslash{}resp \$f \textbackslash{}ast\par
\textbf{GPT-6.1 Sol.} nothing reported here\par
a fix from one model only\par
\par\noindent\hypertarget{DP0307}{}\textbf{DP0307 -- Vbis, p. 128 (conflict)}\par
augmentation plate \textbackslash{}overline\{XfX'\} \ensuremath{\to} C\_S\textasciicircum{}\{\textbackslash{}overline\{\ensuremath{\Gamma}\ensuremath{\times}D\}\} | But de l'augmentation t\par
\textbf{Claude Opus 5.5.} Place: augmentation plate \textbackslash{}overline\{XfX'\} \ensuremath{\to} C\_S\textasciicircum{}\{\textbackslash{}overline\{\ensuremath{\Gamma}\ensuremath{\times}D\}\}; Reported print: C\_S\textasciicircum{}\{\textbackslash{}overline\{\ensuremath{\Gamma}\ensuremath{\times}D\}\}; Source TeX: C\textasciicircum{}\{\textbackslash{}bar\{\textbackslash{}Gamma \textbackslash{}times D\}\}\_S; Proposed correction: C\textasciicircum{}\{\textbackslash{}bar\{I \textbackslash{}times D\}\}\_S; Note: The print has \ensuremath{\Gamma}\ensuremath{\times}D, presumably for the category \textbackslash{}overline\{I\ensuremath{\times}D\} just defined; the TeX keeps the printed reading.\par
\textbf{GPT-6.1 Sol.} Place: But de l'augmentation t; Reported print: C\textasciicircum{}\{\textbackslash{}bar\{I\textbackslash{}times D\}\}\_S; Source TeX: C\textasciicircum{}\{\textbackslash{}bar\{\textbackslash{}Gamma \textbackslash{}times D\}\}\_S; Note: Le premier facteur de l'exposant est I, reprenant la catégorie définie dans la phrase précédente, et non Gamma.; Proposed edit: C\textasciicircum{}\{\textbackslash{}bar\{\textbackslash{}Gamma \textbackslash{}times D\}\}\_S =\allowbreak > C\textasciicircum{}\{\textbackslash{}bar\{I \textbackslash{}times D\}\}\_S\par
the two models differ at this place: one model fixes the TeX, the other calls it a misprint of the print\par
\par\noindent\hypertarget{DP0308}{}\textbf{DP0308 -- Vbis, p. 129 (one\_model\_fix)}\par
Lemme 3.1.5.1, first line\par
\textbf{Claude Opus 5.5.} Place: Lemme 3.1.5.1, first line; Reported print: un morphisme plat, Y \ensuremath{\to}u C\_S\textasciicircum{}\ensuremath{\Delta}; Source TeX: un morphisme plat \$Y \textbackslash{}xrightarrow\{u\}; Note: The comma separating the two hypotheses is missing.; Proposed edit: un morphisme plat \$Y =\allowbreak > un morphisme plat, \$Y\par
\textbf{GPT-6.1 Sol.} nothing reported here\par
a fix from one model only\par
\par\noindent\hypertarget{DP0309}{}\textbf{DP0309 -- Vbis, p. 130 (conflict)}\par
(3.2.2), « soit F un objet de Mod(E\textasciicircum{}o\_S, O\_S) » | Paragraphe (3.2.2), catégorie à laquelle appartient F\par
\textbf{Claude Opus 5.5.} Place: (3.2.2), « soit F un objet de Mod(E\textasciicircum{}o\_S, O\_S) »; Reported print: Mod(E\textasciicircum{}o\_S, O\_S); Source TeX: \textbackslash{}Mod(\textbackslash{}mathcal\{E\}\textasciicircum{}0\_Y, \textbackslash{}mathcal\{O\}\_\{Y\}); Note: The print has index S, not Y (possibly an editorial correction, but not marked by \textbackslash{}sisi).; Proposed edit: \textbackslash{}Mod(\textbackslash{}mathcal\{E\}\textasciicircum{}0\_Y, \textbackslash{}mathcal\{O\}\_\{Y\}) =\allowbreak > \textbackslash{}Mod(\textbackslash{}mathcal\{E\}\textasciicircum{}0\_S, \textbackslash{}mathcal\{O\}\_\{S\})\par
\textbf{GPT-6.1 Sol.} Place: Paragraphe (3.2.2), catégorie à laquelle appartient F; Reported print: \textbackslash{}Mod(\textbackslash{}mathcal\{E\}\textasciicircum{}\{\textbackslash{}circ\}\_S,\textbackslash{}mathcal\{O\}\_S); Source TeX: \textbackslash{}Mod(\textbackslash{}mathcal\{E\}\textasciicircum{}0\_Y, \textbackslash{}mathcal\{O\}\_\{Y\}); Proposed correction: \textbackslash{}Mod(\textbackslash{}mathcal\{E\}\textasciicircum{}0\_Y, \textbackslash{}mathcal\{O\}\_\{Y\}); Note: Le tirage porte deux indices S; F doit être sur Y pour les foncteurs de la formule suivante. Le TeX corrige silencieusement les deux indices en Y.\par
the two models differ at this place: one model fixes the TeX, the other calls it a misprint of the print\par
\par\noindent\hypertarget{DP0310}{}\textbf{DP0310 -- Vbis, p. 130 (conflict)}\par
diagram (3.2.2.2), bottom node | Diagramme (3.2.2.2), étiquettes des cinq flèches | Diagramme (3.2.2.2), sommet inférieur\par
\textbf{Claude Opus 5.5.} Place: diagram (3.2.2.2), bottom node; Reported print: f'*(g'\_*(F)); Source TeX: f'\{\textbackslash{}ast\}(g'\_\{\textbackslash{}ast\}(F)).; Note: The superscript \textasciicircum{} is missing, so the star prints on the line; the print has no period after the node.; Proposed edit: \&f'\{\textbackslash{}ast\}(g'\_\{\textbackslash{}ast\}(F)). =\allowbreak > \&f'\textasciicircum{}\{\textbackslash{}ast\}(g'\_\{\textbackslash{}ast\}(F))\par
\textbf{GPT-6.1 Sol.} Place: Diagramme (3.2.2.2), étiquettes des cinq flèches; Reported print: Les cinq flèches sont sans étiquette.; Source TeX: Étiquettes f\textasciicircum{}\textbackslash{}circ, g\textasciicircum{}\textbackslash{}circ, g'\textasciicircum{}\textbackslash{}circ, \textbackslash{}Id et f'\textasciicircum{}\textbackslash{}circ.; Note: Ces étiquettes sont des ajouts éditoriaux au diagramme; elles ne figurent pas dans le tirage.; Proposed edit: f\textasciicircum{}\{\textbackslash{}ast\}(F) \textbackslash{}ar[r]\textasciicircum{}\{f\textasciicircum{}\textbackslash{}circ\} \textbackslash{}ar[ld]\_\{g\textasciicircum{}\textbackslash{}circ\}\& F\textbackslash{}ar[d]\textasciicircum{}\{g'\textasciicircum{}\textbackslash{}circ\} \textbackslash{}\textbackslash{} g\_\{\textbackslash{}ast\}(f\textasciicircum{}\{\textbackslash{}ast\}(F))  \textbackslash{}ar[rd]\_\{\textbackslash{}Id\}\&             \& g'\_\{\textbackslash{}ast\}(F)\textbackslash{}\textbackslash{}                       \&f'\{\textbackslash{}ast\}(g'\_\{\textbackslash{}ast\}(F)). \textbackslash{}ar[ru]\_\{f'\textasciicircum{}\textbackslash{}circ\} =\allowbreak > f\textasciicircum{}\{\textbackslash{}ast\}(F) \textbackslash{}ar[r] \textbackslash{}ar[ld]\& F\textbackslash{}ar[d] \textbackslash{}\textbackslash{} g\_\{\textbackslash{}ast\}(f\textasciicircum{}\{\textbackslash{}ast\}(F))  \textbackslash{}ar[rd]\&             \& g'\_\{\textbackslash{}ast\}(F)\textbackslash{}\textbackslash{}                       \&f'\{\textbackslash{}ast\}(g'\_\{\textbackslash{}ast\}(F)). \textbackslash{}ar[ru] || Place: Diagramme (3.2.2.2), sommet inférieur; Reported print: f'\textasciicircum{}\{\textbackslash{}ast\}(g'\_\{\textbackslash{}ast\}(F)); Source TeX: f'\{\textbackslash{}ast\}(g'\_\{\textbackslash{}ast\}(F)); Note: L'astérisque est en exposant dans le tirage; le caractère \textasciicircum{} manque dans le TeX.; Proposed edit: f'\{\textbackslash{}ast\}(g'\_\{\textbackslash{}ast\}(F)) =\allowbreak > f'\textasciicircum{}\{\textbackslash{}ast\}(g'\_\{\textbackslash{}ast\}(F))\par
the two models differ at this place\par
\par\noindent\hypertarget{DP0311}{}\textbf{DP0311 -- Vbis, p. 130 (doubt)}\par
diagram (3.2.2.2), arrow labels\par
\textbf{Claude Opus 5.5.} Place: diagram (3.2.2.2), arrow labels; Reported print: arrows without labels; Source TeX: \textbackslash{}ar[r]\textasciicircum{}\{f\textasciicircum{}\textbackslash{}circ\} ... \textbackslash{}ar[ld]\_\{g\textasciicircum{}\textbackslash{}circ\} ... \textbackslash{}ar[d]\textasciicircum{}\{g'\textasciicircum{}\textbackslash{}circ\} ... \textbackslash{}ar[rd]\_\{\textbackslash{}Id\} ... \textbackslash{}ar[ru]\_\{f'\textasciicircum{}\textbackslash{}circ\}; Note: The arrow labels are the editors' addition, announced in their N.D.E.; the print has none. Left as an editorial device.\par
\textbf{GPT-6.1 Sol.} nothing reported here\par
a model flags this place without giving a fix\par
\par\noindent\hypertarget{DP0312}{}\textbf{DP0312 -- Vbis, p. 131 (one\_model\_fix)}\par
Formula (3.2.2.3); recurring simplification of derived-functor notation\par
\textbf{Claude Opus 5.5.} nothing reported here\par
\textbf{GPT-6.1 Sol.} Place: Formula (3.2.2.3); recurring simplification of derived-functor notation; Reported print: \textbackslash{}xi\_G : \textbackslash{}mathbf\{L\}\textasciicircum{}\{+\}\_G g\_\{\textbackslash{}ast\}\textbackslash{}circ\textbackslash{}RR\textasciicircum{}\{+\}\_G f\textasciicircum{}\{\textbackslash{}ast\}\textbackslash{}longrightarrow\textbackslash{}RR\textasciicircum{}\{+\}\_G \{f'\}\textasciicircum{}\{\textbackslash{}ast\}\textbackslash{}circ\textbackslash{}mathbf\{L\}\textasciicircum{}\{+\}\_G \{g'\}\_\{\textbackslash{}ast\}; Source TeX: \textbackslash{}xi\_G : g\_\{G\textbackslash{}ast\}\textbackslash{}circ\textbackslash{}RR\textasciicircum{}\{+\}\_G f\textasciicircum{}\{\textbackslash{}ast\}\textbackslash{}longrightarrow\textbackslash{}RR\textasciicircum{}\{+\}\_G \{f'\}\textasciicircum{}\{\textbackslash{}ast\}\textbackslash{}circ \{g'\}\_\{G\textbackslash{}ast\}; Note: Recurring simplification, reported once with a representative fix. Three examples: (3.2.2.3) replaces printed L\_G\textasciicircum{}+g\_* by g\_\{G*\}; the formula for eta\_G\textasciicircum{}f replaces L\textasciicircum{}+\textbackslash{}bar u'\textasciicircum{}* by \textbackslash{}bar u'\textasciicircum{}*; diagram (3.2.3.2), p.132, replaces R\_G\textasciicircum{}+\textbackslash{}bar u'\_* by R\textasciicircum{}+\textbackslash{}bar u'\_*. Restore the omitted L\textasciicircum{}+ throughout these formulas on pp.131–132, and the G subscripts on both R\textasciicircum{}+ in (3.2.3.2). The first arguments of the existing \textbackslash{}sisi commands otherwise agree with the printed R\textasciicircum{}+.; Proposed edit: g\_\{G\textbackslash{}ast\} \textbackslash{}circ  \textbackslash{}sisi\{\textbackslash{}RR\textasciicircum{}+\}\{R\}\_G f\textasciicircum{}\{\textbackslash{}ast\} \textbackslash{}longrightarrow  \textbackslash{}sisi\{\textbackslash{}RR\textasciicircum{}+\}\{R\}\_G \{f'\}\textasciicircum{}\{\textbackslash{}ast\} \textbackslash{}circ          \{g'\}\_\{G\textbackslash{}ast\} =\allowbreak > \textbackslash{}mathbf\{L\}\textasciicircum{}\{+\}\_G g\_\{\textbackslash{}ast\} \textbackslash{}circ  \textbackslash{}sisi\{\textbackslash{}RR\textasciicircum{}+\}\{R\}\_G f\textasciicircum{}\{\textbackslash{}ast\} \textbackslash{}longrightarrow  \textbackslash{}sisi\{\textbackslash{}RR\textasciicircum{}+\}\{R\}\_G \{f'\}\textasciicircum{}\{\textbackslash{}ast\} \textbackslash{}circ          \textbackslash{}mathbf\{L\}\textasciicircum{}\{+\}\_G \{g'\}\_\{\textbackslash{}ast\}\par
a fix from one model only\par
\par\noindent\hypertarget{DP0313}{}\textbf{DP0313 -- Vbis, p. 132 (conflict)}\par
diagram (3.2.3.2), right ascending arrow | Diagram (3.2.3.2), ascending right-hand diagonal\par
\textbf{Claude Opus 5.5.} Place: diagram (3.2.3.2), right ascending arrow; Reported print: label \ensuremath{\alpha}(L\textasciicircum{}+\_G h\_*(K\ensuremath{\cdot})) on the right arrow; Source TeX: \textbackslash{}ar[ru]\&.; Note: The label of the right arrow is missing in the TeX. The new text follows the edition's convention of dropping L\textasciicircum{}+\_G; in the print's notation it reads \ensuremath{\alpha}(L\textasciicircum{}+\_G h\_*(K\ensuremath{\cdot})).; Proposed edit: \textbackslash{}ar[ru]\&. =\allowbreak > \textbackslash{}ar[ru]\_\{\textbackslash{}alpha(h\_\{c\textbackslash{}ast\}(K\textasciicircum{}\{\textbackslash{}bigdot\}))\}\&.\par
\textbf{GPT-6.1 Sol.} Place: Diagram (3.2.3.2), ascending right-hand diagonal; Reported print: \textbackslash{}alpha(\textbackslash{}mathbf\{L\}\textasciicircum{}\{+\}\_G h\_\{\textbackslash{}ast\}(K\textasciicircum{}\{\textbackslash{}bigdot\})); Source TeX: The diagonal arrow has no label.; Note: The printed right-hand diagonal is labelled alpha(L\_G\textasciicircum{}+h\_*(K)); this label is entirely absent from the TeX.; Proposed edit: \textbackslash{}ar[ru]\&. =\allowbreak > \textbackslash{}ar[ru]\_\{\textbackslash{}alpha(\textbackslash{}mathbf\{L\}\textasciicircum{}\{+\}\_G h\_\{\textbackslash{}ast\}(K\textasciicircum{}\{\textbackslash{}bigdot\}))\}\&.\par
the two models differ at this place\par
\par\noindent\hypertarget{DP0314}{}\textbf{DP0314 -- Vbis, p. 133 (one\_model\_fix)}\par
Functor applied to t, immediately after the opening factorization\par
\textbf{Claude Opus 5.5.} nothing reported here\par
\textbf{GPT-6.1 Sol.} Place: Functor applied to t, immediately after the opening factorization; Reported print: \textbackslash{}epsilon\_\{Z\textasciicircum{}\{\textbackslash{}bigdot\}\textbackslash{}ast\}; Source TeX: \textbackslash{}epsilon\_\{Z\textasciicircum{}\{\textbackslash{}bigdot\textbackslash{}ast\}\}; Note: The direct-image asterisk follows Z\textasciicircum{}\{\textbackslash{}bigdot\} in epsilon's subscript. The TeX instead puts it in the exponent of Z alongside the complex dot.; Proposed edit: \textbackslash{}epsilon\_\{Z\textasciicircum{}\{\textbackslash{}bigdot\textbackslash{}ast\}\} =\allowbreak > \textbackslash{}epsilon\_\{Z\textasciicircum{}\{\textbackslash{}bigdot\}\textbackslash{}ast\}\par
a fix from one model only\par
\par\noindent\hypertarget{DP0315}{}\textbf{DP0315 -- Vbis, p. 137 (conflict)}\par
end of the proof of Lemme 3.3.3.3, last line of the page | Fin de page, après les deux limites projectives\par
\textbf{Claude Opus 5.5.} Place: end of the proof of Lemme 3.3.3.3, last line of the page; Reported print: (resp. lim\_\{\ensuremath{\Delta}\textasciicircum{}+\_\{n+1[p]\}\} X'\_q) . Or pour toute flèche ι : i \ensuremath{\longrightarrow} j de \ensuremath{\Delta}\textasciicircum{}+\_\{n+1[p]\} on a un diagramme commutatif à lignes exactes :; Source TeX: (\textbackslash{}resp \$\textbackslash{}varprojlim\_\{\textbackslash{}substack\{\textbackslash{}Delta\textasciicircum{}+\textbackslash{}\textbackslash{} n+1[p]\}\} \{X'\}\_q\$).; Note: The edition rewrote the proof of 3.3.3.3. The printed sentence survives only commented out in the TeX of p. 138, so it is not printed.; Proposed edit: \{X'\}\_q\$). =\allowbreak > \{X'\}\_q\$). Or pour toute flèche \$\textbackslash{}iota : i \textbackslash{}to j\$ de \$\textbackslash{}Delta\textasciicircum{}+\_\{n+1[p]\}\$ on a un diagramme commutatif à lignes exactes:\par
\textbf{GPT-6.1 Sol.} Place: Fin de page, après les deux limites projectives; Reported print: Or pour toute flèche \$\textbackslash{}iota : i \textbackslash{}to j\$ de \$\textbackslash{}Delta\textasciicircum{}+\_\{n+1[p]\}\$ on a un diagramme commutatif à lignes exactes :; Source TeX: Le texte actif s'arrête après la seconde limite projective.; Note: Cette phrase termine bien la page imprimée. Elle ne subsiste que dans les commentaires du fragment p.138 : elle manque au texte actif.; Proposed edit: \{X'\}\_q\$). =\allowbreak > \{X'\}\_q\$). Or pour toute flèche \$\textbackslash{}iota : i \textbackslash{}to j\$ de \$\textbackslash{}Delta\textasciicircum{}+\_\{n+1[p]\}\$ on a un diagramme commutatif à lignes exactes :\par
the two models differ at this place\par
\par\noindent\hypertarget{DP0316}{}\textbf{DP0316 -- Vbis, p. 138 (conflict)}\par
top of the page: rewritten proof 'On peut supposer p>n+1 ... revient au même)' | Début de page, preuve du lemme 3.3.3.3\par
\textbf{Claude Opus 5.5.} Place: top of the page: rewritten proof 'On peut supposer p>n+1 ... revient au même)'; Reported print: (print has the diagram and 'dont le carré de gauche est cartésien ...', no such text); Source TeX: On peut supposer \$p>n+1\$. On écrit alors \$\$\textbackslash{}cosq\_\{n+1\}(X)\_p=\allowbreak ... revient au même)\textbackslash{}nde\{La preuve du lemme ... a été légèrement réécrite\} .; Note: The whole uncommented block from 'On peut supposer \$p>n+1\$.' through 'revient au même)' before the \textbackslash{}nde is an editorial rewrite that the print does not have. It should be deleted up to 'revient au même)'; the \textbackslash{}nde itself is not reported. The fix shown removes only the first words: the coordinator must delete the rest of the block.; Proposed edit: On peut supposer \$p>n+1\$. On écrit alors =\allowbreak > \par
\textbf{GPT-6.1 Sol.} Place: Début de page, preuve du lemme 3.3.3.3; Reported print: Diagramme reliant K'\_\textbackslash{}iota, K\_\textbackslash{}iota, les deux produits, X'\_i et X\_i, avec branches par X'\_j et X\_j ; puis « dont le carré de gauche est cartésien (ou bien i =\allowbreak  n+1 =\allowbreak  j et \textbackslash{}iota =\allowbreak  id, ou bien i < n+1 et f\_i est un isomorphisme) ».; Source TeX: Nouvelle preuve allant de « On peut supposer \$p>n+1\$ » à « revient au même) », utilisant \textbackslash{}Pi X, \textbackslash{}tilde X et deux carrés de noyaux.; Note: Il s'agit d'une réécriture de contenu, non d'un simple changement de présentation. Le remplacement rétablit le diagramme et son commentaire imprimés. Les deux indices fautifs de l'imprimé sont conservés dans les premiers arguments des \textbackslash{}sisi.; Proposed edit: On peut supposer \$p>n+1\$. On écrit alors \$\$\textbackslash{}cosq\_\{n+1\}(X)\_p=\allowbreak \textbackslash{}varprojlim\_\{\textbackslash{}substack\{[q]\textbackslash{}rightarrow[p]\textbackslash{}\textbackslash{} q\textbackslash{}leq n+1\}\}X\_q\$\$ comme le noyau de la double flèche  \$\$\textbackslash{}Pi X=\allowbreak \textbackslash{}prod\_\{\textbackslash{}substack\{[q]\textbackslash{}rightarrow[p]\textbackslash{}\textbackslash{} q\textbackslash{}leq n+1\}\}X\_q\textbackslash{}rightrightarrows \textbackslash{}prod\_\{\textbackslash{}substack\{[i]\textbackslash{}stackrel\{\textbackslash{}alpha\}\{\textbackslash{}hookrightarrow\}[j]\textbackslash{}\textbackslash{} \textbackslash{}searrow\textbackslash{}swarrow\textbackslash{}\textbackslash{} [p] \textbackslash{}\textbackslash{}j\textbackslash{}leq n+1\}\}X\_i\$\$ où la composante \$\textbackslash{}alpha\_X\$ d'indice \$\textbackslash{}alpha\textbackslash{}in\textbackslash{}Hom\_\{[p]\}([i],[j])\$ de la double flèche est la double flèche formée d'une part du morphisme \$\$\textbackslash{}Pi X\textbackslash{}rightarrow X\_i\$\$ de projection d'indice \$[i]\textbackslash{}rightarrow[p]\$ et, d'autre part, du morphisme \$\$\textbackslash{}Pi X\textbackslash{}rightarrow X\_j\textbackslash{}rightarrow X\_i,\$\$ composé de la projection d'indice \$[j]\textbackslash{}rightarrow[p]\$ et de \$\textbackslash{}alpha\textbackslash{}in\textbackslash{}Hom(X\_j,X\_i)\$.  On observe alors que le diagramme \$\$\textbackslash{}xymatrix\{\textbackslash{}cosq\_\{n+1\}(X)\_p\textbackslash{}ar[r]\textbackslash{}ar[d]\&\textbackslash{}Pi X\textbackslash{}ar[d]\textbackslash{}\textbackslash{} \textbackslash{}cosq\_\{n+1\}(\{\textbackslash{}tilde X\})\_p\textbackslash{}ar[r]\&\textbackslash{}Pi \{\textbackslash{}tilde X\}\}\$\$ est cartésien.  En effet, d'après la description précédente, on doit prouver que pour toute \textbackslash{}textit\{injection croissante\} \$\textbackslash{}alpha\textbackslash{}in\textbackslash{}Hom\_\{[p]\}([i],[j])\$, le diagramme \$\$\textbackslash{}xymatrix\{\textbackslash{}Ker(\textbackslash{}alpha\_X:\textbackslash{} \textbackslash{}Pi X\textbackslash{}rightrightarrows X\_i)\textbackslash{}ar[r]\textbackslash{}ar[d]\&\textbackslash{}Pi X\textbackslash{}ar[d]\textbackslash{}\textbackslash{} \textbackslash{}Ker(\textbackslash{}alpha\_\{\{\textbackslash{}tilde X\}\}:\textbackslash{} \textbackslash{}Pi \{\textbackslash{}tilde X\}\textbackslash{}rightrightarrows \{\textbackslash{}tilde X\}\_i)\textbackslash{}ar[r]\&\textbackslash{}Pi \{\textbackslash{}tilde X\}\}\$\$ est cartésien. Mais si \$i=\allowbreak n+1\$, alors \$\textbackslash{}alpha: [i]\textbackslash{}hookrightarrow [j] \$ est l'identité de \$[n+1]\$ : les doubles flèches n'imposent alors pas de condition. \% de sorte que les doubles flèches \$\textbackslash{}alpha\_\{X\}\$ et \%\$\textbackslash{}alpha\_\{\{\textbackslash{}tilde X\}\}\$ sont les couples formés des projections sur \$X\_i\$ et \%\$\{\textbackslash{}tilde X\}\_i\$ prouvant que Les noyaux respectifs sont \$\textbackslash{}Pi X\$ et \$\textbackslash{}Pi \{\textbackslash{}tilde X\}\$, ce qui traite ce cas. Sinon, on a \$i<n+1\$ et \$X\_i\textbackslash{}rightarrow \{\textbackslash{}tilde X\}\_i\$ est un isomorphisme, de sorte que le diagramme est cartésien aussi dans ce cas (tester une égalité dans \$X\_i\$ ou \$\{\textbackslash{}tilde X\}\_i\$ revient au même) =\allowbreak > \textbackslash{}[ \textbackslash{}xymatrix@R=\allowbreak .5cm@C=\allowbreak 1.5cm\{ K'\_\{\textbackslash{}iota\} \textbackslash{}ar[r]\textbackslash{}ar[ddd]\_\{\textbackslash{}alpha\_\{\textbackslash{}iota\}\} \& \textbackslash{}prod\_\{r\textbackslash{}in\textbackslash{}operatorname\{ob\}(\textbackslash{}Delta\textasciicircum{}+\_\{n+1[p]\})\} \textbackslash{}sisi\{X'\_2\}\{X'\_r\} \textbackslash{}ar[rr]\textbackslash{}ar[ddd]\textbackslash{}ar[dr] \&\& X'\_i \textbackslash{}ar[ddd]\textasciicircum{}\{f\_i\}\textbackslash{}\textbackslash{} \&\& X'\_j \textbackslash{}ar[ur] \&\textbackslash{}\textbackslash{} \&\& X\_j \textbackslash{}ar[dr] \&\textbackslash{}\textbackslash{} K\_\{\textbackslash{}iota\} \textbackslash{}ar[r] \& \textbackslash{}prod\_\{\textbackslash{}operatorname\{ob\}(\textbackslash{}Delta\textasciicircum{}+\_\{n+1[p]\})\} \textbackslash{}sisi\{X\_2\}\{X\_r\} \textbackslash{}ar[rr]\textbackslash{}ar[ur] \&\& X\_i \} \textbackslash{}] dont le carré de gauche est cartésien (ou bien \$i =\allowbreak  n + 1 =\allowbreak j\$ et \$\textbackslash{}iota =\allowbreak  \textbackslash{}mathrm\{id\}\$, ou bien \$i < n+1\$ et \$f\_i\$ est un isomorphisme)\par
the two models differ at this place\par
\par\noindent\hypertarget{DP0317}{}\textbf{DP0317 -- Vbis, p. 138 (one\_model\_fix)}\par
diagram and the sentence after it, before 'Utilisant alors (3.3.1)'\par
\textbf{Claude Opus 5.5.} Place: diagram and the sentence after it, before 'Utilisant alors (3.3.1)'; Reported print: diagram K'\_ι \ensuremath{\to} \ensuremath{\Pi}\_\{r\ensuremath{\in}ob(\ensuremath{\Delta}\textasciicircum{}+\_\{n+1[p]\})\} X'\_2 \ensuremath{\rightrightarrows} X'\_i (second arrow through X'\_j), \ensuremath{\alpha}\_ι and f\_i vertical, bottom row K\_ι \ensuremath{\to} \ensuremath{\Pi} X\_2 \ensuremath{\rightrightarrows} X\_i (through X\_j); then 'dont le carré de gauche est cartésien (ou bien i =\allowbreak  n+1 =\allowbreak  j et ι =\allowbreak  id , ou bien i < n+1 et f\_i est un isomorphisme).'; Source TeX: (only commented out); Note: The printed diagram (the index 2 on X stands so in the print, probably for r) and the sentence 'dont le carré ...' are only in the comment. The sketched xymatrix may be replaced by the commented one, uncommented and with a\_ι corrected to \ensuremath{\alpha}\_ι.; Proposed edit: Utilisant alors (\textbackslash{}hyperref[Vbis.prop3.3.1]\{3.3.1\}) on voit que \$\textbackslash{}alpha\_\textbackslash{}iota\$ est un morphisme =\allowbreak > \textbackslash{}[ \textbackslash{}xymatrix@C=\allowbreak 1.5cm\{K'\_\{\textbackslash{}iota\} \textbackslash{}ar[r]\textbackslash{}ar[d]\_\{\textbackslash{}alpha\_\{\textbackslash{}iota\}\} \& \textbackslash{}displaystyle\textbackslash{}prod\_\{r\textbackslash{}in\textbackslash{}ob(\textbackslash{}Delta\textasciicircum{}+\_\{n+1[p]\})\} X'\_2 \textbackslash{}ar[r]\textbackslash{}ar[d]\textbackslash{}ar[rd] \& X'\_i \textbackslash{}ar[d]\textasciicircum{}\{f\_i\}\textbackslash{}\textbackslash{} K\_\{\textbackslash{}iota\} \textbackslash{}ar[r] \& \textbackslash{}displaystyle\textbackslash{}prod\_\{\textbackslash{}ob(\textbackslash{}Delta\textasciicircum{}+\_\{n+1[p]\})\} X\_2 \textbackslash{}ar[r] \& X\_i\} \textbackslash{}] (les secondes flèches des lignes passant par \$X'\_j\$, \textbackslash{}resp \$X\_j\$) dont le carré de gauche est cartésien (ou bien \$i =\allowbreak  n+1 =\allowbreak  j\$ et \$\textbackslash{}iota =\allowbreak  \textbackslash{}id\$, ou bien \$i < n+1\$ et \$f\_i\$ est un isomorphisme). Utilisant alors (\textbackslash{}hyperref[Vbis.prop3.3.1]\{3.3.1\}) on voit que \$\textbackslash{}alpha\_\textbackslash{}iota\$ est un morphisme\par
\textbf{GPT-6.1 Sol.} nothing reported here\par
a fix from one model only\par
\par\noindent\hypertarget{DP0318}{}\textbf{DP0318 -- Vbis, p. 138 (conflict)}\par
after 'Utilisant alors (3.3.1) ... universelle.' | Après « un morphisme de G-1-descente cohomologique universelle »\par
\textbf{Claude Opus 5.5.} Place: after 'Utilisant alors (3.3.1) ... universelle.'; Reported print: On achève la démonstration en remarquant que f\_p s'identifie au morphisme canonique \ensuremath{\cap}\_ι K'\_ι —(\ensuremath{\times} \ensuremath{\alpha}\_ι)\ensuremath{\to} \ensuremath{\cap}\_ι K\_ι .; Source TeX: cohomologique universelle.Soit \$[X'/X]\textasciicircum{}p\$; Note: The final sentence of the printed proof is dropped; it survives only commented out.; Proposed edit: universelle.Soit \$[X'/X]\textasciicircum{}p\$ =\allowbreak > universelle. On achève la démonstration en remarquant que \$f\_p\$ s'identifie au morphisme canonique \$\textbackslash{}bigcap\_\{\textbackslash{}iota\} K'\_\textbackslash{}iota \textbackslash{}xrightarrow\{\textbackslash{}times \textbackslash{}alpha\_\textbackslash{}iota\} \textbackslash{}bigcap\_\{\textbackslash{}iota\} K\_\textbackslash{}iota\$.  Soit \$[X'/X]\textasciicircum{}p\$\par
\textbf{GPT-6.1 Sol.} Place: Après « un morphisme de G-1-descente cohomologique universelle »; Reported print: On achève la démonstration en remarquant que \$f\_p\$ s'identifie au morphisme canonique \$\textbackslash{}displaystyle\textbackslash{}mathop\{\textbackslash{}cap\}\_\{\textbackslash{}iota\} K'\_\textbackslash{}iota \textbackslash{}xrightarrow\{X \textbackslash{}alpha\_\textbackslash{}iota\} \textbackslash{}displaystyle\textbackslash{}mathop\{\textbackslash{}cap\}\_\{\textbackslash{}iota\} K\_\textbackslash{}iota\$.; Source TeX: universelle.Soit \$[X'/X]\textasciicircum{}p\$; Note: La dernière phrase de la preuve du lemme 3.3.3.3 manque entièrement au texte actif.; Proposed edit: universelle.Soit \$[X'/X]\textasciicircum{}p\$ =\allowbreak > universelle. On achève la démonstration en remarquant que \$f\_p\$ s'identifie au morphisme canonique \$\textbackslash{}displaystyle\textbackslash{}mathop\{\textbackslash{}cap\}\_\{\textbackslash{}iota\} K'\_\textbackslash{}iota \textbackslash{}xrightarrow\{X \textbackslash{}alpha\_\textbackslash{}iota\} \textbackslash{}displaystyle\textbackslash{}mathop\{\textbackslash{}cap\}\_\{\textbackslash{}iota\} K\_\textbackslash{}iota\$.  Soit \$[X'/X]\textasciicircum{}p\$\par
the two models differ at this place\par
\par\noindent\hypertarget{DP0319}{}\textbf{DP0319 -- Vbis, p. 141 (cut)}\par
Fin de la coupe : titre du lemme 4.1.3\par
\textbf{Claude Opus 5.5.} nothing reported here\par
\textbf{GPT-6.1 Sol.} Place: Fin de la coupe : titre du lemme 4.1.3; Reported print: La page finit par « résulte du lemme suivant. » Le titre « Lemme (4.1.3) » ouvre la page 142.; Source TeX: \textbackslash{}begin\{enonce*\}\{Lemme 4.1.3\}\textbackslash{}etiquette[4.1.3]\{Vbis.lem4.1.3\}; Note: La ligne 2882 appartient à la page 142, sans que le commentaire de coupe le signale. Problème de coupe seulement : ne pas déplacer le texte dans la source.\par
a model reports that the automatic cut gives lines of this page to its neighbour, or the neighbour's lines to this page: check the lines it names against the scan of the page that prints them (that page's checker did not have their TeX)\par
\par\noindent\hypertarget{DP0320}{}\textbf{DP0320 -- Vbis, p. 143 (conflict)}\par
display in the proof of 4.1.3, last '=\allowbreak ' of the first line | Preuve de 4.1.3, égalité donnant H\textasciicircum{}i \ensuremath{\Gamma}(K, F*)\par
\textbf{Claude Opus 5.5.} Place: display in the proof of 4.1.3, last '=\allowbreak ' of the first line; Reported print: =\allowbreak  H\textasciicircum{}i \ensuremath{\Gamma}(K,F*) with '4.1.3(i=\allowbreak 0)' written under the '=\allowbreak '; Source TeX: \textbackslash{}stackrel\{\textbackslash{}ndemark\}\{=\allowbreak \} H\textasciicircum{}i \textbackslash{}Gamma (K, F\textasciicircum{}\{\textbackslash{}ast\}); Note: The print labels this equality '4.1.3(i=\allowbreak 0)', just as the next one is labelled (4.1.5). The TeX drops the label and keeps only an editor's note mark.; Proposed edit: \textbackslash{}stackrel\{\textbackslash{}ndemark\}\{=\allowbreak \} H\textasciicircum{}i \textbackslash{}Gamma (K =\allowbreak > \textbackslash{}stackrel\{\textbackslash{}text\{4.1.3(\$i=\allowbreak 0\$)\}\textbackslash{}ndemark\}\{=\allowbreak \} H\textasciicircum{}i \textbackslash{}Gamma (K\par
\textbf{GPT-6.1 Sol.} Place: Preuve de 4.1.3, égalité donnant H\textasciicircum{}i \ensuremath{\Gamma}(K, F*); Reported print: L'égalité porte la justification « 4.1.3 (i=\allowbreak 0) ».; Source TeX: \textbackslash{}stackrel\{\textbackslash{}ndemark\}\{=\allowbreak \} H\textasciicircum{}i \textbackslash{}Gamma (K, F\textasciicircum{}\{\textbackslash{}ast\}); Note: La justification imprimée manque dans la formule. La marque de note éditoriale est conservée par l'édition proposée.; Proposed edit: \textbackslash{}stackrel\{\textbackslash{}ndemark\}\{=\allowbreak \} H\textasciicircum{}i \textbackslash{}Gamma (K, F\textasciicircum{}\{\textbackslash{}ast\}) =\allowbreak > \textbackslash{}underset\{\textbackslash{}text\{4.1.3 \}(i=\allowbreak 0)\}\{\textbackslash{}stackrel\{\textbackslash{}ndemark\}\{=\allowbreak \}\} H\textasciicircum{}i \textbackslash{}Gamma (K, F\textasciicircum{}\{\textbackslash{}ast\})\par
the two models differ at this place\par
\par\noindent\hypertarget{DP0321}{}\textbf{DP0321 -- Vbis, p. 155 (one\_model\_fix)}\par
proof of 5.2.3, reference\par
\textbf{Claude Opus 5.5.} Place: proof of 5.2.3, reference; Reported print: [[6]. Th. 1']; Source TeX: [\textbackslash{}cite\{SGA4-05bis:6\}, Th, \$1'\$]; Note: The abbreviation is « Th. » in the print, « Th, » in the TeX.; Proposed edit: Th, \$1'\$ =\allowbreak > Th. \$1'\$\par
\textbf{GPT-6.1 Sol.} nothing reported here\par
a fix from one model only\par
\par\noindent\hypertarget{DP0322}{}\textbf{DP0322 -- Vbis, p. 157 (one\_model\_fix)}\par
proof of 5.2.8, morphisme fonctoriel\par
\textbf{Claude Opus 5.5.} Place: proof of 5.2.8, morphisme fonctoriel; Reported print: ℍ'\_DR\ensuremath{\circ}\ensuremath{\bar{\pi}} \ensuremath{\longrightarrow} H'\_DR (target in ordinary H); Source TeX: \textbackslash{}to \textbackslash{}HH\textasciicircum{}\{\textbackslash{}bigdot\}\_\{DR\}\$, de sorte; Note: The target of the arrow is printed in ordinary H, not double-struck; it matches the top-left vertex of the diagram.; Proposed edit: \textbackslash{}to \textbackslash{}HH\textasciicircum{}\{\textbackslash{}bigdot\}\_\{DR\}\$, de sorte =\allowbreak > \textbackslash{}to H\textasciicircum{}\{\textbackslash{}bigdot\}\_\{DR\}\$, de sorte\par
\textbf{GPT-6.1 Sol.} nothing reported here\par
a fix from one model only\par
\par\noindent\hypertarget{DP0323}{}\textbf{DP0323 -- Vbis, p. 157 (one\_model\_fix)}\par
diagram, top-right vertex\par
\textbf{Claude Opus 5.5.} Place: diagram, top-right vertex; Reported print: ℍ'\_DR \ensuremath{\circ} \ensuremath{\bar{\pi}} (double-struck H); Source TeX: \& H\textasciicircum{}\{\textbackslash{}bigdot\}\_\{DR\} \textbackslash{}circ \textbackslash{}bar\{\textbackslash{}pi\} \textbackslash{}ar[l]; Note: The print uses double-struck ℍ here, as in the text.; Proposed edit: \& H\textasciicircum{}\{\textbackslash{}bigdot\}\_\{DR\} \textbackslash{}circ \textbackslash{}bar\{\textbackslash{}pi\} \textbackslash{}ar[l] =\allowbreak > \& \textbackslash{}HH\textasciicircum{}\{\textbackslash{}bigdot\}\_\{DR\} \textbackslash{}circ \textbackslash{}bar\{\textbackslash{}pi\} \textbackslash{}ar[l]\par
\textbf{GPT-6.1 Sol.} nothing reported here\par
a fix from one model only\par
\par\noindent\hypertarget{DP0324}{}\textbf{DP0324 -- Vbis, p. 162 (conflict)}\par
[6] Grothendieck | Bibliography [6], publication information\par
\textbf{Claude Opus 5.5.} Place: [6] Grothendieck; Reported print: Publications de l'I.H.E.S. n° 29.; Source TeX: Publications de 1' I.H.E.S \$\{\textbackslash{}rm n\textasciicircum{}\{\textbackslash{}circ\}\}\$ 29.; Note: The TeX has the digit 1 where the print has the letter l of « l' ». The TeX also lacks the final period of « I.H.E.S. ».; Proposed edit: Publications de 1' I.H.E.S =\allowbreak > Publications de l'I.H.E.S.\par
\textbf{GPT-6.1 Sol.} Place: Bibliography [6], publication information; Reported print: Publications de l'I.H.E.S.; Source TeX: Publications de 1' I.H.E.S; Note: The French article l' has been transcribed with the digit 1.; Proposed edit: de 1' I.H.E.S =\allowbreak > de l' I.H.E.S\par
the two models differ at this place\par
\par\noindent\hypertarget{DP0325}{}\textbf{DP0325 -- VI, p. 165 (conflict)}\par
§0, 4th paragraph, first line | Alinéa « Il est sans doute utile… », accord avec « notions de finitude »\par
\textbf{Claude Opus 5.5.} Place: §0, 4th paragraph, first line; Reported print: les notions de finitude utilisés ici; Source TeX: les notions de finitude utilisés ici; Proposed correction: utilisées; Note: The agreement with 'notions' requires 'utilisées'. The TeX keeps the printed reading.\par
\textbf{GPT-6.1 Sol.} Place: Alinéa « Il est sans doute utile… », accord avec « notions de finitude »; Reported print: notions de finitude utilisées ici; Source TeX: notions de finitude utilisés ici; Note: L’imprimé porte le féminin pluriel « utilisées ».; Proposed edit: utilisés ici =\allowbreak > utilisées ici\par
the two models differ at this place: one model fixes the TeX, the other calls it a misprint of the print\par
\par\noindent\hypertarget{DP0326}{}\textbf{DP0326 -- VI, p. 165 (cut)}\par
Début du fragment, avant « 0. Introduction »\par
\textbf{Claude Opus 5.5.} nothing reported here\par
\textbf{GPT-6.1 Sol.} Place: Début du fragment, avant « 0. Introduction »; Reported print: La page commence par « 0. Introduction », suivi de « Dans le présent exposé… ».; Source TeX: \textbackslash{}chapter\{Conditions de finitude. Topos et sites fibrés. Applications aux questions de passage à la limite.\}; Note: Le titre d’exposé et les noms d’auteurs précédant \textbackslash{}section\{Introduction\} ne figurent pas sur cette page. Ce débordement initial n’est pas annoncé dans le commentaire de découpage. La page 165 commence à \textbackslash{}section\{Introduction\}; rectifier le découpage sans déplacer le texte.\par
a model reports that the automatic cut gives lines of this page to its neighbour, or the neighbour's lines to this page: check the lines it names against the scan of the page that prints them (that page's checker did not have their TeX)\par
\par\noindent\hypertarget{DP0327}{}\textbf{DP0327 -- VI, p. 166 (one\_model\_fix)}\par
'(VII 3.3 et 5.7)'\par
\textbf{Claude Opus 5.5.} Place: '(VII 3.3 et 5.7)'; Reported print: (VII 3.3 et 5.7); Source TeX: \textbackslash{}hyperref[VI.prop5.7]\{5.7\}; Note: The print means VII 5.7, but the label points to VI 5.7, a wrong link target. The printed text is the same. Replace with plain text, or with the correct VII label if it exists.; Proposed edit: \textbackslash{}hyperref[VI.prop5.7]\{5.7\} =\allowbreak > 5.7\par
\textbf{GPT-6.1 Sol.} nothing reported here\par
a fix from one model only\par
\par\noindent\hypertarget{DP0328}{}\textbf{DP0328 -- VI, p. 168 (one\_model\_fix)}\par
'donc (II 4.4) la famille des X\_ij \ensuremath{\to} X est couvrante'\par
\textbf{Claude Opus 5.5.} Place: 'donc (II 4.4) la famille des X\_ij \ensuremath{\to} X est couvrante'; Reported print: (II 4.4); Source TeX: (II \textbackslash{}hyperref[VI.cor4.4]\{4.4\}); Note: The print cites II 4.4; the label VI.cor4.4 links into Exposé VI. Wrong link target, same printed text.; Proposed edit: (II \textbackslash{}hyperref[VI.cor4.4]\{4.4\}) =\allowbreak > (II 4.4)\par
\textbf{GPT-6.1 Sol.} nothing reported here\par
a fix from one model only\par
\par\noindent\hypertarget{DP0329}{}\textbf{DP0329 -- VI, p. 168 (conflict)}\par
'il résulte de nouveau de loc. cit. que les morphismes correspondants' | Fin de la preuve de la proposition 1.2, après l'extraction d'une sous-famille finie\par
\textbf{Claude Opus 5.5.} Place: 'il résulte de nouveau de loc. cit. que les morphismes correspondants'; Reported print: il résulte de nouveau de loc. cit. que les morphismes correspondants \ensuremath{\varepsilon}(X\_ij) \ensuremath{\to} \ensuremath{\varepsilon}(X) forment; Source TeX: il résulte de nouveau une famille épimorphismes correspondants; Note: The TeX garbles the clause: 'de loc. cit. que les morphismes' is replaced by 'une famille épimorphismes'.; Proposed edit: nouveau une famille épimorphismes correspondants =\allowbreak > nouveau de \textbackslash{}loccit\{\} que les morphismes correspondants\par
\textbf{GPT-6.1 Sol.} Place: Fin de la preuve de la proposition 1.2, après l'extraction d'une sous-famille finie; Reported print: il résulte de nouveau de loc. cit. que les morphismes correspondants; Source TeX: il résulte de nouveau une famille épimorphismes correspondants; Note: Le passage est corrompu : l'imprimé contient le renvoi « de loc. cit. » suivi de « que les morphismes ».; Proposed edit: nouveau une famille épimorphismes =\allowbreak > nouveau de \textbackslash{}loccit que les morphismes\par
the two models differ at this place\par
\par\noindent\hypertarget{DP0330}{}\textbf{DP0330 -- VI, p. 169 (agreed\_fix\_held)}\par
Remarque 1.5.1, 'topos induit E\_\{/X\},' | Remarque 1.5.1, topos induit\par
\textbf{Claude Opus 5.5.} Place: Remarque 1.5.1, 'topos induit E\_\{/X\},'; Reported print: E\_\{/X\} followed by a mark at subscript level (apparently a comma), then 'et signifie'; Source TeX: \$E\_\{/X'\}\$; Note: The mark after X at subscript level looks like a lowered comma rather than a prime; the sense requires E\_\{/X\}. Check the scan.; Proposed edit: \$E\_\{/X'\}\$ =\allowbreak > \$E\_\{/X\}\$,\par
\textbf{GPT-6.1 Sol.} Place: Remarque 1.5.1, topos induit; Reported print: \$E\_\{/X\}\$, et signifie; Source TeX: \$E\_\{/X'\}\$ et signifie; Note: La marque suivant X est la virgule de la phrase, non une prime de l'indice.; Proposed edit: \$E\_\{/X'\}\$ et =\allowbreak > \$E\_\{/X\}\$, et\par
both models give the same fix, held: a model marks it low confidence\par
\par\noindent\hypertarget{DP0331}{}\textbf{DP0331 -- VI, p. 169 (one\_model\_fix)}\par
proof of 1.4, '(II 4.5.1)'\par
\textbf{Claude Opus 5.5.} Place: proof of 1.4, '(II 4.5.1)'; Reported print: (II 4.5.1); Source TeX: \textbackslash{}hyperref[VI.lem4.5.1]\{4.5.1\}; Note: The print cites II 4.5.1; the label VI.lem4.5.1 links into Exposé VI. Wrong link target, same printed text.; Proposed edit: \textbackslash{}hyperref[VI.lem4.5.1]\{4.5.1\} =\allowbreak > 4.5.1\par
\textbf{GPT-6.1 Sol.} nothing reported here\par
a fix from one model only\par
\par\noindent\hypertarget{DP0332}{}\textbf{DP0332 -- VI, p. 173 (agreed\_fix\_held)}\par
end of page, proof of (iii) | Last paragraph, explanation of the conclusion that f is quasi-separated\par
\textbf{Claude Opus 5.5.} Place: end of page, proof of (iii); Reported print: et que la conclusion f quasi-séparé signifie que \ensuremath{\Delta}\_\{X/Y\} est quasi-compact, ce qui résulte de (iv); Source TeX: the clause « et que la conclusion \$f\$ quasi-séparé signifie que \$\textbackslash{}Delta\_\{X/Y\}\$ est quasi-compact » is written twice; Note: The clause appears once in the print; the TeX repeats it.; Proposed edit: quasi-compact, et que la conclusion \$f\$ quasi-séparé signifie que \$\textbackslash{}Delta\_\{X/Y\}\$ est quasi-compact, ce qui =\allowbreak > quasi-compact, ce qui\par
\textbf{GPT-6.1 Sol.} Place: Last paragraph, explanation of the conclusion that f is quasi-separated; Reported print: The clause « et que la conclusion f quasi-séparé signifie que \ensuremath{\Delta}\_\{X/Y\} est quasi-compact » occurs once.; Source TeX: The clause occurs twice consecutively.; Note: Remove the second occurrence of the clause.; Proposed edit: , et que la conclusion \$f\$ quasi-séparé signifie que \$\textbackslash{}Delta\_\{X/Y\}\$ est quasi-compact =\allowbreak > \par
both models give the same fix, held: a deletion in the first or last three lines of the page (the page cut may be off)\par
\par\noindent\hypertarget{DP0333}{}\textbf{DP0333 -- VI, p. 174 (agreed\_fix\_held)}\par
Remarques 1.9.1, end | Remarques 1.9.1, sentence asserting validity of 1.8 and 1.8.1\par
\textbf{Claude Opus 5.5.} Place: Remarques 1.9.1, end; Reported print: 1.8 et 1.8.1 sont valables; Source TeX: \{1.8.1\}. sont valables; Note: The TeX has a stray full stop after 1.8.1 that the print does not have.; Proposed edit: \{1.8.1\}. sont =\allowbreak > \{1.8.1\} sont\par
\textbf{GPT-6.1 Sol.} Place: Remarques 1.9.1, sentence asserting validity of 1.8 and 1.8.1; Reported print: Alors 1.8 et 1.8.1 sont valables; Source TeX: Alors 1.8 et 1.8.1. sont valables; Note: The TeX inserts a full stop between the subject and « sont valables »; the print has no sentence break.; Proposed edit: \textbackslash{}hyperref[VI.cor1.8.1]\{1.8.1\}. sont =\allowbreak > \textbackslash{}hyperref[VI.cor1.8.1]\{1.8.1\} sont\par
both models give the same fix, held: a model marks it low confidence\par
\par\noindent\hypertarget{DP0334}{}\textbf{DP0334 -- VI, p. 179 (conflict)}\par
Corollaire 1.15 | Corollary 1.15, hypothesis on the indexing set\par
\textbf{Claude Opus 5.5.} Place: Corollaire 1.15; Reported print: avec I \ensuremath{\in} U (U underlined); Source TeX: avec \$I \textbackslash{}in U\$; Note: The underlined U is the universe, which the edition elsewhere writes in script. Check that \textbackslash{}Uu exists in macros.tex; if not, use \textbackslash{}mathcal\{U\}.; Proposed edit: \$I \textbackslash{}in U\$ =\allowbreak > \$I \textbackslash{}in \textbackslash{}Uu\$\par
\textbf{GPT-6.1 Sol.} Place: Corollary 1.15, hypothesis on the indexing set; Reported print: avec I \ensuremath{\in} U ; alors X est quasi-séparé (resp. cohérent) si et seulement si les X\_i le sont; Source TeX: avec \$I \textbackslash{}in U\$; alors \$X =\allowbreak  \textbackslash{}amalg\_i X\_i\$ est quasi-séparé (\textbackslash{}emph\{\textbackslash{}resp\} cohérent) si et seulement si les \$X\_i\$ le sont.; Proposed correction: avec \$I \textbackslash{}in U\$ (\textbackslash{}resp \$I\$ fini); Note: The print omits a finiteness qualification for the coherent case, and the TeX retains that omission. An arbitrary coproduct of coherent objects need not be quasi-compact. The stated equivalence is valid with I finite in the resp. case; more generally, only finitely many summands may be non-initial.\par
the two models differ at this place: one model fixes the TeX, the other calls it a misprint of the print\par
\par\noindent\hypertarget{DP0335}{}\textbf{DP0335 -- VI, p. 182 (agreed\_fix\_held)}\par
heading of 1.20.1 | Heading of 1.20.1\par
\textbf{Claude Opus 5.5.} Place: heading of 1.20.1; Reported print: Remarques 1.20.1.; Source TeX: \{Remarque 1.20.1\}; Note: The print reads 'Remarques' (plural; the ending 'es' looks overtyped). The TeX has the singular.; Proposed edit: \{Remarque 1.20.1\} =\allowbreak > \{Remarques 1.20.1\}\par
\textbf{GPT-6.1 Sol.} Place: Heading of 1.20.1; Reported print: Remarques 1.20.1.; Source TeX: Remarque 1.20.1; Note: The printed heading is plural.; Proposed edit: Remarque 1.20.1 =\allowbreak > Remarques 1.20.1\par
both models give the same fix, held: a model marks it low confidence\par
\par\noindent\hypertarget{DP0336}{}\textbf{DP0336 -- VI, p. 185 (agreed\_fix\_held)}\par
site X\_ét | 1.22.1, notation for the site of étale schemes\par
\textbf{Claude Opus 5.5.} Place: site X\_ét; Reported print: du site X\_ét des schémas étales sur X; Source TeX: \$X\_\{\textbackslash{}text\{et\}\}\$ des schémas; Note: On this page the printed subscript carries an accent (ét). On p. 184 it is printed 'et' without one.; Proposed edit: \$X\_\{\textbackslash{}text\{et\}\}\$ des schémas =\allowbreak > \$X\_\{\textbackslash{}text\{ét\}\}\$ des schémas\par
\textbf{GPT-6.1 Sol.} Place: 1.22.1, notation for the site of étale schemes; Reported print: \$X\_\{\textbackslash{}text\{ét\}\}\$; Source TeX: \$X\_\{\textbackslash{}text\{et\}\}\$; Note: The printed subscript has an acute accent on e.; Proposed edit: \$X\_\{\textbackslash{}text\{et\}\}\$ =\allowbreak > \$X\_\{\textbackslash{}text\{ét\}\}\$\par
both models give the same fix, held: a model marks it low confidence\par
\par\noindent\hypertarget{DP0337}{}\textbf{DP0337 -- VI, p. 185 (cut)}\par
end of page TeX, '\textbackslash{}item \textbackslash{}underline\{Nécessité\}.'\par
\textbf{Claude Opus 5.5.} Place: end of page TeX, '\textbackslash{}item \textbackslash{}underline\{Nécessité\}.'; Reported print: page ends with '... l'application (1.23.1) soit bijective.'; Source TeX: \textbackslash{}begin\{enumerate\} ... \textbackslash{}item \textbackslash{}underline\{Nécessité\}.; Note: Printed page 185 ends after the statement of Theorem 1.23. The proof's '(i) Nécessité.' that closes the page's TeX belongs to p. 186, but the comment line does not say so.\par
\textbf{GPT-6.1 Sol.} nothing reported here\par
a model reports that the automatic cut gives lines of this page to its neighbour, or the neighbour's lines to this page: check the lines it names against the scan of the page that prints them (that page's checker did not have their TeX)\par
\par\noindent\hypertarget{DP0338}{}\textbf{DP0338 -- VI, p. 185 (cut)}\par
End of the supplied TeX, after Theorem 1.23\par
\textbf{Claude Opus 5.5.} nothing reported here\par
\textbf{GPT-6.1 Sol.} Place: End of the supplied TeX, after Theorem 1.23; Reported print: The page ends after « l'application (1.23.1) soit bijective. »; Source TeX: \textbackslash{}item \textbackslash{}underline\{Nécessité\}.; Note: The proof-enumeration block following \textbackslash{}end\{enonce*\}, including « (i) Nécessité. », is not printed on p. 185 and belongs to the following page. This overrun is not covered by the cut's opening comment.\par
a model reports that the automatic cut gives lines of this page to its neighbour, or the neighbour's lines to this page: check the lines it names against the scan of the page that prints them (that page's checker did not have their TeX)\par
\par\noindent\hypertarget{DP0339}{}\textbf{DP0339 -- VI, p. 186 (cut)}\par
start of page\par
\textbf{Claude Opus 5.5.} Place: start of page; Reported print: (i)  Nécessité. Supposons X quasi-compact, ...; Source TeX: Supposons\textbackslash{}pageoriginale \$X\$ quasi-compact, ...; Note: The page prints the item label (i) and the underlined word « Nécessité. » before « Supposons ». This text is not in this page's cut, so it must sit in the preceding source line (776). The page really begins at « (i) Nécessité. » (source line 776), not at line 777.\par
\textbf{GPT-6.1 Sol.} nothing reported here\par
a model reports that the automatic cut gives lines of this page to its neighbour, or the neighbour's lines to this page: check the lines it names against the scan of the page that prints them (that page's checker did not have their TeX)\par
\par\noindent\hypertarget{DP0340}{}\textbf{DP0340 -- VI, p. 186 (cut)}\par
Opening of the page, before « Supposons »\par
\textbf{Claude Opus 5.5.} nothing reported here\par
\textbf{GPT-6.1 Sol.} Place: Opening of the page, before « Supposons »; Reported print: (i) Nécessité. Supposons; Source TeX: Supposons\textbackslash{}pageoriginale; Note: The printed page begins with the item label and « Nécessité. », absent from the supplied cut. Check the source lines immediately before line 777 and include this opening in the p. 186 slice; do not insert a duplicate heading.\par
a model reports that the automatic cut gives lines of this page to its neighbour, or the neighbour's lines to this page: check the lines it names against the scan of the page that prints them (that page's checker did not have their TeX)\par
\par\noindent\hypertarget{DP0341}{}\textbf{DP0341 -- VI, p. 190 (conflict)}\par
proof of b), sentence « En vertu de a) » | Proof of Corollaire 1.24.2 b), after R =\allowbreak  lim R\_i\par
\textbf{Claude Opus 5.5.} Place: proof of b), sentence « En vertu de a) »; Reported print: on a R =\allowbreak  lim R\_i, limite inductive filtrante avec les R\_i dans E\_PF, donc X\_o est limite inductive filtrante des Coker(R\_i \ensuremath{\rightrightarrows} X\_o) =\allowbreak  X\_i.; Source TeX: limite inductive filtrante des coker \$(R\_i \textbackslash{}rightrightarrows X\_0) =\allowbreak  X\_i\$; Note: The TeX drops a clause: « avec les R\_i dans E\_PF, donc X\_o est limite inductive filtrante ». The print also has « Coker », not « coker ».; Proposed edit: filtrante des coker \$(R\_i \textbackslash{}rightrightarrows X\_0) =\allowbreak  X\_i\$ =\allowbreak > filtrante avec les \$R\_i\$ dans \$E\_\{PF\}\$, donc \$X\_0\$ est limite inductive filtrante des \$\textbackslash{}Coker(R\_i \textbackslash{}rightrightarrows X\_0) =\allowbreak  X\_i\$\par
\textbf{GPT-6.1 Sol.} Place: Proof of Corollaire 1.24.2 b), after R =\allowbreak  lim R\_i; Reported print: limite inductive filtrante avec les \$R\_i\$ dans \$E\_\{PF\}\$, donc \$X\_0\$ est limite inductive filtrante des \$\textbackslash{}Coker(R\_i \textbackslash{}rightrightarrows X\_0) =\allowbreak  X\_i\$; Source TeX: limite inductive filtrante des coker \$(R\_i \textbackslash{}rightrightarrows X\_0) =\allowbreak  X\_i\$; Note: A whole clause is omitted, making the TeX describe R itself as the limit of the cokernels. The insertion reproduces the printed X\_0; that printed misprint is reported separately.; Proposed edit: des coker \$(R\_i \textbackslash{}rightrightarrows X\_0) =\allowbreak  X\_i\$ =\allowbreak > avec les \$R\_i\$ dans \$E\_\{PF\}\$, donc \$X\_0\$ est limite inductive filtrante des \$\textbackslash{}Coker (R\_i \textbackslash{}rightrightarrows X\_0) =\allowbreak  X\_i\$\par
the two models differ at this place\par
\par\noindent\hypertarget{DP0342}{}\textbf{DP0342 -- VI, p. 210 (one\_model\_fix)}\par
Définition 2.9.1\par
\textbf{Claude Opus 5.5.} Place: Définition 2.9.1; Reported print: de 1.25, i.e. si la sous-catégorie; Source TeX: de \textbackslash{}hyperref[VI.cor1.25]\{1.25\}. \textbackslash{}ie si la; Note: The print has a comma before 'i.e.', not a full stop.; Proposed edit: \{1.25\}. \textbackslash{}ie si la =\allowbreak > \{1.25\}, \textbackslash{}ie si la\par
\textbf{GPT-6.1 Sol.} nothing reported here\par
a fix from one model only\par
\par\noindent\hypertarget{DP0343}{}\textbf{DP0343 -- VI, p. 215 (conflict)}\par
2.14, définition de C | 2.14, sous-catégorie C des objets noethériens\par
\textbf{Claude Opus 5.5.} Place: 2.14, définition de C; Reported print: objets noethériens (=\allowbreak  cohérents) de E; Source TeX: objets noethériens (cohérents) de \$E\$; Note: The =\allowbreak  sign inside the parenthesis is missing from the TeX.; Proposed edit: (cohérents) de \$E\$, alors =\allowbreak > (=\allowbreak  cohérents) de \$E\$, alors\par
\textbf{GPT-6.1 Sol.} Place: 2.14, sous-catégorie C des objets noethériens; Reported print: noethériens (=\allowbreak  cohérents) de E; Source TeX: noethériens (cohérents) de \$E\$; Note: Le signe d’égalité imprimé dans la parenthèse manque.; Proposed edit: (cohérents) de \$E\$ =\allowbreak > (\$=\allowbreak \$ cohérents) de \$E\$\par
the two models differ at this place\par
\par\noindent\hypertarget{DP0344}{}\textbf{DP0344 -- VI, p. 217 (conflict)}\par
Exercice 2.16, preamble | Exercice 2.16, introductory paragraph, induced topos\par
\textbf{Claude Opus 5.5.} Place: Exercice 2.16, preamble; Reported print: le topos induit E\_\{/X'\} est canoniquement équivalent; Source TeX: le topos induit \$E\_\{/X\}\$, est; Note: The prime on X' is lost: the TeX has a comma after the formula instead of the prime inside it.; Proposed edit: \$E\_\{/X\}\$, =\allowbreak > \$E\_\{/X'\}\$\par
\textbf{GPT-6.1 Sol.} Place: Exercice 2.16, introductory paragraph, induced topos; Reported print: E\_\{/X'\}; Source TeX: E\_\{/X\}; Note: The print has a prime: the slice is over X', not X.; Proposed edit: E\_\{/X\} =\allowbreak > E\_\{/X'\}\par
the two models differ at this place\par
\par\noindent\hypertarget{DP0345}{}\textbf{DP0345 -- VI, p. 224 (conflict)}\par
Exercice 2.18 b), fin | Exercice 2.18 b), index of Sup | Exercice 2.18 b), formula defining the summands\par
\textbf{Claude Opus 5.5.} Place: Exercice 2.18 b), fin; Reported print: X est somme des X(j)=\allowbreak  Sup\_\{i\ensuremath{\in}J\} X\_i; Source TeX: X(i) =\allowbreak  \textbackslash{}displaystyle; Note: The print indexes by j (the classes of J =\allowbreak  I/R), not i.; Proposed edit: X(i) =\allowbreak  \textbackslash{}displaystyle =\allowbreak > X(j) =\allowbreak  \textbackslash{}displaystyle || Place: Exercice 2.18 b), index of Sup; Reported print: i\ensuremath{\in} J (faint, broken glyph); Source TeX: \textbackslash{}mathop\{\textbackslash{}Sup\}\_\{i \textbackslash{}in I\}\}; Note: The glyph under Sup is broken but is not I; it reads J (or j, the meaning being i \ensuremath{\in} j). Report together with the X(j) fix.; Proposed edit: \textbackslash{}mathop\{\textbackslash{}Sup\}\_\{i \textbackslash{}in I\}\} =\allowbreak > \textbackslash{}mathop\{\textbackslash{}Sup\}\_\{i \textbackslash{}in J\}\}\par
\textbf{GPT-6.1 Sol.} Place: Exercice 2.18 b), formula defining the summands; Reported print: \$X(j) =\allowbreak  \textbackslash{}mathop\{\textbackslash{}Sup\}\_\{i \textbackslash{}in j\} X\_i\$; Source TeX: X(i) =\allowbreak  \textbackslash{}displaystyle\{\textbackslash{}mathop\{\textbackslash{}Sup\}\_\{i \textbackslash{}in I\}\} X\_i; Note: The print labels the equivalence class j in both X(j) and the lower limit i\ensuremath{\in}j. The foot of the lower-index j is faint/broken, but its remaining shape is readable.; Proposed edit: X(i) =\allowbreak  \textbackslash{}displaystyle\{\textbackslash{}mathop\{\textbackslash{}Sup\}\_\{i \textbackslash{}in I\}\} =\allowbreak > X(j) =\allowbreak  \textbackslash{}displaystyle\{\textbackslash{}mathop\{\textbackslash{}Sup\}\_\{i \textbackslash{}in j\}\}\par
the two models differ at this place\par
\par\noindent\hypertarget{DP0346}{}\textbf{DP0346 -- VI, p. 234 (one\_model\_fix)}\par
3.11 e), 1.25. (Consulter\par
\textbf{Claude Opus 5.5.} Place: 3.11 e), 1.25. (Consulter; Reported print: les conditions équivalentes de 1.25. (Consulter 1.28 b.); Source TeX: \textbackslash{}hyperref[VI.cor1.25]\{1.25\} (Consulter; Note: The print has a full stop after 1.25.; Proposed edit: \textbackslash{}hyperref[VI.cor1.25]\{1.25\} (Consulter =\allowbreak > \textbackslash{}hyperref[VI.cor1.25]\{1.25\}. (Consulter\par
\textbf{GPT-6.1 Sol.} nothing reported here\par
a fix from one model only\par
\par\noindent\hypertarget{DP0347}{}\textbf{DP0347 -- VI, p. 235 (conflict)}\par
3.12 c), (Utiliser 3.2 et 2.17 g).) | Exercice 3.12 c), indication entre parenthèses\par
\textbf{Claude Opus 5.5.} Place: 3.12 c), (Utiliser 3.2 et 2.17 g).); Reported print: (Utiliser 3.2 et 2.17 g).); Source TeX: 2,17 g).); Note: The print reads 2.17 with a point; the TeX has a comma.; Proposed edit: 2,17 g) =\allowbreak > \textbackslash{}hyperref[VI.exer2.17]\{2.17\} g)\par
\textbf{GPT-6.1 Sol.} Place: Exercice 3.12 c), indication entre parenthèses; Reported print: 2.17 g); Source TeX: 2,17 g); Note: Le numéro du renvoi imprimé est 2.17, avec un point.; Proposed edit: 2,17 g) =\allowbreak > 2.17 g)\par
the two models differ at this place\par
\par\noindent\hypertarget{DP0348}{}\textbf{DP0348 -- VI, p. 241 (cut)}\par
last line of the page (end of b), source line 2630)\par
\textbf{Claude Opus 5.5.} Place: last line of the page (end of b), source line 2630); Reported print: p. 241 ends '... car U est maintenant cohérent,'; the next word 'e' opens p. 242 ('e l'étant et j: U \ensuremath{\to} e ...'); Source TeX: ... car \$U\$ est maintenant cohérent, \$e\$; Note: The TeX of p. 241 ends with '\$e\$', which the print sets as the first word of p. 242. The comment line does not flag this line as shared. Page 242 really begins with '\$e\$ l'étant'.\par
\textbf{GPT-6.1 Sol.} nothing reported here\par
a model reports that the automatic cut gives lines of this page to its neighbour, or the neighbour's lines to this page: check the lines it names against the scan of the page that prints them (that page's checker did not have their TeX)\par
\par\noindent\hypertarget{DP0349}{}\textbf{DP0349 -- VI, p. 241 (cut)}\par
Fin de page, démonstration de 4.6 b)\par
\textbf{Claude Opus 5.5.} nothing reported here\par
\textbf{GPT-6.1 Sol.} Place: Fin de page, démonstration de 4.6 b); Reported print: La page finit par « car U est maintenant cohérent, ». La page 242 commence par « e l'étant ».; Source TeX: cohérent, \$e\$; Note: Le \$e\$ final appartient à la page 242. Le partage de la ligne source 2630 entre les deux pages n'est pas signalé dans les commentaires de coupe. Corriger l'extraction, sans modifier le texte.\par
a model reports that the automatic cut gives lines of this page to its neighbour, or the neighbour's lines to this page: check the lines it names against the scan of the page that prints them (that page's checker did not have their TeX)\par
\par\noindent\hypertarget{DP0350}{}\textbf{DP0350 -- VI, p. 246 (one\_model\_fix)}\par
"on sait (II 6.10)"\par
\textbf{Claude Opus 5.5.} Place: "on sait (II 6.10)"; Reported print: (II 6.10); Source TeX: (II \textbackslash{}hyperref[VI.subsec6.10]\{6.10\}); Note: The print refers to II 6.10, but the TeX links to label VI.subsec6.10, a paragraph of exposé VI. The printed number is the same; the link target is wrong.; Proposed edit: (II \textbackslash{}hyperref[VI.subsec6.10]\{6.10\}) =\allowbreak > (II 6.10)\par
\textbf{GPT-6.1 Sol.} nothing reported here\par
a fix from one model only\par
\par\noindent\hypertarget{DP0351}{}\textbf{DP0351 -- VI, p. 252 (one\_model\_fix)}\par
c), après la définition de g(X'')\par
\textbf{Claude Opus 5.5.} nothing reported here\par
\textbf{GPT-6.1 Sol.} Place: c), après la définition de g(X''); Reported print: dont les composantes sont dans E' coh; Source TeX: dont les composants; Note: L'imprimé emploie le féminin « composantes ».; Proposed edit: les composants =\allowbreak > les composantes\par
a fix from one model only\par
\par\noindent\hypertarget{DP0352}{}\textbf{DP0352 -- VI, p. 258 (conflict)}\par
Corollaire 4.14, f ↦ f\_0 =\allowbreak  f|E'\_coh | Corollaire 4.14, définition de f₀\par
\textbf{Claude Opus 5.5.} Place: Corollaire 4.14, f ↦ f\_0 =\allowbreak  f|E'\_coh; Reported print: f|E'\_coh (vertical bar, restriction); Source TeX: f / E'\_\{\textbackslash{}coh\}; Note: The print uses a restriction bar. The TeX has a slash, which reads as a localization.; Proposed edit: f / E'\_\{\textbackslash{}coh\} =\allowbreak > f | E'\_\{\textbackslash{}coh\}\par
\textbf{GPT-6.1 Sol.} Place: Corollaire 4.14, définition de f₀; Reported print: f\_0 =\allowbreak  f\textbackslash{}vert E'\_\{\textbackslash{}coh\}; Source TeX: f\_0 =\allowbreak  f / E'\_\{\textbackslash{}coh\}; Note: Le signe imprimé de restriction est une barre verticale, non une barre oblique.; Proposed edit: f / E'\_\{\textbackslash{}coh\} =\allowbreak > f \textbackslash{}vert E'\_\{\textbackslash{}coh\}\par
the two models differ at this place\par
\par\noindent\hypertarget{DP0353}{}\textbf{DP0353 -- VI, p. 260 (cut)}\par
end of page\par
\textbf{Claude Opus 5.5.} Place: end of page; Reported print: page ends at 'Ceci résulte du lemme suivant appliqué au cas où C =\allowbreak  E'\_coh :'; Source TeX: \textbackslash{}begin\{enonce*\}\{Lemme 5.4\}\textbackslash{}etiquette[5.4]\{VI.lem5.4\}; Note: The heading 'Lemme 5.4' and the whole statement are printed on p. 261, not p. 260. The TeX cut includes the heading line before the flagged last line.\par
\textbf{GPT-6.1 Sol.} nothing reported here\par
a model reports that the automatic cut gives lines of this page to its neighbour, or the neighbour's lines to this page: check the lines it names against the scan of the page that prints them (that page's checker did not have their TeX)\par
\par\noindent\hypertarget{DP0354}{}\textbf{DP0354 -- VI, p. 260 (cut)}\par
Fin de la coupe, titre du lemme 5.4\par
\textbf{Claude Opus 5.5.} nothing reported here\par
\textbf{GPT-6.1 Sol.} Place: Fin de la coupe, titre du lemme 5.4; Reported print: La page finit après « appliqué au cas où C =\allowbreak  E'\_\{\textbackslash{}coh\} : ».; Source TeX: \textbackslash{}begin\{enonce*\}\{Lemme 5.4\}\textbackslash{}etiquette[5.4]\{VI.lem5.4\}; Note: Le titre du lemme, à la ligne source 3289, appartient lui aussi à la page 261. Le commentaire de coupe ne réserve à cette page voisine que la ligne 3290. Rectifier les bornes de la coupe, sans déplacer le texte du source.\par
a model reports that the automatic cut gives lines of this page to its neighbour, or the neighbour's lines to this page: check the lines it names against the scan of the page that prints them (that page's checker did not have their TeX)\par
\par\noindent\hypertarget{DP0355}{}\textbf{DP0355 -- VI, p. 261 (conflict)}\par
Corollaire 5.5, statement | Corollaire 5.5, second foncteur de l'énoncé\par
\textbf{Claude Opus 5.5.} Place: Corollaire 5.5, statement; Reported print: H\textasciicircum{}q\_Φ(E,-) et H\textasciicircum{}q\_Φ with the H underlined (the sheaf); Source TeX: et \$H\textasciicircum{}q\_\{\textbackslash{}Phi\}\$; Note: The second functor is the sheaf H\textasciicircum{}q\_Φ (underlined in the print, cf. V 6.13). The TeX writes it like the group H\textasciicircum{}q\_Φ; it has \ensuremath{\mathcal{H}}\textasciicircum{}q\_Φ at the end of the proof.; Proposed edit: et \$H\textasciicircum{}q\_\{\textbackslash{}Phi\}\$ =\allowbreak > et \$\ensuremath{\mathcal{H}}\textasciicircum{}q\_\{\textbackslash{}Phi\}\$\par
\textbf{GPT-6.1 Sol.} Place: Corollaire 5.5, second foncteur de l'énoncé; Reported print: H\textasciicircum{}q\_\textbackslash{}Phi(E,-) et \textbackslash{}underline\{H\}\textasciicircum{}q\_\textbackslash{}Phi; Source TeX: \$H\textasciicircum{}q\_\{\textbackslash{}Phi\} (E, -)\$ et \$H\textasciicircum{}q\_\{\textbackslash{}Phi\}\$; Note: Le second H est souligné pour désigner le faisceau de cohomologie, distinct du groupe désigné par le premier H.; Proposed edit: et \$H\textasciicircum{}q\_\{\textbackslash{}Phi\}\$ =\allowbreak > et \$\textbackslash{}Hh\textasciicircum{}q\_\{\textbackslash{}Phi\}\$\par
the two models differ at this place\par
\par\noindent\hypertarget{DP0356}{}\textbf{DP0356 -- VI, p. 261 (cut)}\par
Début de la page, titre du lemme 5.4\par
\textbf{Claude Opus 5.5.} nothing reported here\par
\textbf{GPT-6.1 Sol.} Place: Début de la page, titre du lemme 5.4; Reported print: Lemme 5.4 : Soient E un topos localement cohérent...; Source TeX: Soient \$E\$\textbackslash{}pageoriginale un topos localement cohérent; Note: La page commence au titre « Lemme 5.4 ». L'ouverture de cet énoncé précède le début du fragment fourni et doit être récupérée dans la découpe voisine, sans ajouter un second titre au fichier source.\par
a model reports that the automatic cut gives lines of this page to its neighbour, or the neighbour's lines to this page: check the lines it names against the scan of the page that prints them (that page's checker did not have their TeX)\par
\par\noindent\hypertarget{DP0357}{}\textbf{DP0357 -- VI, p. 262 (conflict)}\par
5.9.1, first line | 5.9.1, second line | 5.9, les deux lignes de (5.9.1)\par
\textbf{Claude Opus 5.5.} Place: 5.9.1, first line; Reported print: Ext\textasciicircum{}p\_A(X;G,H\textasciicircum{}q\_Z F) with the H underlined; Source TeX: \textbackslash{}Ext\textasciicircum{}p\_A (X; G, H\textasciicircum{}q\_Z F); Note: The underlined H is the sheaf H\textasciicircum{}q\_Z, not the group.; Proposed edit: \textbackslash{}Ext\textasciicircum{}p\_A (X; G, H\textasciicircum{}q\_Z F) =\allowbreak > \textbackslash{}Ext\textasciicircum{}p\_A (X; G, \ensuremath{\mathcal{H}}\textasciicircum{}q\_Z F) || Place: 5.9.1, second line; Reported print: \ensuremath{\mathcal{E}}xt\textasciicircum{}p\_A(G,H\textasciicircum{}q\_Z F) with the H underlined; Source TeX: \textbackslash{}Eext\textasciicircum{}p\_A (G, H\textasciicircum{}q\_Z F); Note: The underlined H is the sheaf H\textasciicircum{}q\_Z.; Proposed edit: \textbackslash{}Eext\textasciicircum{}p\_A (G, H\textasciicircum{}q\_Z F) =\allowbreak > \textbackslash{}Eext\textasciicircum{}p\_A (G, \ensuremath{\mathcal{H}}\textasciicircum{}q\_Z F)\par
\textbf{GPT-6.1 Sol.} Place: 5.9, les deux lignes de (5.9.1); Reported print: \textbackslash{}underline\{H\}\textasciicircum{}q\_Z F dans les deux termes E\_2; Source TeX: H\textasciicircum{}q\_Z F dans les deux termes E\_2; Note: Dans les deux lignes, le H souligné désigne le faisceau de cohomologie à support, non un groupe de cohomologie.; Proposed edit: H\textasciicircum{}q\_Z F) \textbackslash{}Rightarrow \textbackslash{}Ext\textasciicircum{}\{p+q\}\_\{A, \textbackslash{}Phi\} (X; G, F)\{\textbackslash{}quoi\},\textbackslash{}\textbackslash{} ''E\textasciicircum{}\{pq\}\_2 =\allowbreak  \textbackslash{}varinjlim\_\{Z \textbackslash{}in \textbackslash{}Phi\} \textbackslash{}Eext\textasciicircum{}p\_A (G, H\textasciicircum{}q\_Z F) =\allowbreak > \textbackslash{}Hh\textasciicircum{}q\_Z F) \textbackslash{}Rightarrow \textbackslash{}Ext\textasciicircum{}\{p+q\}\_\{A, \textbackslash{}Phi\} (X; G, F)\{\textbackslash{}quoi\},\textbackslash{}\textbackslash{} ''E\textasciicircum{}\{pq\}\_2 =\allowbreak  \textbackslash{}varinjlim\_\{Z \textbackslash{}in \textbackslash{}Phi\} \textbackslash{}Eext\textasciicircum{}p\_A (G, \textbackslash{}Hh\textasciicircum{}q\_Z F)\par
the two models differ at this place\par
\par\noindent\hypertarget{DP0358}{}\textbf{DP0358 -- VI, p. 263 (conflict)}\par
5.9.2, first line | 5.9.2, second line | Début de page, les deux lignes de (5.9.2)\par
\textbf{Claude Opus 5.5.} Place: 5.9.2, first line; Reported print: Ext\textasciicircum{}p\_A(X;G,H\textasciicircum{}q\_Φ F) with the H underlined; Source TeX: \textbackslash{}Ext\textasciicircum{}p\_A (X; G, H\textasciicircum{}q\_\textbackslash{}Phi F); Note: The underlined H is the sheaf H\textasciicircum{}q\_Φ.; Proposed edit: \textbackslash{}Ext\textasciicircum{}p\_A (X; G, H\textasciicircum{}q\_\textbackslash{}Phi F) =\allowbreak > \textbackslash{}Ext\textasciicircum{}p\_A (X; G, \ensuremath{\mathcal{H}}\textasciicircum{}q\_\textbackslash{}Phi F) || Place: 5.9.2, second line; Reported print: \ensuremath{\mathcal{E}}xt\textasciicircum{}p\_A(G;H\textasciicircum{}q\_Φ F) with the H underlined; Source TeX: \textbackslash{}Eext\textasciicircum{}p\_A (G, H\textasciicircum{}q\_\textbackslash{}Phi F); Note: The underlined H is the sheaf H\textasciicircum{}q\_Φ.; Proposed edit: \textbackslash{}Eext\textasciicircum{}p\_A (G, H\textasciicircum{}q\_\textbackslash{}Phi F) =\allowbreak > \textbackslash{}Eext\textasciicircum{}p\_A (G, \ensuremath{\mathcal{H}}\textasciicircum{}q\_\textbackslash{}Phi F)\par
\textbf{GPT-6.1 Sol.} Place: Début de page, les deux lignes de (5.9.2); Reported print: \textbackslash{}underline\{H\}\textasciicircum{}q\_\textbackslash{}Phi F dans les deux isomorphismes; Source TeX: H\textasciicircum{}q\_\textbackslash{}Phi F dans les deux isomorphismes; Note: Les deux H sont soulignés dans le scan : il s'agit des faisceaux de cohomologie à support.; Proposed edit: H\textasciicircum{}q\_\textbackslash{}Phi F)\&\textbackslash{}quav \textbackslash{}frac\{p(p-1)\}\{2\} \textbackslash{}leq r - 1\{\textbackslash{}quoi\},\textbackslash{}\textbackslash{} ''E\textasciicircum{}\{pq\}\_2 \textbackslash{}simeq \textbackslash{}Eext\textasciicircum{}p\_A (G, H\textasciicircum{}q\_\textbackslash{}Phi F) =\allowbreak > \textbackslash{}Hh\textasciicircum{}q\_\textbackslash{}Phi F)\&\textbackslash{}quav \textbackslash{}frac\{p(p-1)\}\{2\} \textbackslash{}leq r - 1\{\textbackslash{}quoi\},\textbackslash{}\textbackslash{} ''E\textasciicircum{}\{pq\}\_2 \textbackslash{}simeq \textbackslash{}Eext\textasciicircum{}p\_A (G, \textbackslash{}Hh\textasciicircum{}q\_\textbackslash{}Phi F)\par
the two models differ at this place\par
\par\noindent\hypertarget{DP0359}{}\textbf{DP0359 -- VI, p. 265 (cut)}\par
last line\par
\textbf{Claude Opus 5.5.} Place: last line; Reported print: the page ends with '… à valeur dans'; Source TeX: à valeur dans la; Note: The word 'la' at the end of the cut belongs to p. 266. The comment does not mark the last line as shared.\par
\textbf{GPT-6.1 Sol.} nothing reported here\par
a model reports that the automatic cut gives lines of this page to its neighbour, or the neighbour's lines to this page: check the lines it names against the scan of the page that prints them (that page's checker did not have their TeX)\par
\par\noindent\hypertarget{DP0360}{}\textbf{DP0360 -- VI, p. 265 (cut)}\par
Dernière ligne de la page\par
\textbf{Claude Opus 5.5.} nothing reported here\par
\textbf{GPT-6.1 Sol.} Place: Dernière ligne de la page; Reported print: à valeur dans; Source TeX: à valeur dans la; Note: La page 265 se termine après « dans ». L'article « la », inclus dans le fragment sans débordement final annoncé, appartient à la page 266. Ne pas le supprimer du fichier source.\par
a model reports that the automatic cut gives lines of this page to its neighbour, or the neighbour's lines to this page: check the lines it names against the scan of the page that prints them (that page's checker did not have their TeX)\par
\par\noindent\hypertarget{DP0361}{}\textbf{DP0361 -- VI, p. 266 (cut)}\par
Beginning of page\par
\textbf{Claude Opus 5.5.} nothing reported here\par
\textbf{GPT-6.1 Sol.} Place: Beginning of page; Reported print: la catégorie des ensembles appartenant à l'univers; Source TeX: catégorie\textbackslash{}pageoriginale des ensembles appartenant à l'univers:; Note: The supplied cut omits the printed initial word « la ». Check the end of source line 3556: the article may already be there. The page begins with « la catégorie »; verify the extraction boundary rather than inserting a potentially duplicate article.\par
a model reports that the automatic cut gives lines of this page to its neighbour, or the neighbour's lines to this page: check the lines it names against the scan of the page that prints them (that page's checker did not have their TeX)\par
\par\noindent\hypertarget{DP0362}{}\textbf{DP0362 -- VI, p. 267 (one\_model\_fix)}\par
On notera toute-fois\par
\textbf{Claude Opus 5.5.} Place: On notera toute-fois; Reported print: toute-/fois (hyphenation at a line break); Source TeX: toute-fois; Note: The line-break hyphen has been kept as part of the word.; Proposed edit: toute-fois =\allowbreak > toutefois\par
\textbf{GPT-6.1 Sol.} nothing reported here\par
a fix from one model only\par
\par\noindent\hypertarget{DP0363}{}\textbf{DP0363 -- VI, p. 270 (one\_model\_fix)}\par
Exercice 6.6 c)\par
\textbf{Claude Opus 5.5.} Place: Exercice 6.6 c); Reported print: (II 5.2); Source TeX: (II \textbackslash{}hyperref[VI.cor5.2]\{5.2\}); Note: The reference is to Exposé II 5.2, but the link points to label VI.cor5.2, which is in this exposé. The printed text is the same.; Proposed edit: (II \textbackslash{}hyperref[VI.cor5.2]\{5.2\}) =\allowbreak > (II 5.2)\par
\textbf{GPT-6.1 Sol.} nothing reported here\par
a fix from one model only\par
\par\noindent\hypertarget{DP0364}{}\textbf{DP0364 -- VI, p. 270 (one\_model\_fix)}\par
Exercise 6.6 e), sheafification formula\par
\textbf{Claude Opus 5.5.} nothing reported here\par
\textbf{GPT-6.1 Sol.} Place: Exercise 6.6 e), sheafification formula; Reported print: \textbackslash{}sheaf\{a\}G(X) \textbackslash{}xrightarrow\{\textbackslash{}sim\} \textbackslash{}varinjlim\_\{S(X)\} G(\textbackslash{}cdot); Source TeX: \textbackslash{}frac\{\textbackslash{}lim\}\{S(X)\}; Note: The print has an indexed inductive-limit sign, not a fraction with denominator S(X). The right end of the stroke beneath « lim » is weak in this scan.; Proposed edit: \textbackslash{}frac\{\textbackslash{}lim\}\{S(X)\} =\allowbreak > \textbackslash{}varinjlim\_\{S(X)\}\par
a fix from one model only\par
\par\noindent\hypertarget{DP0365}{}\textbf{DP0365 -- VI, p. 272 (one\_model\_fix)}\par
Exercice 6.8, 1), foncteur canonique ; convention également présente en 6.7 et 6.10–6.11\par
\textbf{Claude Opus 5.5.} nothing reported here\par
\textbf{GPT-6.1 Sol.} Place: Exercice 6.8, 1), foncteur canonique ; convention également présente en 6.7 et 6.10–6.11; Reported print: \textbackslash{}underrightarrow\{\textbackslash{}mathrm\{Lim\}\}\_\{\textbackslash{}Ee\textasciicircum{}\{\textbackslash{}circ\}\} \textbackslash{}Gg \textbackslash{}to \textbackslash{}varinjlim\_\{\textbackslash{}Ee\textasciicircum{}\{\textbackslash{}circ\}\} g; Source TeX: \textbackslash{}varinjlim\_\{\textbackslash{}Ee\textasciicircum{}\{\textbackslash{}circ\}\} \textbackslash{}Gg \textbackslash{}to \textbackslash{}varinjlim\_\{\textbackslash{}Ee\textasciicircum{}\{\textbackslash{}circ\}\} g; Note: Le print distingue Lim, pour la limite de catégories fibrées, de lim, pour la limite ordinaire. Ce contraste est explicite ici et à la fin de 6.11. Le TeX efface cette distinction. Convention signalée une seule fois ; l'édition proposée est représentative.; Proposed edit: \$\textbackslash{}varinjlim\_\{\textbackslash{}Ee\textasciicircum{}\{\textbackslash{}circ\}\} \textbackslash{}Gg \textbackslash{}to   \textbackslash{}varinjlim\_\{\textbackslash{}Ee\textasciicircum{}\{\textbackslash{}circ\}\} g\$. =\allowbreak > \$\textbackslash{}mathop\{\textbackslash{}underrightarrow\{\textbackslash{}mathrm\{Lim\}\}\}\_\{\textbackslash{}Ee\textasciicircum{}\{\textbackslash{}circ\}\} \textbackslash{}Gg \textbackslash{}to   \textbackslash{}varinjlim\_\{\textbackslash{}Ee\textasciicircum{}\{\textbackslash{}circ\}\} g\$.\par
a fix from one model only\par
\par\noindent\hypertarget{DP0366}{}\textbf{DP0366 -- VI, p. 280 (one\_model\_fix)}\par
7.2.5\par
\textbf{Claude Opus 5.5.} Place: 7.2.5; Reported print: (V-U-Top\_\{/I\}) \ensuremath{\to} (V-U-Site\_\{/I\}) with the U underlined; Source TeX: \textbackslash{}text\{-\}U\textbackslash{}text\{-\}\textbackslash{}Site\_\{/I\}; Note: The U is underlined (the universe) in the print, as in the Top side.; Proposed edit: \textbackslash{}text\{-\}U\textbackslash{}text\{-\} =\allowbreak > \textbackslash{}text\{-\}\ensuremath{𝒰}\textbackslash{}text\{-\}\par
\textbf{GPT-6.1 Sol.} nothing reported here\par
a fix from one model only\par
\par\noindent\hypertarget{DP0367}{}\textbf{DP0367 -- VI, p. 283 (conflict)}\par
(7.3.1.4) | (7.3.1.4), source and target\par
\textbf{Claude Opus 5.5.} Place: (7.3.1.4); Reported print: loc(m) : Fl(E)\_1 \ensuremath{\to} Fl(E)\_2; Source TeX: \textbackslash{}loc(m) : F1(E)\_1 \textbackslash{}to F1 (E)\_2; Note: The TeX has the digit 1 (F1) instead of the letters Fl, as written everywhere else on these pages.; Proposed edit: \textbackslash{}loc(m) : F1(E)\_1 \textbackslash{}to F1 (E)\_2 =\allowbreak > \textbackslash{}loc(m) : \textbackslash{}text\{Fl\}(E)\_1 \textbackslash{}to \textbackslash{}text\{Fl\} (E)\_2\par
\textbf{GPT-6.1 Sol.} Place: (7.3.1.4), source and target; Reported print: \textbackslash{}loc(m) : \textbackslash{}text\{Fl\}(E)\_1 \textbackslash{}to \textbackslash{}text\{Fl\}(E)\_2; Source TeX: \textbackslash{}loc(m) : F1(E)\_1 \textbackslash{}to F1 (E)\_2; Note: Both category names contain lowercase l in the print, not the digit 1.; Proposed edit: F1(E)\_1 \textbackslash{}to F1 (E)\_2 =\allowbreak > \textbackslash{}text\{Fl\}(E)\_1 \textbackslash{}to \textbackslash{}text\{Fl\}(E)\_2\par
the two models differ at this place\par
\par\noindent\hypertarget{DP0368}{}\textbf{DP0368 -- VI, p. 287 (one\_model\_fix)}\par
i) \ensuremath{\Rightarrow} ii), reference\par
\textbf{Claude Opus 5.5.} Place: i) \ensuremath{\Rightarrow} ii), reference; Reported print: (II 5.2); Source TeX: (II \textbackslash{}hyperref[VI.cor5.2]\{5.2\}); Note: The print refers to exposé II 5.2, but the link targets VI 5.2.; Proposed edit: (II \textbackslash{}hyperref[VI.cor5.2]\{5.2\}) =\allowbreak > (II 5.2)\par
\textbf{GPT-6.1 Sol.} nothing reported here\par
a fix from one model only\par
\par\noindent\hypertarget{DP0369}{}\textbf{DP0369 -- VI, p. 287 (one\_model\_fix)}\par
Dans cet isomorphisme ... n' : Y \ensuremath{\rightrightarrows} f\textasciicircum{}*(R)\par
\textbf{Claude Opus 5.5.} Place: Dans cet isomorphisme ... n' : Y \ensuremath{\rightrightarrows} f\textasciicircum{}*(R); Reported print: f\textasciicircum{}*(R); Source TeX: f\textasciicircum{}\{\textbackslash{}ast\} R; Note: The print has parentheses around R.; Proposed edit: f\textasciicircum{}\{\textbackslash{}ast\} R\$ =\allowbreak > f\textasciicircum{}\{\textbackslash{}ast\}(R)\$\par
\textbf{GPT-6.1 Sol.} nothing reported here\par
a fix from one model only\par
\par\noindent\hypertarget{DP0370}{}\textbf{DP0370 -- VI, p. 288 (one\_model\_fix)}\par
(II 5.3 ii)) and (II 5.2)\par
\textbf{Claude Opus 5.5.} Place: (II 5.3 ii)) and (II 5.2); Reported print: (II 5.3 ii)) ... (II 5.2); Source TeX: \textbackslash{}hyperref[VI.cor5.3]\{5.3\}, \textbackslash{}hyperref[VI.cor5.2]\{5.2\}; Note: These are references to exposé II but link to VI labels. The same applies to '(II \textbackslash{}hyperref[VI.cor5.2]\{5.2\})' later on this page; replace it by '(II 5.2)'.; Proposed edit: (II \textbackslash{}hyperref[VI.cor5.3]\{5.3\} ii)) =\allowbreak > (II 5.3 ii))\par
\textbf{GPT-6.1 Sol.} nothing reported here\par
a fix from one model only\par
\par\noindent\hypertarget{DP0371}{}\textbf{DP0371 -- VI, p. 289 (one\_model\_fix)}\par
7.4.5, avec les notations de\par
\textbf{Claude Opus 5.5.} Place: 7.4.5, avec les notations de; Reported print: (avec les notations de 7.4.3 4)); Source TeX: (avec les notations de 7.4.3.4); Note: The print cites item 4) of 7.4.3, not a number 7.4.3.4.; Proposed edit: notations de 7.4.3.4) =\allowbreak > notations de 7.4.3 4))\par
\textbf{GPT-6.1 Sol.} nothing reported here\par
a fix from one model only\par
\par\noindent\hypertarget{DP0372}{}\textbf{DP0372 -- VI, p. 290 (one\_model\_fix)}\par
une application fonctorielle en Y\par
\textbf{Claude Opus 5.5.} Place: une application fonctorielle en Y; Reported print: en Y :; Source TeX: en \$Y\$;; Note: The print has a colon introducing (7.4.5.6).; Proposed edit: fonctorielle en \$Y\$; =\allowbreak > fonctorielle en \$Y\$:\par
\textbf{GPT-6.1 Sol.} nothing reported here\par
a fix from one model only\par
\par\noindent\hypertarget{DP0373}{}\textbf{DP0373 -- VI, p. 291 (conflict)}\par
7.4.6, family of morphisms \ensuremath{\gamma}\_f F | 7.4.6, displayed family of morphisms\par
\textbf{Claude Opus 5.5.} Place: 7.4.6, family of morphisms \ensuremath{\gamma}\_f F; Reported print: f \ensuremath{\in} Fl(I); Source TeX: f\textbackslash{}in F1(I); Note: The print means Fl(I), the set of arrows of I; in this typewriter face l and 1 are hard to tell apart. The TeX gives the digit 1.; Proposed edit: f\textbackslash{}in F1(I) =\allowbreak > f\textbackslash{}in \textbackslash{}mathrm\{Fl\}(I)\par
\textbf{GPT-6.1 Sol.} Place: 7.4.6, displayed family of morphisms; Reported print: f\textbackslash{}in Fl(I); Source TeX: f\textbackslash{}in F1(I); Note: The second character of Fl is a lower-case l, not the numeral 1.; Proposed edit: f\textbackslash{}in F1(I) =\allowbreak > f\textbackslash{}in Fl(I)\par
the two models differ at this place\par
\par\noindent\hypertarget{DP0374}{}\textbf{DP0374 -- VI, p. 292 (one\_model\_fix)}\par
(7.4.9.2), lower horizontal arrow; printed on page 293\par
\textbf{Claude Opus 5.5.} nothing reported here\par
\textbf{GPT-6.1 Sol.} Place: (7.4.9.2), lower horizontal arrow; printed on page 293; Reported print: \textbackslash{}Theta\_D; Source TeX: \textbackslash{}theta\_D; Note: The printed arrow label uses capital theta.; Proposed edit: \textbackslash{}theta\_D =\allowbreak > \textbackslash{}Theta\_D\par
a fix from one model only\par
\par\noindent\hypertarget{DP0375}{}\textbf{DP0375 -- VI, p. 292 (one\_model\_fix)}\par
(7.4.9.2), lower-right category; printed on page 293\par
\textbf{Claude Opus 5.5.} nothing reported here\par
\textbf{GPT-6.1 Sol.} Place: (7.4.9.2), lower-right category; printed on page 293; Reported print: \textbackslash{}Hhom\_\{I\textasciicircum{}\{\textbackslash{}circ\}\}(I\textasciicircum{}\{\textbackslash{}circ\},(D\textasciicircum{}\{\textbackslash{}sim/I\})'); Source TeX: \textbackslash{}Hhom\_I\textasciicircum{}\{\textbackslash{}circ\}(I\textasciicircum{}\{\textbackslash{}circ\},(D\textasciicircum{}\{\textbackslash{}sim/I\})'); Note: The circle belongs to the subscript I, not to Hom.; Proposed edit: \textbackslash{}Hhom\_I\textasciicircum{}\{\textbackslash{}circ\}(I\textasciicircum{}\{\textbackslash{}circ\},(D\textasciicircum{}\{\textbackslash{}sim/I\})') =\allowbreak > \textbackslash{}Hhom\_\{I\textasciicircum{}\{\textbackslash{}circ\}\}(I\textasciicircum{}\{\textbackslash{}circ\},(D\textasciicircum{}\{\textbackslash{}sim/I\})')\par
a fix from one model only\par
\par\noindent\hypertarget{DP0376}{}\textbf{DP0376 -- VI, p. 296 (cut)}\par
End-of-page cut, source line 5151\par
\textbf{Claude Opus 5.5.} nothing reported here\par
\textbf{GPT-6.1 Sol.} Place: End-of-page cut, source line 5151; Reported print: Page 296 ends with « ce que nous supposerons désormais. »; the heading 7.4.13.5 starts on page 297.; Source TeX: \textbackslash{}paragraph\{\}\textbackslash{}etiquette[7.4.13.5]\{VI.par7.4.13.5\}; Note: The heading-number line 5151 also belongs to page 297, whereas the cut comment only identifies line 5152 as overlapping that page. No content edit is proposed.\par
a model reports that the automatic cut gives lines of this page to its neighbour, or the neighbour's lines to this page: check the lines it names against the scan of the page that prints them (that page's checker did not have their TeX)\par
\par\noindent\hypertarget{DP0377}{}\textbf{DP0377 -- VI, p. 300 (one\_model\_fix)}\par
8.1.3.2, m\_i : D ... D\ensuremath{\times}\{i\}\par
\textbf{Claude Opus 5.5.} Place: 8.1.3.2, m\_i : D ... D\ensuremath{\times}\{i\}; Reported print: D \ensuremath{\simeq} D \ensuremath{\times}\{i\}; Source TeX: D =\allowbreak  D \textbackslash{}times \textbackslash{}\{i\textbackslash{}\}; Note: The print appears to have an isomorphism sign (\ensuremath{\simeq}), as for \ensuremath{\mu}\_i above, not '=\allowbreak '.; Proposed edit: D =\allowbreak  D \textbackslash{}times \textbackslash{}\{i\textbackslash{}\} =\allowbreak > D \textbackslash{}simeq D \textbackslash{}times \textbackslash{}\{i\textbackslash{}\}\par
\textbf{GPT-6.1 Sol.} nothing reported here\par
a fix from one model only\par
\par\noindent\hypertarget{DP0378}{}\textbf{DP0378 -- VI, p. 300 (cut)}\par
End-of-page cut, source lines 5350–5351\par
\textbf{Claude Opus 5.5.} nothing reported here\par
\textbf{GPT-6.1 Sol.} Place: End-of-page cut, source lines 5350–5351; Reported print: The page ends after « un morphisme de topos ».; Source TeX: morphisme de topos \$n : D \textbackslash{}to \textbackslash{}underleftarrow\{F\}\$\textbackslash{}pageoriginale et une famille d'isomorphismes:; Note: The formula n:D\ensuremath{\to}F with the left arrow underneath F begins page 301. Its beginning is already on source line 5350, whereas the cut comment only permits line 5351 to overlap page 301. No text should be moved as a transcription edit.\par
a model reports that the automatic cut gives lines of this page to its neighbour, or the neighbour's lines to this page: check the lines it names against the scan of the page that prints them (that page's checker did not have their TeX)\par
\par\noindent\hypertarget{DP0379}{}\textbf{DP0379 -- VI, p. 301 (one\_model\_fix)}\par
Réciproquement…, reference after b\_f\par
\textbf{Claude Opus 5.5.} Place: Réciproquement…, reference after b\_f; Reported print: (8.1.3.1),; Source TeX: \textbackslash{}eqhyperref[VI.eq8.1.3.1]\{8.1.3.1.\},; Note: The TeX adds a stray period inside the formula reference.; Proposed edit: \{8.1.3.1.\} =\allowbreak > \{8.1.3.1\}\par
\textbf{GPT-6.1 Sol.} nothing reported here\par
a fix from one model only\par
\par\noindent\hypertarget{DP0380}{}\textbf{DP0380 -- VI, p. 301 (cut)}\par
Début de la page\par
\textbf{Claude Opus 5.5.} nothing reported here\par
\textbf{GPT-6.1 Sol.} Place: Début de la page; Reported print: de topos \$n : D \textbackslash{}to \textbackslash{}underleftarrow\{F\}\$ et une famille d'isomorphismes; Source TeX: \textbackslash{}underleftarrow\{F\}\$\textbackslash{}pageoriginale et une famille d'isomorphismes:; Note: La page commence par « de topos n : D \ensuremath{\to} F », et non au seul F. La coupe omet donc du texte appartenant à cette page, en amont de la première ligne fournie.\par
a model reports that the automatic cut gives lines of this page to its neighbour, or the neighbour's lines to this page: check the lines it names against the scan of the page that prints them (that page's checker did not have their TeX)\par
\par\noindent\hypertarget{DP0381}{}\textbf{DP0381 -- VI, p. 302 (one\_model\_fix)}\par
Entre (8.1.6.1) et (8.1.6.2)\par
\textbf{Claude Opus 5.5.} nothing reported here\par
\textbf{GPT-6.1 Sol.} Place: Entre (8.1.6.1) et (8.1.6.2); Reported print: (7.2.2); Source TeX: \textbackslash{}eqhyperref[VI.eq7.2.2.2]\{7.2.2.2\}; Note: Le renvoi imprimé est 7.2.2, sans dernier « .2 ».; Proposed edit: \textbackslash{}eqhyperref[VI.eq7.2.2.2]\{7.2.2.2\} =\allowbreak > (\textbackslash{}hyperref[VI.subsubsec7.2.2]\{7.2.2\})\par
a fix from one model only\par
\par\noindent\hypertarget{DP0382}{}\textbf{DP0382 -- VI, p. 303 (one\_model\_fix)}\par
first line\par
\textbf{Claude Opus 5.5.} Place: first line; Reported print: par définition de la Limite inductive (6.3); Source TeX: la limite inductive (; Note: The print capitalizes 'Limite', the defined notion; compare 'la Limite inductive' in 8.2.2.; Proposed edit: la limite inductive ( =\allowbreak > la Limite inductive (\par
\textbf{GPT-6.1 Sol.} nothing reported here\par
a fix from one model only\par
\par\noindent\hypertarget{DP0383}{}\textbf{DP0383 -- VI, p. 305 (conflict)}\par
«On a donc pr\_1\ensuremath{\circ}(n\ensuremath{\pi},p)\textasciicircum{}\{\textasciitilde{}/I\} =\allowbreak  n'» | Égalité commençant par « On a donc »\par
\textbf{Claude Opus 5.5.} Place: «On a donc pr\_1\ensuremath{\circ}(n\ensuremath{\pi},p)\textasciicircum{}\{\textasciitilde{}/I\} =\allowbreak  n'»; Reported print: pr\_1 \ensuremath{\circ} (n\ensuremath{\pi},p)\textasciicircum{}\{\textasciitilde{}/I\}; Source TeX: \textbackslash{}pr\_\{1\textasciicircum{}\{\textbackslash{}circ\}\}(n\textbackslash{}pi, p)\textasciicircum{}\{\textbackslash{}sim/I\}; Note: The composition sign has become a degree sign inside the subscript.; Proposed edit: \textbackslash{}pr\_\{1\textasciicircum{}\{\textbackslash{}circ\}\}(n\textbackslash{}pi =\allowbreak > \textbackslash{}pr\_1\textbackslash{}circ(n\textbackslash{}pi\par
\textbf{GPT-6.1 Sol.} Place: Égalité commençant par « On a donc »; Reported print: \$\textbackslash{}pr\_1 \textbackslash{}circ (n\textbackslash{}pi,p)\textasciicircum{}\{\textbackslash{}sim/I\}\$; Source TeX: \textbackslash{}pr\_\{1\textasciicircum{}\{\textbackslash{}circ\}\}(n\textbackslash{}pi, p)\textasciicircum{}\{\textbackslash{}sim/I\}; Note: Le cercle imprimé est le signe de composition après pr₁, non un exposant attaché à son indice.; Proposed edit: \textbackslash{}pr\_\{1\textasciicircum{}\{\textbackslash{}circ\}\}(n\textbackslash{}pi, p)\textasciicircum{}\{\textbackslash{}sim/I\} =\allowbreak > \textbackslash{}pr\_1 \textbackslash{}circ (n\textbackslash{}pi, p)\textasciicircum{}\{\textbackslash{}sim/I\}\par
the two models differ at this place\par
\par\noindent\hypertarget{DP0384}{}\textbf{DP0384 -- VI, p. 305 (cut)}\par
last line\par
\textbf{Claude Opus 5.5.} Place: last line; Reported print: D'après ce; Source TeX: D'après ce qui; Note: 'qui' is printed on p. 306; this page ends at 'D'après ce'.\par
\textbf{GPT-6.1 Sol.} nothing reported here\par
a model reports that the automatic cut gives lines of this page to its neighbour, or the neighbour's lines to this page: check the lines it names against the scan of the page that prints them (that page's checker did not have their TeX)\par
\par\noindent\hypertarget{DP0385}{}\textbf{DP0385 -- VI, p. 305 (cut)}\par
Fin de la page\par
\textbf{Claude Opus 5.5.} nothing reported here\par
\textbf{GPT-6.1 Sol.} Place: Fin de la page; Reported print: D'après ce; Source TeX: D'après ce qui; Note: La page imprimée s'arrête après « ce ». Le mot « qui » appartient à la page 306, sans que ce débordement soit annoncé dans le commentaire de coupe.\par
a model reports that the automatic cut gives lines of this page to its neighbour, or the neighbour's lines to this page: check the lines it names against the scan of the page that prints them (that page's checker did not have their TeX)\par
\par\noindent\hypertarget{DP0386}{}\textbf{DP0386 -- VI, p. 306 (cut)}\par
Début de la page\par
\textbf{Claude Opus 5.5.} nothing reported here\par
\textbf{GPT-6.1 Sol.} Place: Début de la page; Reported print: qui précède, le foncteur; Source TeX: précède,\textbackslash{}pageoriginale le foncteur; Note: La page imprimée commence par « qui précède ». La coupe commence un mot trop tard : vérifier « qui » dans la ligne source précédente, sans l'ajouter en double au texte.\par
a model reports that the automatic cut gives lines of this page to its neighbour, or the neighbour's lines to this page: check the lines it names against the scan of the page that prints them (that page's checker did not have their TeX)\par
\par\noindent\hypertarget{DP0387}{}\textbf{DP0387 -- VI, p. 306 (cut)}\par
Fin de la coupe, titre 8.2.6\par
\textbf{Claude Opus 5.5.} nothing reported here\par
\textbf{GPT-6.1 Sol.} Place: Fin de la coupe, titre 8.2.6; Reported print: La page finit par « du topos fibré associé. » ; le titre 8.2.6 ouvre la page 307.; Source TeX: \textbackslash{}subsubsection\{Démonstration du théorème : Réduction au cas d'une   petite catégorie d'indices\}\textbackslash{}etiquette[8.2.6]\{VI.subsubsec8.2.6\}; Note: Ce titre appartient à la page 307, contrairement à la coupe fournie, dont le commentaire n'autorise pas ce débordement. Problème signalé une seule fois pour les deux pages.\par
a model reports that the automatic cut gives lines of this page to its neighbour, or the neighbour's lines to this page: check the lines it names against the scan of the page that prints them (that page's checker did not have their TeX)\par
\par\noindent\hypertarget{DP0388}{}\textbf{DP0388 -- VI, p. 310 (conflict)}\par
proof of 8.2.9, \ensuremath{\mu}\_\{j*\} and \ensuremath{\mu}\_j | 8.2.12, domaine du foncteur \ensuremath{\mu}\_\{j*\}\par
\textbf{Claude Opus 5.5.} Place: proof of 8.2.9, \ensuremath{\mu}\_\{j*\} and \ensuremath{\mu}\_j; Reported print: \ensuremath{\mu}\_\{j*\} : (Lim\_\{I°\} F\_j)\textasciicircum{}\ensuremath{∿} \ensuremath{\longrightarrow} F\_j … \ensuremath{\mu}\_j : (Lim\_\{I°\} F\_j)\textasciicircum{}\ensuremath{∿} \ensuremath{\longrightarrow} F\_j; Source TeX: \textbackslash{}mu\_\{j\_\textbackslash{}ast\} : (\textbackslash{}varinjlim\_\{I\textasciicircum{}\{\textbackslash{}circ\}\} F\_i)\textasciicircum{}\{\textbackslash{}sim\} \textbackslash{}to F\_j; Proposed correction: (\textbackslash{}varinjlim\_\{I\textasciicircum{}\{\textbackslash{}circ\}\} F\_i)\textasciicircum{}\{\textbackslash{}sim\}; Note: Both times the print writes Lim F\_j (wrong index) inside the limit. The TeX silently corrects it to F\_i.\par
\textbf{GPT-6.1 Sol.} Place: 8.2.12, domaine du foncteur \ensuremath{\mu}\_\{j*\}; Reported print: \textbackslash{}mu\_\{j\_\textbackslash{}ast\} : (\textbackslash{}varinjlim\_\{I\textasciicircum{}\{\textbackslash{}circ\}\} F\_j)\textasciicircum{}\{\textbackslash{}sim\} \textbackslash{}to F\_j; Source TeX: \textbackslash{}mu\_\{j\_\textbackslash{}ast\} : (\textbackslash{}varinjlim\_\{I\textasciicircum{}\{\textbackslash{}circ\}\} F\_i)\textasciicircum{}\{\textbackslash{}sim\} \textbackslash{}to F\_j; Note: L'imprimé emploie l'indice muet j sous cette limite ; le TeX le remplace par i.; Proposed edit: \textbackslash{}mu\_\{j\_\textbackslash{}ast\} : (\textbackslash{}varinjlim\_\{I\textasciicircum{}\{\textbackslash{}circ\}\} F\_i) =\allowbreak > \textbackslash{}mu\_\{j\_\textbackslash{}ast\} : (\textbackslash{}varinjlim\_\{I\textasciicircum{}\{\textbackslash{}circ\}\} F\_j)\par
the two models differ at this place: one model fixes the TeX, the other calls it a misprint of the print\par
\par\noindent\hypertarget{DP0389}{}\textbf{DP0389 -- VI, p. 310 (one\_model\_fix)}\par
8.2.12, domaine du morphisme \ensuremath{\mu}\_j\par
\textbf{Claude Opus 5.5.} nothing reported here\par
\textbf{GPT-6.1 Sol.} Place: 8.2.12, domaine du morphisme \ensuremath{\mu}\_j; Reported print: \textbackslash{}mu\_j : (\textbackslash{}varinjlim\_\{I\textasciicircum{}\{\textbackslash{}circ\}\} F\_j)\textasciicircum{}\{\textbackslash{}sim\} \textbackslash{}to F\_j; Source TeX: \textbackslash{}mu\_j : (\textbackslash{}varinjlim\_\{I\textasciicircum{}\{\textbackslash{}circ\}\} F\_i)\textasciicircum{}\{\textbackslash{}sim\} \textbackslash{}to F\_j; Note: Ici encore, l'indice muet imprimé sous la limite est j, non i.; Proposed edit: \textbackslash{}mu\_j : (\textbackslash{}varinjlim\_\{I\textasciicircum{}\{\textbackslash{}circ\}\} F\_i) =\allowbreak > \textbackslash{}mu\_j : (\textbackslash{}varinjlim\_\{I\textasciicircum{}\{\textbackslash{}circ\}\} F\_j)\par
a fix from one model only\par
\par\noindent\hypertarget{DP0390}{}\textbf{DP0390 -- VI, p. 313 (one\_model\_fix)}\par
first display, index diagram of lim\par
\textbf{Claude Opus 5.5.} Place: first display, index diagram of lim; Reported print: \ensuremath{\ell}; Source TeX: l \textbackslash{}ar[ru]\textasciicircum{}\{f\}; Note: The print has script \ensuremath{\ell}, as in F\_\ensuremath{\ell}.; Proposed edit: l \textbackslash{}ar[ru]\textasciicircum{}\{f\} =\allowbreak > \textbackslash{}ell \textbackslash{}ar[ru]\textasciicircum{}\{f\}\par
\textbf{GPT-6.1 Sol.} nothing reported here\par
a fix from one model only\par
\par\noindent\hypertarget{DP0391}{}\textbf{DP0391 -- VI, p. 315 (one\_model\_fix)}\par
'Comme les familles', target of arrow\par
\textbf{Claude Opus 5.5.} Place: 'Comme les familles', target of arrow; Reported print: f*\_\ensuremath{\alpha}(Z\_\ensuremath{\alpha})); Source TeX: f\textasciicircum{}\{\textbackslash{}ast\}(Z\_\{\textbackslash{}alpha\}))\$ appartiennent; Note: The index \ensuremath{\alpha} on f* is missing in the TeX.; Proposed edit: f\textasciicircum{}\{\textbackslash{}ast\}(Z\_\{\textbackslash{}alpha\}))\$ appartiennent =\allowbreak > f\textasciicircum{}\{\textbackslash{}ast\}\_\{\textbackslash{}alpha\}(Z\_\{\textbackslash{}alpha\}))\$ appartiennent\par
\textbf{GPT-6.1 Sol.} nothing reported here\par
a fix from one model only\par
\par\noindent\hypertarget{DP0392}{}\textbf{DP0392 -- VI, p. 315 (one\_model\_fix)}\par
Dernière ligne, famille que l'on remplace\par
\textbf{Claude Opus 5.5.} nothing reported here\par
\textbf{GPT-6.1 Sol.} Place: Dernière ligne, famille que l'on remplace; Reported print: (Y\_\{\textbackslash{}alpha\}\textbackslash{}to Y)\_\{\textbackslash{}alpha\textbackslash{}in A\}; Source TeX: \$(Y\_\{\textbackslash{}alpha\} \textbackslash{}to Y) \textbackslash{}alpha \textbackslash{}in A\$; Note: Dans l'imprimé, \ensuremath{\alpha} \ensuremath{\in} A est l'indice de la famille ; le TeX le laisse sur la ligne principale.; Proposed edit: \$(Y\_\{\textbackslash{}alpha\} \textbackslash{}to Y) \textbackslash{}alpha \textbackslash{}in A\$ =\allowbreak > \$(Y\_\{\textbackslash{}alpha\} \textbackslash{}to Y)\_\{\textbackslash{}alpha \textbackslash{}in A\}\$\par
a fix from one model only\par
\par\noindent\hypertarget{DP0393}{}\textbf{DP0393 -- VI, p. 317 (one\_model\_fix)}\par
8.3.12, 'et \ensuremath{\phi}\_\ensuremath{\alpha} =\allowbreak  id\_\{\ensuremath{\pi}(Y\_\ensuremath{\alpha})\}'\par
\textbf{Claude Opus 5.5.} Place: 8.3.12, 'et \ensuremath{\phi}\_\ensuremath{\alpha} =\allowbreak  id\_\{\ensuremath{\pi}(Y\_\ensuremath{\alpha})\}'; Reported print: \ensuremath{\phi}\_\ensuremath{\alpha} (lower-case phi); Source TeX: \textbackslash{}Phi\_\{\textbackslash{}alpha\}; Note: The print appears to have a lower-case phi; the TeX has a capital Phi.; Proposed edit: \$\textbackslash{}Phi\_\{\textbackslash{}alpha\} =\allowbreak  \textbackslash{}id =\allowbreak > \$\textbackslash{}phi\_\{\textbackslash{}alpha\} =\allowbreak  \textbackslash{}id\par
\textbf{GPT-6.1 Sol.} nothing reported here\par
a fix from one model only\par
\par\noindent\hypertarget{DP0394}{}\textbf{DP0394 -- VI, p. 318 (conflict)}\par
'représentables dans Lim\_\{I°,t\} (F\_i)\_coh' | Application du lemme 8.3.4, catégorie où les produits sont représentables\par
\textbf{Claude Opus 5.5.} Place: 'représentables dans Lim\_\{I°,t\} (F\_i)\_coh'; Reported print: index I°,t; Source TeX: \textbackslash{}varprojlim\_\{I\textasciicircum{}\{\textbackslash{}circ\},F\}; Note: The index appears to be 'I°,t' (small letter), not 'I°,F'. The print is hard to read here.; Proposed edit: \textbackslash{}varprojlim\_\{I\textasciicircum{}\{\textbackslash{}circ\},F\} =\allowbreak > \textbackslash{}varprojlim\_\{I\textasciicircum{}\{\textbackslash{}circ\},t\}\par
\textbf{GPT-6.1 Sol.} Place: Application du lemme 8.3.4, catégorie où les produits sont représentables; Reported print: Lim muni d'une flèche vers la droite, devant (F\_i)\_\{\textbackslash{}coh\}; Source TeX: \textbackslash{}varprojlim\_\{I\textasciicircum{}\{\textbackslash{}circ\},F\} (F\_i)\_\{\textbackslash{}coh\}; Note: La flèche imprimée sous Lim est dirigée vers la droite : il s'agit du signe de limite inductive.; Proposed edit: \textbackslash{}varprojlim =\allowbreak > \textbackslash{}varinjlim\par
the two models differ at this place\par
\par\noindent\hypertarget{DP0395}{}\textbf{DP0395 -- VI, p. 320 (one\_model\_fix)}\par
'l'exercice I 8.2.8'\par
\textbf{Claude Opus 5.5.} Place: 'l'exercice I 8.2.8'; Reported print: I 8.2.8; Source TeX: I \textbackslash{}hyperref[VI.subsubsec8.2.8]\{8.2.8\}; Note: The print refers to exercise 8.2.8 of exposé I. The link points to VI 8.2.8, which is the wrong target.; Proposed edit: l'exercice I \textbackslash{}hyperref[VI.subsubsec8.2.8]\{8.2.8\} =\allowbreak > l'exercice I 8.2.8\par
\textbf{GPT-6.1 Sol.} nothing reported here\par
a fix from one model only\par
\par\noindent\hypertarget{DP0396}{}\textbf{DP0396 -- VI, p. 320 (cut)}\par
end of TeX: the diagram with can. and \ensuremath{\phi}\par
\textbf{Claude Opus 5.5.} Place: end of TeX: the diagram with can. and \ensuremath{\phi}; Reported print: page 320 ends with 'qui s'insère dans un diagramme commutatif à isomorphisme près'; Source TeX: \textbackslash{}[ \textbackslash{}xymatrix@R=\allowbreak 1.5cm\{ ... \} \textbackslash{}]; Note: The triangular diagram is not printed on p. 320. Page 321 begins with this diagram, before 'où \ensuremath{\phi} est le morphisme canonique'.\par
\textbf{GPT-6.1 Sol.} nothing reported here\par
a model reports that the automatic cut gives lines of this page to its neighbour, or the neighbour's lines to this page: check the lines it names against the scan of the page that prints them (that page's checker did not have their TeX)\par
\par\noindent\hypertarget{DP0397}{}\textbf{DP0397 -- VI, p. 320 (cut)}\par
Fin de page, après « à isomorphisme près »\par
\textbf{Claude Opus 5.5.} nothing reported here\par
\textbf{GPT-6.1 Sol.} Place: Fin de page, après « à isomorphisme près »; Reported print: La page finit par « qui s'insère dans un diagramme commutatif à isomorphisme près » ; aucun diagramme ne figure sur cette page.; Source TeX: \textbackslash{}xymatrix@R=\allowbreak 1.5cm\{ \textbackslash{}Loc\_\{\textbackslash{}theta(m)\} (\textbackslash{}Loc\_p(E))  \textbackslash{}ar[rr]\textasciicircum{}\{\textbackslash{}can\textbackslash{}bigdot\} \textbackslash{}ar[rd]\_\{\textbackslash{}phi\}  \&\&\textbackslash{}Loc\_\{p'\}(E)\textbackslash{}ar[ld]\textasciicircum{}\{m\}; Note: Tout le diagramme suivant appartient à la page 321, et non seulement les mots de la dernière ligne signalés par le commentaire de coupe. La page 320 s'arrête avant le display. Faire contrôler ce diagramme avec 321, sans le déplacer dans le texte source.\par
a model reports that the automatic cut gives lines of this page to its neighbour, or the neighbour's lines to this page: check the lines it names against the scan of the page that prints them (that page's checker did not have their TeX)\par
\par\noindent\hypertarget{DP0398}{}\textbf{DP0398 -- VI, p. 321 (one\_model\_fix)}\par
8.4.6, after \ensuremath{\theta}\_p(id\_p)\par
\textbf{Claude Opus 5.5.} Place: 8.4.6, after \ensuremath{\theta}\_p(id\_p); Reported print: \ensuremath{\theta}\_p(id\_p) (8.4.3.1). L'image; Source TeX: \textbackslash{}eqhyperref[VI.eq8.4.3.1]\{8.4.3.1\}).; Note: \textbackslash{}eqhyperref already supplies the parentheses, so the TeX prints '(8.4.3.1))' with an extra closing parenthesis.; Proposed edit: \textbackslash{}eqhyperref[VI.eq8.4.3.1]\{8.4.3.1\}). =\allowbreak > \textbackslash{}eqhyperref[VI.eq8.4.3.1]\{8.4.3.1\}.\par
\textbf{GPT-6.1 Sol.} nothing reported here\par
a fix from one model only\par
\par\noindent\hypertarget{DP0399}{}\textbf{DP0399 -- VI, p. 322 (one\_model\_fix)}\par
8.5.1, before (8.5.1.4)\par
\textbf{Claude Opus 5.5.} Place: 8.5.1, before (8.5.1.4); Reported print: D'après 7.4.3,4,; Source TeX: D'après 7.4.3.4,; Note: The print writes '7.4.3,4' with a comma, like '7.4.3,3' just above, which the TeX keeps.; Proposed edit: D'après 7.4.3.4 =\allowbreak > D'après 7.4.3,4\par
\textbf{GPT-6.1 Sol.} nothing reported here\par
a fix from one model only\par
\par\noindent\hypertarget{DP0400}{}\textbf{DP0400 -- VI, p. 329 (one\_model\_fix)}\par
proof of 8.5.9, « Soit alors Φ : I \ensuremath{\to} Fl(I) » and « Φ° »\par
\textbf{Claude Opus 5.5.} Place: proof of 8.5.9, « Soit alors Φ : I \ensuremath{\to} Fl(I) » and « Φ° »; Reported print: Φ (capital); Source TeX: \textbackslash{}phi; Note: The typed letter looks like a capital Φ, larger than the \ensuremath{\varphi}\_f of p. 330. If so, the following \$\textbackslash{}phi\textasciicircum{}\{\textbackslash{}circ\}\$ must also become \$\textbackslash{}Phi\textasciicircum{}\{\textbackslash{}circ\}\$.; Proposed edit: Soit alors \$\textbackslash{}phi : I =\allowbreak > Soit alors \$\textbackslash{}Phi : I\par
\textbf{GPT-6.1 Sol.} nothing reported here\par
a fix from one model only\par
\par\noindent\hypertarget{DP0401}{}\textbf{DP0401 -- VI, p. 330 (cut)}\par
Fin de page, après « Alors le morphisme de topos annelé »\par
\textbf{Claude Opus 5.5.} nothing reported here\par
\textbf{GPT-6.1 Sol.} Place: Fin de page, après « Alors le morphisme de topos annelé »; Reported print: La page finit par « Alors le morphisme de topos annelé », sans la formule définissant Q.; Source TeX: \$Q : (\textbackslash{}limtop\{I\} F\_i, \textbackslash{}varinjlim\_\{I\textasciicircum{}\{\textbackslash{}circ\}\} \textbackslash{}mu\textasciicircum{}\{\textbackslash{}ast\}\_i (A\_i)) \textbackslash{}to (\textbackslash{}Hom\_\{I\textasciicircum{}\{\textbackslash{}circ\}\} (I\textasciicircum{}\{\textbackslash{}circ\}, F'), A)\$\textbackslash{}pageoriginale; Note: La formule commençant par Q appartient entièrement à 331. La réserve du commentaire ne vise que la dernière ligne, alors que les deux lignes précédentes sont également hors de 330. Corriger le découpage des extraits, sans déplacer le texte dans la source.\par
a model reports that the automatic cut gives lines of this page to its neighbour, or the neighbour's lines to this page: check the lines it names against the scan of the page that prints them (that page's checker did not have their TeX)\par
\par\noindent\hypertarget{DP0402}{}\textbf{DP0402 -- VI, p. 331 (cut)}\par
start of page\par
\textbf{Claude Opus 5.5.} Place: start of page; Reported print: Q : (\textbackslash{}underleftarrow\{Limtop\}\_I F\_i , \textbackslash{}varinjlim\_\{I°\} \ensuremath{\mu}\_i\textasciicircum{}*(A\_i)) \ensuremath{\longrightarrow} (Hom\_\{I°\}(I°,F'),A) (8.5.1) est plat à droite; Source TeX: (I\textasciicircum{}\{\textbackslash{}circ\}, F'), A)\$\textbackslash{}pageoriginale (\textbackslash{}hyperref[VI.subsubsec8.5.1]\{8.5.1\}) est; Note: The printed page begins with the whole formula 'Q : (Limtop\_I F\_i, lim\_\{I°\} \ensuremath{\mu}\_i\textasciicircum{}*(A\_i)) \ensuremath{\to} (Hom\_\{I°\}(I°,F'),A)'. Its first part, up to 'Hom\_\{I°\}(', is in the previous source line (6897), which was cut to p. 330. Page 331 really begins at 'Q :' in that line.\par
\textbf{GPT-6.1 Sol.} nothing reported here\par
a model reports that the automatic cut gives lines of this page to its neighbour, or the neighbour's lines to this page: check the lines it names against the scan of the page that prints them (that page's checker did not have their TeX)\par
\par\noindent\hypertarget{DP0403}{}\textbf{DP0403 -- VI, p. 331 (cut)}\par
Début de page, formule définissant Q\par
\textbf{Claude Opus 5.5.} nothing reported here\par
\textbf{GPT-6.1 Sol.} Place: Début de page, formule définissant Q; Reported print: La page commence par la formule entière Q : (\textbackslash{}limtop\{I\} F\_i, \textbackslash{}varinjlim\_\{I\textasciicircum{}\{\textbackslash{}circ\}\} \textbackslash{}mu\_i\textasciicircum{}\{\textbackslash{}ast\}(A\_i)) \textbackslash{}to (\textbackslash{}operatorname\{Hom\}\_\{I\textasciicircum{}\{\textbackslash{}circ\}\}(I\textasciicircum{}\{\textbackslash{}circ\},F'),A), suivie de (8.5.1).; Source TeX: (I\textasciicircum{}\{\textbackslash{}circ\}, F'), A)\$\textbackslash{}pageoriginale; Note: La coupe ne contient que la fin de la formule qui ouvre la page imprimée. Son début se trouve en amont des lignes fournies, vraisemblablement dans la coupe de p. 330. Il faut rectifier la coupe, sans ajouter cette formule au source.\par
a model reports that the automatic cut gives lines of this page to its neighbour, or the neighbour's lines to this page: check the lines it names against the scan of the page that prints them (that page's checker did not have their TeX)\par
\par\noindent\hypertarget{DP0404}{}\textbf{DP0404 -- VI, p. 335 (cut)}\par
start of page\par
\textbf{Claude Opus 5.5.} Place: start of page; Reported print: (page 335 begins with 'Corollaire 8.7.7'); Source TeX: (\textbackslash{}hyperref[VI.cor8.5.4]\{8.5.4\}).; Note: The first line '(8.5.4).' is printed at the end of p. 334 ('compte tenu de (8.5.4).'). It belongs to p. 334, and p. 335 begins at Corollaire 8.7.7. The comment line does not flag this.\par
\textbf{GPT-6.1 Sol.} nothing reported here\par
a model reports that the automatic cut gives lines of this page to its neighbour, or the neighbour's lines to this page: check the lines it names against the scan of the page that prints them (that page's checker did not have their TeX)\par
\par\noindent\hypertarget{DP0405}{}\textbf{DP0405 -- VI, p. 337 (one\_model\_fix)}\par
Lemme 9.1, name of the sheaf\par
\textbf{Claude Opus 5.5.} nothing reported here\par
\textbf{GPT-6.1 Sol.} Place: Lemme 9.1, name of the sheaf; Reported print: \textbackslash{}Ff; Source TeX: \textbackslash{}mathfrak\{F\}; Note: The print uses the same handwritten script F as elsewhere in this argument, not a distinct Fraktur F.; Proposed edit: \textbackslash{}mathfrak\{F\} =\allowbreak > \textbackslash{}Ff\par
a fix from one model only\par
\par\noindent\hypertarget{DP0406}{}\textbf{DP0406 -- VII, p. 342 (cut)}\par
Beginning of the TeX cut, before the introductory paragraph\par
\textbf{Claude Opus 5.5.} nothing reported here\par
\textbf{GPT-6.1 Sol.} Place: Beginning of the TeX cut, before the introductory paragraph; Reported print: The page begins with « Dans le présent exposé et le suivant… »; no title or author block is printed here.; Source TeX: \textbackslash{}chapter\{Site et topos étales d'un schéma\} followed by the author block; Note: The cut starts too early: the title and author material precedes this printed page. Page 342 begins at Dans\textbackslash{}pageoriginale.\par
a model reports that the automatic cut gives lines of this page to its neighbour, or the neighbour's lines to this page: check the lines it names against the scan of the page that prints them (that page's checker did not have their TeX)\par
\par\noindent\hypertarget{DP0407}{}\textbf{DP0407 -- VII, p. 352 (one\_model\_fix)}\par
u\_* display\par
\textbf{Claude Opus 5.5.} Place: u\_* display; Reported print: u\_* : \textbackslash{}widetilde\{Sch\_\{/X\}\} \ensuremath{\to} \textbackslash{}widetilde\{X\_ét\} (no parentheses around Sch); Source TeX: \textbackslash{}widetilde\{(\textbackslash{}Sch)\_\{/X\}\}; Note: The print writes Sch\_\{/X\} under the tilde, as in the u\textasciicircum{}* display below.; Proposed edit: \textbackslash{}widetilde\{(\textbackslash{}Sch)\_\{/X\}\} =\allowbreak > \textbackslash{}widetilde\{\textbackslash{}Sch\_\{/X\}\}\par
\textbf{GPT-6.1 Sol.} nothing reported here\par
a fix from one model only\par
\par\noindent\hypertarget{DP0408}{}\textbf{DP0408 -- VII, p. 353 (conflict)}\par
second isomorphism of Prop. 4.1 (ii) | Proposition 4.1, second displayed cohomology isomorphism\par
\textbf{Claude Opus 5.5.} Place: second isomorphism of Prop. 4.1 (ii); Reported print: H\textasciicircum{}*(X\_ét,G) \ensuremath{⥲} H\textasciicircum{}*((Sch)\_\{/X\}, u\textasciicircum{}*(G)) .; Source TeX: second line repeats the first: H\textasciicircum{}\{\textbackslash{}ast\} (X\_\{\textbackslash{}text\{\textbackslash{}rm ét\}\}, F \textbackslash{}circ u) \textbackslash{}xrightarrow\{\textbackslash{}sim\} H\textasciicircum{}\{\textbackslash{}ast\}((\textbackslash{}Sch)\_\{/X\}, F); Note: The second formula is a copy of the first; the print has G on the left and u\textasciicircum{}*(G) on the right.; Proposed edit: H\textasciicircum{}\{\textbackslash{}ast\} (X\_\{\textbackslash{}text\{\textbackslash{}rm ét\}\}, F \textbackslash{}circ u) \textbackslash{}xrightarrow\{\textbackslash{}sim\} H\textasciicircum{}\{\textbackslash{}ast\}((\textbackslash{}Sch)\_\{/X\}, F)\{\textbackslash{}quoi\}. =\allowbreak > H\textasciicircum{}\{\textbackslash{}ast\} (X\_\{\textbackslash{}text\{\textbackslash{}rm ét\}\}, G) \textbackslash{}xrightarrow\{\textbackslash{}sim\} H\textasciicircum{}\{\textbackslash{}ast\}((\textbackslash{}Sch)\_\{/X\}, u\textasciicircum{}\{\textbackslash{}ast\}(G))\{\textbackslash{}quoi\}.\par
\textbf{GPT-6.1 Sol.} Place: Proposition 4.1, second displayed cohomology isomorphism; Reported print: H\textasciicircum{}\{\textbackslash{}ast\}(X\_\{\textbackslash{}text\{ét\}\},G) \textbackslash{}xrightarrow\{\textbackslash{}sim\} H\textasciicircum{}\{\textbackslash{}ast\}((\textbackslash{}Sch)\_\{/X\},u\textasciicircum{}\{\textbackslash{}ast\}(G)); Source TeX: H\textasciicircum{}\{\textbackslash{}ast\} (X\_\{\textbackslash{}text\{\textbackslash{}rm ét\}\}, F \textbackslash{}circ u) \textbackslash{}xrightarrow\{\textbackslash{}sim\} H\textasciicircum{}\{\textbackslash{}ast\}((\textbackslash{}Sch)\_\{/X\}, F)\{\textbackslash{}quoi\}.; Note: The TeX duplicates the first formula. The printed second formula has G on the left and u\textasciicircum{}\{\textbackslash{}ast\}(G) on the right.; Proposed edit: H\textasciicircum{}\{\textbackslash{}ast\} (X\_\{\textbackslash{}text\{\textbackslash{}rm ét\}\}, F \textbackslash{}circ u) \textbackslash{}xrightarrow\{\textbackslash{}sim\} H\textasciicircum{}\{\textbackslash{}ast\}((\textbackslash{}Sch)\_\{/X\}, F)\{\textbackslash{}quoi\}. =\allowbreak > H\textasciicircum{}\{\textbackslash{}ast\} (X\_\{\textbackslash{}text\{\textbackslash{}rm ét\}\}, G) \textbackslash{}xrightarrow\{\textbackslash{}sim\} H\textasciicircum{}\{\textbackslash{}ast\}((\textbackslash{}Sch)\_\{/X\}, u\textasciicircum{}\{\textbackslash{}ast\} (G))\{\textbackslash{}quoi\}.\par
the two models differ at this place\par
\par\noindent\hypertarget{DP0409}{}\textbf{DP0409 -- VII, p. 353 (one\_model\_fix)}\par
4.2, end\par
\textbf{Claude Opus 5.5.} Place: 4.2, end; Reported print: faisceaux sur (Sch)\_\{/X\} envisagés au N° 2; Source TeX: \$(\textbackslash{}Sch)/X\$; Note: The print sets /X as a subscript, as elsewhere on the page.; Proposed edit: \$(\textbackslash{}Sch)/X\$ =\allowbreak > \$(\textbackslash{}Sch)\_\{/X\}\$\par
\textbf{GPT-6.1 Sol.} nothing reported here\par
a fix from one model only\par
\par\noindent\hypertarget{DP0410}{}\textbf{DP0410 -- VII, p. 353 (one\_model\_fix)}\par
Section 4.2, name of the fpqc topology\par
\textbf{Claude Opus 5.5.} nothing reported here\par
\textbf{GPT-6.1 Sol.} Place: Section 4.2, name of the fpqc topology; Reported print: fidèlement plate quasi-compacte; Source TeX: fidèlement plat quasi-compacte; Note: The print has « plate », agreeing with « topologie ».; Proposed edit: fidèlement plat quasi-compacte =\allowbreak > fidèlement plate quasi-compacte\par
a fix from one model only\par
\par\noindent\hypertarget{DP0411}{}\textbf{DP0411 -- VII, p. 356 (one\_model\_fix)}\par
after (*)\par
\textbf{Claude Opus 5.5.} Place: after (*); Reported print: TDTE I, B, 1.1, (dans FGA; Source TeX: \textbackslash{}hyperref[VII.subsec1.1]\{1.1\}. (dans FGA; Note: The print has a comma after 1.1. Also, the link points to VII 1.1, but this is the external reference TDTE I B 1.1.; Proposed edit: \{1.1\}. (dans FGA =\allowbreak > \{1.1\}, (dans FGA\par
\textbf{GPT-6.1 Sol.} nothing reported here\par
a fix from one model only\par
\par\noindent\hypertarget{DP0412}{}\textbf{DP0412 -- VII, p. 363 (conflict)}\par
line 5 | Démonstration de 5.9, renvoi après F\ensuremath{\infty}\par
\textbf{Claude Opus 5.5.} Place: line 5; Reported print: le faisceau F\_\ensuremath{\infty} de 5.7 est isomorphe; Source TeX: de 5,7 est; Note: The TeX has a comma for the point in the number; the print reads 5.7.; Proposed edit: de 5,7 est =\allowbreak > de \textbackslash{}hyperref[VII.thm5.7]\{5.7\} est\par
\textbf{GPT-6.1 Sol.} Place: Démonstration de 5.9, renvoi après F\ensuremath{\infty}; Reported print: de 5.7; Source TeX: de 5,7; Note: Le numéro imprimé du théorème est 5.7.; Proposed edit: 5,7 =\allowbreak > 5.7\par
the two models differ at this place\par
\par\noindent\hypertarget{DP0413}{}\textbf{DP0413 -- VII, p. 363 (conflict)}\par
paragraph before 5.11 | Paragraphe précédant le corollaire 5.11\par
\textbf{Claude Opus 5.5.} Place: paragraph before 5.11; Reported print: les variantes de 5.8 et 5.9 en termes de foncteurs; Source TeX: et 5,9 en termes; Note: The TeX has a comma for the point in the number; the print reads 5.9.; Proposed edit: et 5,9 en termes =\allowbreak > et \textbackslash{}hyperref[VII.cor5.9]\{5.9\} en termes\par
\textbf{GPT-6.1 Sol.} Place: Paragraphe précédant le corollaire 5.11; Reported print: variantes de 5.8 et 5.9; Source TeX: variantes de \textbackslash{}hyperref[VII.cor5.8]\{5.8\} et 5,9; Note: Le numéro imprimé du corollaire est 5.9.; Proposed edit: 5,9 =\allowbreak > 5.9\par
the two models differ at this place\par
\par\noindent\hypertarget{DP0414}{}\textbf{DP0414 -- VIII, p. 367 (cut)}\par
Début de la coupe, avant la section 1\par
\textbf{Claude Opus 5.5.} nothing reported here\par
\textbf{GPT-6.1 Sol.} Place: Début de la coupe, avant la section 1; Reported print: La page commence par « 1. Invariance topologique du topos étale », sans titre d'exposé ni nom d'auteur.; Source TeX: \textbackslash{}chapter\{Foncteurs fibres, supports, étude cohomologique des morphismes finis\} et bloc « A. Grothendieck »; Note: La coupe comprend un bloc de titre antérieur à cette page, sans que son commentaire initial le signale. Le contenu de p. 367 commence à la section 1. Corriger l'attribution de page, sans supprimer ce bloc du fichier source.\par
a model reports that the automatic cut gives lines of this page to its neighbour, or the neighbour's lines to this page: check the lines it names against the scan of the page that prints them (that page's checker did not have their TeX)\par
\par\noindent\hypertarget{DP0415}{}\textbf{DP0415 -- VIII, p. 370 (one\_model\_fix)}\par
Fin de la proposition 2.1, « topos classifiant »\par
\textbf{Claude Opus 5.5.} nothing reported here\par
\textbf{GPT-6.1 Sol.} Place: Fin de la proposition 2.1, « topos classifiant »; Reported print: B\_\ensuremath{\pi}.) — le point termine la phrase; Source TeX: \$\textbackslash{}Bb\_\{\textbackslash{}pi\textasciicircum{}\{\textbackslash{}bigdot\}\}\$); Note: Le point imprimé est la ponctuation finale, non un exposant de \ensuremath{\pi} dans l'indice de B.; Proposed edit: \$\textbackslash{}Bb\_\{\textbackslash{}pi\textasciicircum{}\{\textbackslash{}bigdot\}\}\$) =\allowbreak > \$\textbackslash{}Bb\_\{\textbackslash{}pi\}\$.)\par
a fix from one model only\par
\par\noindent\hypertarget{DP0416}{}\textbf{DP0416 -- VIII, p. 373 (agreed\_fix\_held)}\par
3.4, display after 'isomorphisme fonctoriel' | 3.4, displayed isomorphism between the fibres\par
\textbf{Claude Opus 5.5.} Place: 3.4, display after 'isomorphisme fonctoriel'; Reported print: F\_ξ \ensuremath{⥲} F\_ξ' (arrow with tilde); Source TeX: F\_\{\textbackslash{}xi\} \textbackslash{}simeq F\_\{\textbackslash{}xi'\}; Note: The print has an isomorphism arrow, not the sign \ensuremath{\simeq}.; Proposed edit: F\_\{\textbackslash{}xi\} \textbackslash{}simeq F\_\{\textbackslash{}xi'\}\{\textbackslash{}quoi\}. =\allowbreak > F\_\{\textbackslash{}xi\} \textbackslash{}xrightarrow\{\textbackslash{}sim\} F\_\{\textbackslash{}xi'\}\{\textbackslash{}quoi\}.\par
\textbf{GPT-6.1 Sol.} Place: 3.4, displayed isomorphism between the fibres; Reported print: F\_\{\textbackslash{}xi\} \textbackslash{}xrightarrow\{\textbackslash{}sim\} F\_\{\textbackslash{}xi'\}; Source TeX: F\_\{\textbackslash{}xi\} \textbackslash{}simeq F\_\{\textbackslash{}xi'\}; Note: The print has a right-pointing isomorphism arrow, not a \textbackslash{}simeq sign.; Proposed edit: F\_\{\textbackslash{}xi\} \textbackslash{}simeq F\_\{\textbackslash{}xi'\} =\allowbreak > F\_\{\textbackslash{}xi\} \textbackslash{}xrightarrow\{\textbackslash{}sim\} F\_\{\textbackslash{}xi'\}\par
both models give the same fix, held: a model marks it low confidence\par
\par\noindent\hypertarget{DP0417}{}\textbf{DP0417 -- VIII, p. 373 (one\_model\_fix)}\par
3.4, sentence beginning « Notons d'ailleurs »\par
\textbf{Claude Opus 5.5.} nothing reported here\par
\textbf{GPT-6.1 Sol.} Place: 3.4, sentence beginning « Notons d'ailleurs »; Reported print: que si \$\textbackslash{}xi'\$ est; Source TeX: que si \$\textbackslash{}xi\$ est; Note: The point following « si » has a prime in the print.; Proposed edit: que si \$\textbackslash{}xi\$ est =\allowbreak > que si \$\textbackslash{}xi'\$ est\par
a fix from one model only\par
\par\noindent\hypertarget{DP0418}{}\textbf{DP0418 -- VIII, p. 376 (one\_model\_fix)}\par
first sentence, target of \ensuremath{\phi}\par
\textbf{Claude Opus 5.5.} Place: first sentence, target of \ensuremath{\phi}; Reported print: \ensuremath{\phi} : \textbackslash{}widetilde\{X\_\{ét\}\} \ensuremath{\to} \textbackslash{}widehat\{ξ\_\{ét\}\} (the ξ has an arc, like the hat of î : \ensuremath{\tilde{X}} \ensuremath{\to} \ensuremath{\hat{X}} two lines below); Source TeX: \textbackslash{}widetilde\{X\_\{\textbackslash{}text\{ét\}\}\} \textbackslash{}to \textbackslash{}widetilde\{\textbackslash{}xi\_\{\textbackslash{}text\{ét\}\}\}\$ le foncteur « image réciproque »; Note: The print appears to put a hat on ξ\_ét here, not a tilde. This may be a misprint of the print, which the TeX would then correct silently; check it.; Proposed edit: \textbackslash{}widetilde\{X\_\{\textbackslash{}text\{ét\}\}\} \textbackslash{}to \textbackslash{}widetilde\{\textbackslash{}xi\_\{\textbackslash{}text\{ét\}\}\}\$ le foncteur «~image réciproque~» =\allowbreak > \textbackslash{}widetilde\{X\_\{\textbackslash{}text\{ét\}\}\} \textbackslash{}to \textbackslash{}widehat\{\textbackslash{}xi\_\{\textbackslash{}text\{ét\}\}\}\$ le foncteur « image réciproque »\par
\textbf{GPT-6.1 Sol.} nothing reported here\par
a fix from one model only\par
\par\noindent\hypertarget{DP0419}{}\textbf{DP0419 -- VIII, p. 377 (one\_model\_fix)}\par
3.11, introduction of the morphism i\_x\par
\textbf{Claude Opus 5.5.} nothing reported here\par
\textbf{GPT-6.1 Sol.} Place: 3.11, introduction of the morphism i\_x; Reported print: pour le morphisme structurel habituel; Source TeX: pour le morphisme structural habituel; Note: The printed adjective is « structurel », not « structural ».; Proposed edit: morphisme structural habituel =\allowbreak > morphisme structurel habituel\par
a fix from one model only\par
\par\noindent\hypertarget{DP0420}{}\textbf{DP0420 -- VIII, p. 382 (conflict)}\par
after (*), 'D'après la transitivité…' | Between (*) and (**), after the local ring at x'\par
\textbf{Claude Opus 5.5.} Place: after (*), 'D'après la transitivité…'; Reported print: par O\_\{X',x'\}, ou x' est l'image de ξ; Source TeX: par \$\textbackslash{}Oo\_\{X', x'\}\$, ou \$x'\$ est l'image; Proposed correction: où \$x'\$ est l'image; Note: The print has 'ou' for 'où'; the TeX keeps it.\par
\textbf{GPT-6.1 Sol.} Place: Between (*) and (**), after the local ring at x'; Reported print: où x' est l'image; Source TeX: ou \$x'\$ est l'image; Note: The print has the grave accent on « où ».; Proposed edit: ou \$x'\$ est l'image =\allowbreak > où \$x'\$ est l'image\par
the two models differ at this place: one model fixes the TeX, the other calls it a misprint of the print\par
\par\noindent\hypertarget{DP0421}{}\textbf{DP0421 -- VIII, p. 384 (conflict)}\par
(***), index of the limit | (***), last condition under the projective limit\par
\textbf{Claude Opus 5.5.} Place: (***), index of the limit; Reported print: X' ξ-ponctué (ξ on the line, not as subscript); Source TeX: X'\_\{\textbackslash{}xi\}\textbackslash{}mbox\{\textbackslash{}rm - ponctué\}; Note: The print means 'X' ξ-ponctué'; ξ is not a subscript of X'.; Proposed edit: X'\_\{\textbackslash{}xi\}\textbackslash{}mbox\{\textbackslash{}rm - ponctué\} =\allowbreak > X'\textbackslash{} \textbackslash{}xi\textbackslash{}mbox\{\textbackslash{}rm -ponctué\}\par
\textbf{GPT-6.1 Sol.} Place: (***), last condition under the projective limit; Reported print: X' ξ-ponctué; Source TeX: X'\_\{\textbackslash{}xi\}\textbackslash{}mbox\{\textbackslash{}rm - ponctué\}; Note: The condition says that X' is ξ-punctured; ξ belongs to the qualifier « ξ-ponctué », rather than indexing X'.; Proposed edit: X'\_\{\textbackslash{}xi\}\textbackslash{}mbox\{\textbackslash{}rm - ponctué\} =\allowbreak > X'\textbackslash{} \textbackslash{}xi\textbackslash{}mbox\{\textbackslash{}rm - ponctué\}\par
the two models differ at this place\par
\par\noindent\hypertarget{DP0422}{}\textbf{DP0422 -- VIII, p. 393 (one\_model\_fix)}\par
6.5, 'la description IV 3.3'\par
\textbf{Claude Opus 5.5.} Place: 6.5, 'la description IV 3.3'; Reported print: IV 3.3; Source TeX: IV \textbackslash{}hyperref[VIII.def3.3]\{3.3\}; Note: The print refers to IV 3.3, but the TeX label points to VIII 3.3, a wrong target.; Proposed edit: description IV \textbackslash{}hyperref[VIII.def3.3]\{3.3\} =\allowbreak > description IV 3.3\par
\textbf{GPT-6.1 Sol.} nothing reported here\par
a fix from one model only\par
\par\noindent\hypertarget{DP0423}{}\textbf{DP0423 -- VIII, p. 393 (conflict)}\par
6.6, first line | 6.6, first sentence\par
\textbf{Claude Opus 5.5.} Place: 6.6, first line; Reported print: des développements qui précèdent; Source TeX: qui précédent; Note: The print has the verb précèdent (grave accent); the TeX has the adjective précédent.; Proposed edit: qui précédent =\allowbreak > qui précèdent\par
\textbf{GPT-6.1 Sol.} Place: 6.6, first sentence; Reported print: qui précédent; Source TeX: qui précédent; Proposed correction: qui précèdent; Note: The verb requires « précèdent », with a grave accent on the second e. The TeX keeps the printed spelling.\par
the two models differ at this place: one model fixes the TeX, the other calls it a misprint of the print\par
\par\noindent\hypertarget{DP0424}{}\textbf{DP0424 -- VIII, p. 393 (conflict)}\par
6.6, 'Ainsi, on a' | 6.6, identification of the local cohomology functor\par
\textbf{Claude Opus 5.5.} Place: 6.6, 'Ainsi, on a'; Reported print: H\textasciicircum{}q\_Y =\allowbreak  R\textasciicircum{}q i\textasciicircum{}!; Source TeX: R\textasciicircum{}\{q\_\{y\}!\}; Note: The TeX garbles R\textasciicircum{}q i\textasciicircum{}! into R\textasciicircum{}\{q\_y!\}.; Proposed edit: R\textasciicircum{}\{q\_\{y\}!\} =\allowbreak > R\textasciicircum{}\{q\}i\textasciicircum{}\{!\}\par
\textbf{GPT-6.1 Sol.} Place: 6.6, identification of the local cohomology functor; Reported print: \$\ensuremath{\mathcal{H}}\textasciicircum{}\{q\}\_Y =\allowbreak  R\textasciicircum{}q i\textasciicircum{}\{!\}\$; Source TeX: \$\ensuremath{\mathcal{H}}\textasciicircum{}\{q\}\_Y  =\allowbreak  R\textasciicircum{}\{q\_\{y\}!\}\$; Note: The print gives the q-th derived functor of i!, not an exponent containing q with subscript y.; Proposed edit: R\textasciicircum{}\{q\_\{y\}!\} =\allowbreak > R\textasciicircum{}q i\textasciicircum{}\{!\}\par
the two models differ at this place\par
\par\noindent\hypertarget{DP0425}{}\textbf{DP0425 -- VIII, p. 393 (one\_model\_fix)}\par
6.5, opening reference to the residual topos\par
\textbf{Claude Opus 5.5.} nothing reported here\par
\textbf{GPT-6.1 Sol.} Place: 6.5, opening reference to the residual topos; Reported print: \$\ensuremath{\mathcal{E}}\_\{\textbackslash{}complement U\}\$; Source TeX: \$\ensuremath{\mathcal{E}}\_\textbackslash{}complement U\$; Note: Both the complement sign and U belong to the subscript in the print.; Proposed edit: \$\ensuremath{\mathcal{E}}\_\textbackslash{}complement U\$ =\allowbreak > \$\ensuremath{\mathcal{E}}\_\{\textbackslash{}complement U\}\$\par
a fix from one model only\par
\par\noindent\hypertarget{DP0426}{}\textbf{DP0426 -- VIII, p. 397 (one\_model\_fix)}\par
7.8, 'du topos E'\par
\textbf{Claude Opus 5.5.} Place: 7.8, 'du topos E'; Reported print: du topos \ensuremath{\underline{E}} (underlined E, the topos letter); Source TeX: topos \$E\$, dont; Note: Elsewhere in this paragraph the underlined E is rendered \ensuremath{\mathcal{E}}; here the print also underlines it.; Proposed edit: topos \$E\$, dont =\allowbreak > topos \$\ensuremath{\mathcal{E}}\$, dont\par
\textbf{GPT-6.1 Sol.} nothing reported here\par
a fix from one model only\par
\par\noindent\hypertarget{DP0427}{}\textbf{DP0427 -- VIII, p. 401 (conflict)}\par
first sentence, 'sous-catégorie pleine de X\_ét' | Première phrase, catégorie à laquelle on restreint \ensuremath{\phi}\par
\textbf{Claude Opus 5.5.} Place: first sentence, 'sous-catégorie pleine de X\_ét'; Reported print: X\_ét without tilde (the topos of the étale sheaves is meant); Source TeX: \textbackslash{}widetilde\{X\_\{\textbackslash{}text\{ét\}\}\}; Proposed correction: \textbackslash{}widetilde\{X\_\{\textbackslash{}text\{ét\}\}\}; Note: The print lacks the tilde here; 'les ouverts de ce topos' requires the topos, and the tilde is printed later on the page ('Soit F \ensuremath{\in} Ob \ensuremath{\tilde{X}}\_ét'). The TeX silently corrects it.\par
\textbf{GPT-6.1 Sol.} Place: Première phrase, catégorie à laquelle on restreint \ensuremath{\phi}; Reported print: \$X\_\{\textbackslash{}text\{ét\}\}\$; Source TeX: \$\textbackslash{}widetilde\{X\_\{\textbackslash{}text\{ét\}\}\}\$; Note: L'imprimé ne porte pas de tilde ici, contrairement à l'occurrence après « Soit F » plus bas.; Proposed edit: \$\textbackslash{}widetilde\{X\_\{\textbackslash{}text\{ét\}\}\}\$ formée des =\allowbreak > \$X\_\{\textbackslash{}text\{ét\}\}\$ formée des\par
the two models differ at this place: one model fixes the TeX, the other calls it a misprint of the print\par
\par\noindent\hypertarget{DP0428}{}\textbf{DP0428 -- VIII, p. 401 (cut)}\par
Coupure entre les pages 401 et 402\par
\textbf{Claude Opus 5.5.} nothing reported here\par
\textbf{GPT-6.1 Sol.} Place: Coupure entre les pages 401 et 402; Reported print: La page 401 se termine par « et un foncteur fibre \$\textbackslash{}mathcal\{E\}\_\{\textbackslash{}xi\}\$. » ; la page 402 commence par « Pour tout faisceau F ».; Source TeX: L'extrait 401 s'arrête à « et un foncteur fibre » ; l'extrait 402 commence par « \$\textbackslash{}mathcal\{E\}\_\{\textbackslash{}xi\}\$. Pour\textbackslash{}pageoriginale ».; Note: Le nom du foncteur et le point appartiennent à la page 401. Ajuster la coupure des extraits, sans ajouter de texte au source.\par
a model reports that the automatic cut gives lines of this page to its neighbour, or the neighbour's lines to this page: check the lines it names against the scan of the page that prints them (that page's checker did not have their TeX)\par
\par\noindent\hypertarget{DP0429}{}\textbf{DP0429 -- VIII, p. 402 (agreed\_fix\_held)}\par
first display | Première formule, expression de la fibre\par
\textbf{Claude Opus 5.5.} Place: first display; Reported print: lim\ensuremath{\to}\_\{C\_\ensuremath{\phi}\} F(X'); Source TeX: \textbackslash{}varinjlim\_\{C\_\{\textbackslash{}varphi\}\} (F) (X'); Note: The print has no parentheses around F.; Proposed edit: \textbackslash{}varinjlim\_\{C\_\{\textbackslash{}varphi\}\} (F) (X') =\allowbreak > \textbackslash{}varinjlim\_\{C\_\{\textbackslash{}varphi\}\} F(X')\par
\textbf{GPT-6.1 Sol.} Place: Première formule, expression de la fibre; Reported print: \$\textbackslash{}varinjlim\_\{C\_\{\textbackslash{}varphi\}\} F(X')\$; Source TeX: \$\textbackslash{}varinjlim\_\{C\_\{\textbackslash{}varphi\}\} (F)(X')\$; Note: L'imprimé n'a pas de parenthèses autour du seul F.; Proposed edit: \textbackslash{}varinjlim\_\{C\_\{\textbackslash{}varphi\}\} (F) (X') =\allowbreak > \textbackslash{}varinjlim\_\{C\_\{\textbackslash{}varphi\}\} F(X')\par
both models give the same fix, held: a model marks it low confidence\par
\par\noindent\hypertarget{DP0430}{}\textbf{DP0430 -- VIII, p. 409 (one\_model\_fix)}\par
Grand diagramme, première ligne, entre les deuxième et troisième colonnes\par
\textbf{Claude Opus 5.5.} nothing reported here\par
\textbf{GPT-6.1 Sol.} Place: Grand diagramme, première ligne, entre les deuxième et troisième colonnes; Reported print: \$S(F')\textbackslash{}rightrightarrows S(F'')\$; Source TeX: S(F') \textbackslash{}ar[r]\textbackslash{}ar[d]; Note: L'imprimé donne deux flèches horizontales parallèles; le TeX n'en donne qu'une.; Proposed edit: S(F') \textbackslash{}ar[r]\textbackslash{}ar[d] =\allowbreak > S(F') \textbackslash{}ar@<3pt>[r] \textbackslash{}ar@<-3pt>[r]\textbackslash{}ar[d]\par
a fix from one model only\par
\par\noindent\hypertarget{DP0431}{}\textbf{DP0431 -- IX, p. 1 (one\_model\_fix)}\par
after the title, before 0. Introduction\par
\textbf{Claude Opus 5.5.} Place: after the title, before 0. Introduction; Reported print: Sommaire : 0. Introduction / 1. Le sorite des faisceaux de torsion / 2. Faisceaux constructibles / 3. Théories de Kummer et d'Artin-Schreier / 4. Cas d'une courbe algébrique / 5. La méthode de la trace.; Source TeX: (absent); Note: The printed table of contents of the exposé is not in the TeX. It is possibly omitted on purpose if the edition generates a table of contents.; Proposed edit: \textbackslash{}section\{Introduction\}\textbackslash{}etiquette[0]\{IX.sec0\}\textbackslash{}pageoriginale =\allowbreak > \textbackslash{}noindent\textbackslash{}emph\{Sommaire\}: 0. Introduction 1. Le sorite des faisceaux de torsion 2. Faisceaux constructibles 3. Théories de Kummer et d'Artin-Schreier 4. Cas d'une courbe algébrique 5. La méthode de la trace.  \textbackslash{}section\{Introduction\}\textbackslash{}etiquette[0]\{IX.sec0\}\textbackslash{}pageoriginale\par
\textbf{GPT-6.1 Sol.} nothing reported here\par
a fix from one model only\par
\par\noindent\hypertarget{DP0432}{}\textbf{DP0432 -- IX, p. 1 (one\_model\_fix)}\par
Bloc d'auteur\par
\textbf{Claude Opus 5.5.} nothing reported here\par
\textbf{GPT-6.1 Sol.} Place: Bloc d'auteur; Reported print: par M. ARTIN; Source TeX: M. Artin; Note: Le mot « par » manque devant le nom de l'auteur.; Proposed edit: \textbackslash{}begin\{center\} M. Artin =\allowbreak > \textbackslash{}begin\{center\} par M. Artin\par
a fix from one model only\par
\par\noindent\hypertarget{DP0433}{}\textbf{DP0433 -- IX, p. 1 (one\_model\_fix)}\par
Entre le bloc d'auteur et l'introduction\par
\textbf{Claude Opus 5.5.} nothing reported here\par
\textbf{GPT-6.1 Sol.} Place: Entre le bloc d'auteur et l'introduction; Reported print: Sommaire : 0. Introduction 1. Le sorite des faisceaux de torsion 2. Faisceaux constructibles 3. Théories de Kummer et d'Artin-Schreier 4. Cas d'une courbe algébrique 5. La méthode de la trace.; Source TeX: \textbackslash{}setcounter\{section\}\{-1\} \textbackslash{}section\{Introduction\}\textbackslash{}etiquette[0]\{IX.sec0\}\textbackslash{}pageoriginale; Note: Le sommaire et ses six entrées sont absents du TeX fourni.; Proposed edit: \textbackslash{}setcounter\{section\}\{-1\} =\allowbreak > \textbackslash{}noindent Sommaire : \textbackslash{}begin\{enumerate\} \textbackslash{}item[0.] Introduction \textbackslash{}item[1.] Le sorite des faisceaux de torsion \textbackslash{}item[2.] Faisceaux constructibles \textbackslash{}item[3.] Théories de Kummer et d'Artin-Schreier \textbackslash{}item[4.] Cas d'une courbe algébrique \textbackslash{}item[5.] La méthode de la trace. \textbackslash{}end\{enumerate\}  \textbackslash{}setcounter\{section\}\{-1\}\par
a fix from one model only\par
\par\noindent\hypertarget{DP0434}{}\textbf{DP0434 -- IX, p. 10 (one\_model\_fix)}\par
proof of 2.4, «d'après la définition de constructibilité»\par
\textbf{Claude Opus 5.5.} Place: proof of 2.4, «d'après la définition de constructibilité»; Reported print: d'après la définition de constructibilité; Source TeX: d'après la définition de la constructibilité; Note: The TeX adds an article «la» that the print does not have.; Proposed edit: la définition de la constructibilité =\allowbreak > la définition de constructibilité\par
\textbf{GPT-6.1 Sol.} nothing reported here\par
a fix from one model only\par
\par\noindent\hypertarget{DP0435}{}\textbf{DP0435 -- IX, p. 19 (one\_model\_fix)}\par
Premier indice du tuple donnant les composantes de s\_i\par
\textbf{Claude Opus 5.5.} nothing reported here\par
\textbf{GPT-6.1 Sol.} Place: Premier indice du tuple donnant les composantes de s\_i; Reported print: s\_\{i1\}; Source TeX: s\_\{il\}; Note: Le premier terme est celui d'indice j=\allowbreak 1, non celui d'indice l. Le dessin du chiffre 1 est proche d'un l dans cette fonte de machine.; Proposed edit: s\_\{il\} =\allowbreak > s\_\{i1\}\par
a fix from one model only\par
\par\noindent\hypertarget{DP0436}{}\textbf{DP0436 -- IX, p. 28 (conflict)}\par
3.3.1, definition of D | Définition du conoyau D, avant le corollaire 3.4\par
\textbf{Claude Opus 5.5.} Place: 3.3.1, definition of D; Reported print: faisceau dès diviseurs de Cartier (accent on 'des' in the typescript); Source TeX: \textbackslash{}emph\{faisceau dés diviseurs de Cartier\}; Proposed correction: faisceau des diviseurs de Cartier; Note: The typescript puts a stray accent on 'des'. The TeX keeps an accented 'dés'. Which accent the print has cannot be read with certainty.\par
\textbf{GPT-6.1 Sol.} Place: Définition du conoyau D, avant le corollaire 3.4; Reported print: faisceau des diviseurs de Cartier; Source TeX: faisceau dés diviseurs de Cartier; Note: Le mot « des » ne porte pas d'accent dans le scan.; Proposed edit: faisceau dés diviseurs de Cartier =\allowbreak > faisceau des diviseurs de Cartier\par
the two models differ at this place: one model fixes the TeX, the other calls it a misprint of the print\par
\par\noindent\hypertarget{DP0437}{}\textbf{DP0437 -- IX, p. 30 (conflict)}\par
Lemme 4.1 (iii) | Lemme 4.1 (iii), second facteur dans la somme directe\par
\textbf{Claude Opus 5.5.} Place: Lemme 4.1 (iii); Reported print: F ≈ \ensuremath{\oplus}\_x i\_\{x*\} i\_x\textasciicircum{}* F; Source TeX: i\_\{x\textasciicircum{}\{\textbackslash{}ast\}\}i\_\{\textbackslash{}ast\}\{\}\textasciicircum{}\{\textbackslash{}ast\}F; Note: The second factor is i\_x\textasciicircum{}* in the print; the TeX has i\_* followed by a lone *.; Proposed edit: i\_\{\textbackslash{}ast\}\{\}\textasciicircum{}\{\textbackslash{}ast\}F =\allowbreak > i\_x\textasciicircum{}\{\textbackslash{}ast\}F\par
\textbf{GPT-6.1 Sol.} Place: Lemme 4.1 (iii), second facteur dans la somme directe; Reported print: i\_x\textasciicircum{}\{\textbackslash{}ast\}F; Source TeX: i\_\{\textbackslash{}ast\}\{\}\textasciicircum{}\{\textbackslash{}ast\}F; Note: Le second i a pour indice x, non une étoile : il s'agit de l'image inverse par i\_x.; Proposed edit: i\_\{\textbackslash{}ast\}\{\}\textasciicircum{}\{\textbackslash{}ast\}F =\allowbreak > i\_x\{\}\textasciicircum{}\{\textbackslash{}ast\}F\par
the two models differ at this place\par
\par\noindent\hypertarget{DP0438}{}\textbf{DP0438 -- IX, p. 34 (one\_model\_fix)}\par
Remark after 4.3.1, Tsen reference\par
\textbf{Claude Opus 5.5.} Place: Remark after 4.3.1, Tsen reference; Reported print: [6, th. 4.3]; Source TeX: [6, th.\textasciitilde{}\textbackslash{}hyperref[IX.subsec4.3]\{4.3\}]; Note: The theorem 4.3 is in reference [6], not in this exposé; the TeX links it to IX 4.3. The printed number is the same.; Proposed edit: th.\textasciitilde{}\textbackslash{}hyperref[IX.subsec4.3]\{4.3\} =\allowbreak > th.\textasciitilde{}4.3\par
\textbf{GPT-6.1 Sol.} nothing reported here\par
a fix from one model only\par
\par\noindent\hypertarget{DP0439}{}\textbf{DP0439 -- IX, p. 40 (one\_model\_fix)}\par
first paragraph: « d'un \ensuremath{\ell}-groupe h »\par
\textbf{Claude Opus 5.5.} Place: first paragraph: « d'un \ensuremath{\ell}-groupe h »; Reported print: d'un \ensuremath{\ell}-groupe; Source TeX: d'un \$\textbackslash{}ell\$ groupe; Note: The print has a hyphen here.; Proposed edit: d'un \$\textbackslash{}ell\$ groupe =\allowbreak > d'un \$\textbackslash{}ell\$-groupe\par
\textbf{GPT-6.1 Sol.} nothing reported here\par
a fix from one model only\par
\par\noindent\hypertarget{DP0440}{}\textbf{DP0440 -- IX, p. 41 (one\_model\_fix)}\par
Opening vanishing statement, first argument of H\textasciicircum{}q\par
\textbf{Claude Opus 5.5.} nothing reported here\par
\textbf{GPT-6.1 Sol.} Place: Opening vanishing statement, first argument of H\textasciicircum{}q; Reported print: H\textasciicircum{}q(Z,\textbackslash{}ZZ/\textbackslash{}ell)=\allowbreak 0; Source TeX: H\textasciicircum{}q(\textbackslash{}ZZ, \textbackslash{}ZZ/\textbackslash{}ell)=\allowbreak 0; Note: The first Z is the ordinary letter naming the closed subscheme, not the double-struck integer symbol used for the coefficients.; Proposed edit: H\textasciicircum{}q(\textbackslash{}ZZ, \textbackslash{}ZZ/\textbackslash{}ell)=\allowbreak 0 =\allowbreak > H\textasciicircum{}q(Z, \textbackslash{}ZZ/\textbackslash{}ell)=\allowbreak 0\par
a fix from one model only\par
\par\noindent\hypertarget{DP0441}{}\textbf{DP0441 -- IX, p. 42 (one\_model\_fix)}\par
Bibliography [3], volume abbreviation\par
\textbf{Claude Opus 5.5.} nothing reported here\par
\textbf{GPT-6.1 Sol.} Place: Bibliography [3], volume abbreviation; Reported print: Hamb. Abh. Bd. 13; Source TeX: Hamb. Abh. Bol. 13; Note: The printed abbreviation is « Bd. »: the bowl and ascender belong to one d, not to separate o and l.; Proposed edit: Hamb. Abh. Bol. 13 =\allowbreak > Hamb. Abh. Bd. 13\par
a fix from one model only\par
\par\noindent\hypertarget{DP0442}{}\textbf{DP0442 -- X, p. 43 (conflict)}\par
after the author line, before § 1 | Entre la signature et la section 1\par
\textbf{Claude Opus 5.5.} Place: after the author line, before § 1; Reported print: Sommaire : 1. Introduction 2. Résultats auxiliaires pour un corps 3. Corps des fractions d'un anneau strictement local 4. Dimension cohomologique : cas \ensuremath{\ell} inversible dans O\_X 5. Dimension cohomologique : cas \ensuremath{\ell} =\allowbreak  p . 6. Dimension cohomologique pour un préschéma de type fini sur Spec Z .; Source TeX: (absent); Note: The printed table of contents is not in the TeX. The edition may replace it by a generated table of contents. Its wording differs from the section titles: 'pour un corps' here, 'sur un corps' in § 2.; Proposed edit: \textbackslash{}section\{Introduction\} =\allowbreak > Sommaire : \textbackslash{}begin\{itemize\} \textbackslash{}item[1.] Introduction \textbackslash{}item[2.] Résultats auxiliaires pour un corps \textbackslash{}item[3.] Corps des fractions d'un anneau strictement local \textbackslash{}item[4.] Dimension cohomologique : cas \$\textbackslash{}ell\$ inversible dans \$\textbackslash{}Oo\_X\$ \textbackslash{}item[5.] Dimension cohomologique : cas \$\textbackslash{}ell=\allowbreak p\$. \textbackslash{}item[6.] Dimension cohomologique pour un préschéma de type fini sur \$\textbackslash{}Spec\textbackslash{}ZZ\$. \textbackslash{}end\{itemize\}  \textbackslash{}section\{Introduction\}\par
\textbf{GPT-6.1 Sol.} Place: Entre la signature et la section 1; Reported print: Sommaire : 1. Introduction ; 2. Résultats auxiliaires pour un corps ; 3. Corps des fractions d'un anneau strictement local ; 4. Dimension cohomologique : cas \textbackslash{}ell inversible dans \textbackslash{}mathcal\{O\}\_X ; 5. Dimension cohomologique : cas \textbackslash{}ell=\allowbreak p. ; 6. Dimension cohomologique pour un préschéma de type fini sur \textbackslash{}Spec\textbackslash{}mathbf\{Z\}.; Source TeX: \textbackslash{}section\{Introduction\}; Note: Le sommaire imprimé, comprenant son intitulé et ses six entrées, est entièrement absent du TeX.; Proposed edit: \textbackslash{}section\{Introduction\} =\allowbreak > \textbackslash{}noindent Sommaire : \textbackslash{}begin\{enumerate\} \textbackslash{}item Introduction \textbackslash{}item Résultats auxiliaires pour un corps \textbackslash{}item Corps des fractions d'un anneau strictement local \textbackslash{}item Dimension cohomologique : cas \$\textbackslash{}ell\$ inversible dans \$\textbackslash{}mathcal\{O\}\_X\$ \textbackslash{}item Dimension cohomologique : cas \$\textbackslash{}ell=\allowbreak p\$. \textbackslash{}item Dimension cohomologique pour un préschéma de type fini sur \$\textbackslash{}Spec\textbackslash{}mathbf\{Z\}\$. \textbackslash{}end\{enumerate\}  \textbackslash{}section\{Introduction\}\par
the two models differ at this place\par
\par\noindent\hypertarget{DP0443}{}\textbf{DP0443 -- X, p. 43 (one\_model\_fix)}\par
Signature sous le titre\par
\textbf{Claude Opus 5.5.} nothing reported here\par
\textbf{GPT-6.1 Sol.} Place: Signature sous le titre; Reported print: par M. ARTIN; Source TeX: M. Artin; Note: Le mot « par » manque devant le nom de l'auteur.; Proposed edit: M. Artin \textbackslash{}end\{center\} =\allowbreak > par M. Artin \textbackslash{}end\{center\}\par
a fix from one model only\par
\par\noindent\hypertarget{DP0444}{}\textbf{DP0444 -- X, p. 45 (one\_model\_fix)}\par
Rappelons que lorsque \ensuremath{\ell} =\allowbreak  car k … (CG II 2.2)\par
\textbf{Claude Opus 5.5.} Place: Rappelons que lorsque \ensuremath{\ell} =\allowbreak  car k … (CG II 2.2); Reported print: (CG II 2.2 ); Source TeX: (CG II \textbackslash{}hyperref[X.thm2.2]\{2.2\}); Note: The reference is to Serre's Cohomologie galoisienne II 2.2. The TeX links it to the label X.thm2.2 of this exposé, which is a wrong link. The printed number is the same.; Proposed edit: \textbackslash{}hyperref[X.thm2.2]\{2.2\}) =\allowbreak > 2.2)\par
\textbf{GPT-6.1 Sol.} nothing reported here\par
a fix from one model only\par
\par\noindent\hypertarget{DP0445}{}\textbf{DP0445 -- X, p. 46 (one\_model\_fix)}\par
Remarque, (CG II 2.2)\par
\textbf{Claude Opus 5.5.} Place: Remarque, (CG II 2.2); Reported print: (CG II 2.2); Source TeX: (CG II \textbackslash{}hyperref[X.thm2.2]\{2.2\}); Note: The printed number is the same, but the link points to X 2.2 (this exposé's theorem). The reference is to Serre, CG II 2.2.; Proposed edit: (CG II \textbackslash{}hyperref[X.thm2.2]\{2.2\}) =\allowbreak > (CG II 2.2)\par
\textbf{GPT-6.1 Sol.} nothing reported here\par
a fix from one model only\par
\par\noindent\hypertarget{DP0446}{}\textbf{DP0446 -- X, p. 49 (one\_model\_fix)}\par
diagramme, top entry\par
\textbf{Claude Opus 5.5.} Place: diagramme, top entry; Reported print: H\textasciicircum{}\{n-1\}(Ḡ,H\textasciicircum{}1(H', \ensuremath{\mu}\_\ensuremath{\ell})); Source TeX: H\textasciicircum{}\{1\}(H',\ensuremath{𝛍}\_l); Note: The index is \ensuremath{\ell} (script ell) in the print; the TeX has the Latin letter l.; Proposed edit: H\textasciicircum{}\{1\}(H',\ensuremath{𝛍}\_l) =\allowbreak > H\textasciicircum{}\{1\}(H',\ensuremath{𝛍}\_\{\textbackslash{}ell\})\par
\textbf{GPT-6.1 Sol.} nothing reported here\par
a fix from one model only\par
\par\noindent\hypertarget{DP0447}{}\textbf{DP0447 -- X, p. 49 (one\_model\_fix)}\par
diagramme, right entry\par
\textbf{Claude Opus 5.5.} Place: diagramme, right entry; Reported print: H\textasciicircum{}\{n+1\}(Ḡ, \ensuremath{\mu}\_\ensuremath{\ell}); Source TeX: H\textasciicircum{}\{n+1\}(\textbackslash{}sisi\{\textbackslash{}bar\{G\}\}\{G\},\ensuremath{𝛍}\_l); Note: The index is \ensuremath{\ell} in the print; the TeX has l.; Proposed edit: \{G\},\ensuremath{𝛍}\_l) =\allowbreak > \{G\},\ensuremath{𝛍}\_\{\textbackslash{}ell\})\par
\textbf{GPT-6.1 Sol.} nothing reported here\par
a fix from one model only\par
\par\noindent\hypertarget{DP0448}{}\textbf{DP0448 -- X, p. 49 (one\_model\_fix)}\par
diagramme, bottom entry\par
\textbf{Claude Opus 5.5.} Place: diagramme, bottom entry; Reported print: H\textasciicircum{}\{n-1\}(Ḡ, H\textasciicircum{}1(H, \ensuremath{\mu}\_\ensuremath{\ell})); Source TeX: H\textasciicircum{}\{1\}(H,\ensuremath{𝛍}\_l); Note: The index is \ensuremath{\ell} in the print; the TeX has l.; Proposed edit: H\textasciicircum{}\{1\}(H,\ensuremath{𝛍}\_l) =\allowbreak > H\textasciicircum{}\{1\}(H,\ensuremath{𝛍}\_\{\textbackslash{}ell\})\par
\textbf{GPT-6.1 Sol.} nothing reported here\par
a fix from one model only\par
\par\noindent\hypertarget{DP0449}{}\textbf{DP0449 -- X, p. 49 (one\_model\_fix)}\par
qui induit la flèche \ensuremath{\varepsilon}\par
\textbf{Claude Opus 5.5.} Place: qui induit la flèche \ensuremath{\varepsilon}; Reported print: \ensuremath{\varepsilon} (same handwritten epsilon as the diagram label); Source TeX: \$\textbackslash{}epsilon\$; Note: The diagram uses \textbackslash{}varepsilon and the text \textbackslash{}epsilon (lunate) for the same printed letter.; Proposed edit: la flèche \$\textbackslash{}epsilon\$ =\allowbreak > la flèche \$\textbackslash{}varepsilon\$\par
\textbf{GPT-6.1 Sol.} nothing reported here\par
a fix from one model only\par
\par\noindent\hypertarget{DP0450}{}\textbf{DP0450 -- X, p. 50 (one\_model\_fix)}\par
Démonstration du corollaire 2.4, choix de x\par
\textbf{Claude Opus 5.5.} nothing reported here\par
\textbf{GPT-6.1 Sol.} Place: Démonstration du corollaire 2.4, choix de x; Reported print: x\textbackslash{}in\textbackslash{}mathfrak\{r\}\_A; Source TeX: \$x\textbackslash{}in \ensuremath{𝔪}\_A\$; Note: Le symbole manuscrit de l'imprimé est un r gothique, non un m.; Proposed edit: \ensuremath{𝔪}\_A =\allowbreak > \textbackslash{}mathfrak\{r\}\_A\par
a fix from one model only\par
\par\noindent\hypertarget{DP0451}{}\textbf{DP0451 -- X, p. 59 (agreed\_fix\_held)}\par
proof of 5.2, « du schéma affine déduit de … » | Démonstration de 5.2, phrase sur la jacobienne, après « déduit de »\par
\textbf{Claude Opus 5.5.} Place: proof of 5.2, « du schéma affine déduit de … »; Reported print: déduit de [glyph other than the preceding Φ; apparently \ensuremath{\phi}]; Source TeX: déduit de \$\textbackslash{}Phi\$; Note: The print's glyph here differs from the Φ just before it. In context the affine-scheme map Φ is deduced from the morphism \ensuremath{\phi}:V\ensuremath{\to}V. The scan is doubtful (JBIG2), so verify against the print.; Proposed edit: déduit de \$\textbackslash{}Phi\$ =\allowbreak > déduit de \$\textbackslash{}varphi\$\par
\textbf{GPT-6.1 Sol.} Place: Démonstration de 5.2, phrase sur la jacobienne, après « déduit de »; Reported print: déduit de \$\textbackslash{}varphi\$; Source TeX: déduit de \$\textbackslash{}Phi\$; Note: Le phi imprimé à cet endroit est minuscule, comme celui du morphisme V\ensuremath{\to}V, contrairement aux capitales Φ voisines.; Proposed edit: déduit de \$\textbackslash{}Phi\$ =\allowbreak > déduit de \$\textbackslash{}varphi\$\par
both models give the same fix, held: a model marks it low confidence\par
\par\noindent\hypertarget{DP0452}{}\textbf{DP0452 -- X, p. 59 (doubt)}\par
proof of 5.2, « il s'ensuit que Φ est surjectif »\par
\textbf{Claude Opus 5.5.} Place: proof of 5.2, « il s'ensuit que Φ est surjectif »; Reported print: Φ apparently with a mark beneath (underlined?); Source TeX: il s'ensuit que \$\textbackslash{}Phi\$ est surjectif; Note: The glyph appears to carry a stroke underneath, possibly marking the induced map on k-points as distinct from the scheme morphism. The scan is too unclear to decide, so no edit is proposed.\par
\textbf{GPT-6.1 Sol.} nothing reported here\par
a model flags this place without giving a fix\par
\par\noindent\hypertarget{DP0453}{}\textbf{DP0453 -- X, p. 63 (agreed\_fix\_held)}\par
'au-dessus du point (\ensuremath{\ell})' | Même phrase, « au-dessus du point »\par
\textbf{Claude Opus 5.5.} Place: 'au-dessus du point (\ensuremath{\ell})'; Reported print: au-dessus du point; Source TeX: au-dessus \textbackslash{}sisi\{de\}\{du\} point; Note: The print already reads 'du'. The \textbackslash{}sisi gives 'de' as the print's reading, which is wrong. The printed output 'du' is right.; Proposed edit: \textbackslash{}sisi\{de\}\{du\} =\allowbreak > du\par
\textbf{GPT-6.1 Sol.} Place: Même phrase, « au-dessus du point »; Reported print: au-dessus du point; Source TeX: au-dessus \textbackslash{}sisi\{de\}\{du\} point; Note: Le premier argument de \textbackslash{}sisi attribue à tort « de » à l'imprimé : la page porte déjà « du ».; Proposed edit: \textbackslash{}sisi\{de\}\{du\} =\allowbreak > du\par
both models give the same fix, held: it changes the editors' apparatus (\textbackslash{}sisi): a \textbackslash{}sisi correction, note, label, cross-reference or citation is the editors' record and is not removed or retargeted on the models' word\par
\par\noindent\hypertarget{DP0454}{}\textbf{DP0454 -- XI, p. 64 (conflict)}\par
head of the exposé, between the author line and section 1 | Mention de l'auteur\par
\textbf{Claude Opus 5.5.} Place: head of the exposé, between the author line and section 1; Reported print: Sommaire : 1. Introduction / 2. Existence de sections hyperplanes assez générales / 3. Construction des bons voisinages / 4. Le théorème de comparaison.; Source TeX: (absent); Note: The printed table of contents of the exposé is not in the TeX. This is probably the edition's convention, since LaTeX generates its own table of contents; the coordinator should decide.; Proposed edit: M. Artin  \textbackslash{}end\{center\} =\allowbreak > M. Artin  \textbackslash{}end\{center\}  \textbackslash{}noindent Sommaire :  1. Introduction  2. Existence de sections hyperplanes assez générales  3. Construction des bons voisinages  4. Le théorème de comparaison.\par
\textbf{GPT-6.1 Sol.} Place: Mention de l'auteur; Reported print: par M. ARTIN; Source TeX: M. Artin; Note: Le mot « par » est omis.; Proposed edit: M. Artin =\allowbreak > par M. Artin\par
the two models differ at this place\par
\par\noindent\hypertarget{DP0455}{}\textbf{DP0455 -- XI, p. 64 (one\_model\_fix)}\par
Entre la mention de l'auteur et la section 1\par
\textbf{Claude Opus 5.5.} nothing reported here\par
\textbf{GPT-6.1 Sol.} Place: Entre la mention de l'auteur et la section 1; Reported print: Sommaire : 1. Introduction 2. Existence de sections hyperplanes assez générales 3. Construction des bons voisinages 4. Le théorème de comparaison.; Source TeX: \textbackslash{}section\{Introduction\}; Note: Le titre du sommaire et ses quatre entrées sont absents.; Proposed edit: \textbackslash{}section\{Introduction\} =\allowbreak > Sommaire : \textbackslash{}begin\{enumerate\} \textbackslash{}item Introduction \textbackslash{}item Existence de sections hyperplanes assez générales \textbackslash{}item Construction des bons voisinages \textbackslash{}item Le théorème de comparaison. \textbackslash{}end\{enumerate\}  \textbackslash{}section\{Introduction\}\par
a fix from one model only\par
\par\noindent\hypertarget{DP0456}{}\textbf{DP0456 -- XI, p. 65 (cut)}\par
Fin de la dernière ligne, raccord avec la page 66\par
\textbf{Claude Opus 5.5.} nothing reported here\par
\textbf{GPT-6.1 Sol.} Place: Fin de la dernière ligne, raccord avec la page 66; Reported print: les hyperplans de degré d passant par P sont paramé-; Source TeX: les hyperplans de degré \$d\$ passant par \$P\$ sont; Note: Le segment fourni s'arrête après « sont », tandis que la page imprime encore « paramé- », début d'un mot partagé avec la page 66. Vérifier le raccord avec le segment suivant ; ne pas ajouter ce fragment au texte.\par
a model reports that the automatic cut gives lines of this page to its neighbour, or the neighbour's lines to this page: check the lines it names against the scan of the page that prints them (that page's checker did not have their TeX)\par
\par\noindent\hypertarget{DP0457}{}\textbf{DP0457 -- XI, p. 68 (conflict)}\par
Soit V'' =\allowbreak  ..., image of V'' | Définition de l'image de V''\par
\textbf{Claude Opus 5.5.} Place: Soit V'' =\allowbreak  ..., image of V''; Reported print: \textbackslash{}bar V'' =\allowbreak  F(V''); Source TeX: \textbackslash{}sisi\{f\}\{F\}(V''); Note: The print reads F(V''); the correct reading is f(V''). The \textbackslash{}sisi arguments are reversed: the print's reading should be the first argument.; Proposed edit: \textbackslash{}sisi\{f\}\{F\}(V'') =\allowbreak > \textbackslash{}sisi\{F\}\{f\}(V'')\par
\textbf{GPT-6.1 Sol.} Place: Définition de l'image de V''; Reported print: \textbackslash{}bar\{V\}\{\}''=\allowbreak F(V''); Source TeX: \textbackslash{}bar\{V\}\{\}''=\allowbreak \textbackslash{}sisi\{f\}\{F\}(V''); Proposed correction: \textbackslash{}bar\{V\}\{\}''=\allowbreak f(V''); Note: L'application vient d'être notée f ; le F majuscule imprimé est une coquille. La lecture TeX à comparer, premier argument de \textbackslash{}sisi, donne déjà la correction f plutôt que le F imprimé.\par
the two models differ at this place: one model fixes the TeX, the other calls it a misprint of the print\par
\par\noindent\hypertarget{DP0458}{}\textbf{DP0458 -- XI, p. 69 (conflict)}\par
proof of 3.3, U' =\allowbreak  U \ensuremath{\cap} ... | Démonstration de 3.3, choix de U'\par
\textbf{Claude Opus 5.5.} Place: proof of 3.3, U' =\allowbreak  U \ensuremath{\cap} ...; Reported print: U' =\allowbreak  U \ensuremath{\cap} V'; Source TeX: \textbackslash{}sisi\{f\textasciicircum{}\{-1\}(V')\}\{V'\}; Note: The print reads V'; the correct reading is f\textasciicircum{}\{-1\}(V'). The \textbackslash{}sisi arguments are reversed.; Proposed edit: \textbackslash{}sisi\{f\textasciicircum{}\{-1\}(V')\}\{V'\} =\allowbreak > \textbackslash{}sisi\{V'\}\{f\textasciicircum{}\{-1\}(V')\}\par
\textbf{GPT-6.1 Sol.} Place: Démonstration de 3.3, choix de U'; Reported print: U'=\allowbreak U\textbackslash{}cap V'; Source TeX: U'=\allowbreak U\textbackslash{}cap \textbackslash{}sisi\{f\textasciicircum{}\{-1\}(V')\}\{V'\}; Proposed correction: U'=\allowbreak U\textbackslash{}cap f\textasciicircum{}\{-1\}(V'); Note: V' est un ouvert de la base V : il faut prendre son image inverse dans U. La lecture TeX à comparer, premier argument de \textbackslash{}sisi, corrige cette coquille au lieu de reproduire l'imprimé.\par
the two models differ at this place: one model fixes the TeX, the other calls it a misprint of the print\par
\par\noindent\hypertarget{DP0459}{}\textbf{DP0459 -- XI, p. 70 (conflict)}\par
condition on H\_0 | Condition de transversalité imposée à H\_0\par
\textbf{Claude Opus 5.5.} Place: condition on H\_0; Reported print: tel que H\_0 coupe X \ensuremath{\cap} L transversalement; Source TeX: \textbackslash{}sisi\{\textbackslash{}bar\{X\}\}\{X\}\textbackslash{}cap L; Note: The print reads X \ensuremath{\cap} L; the correct reading is \ensuremath{\bar{X}} \ensuremath{\cap} L. The \textbackslash{}sisi arguments are reversed.; Proposed edit: \textbackslash{}sisi\{\textbackslash{}bar\{X\}\}\{X\} =\allowbreak > \textbackslash{}sisi\{X\}\{\textbackslash{}bar\{X\}\}\par
\textbf{GPT-6.1 Sol.} Place: Condition de transversalité imposée à H\_0; Reported print: H\_0 coupe X\textbackslash{}cap L transversalement; Source TeX: \$H\_0\$ coupe \$\textbackslash{}sisi\{\textbackslash{}bar\{X\}\}\{X\}\textbackslash{}cap L\$ transversalement; Note: Le X imprimé est sans barre. Le premier argument de \textbackslash{}sisi ne reproduit donc pas le tirage.; Proposed edit: \textbackslash{}sisi\{\textbackslash{}bar\{X\}\}\{X\} =\allowbreak > \textbackslash{}sisi\{X\}\{X\}\par
the two models differ at this place\par
\par\noindent\hypertarget{DP0460}{}\textbf{DP0460 -- XI, p. 70 (conflict)}\par
display y\_i =\allowbreak  \ensuremath{\Sigma} ... | Formule définissant les coordonnées projectives y\_i\par
\textbf{Claude Opus 5.5.} Place: display y\_i =\allowbreak  \ensuremath{\Sigma} ...; Reported print: a\_\{i\_\ensuremath{\nu}\} x\_\ensuremath{\nu}; Source TeX: a\_\{i\_\{\textbackslash{}nu\}\}x\_\{\textbackslash{}nu\}; Proposed correction: a\_\{i\textbackslash{}nu\}x\_\{\textbackslash{}nu\}; Note: The print sets \ensuremath{\nu} as a subscript of i; it should be a\_\{i\ensuremath{\nu}\}, as in the earlier display. The TeX keeps the printed reading.\par
\textbf{GPT-6.1 Sol.} Place: Formule définissant les coordonnées projectives y\_i; Reported print: y\_i=\allowbreak \textbackslash{}sum\textasciicircum{}N\_\{\textbackslash{}nu=\allowbreak 0\}a\_\{i\textbackslash{}nu\}x\_\textbackslash{}nu; Source TeX: y\_i=\allowbreak \textbackslash{}sum\textasciicircum{}N\_\{\textbackslash{}nu =\allowbreak 0\}a\_\{i\_\{\textbackslash{}nu\}\}x\_\{\textbackslash{}nu\}; Note: Le coefficient porte les deux indices i et \ensuremath{\nu}, comme dans l'équation de H\_i. Le \ensuremath{\nu} manuscrit est placé bas ; cette position ne doit pas être transcrite comme un sous-indice de i.; Proposed edit: a\_\{i\_\{\textbackslash{}nu\}\} =\allowbreak > a\_\{i\textbackslash{}nu\}\par
the two models differ at this place: one model fixes the TeX, the other calls it a misprint of the print\par
\par\noindent\hypertarget{DP0461}{}\textbf{DP0461 -- XI, p. 72 (conflict)}\par
first lines, \textbackslash{}sisi after « sur \$\textbackslash{}PP\textasciicircum{}\{n-1\}\$ » (convention of this file, see coordinator) | Premier paragraphe, après « sur P\textasciicircum{}\{n-1\} »\par
\textbf{Claude Opus 5.5.} Place: first lines, \textbackslash{}sisi after « sur \$\textbackslash{}PP\textasciicircum{}\{n-1\}\$ » (convention of this file, see coordinator); Reported print: isomorphiquement sur P\textasciicircum{}\{n-1\} est définie en chaque point de Y; Source TeX: \$\textbackslash{}PP\textasciicircum{}\{n-1\}\$\textbackslash{}sisi\{. De plus, la projection\}\{\} est définie; Note: In this file the \textbackslash{}sisi arguments are reversed: the second argument holds the print's reading and the first holds the editors' correction. The print omits « . De plus, la projection » (a misprint of the print), and A should be empty.; Proposed edit: \textbackslash{}sisi\{. De plus, la projection\}\{\} =\allowbreak > \textbackslash{}sisi\{\}\{. De plus, la projection\}\par
\textbf{GPT-6.1 Sol.} Place: Premier paragraphe, après « sur P\textasciicircum{}\{n-1\} »; Reported print: sur \$\textbackslash{}PP\textasciicircum{}\{n-1\}\$ est définie en chaque point de \$Y\$; Source TeX: sur \$\textbackslash{}PP\textasciicircum{}\{n-1\}\$\textbackslash{}sisi\{. De plus, la projection\}\{\} est définie; Note: L'imprimé enchaîne directement avec « est définie ». Cette restitution ne doit pas figurer dans le premier argument de \textbackslash{}sisi, censé donner la lecture imprimée ; le second argument n'est pas évalué.; Proposed edit: \textbackslash{}sisi\{. De plus, la projection\}\{\} =\allowbreak > \par
the two models differ at this place\par
\par\noindent\hypertarget{DP0462}{}\textbf{DP0462 -- XI, p. 72 (one\_model\_fix)}\par
4.0, condition définissant un isomorphisme local\par
\textbf{Claude Opus 5.5.} nothing reported here\par
\textbf{GPT-6.1 Sol.} Place: 4.0, condition définissant un isomorphisme local; Reported print: Pour chaque \$x\textbackslash{}in U\$; Source TeX: \textbackslash{}sisi\{pour\}\{Pour\} chaque \$x\textbackslash{}in U\$; Note: L'imprimé porte une majuscule à « Pour » ; le premier argument donne à tort une minuscule.; Proposed edit: \textbackslash{}sisi\{pour\}\{Pour\} =\allowbreak > Pour\par
a fix from one model only\par
\par\noindent\hypertarget{DP0463}{}\textbf{DP0463 -- XI, p. 74 (one\_model\_fix)}\par
Premier paragraphe, après la référence à Satz 8\par
\textbf{Claude Opus 5.5.} nothing reported here\par
\textbf{GPT-6.1 Sol.} Place: Premier paragraphe, après la référence à Satz 8; Reported print: D'après … s'étend à un « revêtement analytique »; Source TeX: D'après … \textbackslash{}sisi\{\$\textbackslash{}eta'\$\}\{\} s'étend à un; Note: Aucun \ensuremath{\eta}′ n'est imprimé entre la référence et « s'étend ». Le sujet restitué est placé à tort dans le premier argument de \textbackslash{}sisi.; Proposed edit: \textbackslash{}sisi\{\$\textbackslash{}eta'\$\}\{\} =\allowbreak > \par
a fix from one model only\par
\par\noindent\hypertarget{DP0464}{}\textbf{DP0464 -- XI, p. 74 (one\_model\_fix)}\par
Premier paragraphe, « que \ensuremath{\eta} est induite »\par
\textbf{Claude Opus 5.5.} nothing reported here\par
\textbf{GPT-6.1 Sol.} Place: Premier paragraphe, « que \ensuremath{\eta} est induite »; Reported print: que \$\textbackslash{}eta\$ est induite par un morphisme; Source TeX: que \$\textbackslash{}eta\$ est \textbackslash{}sisi\{induit\}\{induite\} par un morphisme; Note: La lecture imprimée est « induite », non « induit ». Le contrôle porte sur le premier argument.; Proposed edit: \textbackslash{}sisi\{induit\}\{induite\} =\allowbreak > induite\par
a fix from one model only\par
\par\noindent\hypertarget{DP0465}{}\textbf{DP0465 -- XI, p. 75 (one\_model\_fix)}\par
Théorème 4.4 (iii), second affichage\par
\textbf{Claude Opus 5.5.} nothing reported here\par
\textbf{GPT-6.1 Sol.} Place: Théorème 4.4 (iii), second affichage; Reported print: H\textasciicircum{}q(X\_\{\textbackslash{}et\},\textbackslash{}ZZ/n)\textbackslash{}xrightarrow\{\textbackslash{}sim\} H\textasciicircum{}q(X\_\{\textbackslash{}cl\},\textbackslash{}ZZ/n); Source TeX: H\textasciicircum{}q(X\_\{\textbackslash{}et\}, \textbackslash{}ZZ/\textbackslash{}sisi\{n\textbackslash{}ZZ\}\{n\})\textbackslash{}xrightarrow\{\textbackslash{}sim\} H\textasciicircum{}q(X\_\{\textbackslash{}cl\},\textbackslash{}ZZ/\textbackslash{}sisi\{n\textbackslash{}ZZ\}\{n\}); Note: Dans les deux coefficients, l'imprimé donne le dénominateur n, sans second Z. Les premiers arguments des deux \textbackslash{}sisi ne correspondent donc pas au scan.; Proposed edit: H\textasciicircum{}q(X\_\{\textbackslash{}et\}, \textbackslash{}ZZ/\textbackslash{}sisi\{n\textbackslash{}ZZ\}\{n\})\textbackslash{}xrightarrow\{\textbackslash{}sim\} H\textasciicircum{}q(X\_\{\textbackslash{}cl\},\textbackslash{}ZZ/\textbackslash{}sisi\{n\textbackslash{}ZZ\}\{n\}) =\allowbreak > H\textasciicircum{}q(X\_\{\textbackslash{}et\}, \textbackslash{}ZZ/n)\textbackslash{}xrightarrow\{\textbackslash{}sim\} H\textasciicircum{}q(X\_\{\textbackslash{}cl\},\textbackslash{}ZZ/n)\par
a fix from one model only\par
\par\noindent\hypertarget{DP0466}{}\textbf{DP0466 -- XI, p. 75 (one\_model\_fix)}\par
Phrase introduisant le lemme 4.5\par
\textbf{Claude Opus 5.5.} nothing reported here\par
\textbf{GPT-6.1 Sol.} Place: Phrase introduisant le lemme 4.5; Reported print: on doit donc démontrer le suivant; Source TeX: on doit donc démontrer le \textbackslash{}sisi\{lemme\}\{\} suivant; Note: Le mot « lemme » ne figure pas dans cette phrase imprimée ; il est ajouté dans le premier argument de \textbackslash{}sisi.; Proposed edit: \textbackslash{}sisi\{lemme\}\{\} =\allowbreak > \par
a fix from one model only\par
\par\noindent\hypertarget{DP0467}{}\textbf{DP0467 -- XI, p. 75 (one\_model\_fix)}\par
Démonstration du lemme, première phrase\par
\textbf{Claude Opus 5.5.} nothing reported here\par
\textbf{GPT-6.1 Sol.} Place: Démonstration du lemme, première phrase; Reported print: pour la topologie étale, et on peut donc supposer; Source TeX: pour la topologie étale, \textbackslash{}sisi\{\}\{et\} on peut donc supposer; Note: La conjonction « et » est bien imprimée ; le premier argument est à tort vide.; Proposed edit: \textbackslash{}sisi\{\}\{et\} =\allowbreak > et\par
a fix from one model only\par
\par\noindent\hypertarget{DP0468}{}\textbf{DP0468 -- XI, p. 75 (one\_model\_fix)}\par
Démonstration du lemme, réduction au faisceau constant cyclique\par
\textbf{Claude Opus 5.5.} nothing reported here\par
\textbf{GPT-6.1 Sol.} Place: Démonstration du lemme, réduction au faisceau constant cyclique; Reported print: \$F=\allowbreak \textbackslash{}ZZ/n\$; Source TeX: \$F=\allowbreak \textbackslash{}ZZ/\textbackslash{}sisi\{n\textbackslash{}ZZ\}\{n\}\$; Note: L'imprimé donne n au dénominateur, non nZ ; le premier argument de \textbackslash{}sisi est erroné.; Proposed edit: \$F=\allowbreak \textbackslash{}ZZ/\textbackslash{}sisi\{n\textbackslash{}ZZ\}\{n\}\$ =\allowbreak > \$F=\allowbreak \textbackslash{}ZZ/n\$\par
a fix from one model only\par
\par\noindent\hypertarget{DP0469}{}\textbf{DP0469 -- XI, p. 76 (conflict)}\par
first line: R\textasciicircum{}q j\_\{cl*\}F =\allowbreak  0 si q > 1 | Première phrase, annulation des images directes supérieures\par
\textbf{Claude Opus 5.5.} Place: first line: R\textasciicircum{}q j\_\{cl*\}F =\allowbreak  0 si q > 1; Reported print: R\textasciicircum{}q j\_\{cl*\}F =\allowbreak  0; Source TeX: \$R\textasciicircum{}q\{j\}\_\{\textbackslash{}ast\}F=\allowbreak 0\$; Note: The subscript cl of j\_* is missing in the TeX.; Proposed edit: \$R\textasciicircum{}q\{j\}\_\{\textbackslash{}ast\}F=\allowbreak 0\$ =\allowbreak > \$R\textasciicircum{}qj\_\{\textbackslash{}cl\textbackslash{}!\textbackslash{}ast\}F=\allowbreak 0\$\par
\textbf{GPT-6.1 Sol.} Place: Première phrase, annulation des images directes supérieures; Reported print: R\textasciicircum{}qj\_\{\textbackslash{}cl\textbackslash{}!\textbackslash{}ast\}F=\allowbreak 0; Source TeX: R\textasciicircum{}q\{j\}\_\{\textbackslash{}ast\}F=\allowbreak 0; Note: L'indice cl de j manque dans le TeX.; Proposed edit: R\textasciicircum{}q\{j\}\_\{\textbackslash{}ast\}F=\allowbreak 0 =\allowbreak > R\textasciicircum{}q\{j\}\_\{\textbackslash{}cl\textbackslash{}!\textbackslash{}ast\}F=\allowbreak 0\par
the two models differ at this place\par
\par\noindent\hypertarget{DP0470}{}\textbf{DP0470 -- XI, p. 76 (conflict)}\par
j\_\{cl*\}F est le faisceau constant Z/n sur \ensuremath{\bar{X}} | Quotients d'entiers : même différence répétée cinq fois\par
\textbf{Claude Opus 5.5.} Place: j\_\{cl*\}F est le faisceau constant Z/n sur \ensuremath{\bar{X}}; Reported print: \ensuremath{\mathbf{Z}}/n; Source TeX: \$\textbackslash{}ZZ/\textbackslash{}sisi\{n\textbackslash{}ZZ\}\{n\}\$; Note: \textbackslash{}sisi's first (print) argument gives Z/nZ, but the print has Z/n. The same holds for the three later \textbackslash{}sisi\{n\textbackslash{}ZZ\}\{n\} on this page (R\textasciicircum{}0, R\textasciicircum{}1, R\textasciicircum{}q lines), which all print Z/n. The arguments seem swapped: see coordinator.; Proposed edit: \$\textbackslash{}ZZ/\textbackslash{}sisi\{n\textbackslash{}ZZ\}\{n\}\$ sur =\allowbreak > \$\textbackslash{}ZZ/n\$ sur\par
\textbf{GPT-6.1 Sol.} Place: Quotients d'entiers : même différence répétée cinq fois; Reported print: \textbackslash{}ZZ/n; Source TeX: \textbackslash{}ZZ/\textbackslash{}sisi\{n\textbackslash{}ZZ\}\{n\}, dont la lecture A est \textbackslash{}ZZ/n\textbackslash{}ZZ; Note: Le scan donne toujours n seul au dénominateur. Cette notation abrégée n'est pas une coquille. Le fix est représentatif de la correction à appliquer aux cinq lectures A ; B n'est pas contrôlé.; Proposed edit: le faisceau constant \$\textbackslash{}ZZ/\textbackslash{}sisi\{n\textbackslash{}ZZ\}\{n\}\$ =\allowbreak > le faisceau constant \$\textbackslash{}ZZ/\textbackslash{}sisi\{n\}\{n\}\$\par
the two models differ at this place\par
\par\noindent\hypertarget{DP0471}{}\textbf{DP0471 -- XI, p. 78 (one\_model\_fix)}\par
bibliography [2]\par
\textbf{Claude Opus 5.5.} Place: bibliography [2]; Reported print: GRAUERT H et; Source TeX: GRAUERT H\textbackslash{}sisi\{.\}\{\} et; Note: \textbackslash{}sisi's print argument gives 'H.', but the print has 'H' without a period (arguments swapped).; Proposed edit: GRAUERT H\textbackslash{}sisi\{.\}\{\} et =\allowbreak > GRAUERT H et\par
\textbf{GPT-6.1 Sol.} nothing reported here\par
a fix from one model only\par
\par\noindent\hypertarget{DP0472}{}\textbf{DP0472 -- XI, p. 78 (one\_model\_fix)}\par
bibliography [2]\par
\textbf{Claude Opus 5.5.} Place: bibliography [2]; Reported print: Math. Ann.; Source TeX: Math, Ann.; Note: The print has a period after Math, not a comma.; Proposed edit: Math, Ann. =\allowbreak > Math. Ann.\par
\textbf{GPT-6.1 Sol.} nothing reported here\par
a fix from one model only\par
\par\noindent\hypertarget{DP0473}{}\textbf{DP0473 -- XII, p. 79 (one\_model\_fix)}\par
Ligne d’auteur sous le titre\par
\textbf{Claude Opus 5.5.} nothing reported here\par
\textbf{GPT-6.1 Sol.} Place: Ligne d’auteur sous le titre; Reported print: par M. ARTIN; Source TeX: M. Artin; Note: Le mot « par » précédant le nom de l’auteur manque dans le TeX.; Proposed edit: M. Artin =\allowbreak > par M. Artin\par
a fix from one model only\par
\par\noindent\hypertarget{DP0474}{}\textbf{DP0474 -- XII, p. 80 (one\_model\_fix)}\par
footnote (*)\par
\textbf{Claude Opus 5.5.} Place: footnote (*); Reported print: Comparer, du point de vue \ensuremath{\pi}\_1, avec SGA 3 X (6, Exemples) et (1.6).; Source TeX: Comparer, du point de vue \$\textbackslash{}Pi\_1\$, avec; Note: The print has a lower-case pi (the fundamental group \ensuremath{\pi}\_1). The TeX has a capital \textbackslash{}Pi.; Proposed edit: du point de vue \$\textbackslash{}Pi\_1\$ =\allowbreak > du point de vue \$\textbackslash{}pi\_1\$\par
\textbf{GPT-6.1 Sol.} nothing reported here\par
a fix from one model only\par
\par\noindent\hypertarget{DP0475}{}\textbf{DP0475 -- XII, p. 89 (agreed\_fix\_held)}\par
Corollaire 5.3, first line | Corollary 5.3, opening reference\par
\textbf{Claude Opus 5.5.} Place: Corollaire 5.3, first line; Reported print: Sous les conditions de (5.2); Source TeX: Sous les conditions de (\textbackslash{}hyperrefext[X][X.cor5.2]\{5.2\}); Note: The print refers to 5.2 of this exposé (the corollary just above: f propre, ξ, X\_ξ). The TeX links it to X 5.2 and prints 'X 5.2'.; Proposed edit: Sous les conditions de (\textbackslash{}hyperrefext[X][X.cor5.2]\{5.2\}) =\allowbreak > Sous les conditions de (\textbackslash{}hyperref[XII.cor5.2]\{5.2\})\par
\textbf{GPT-6.1 Sol.} Place: Corollary 5.3, opening reference; Reported print: Sous les conditions de (5.2); Source TeX: Sous les conditions de (\textbackslash{}hyperrefext[X][X.cor5.2]\{5.2\}); Note: The print refers to Corollary 5.2 of the present exposé, not X 5.2. The external reference to X 5.2 in the preceding sentence is correct and should remain unchanged.; Proposed edit: Sous les conditions de (\textbackslash{}hyperrefext[X][X.cor5.2]\{5.2\}) =\allowbreak > Sous les conditions de (\textbackslash{}hyperref[XII.cor5.2]\{5.2\})\par
both models give the same fix, held: it changes the editors' apparatus (\textbackslash{}hyperref, \textbackslash{}hyperrefext, the target of a cross-reference): a \textbackslash{}sisi correction, note, label, cross-reference or citation is the editors' record and is not removed or retargeted on the models' word\par
\par\noindent\hypertarget{DP0476}{}\textbf{DP0476 -- XII, p. 93 (agreed\_fix\_held)}\par
\ensuremath{\pi}\_o throughout (also \ensuremath{\pi}\_1 on p. 94) | 5.6–5.9, notation for components and fundamental groups\par
\textbf{Claude Opus 5.5.} Place: \ensuremath{\pi}\_o throughout (also \ensuremath{\pi}\_1 on p. 94); Reported print: \ensuremath{\pi}\_o(X\_o) (typewriter Greek, apparently small \ensuremath{\pi}); Source TeX: \textbackslash{}Pi\_\{\textbackslash{}circ\}(X\_\{\textbackslash{}circ\}); Note: The print seems to use small \ensuremath{\pi} for \ensuremath{\pi}\_0 and \ensuremath{\pi}\_1, where the TeX has \textbackslash{}Pi throughout. This is a transcription convention; the glyph's case is not fully certain.; Proposed edit: l'ensemble \$\textbackslash{}Pi\_\{\textbackslash{}circ\}(X\_\{\textbackslash{}circ\})\$ des =\allowbreak > l'ensemble \$\textbackslash{}pi\_\{\textbackslash{}circ\}(X\_\{\textbackslash{}circ\})\$ des\par
\textbf{GPT-6.1 Sol.} Place: 5.6–5.9, notation for components and fundamental groups; Reported print: \textbackslash{}pi\_0(X\_0), \textbackslash{}pi\_0(X), \textbackslash{}pi\_1(X\_0,\textbackslash{}bar\{x\}\_0); Source TeX: \textbackslash{}Pi\_\{\textbackslash{}circ\}(X\_\{\textbackslash{}circ\}), \textbackslash{}Pi\_\{\textbackslash{}circ\}(X), \textbackslash{}Pi\_1(X\_\{\textbackslash{}circ\},\textbackslash{}bar\{x\}\_\{\textbackslash{}circ\}); Note: Systematic case difference: the print uses lowercase \ensuremath{\pi}, whereas the TeX uses uppercase \ensuremath{\Pi}. Reported once with a representative edit; it continues on p. 94. The handwritten theta ovals do not establish a Θ/\ensuremath{\theta} distinction safely, so no change to \textbackslash{}Theta is proposed.; Proposed edit: \textbackslash{}Pi\_\{\textbackslash{}circ\}(X\_\{\textbackslash{}circ\})\$ des composantes connexes =\allowbreak > \textbackslash{}pi\_\{\textbackslash{}circ\}(X\_\{\textbackslash{}circ\})\$ des composantes connexes\par
both models give the same fix, held: a model marks it low confidence\par
\par\noindent\hypertarget{DP0477}{}\textbf{DP0477 -- XII, p. 95 (one\_model\_fix)}\par
footnote (*), cf. th. 3.1 dans M. Artin\par
\textbf{Claude Opus 5.5.} Place: footnote (*), cf. th. 3.1 dans M. Artin; Reported print: cf. th. 3.1 dans M. Artin, Algebraic approximation ...; Source TeX: th.\textasciitilde{}\textbackslash{}hyperref[XII.prop3.1]\{3.1\}; Note: 3.1 refers to Artin's paper, not to XII 3.1. The hyperlink to the label XII.prop3.1 is a wrong reference, though the printed text is the same.; Proposed edit: th.\textasciitilde{}\textbackslash{}hyperref[XII.prop3.1]\{3.1\} =\allowbreak > th.\textasciitilde{}3.1\par
\textbf{GPT-6.1 Sol.} nothing reported here\par
a fix from one model only\par
\par\noindent\hypertarget{DP0478}{}\textbf{DP0478 -- XII, p. 104 (one\_model\_fix)}\par
6.5 (i) b''), display\par
\textbf{Claude Opus 5.5.} Place: 6.5 (i) b''), display; Reported print: \ensuremath{\pi}\_o(h') : \ensuremath{\pi}\_o(X') \ensuremath{\to} \ensuremath{\pi}\_o(Y'); Source TeX: \textbackslash{}Pi\_\{\textbackslash{}circ\}(h'):\textbackslash{}Pi\_\{\textbackslash{}circ\}(X')\textbackslash{}longrightarrow \textbackslash{}Pi\_\{\textbackslash{}circ\}(Y'); Note: The typewriter glyph is a small \ensuremath{\pi} (\ensuremath{\pi}\_0), not a capital \ensuremath{\Pi}.; Proposed edit: \textbackslash{}Pi\_\{\textbackslash{}circ\}(h'):\textbackslash{}Pi\_\{\textbackslash{}circ\}(X')\textbackslash{}longrightarrow \textbackslash{}Pi\_\{\textbackslash{}circ\}(Y') =\allowbreak > \textbackslash{}pi\_\{\textbackslash{}circ\}(h'):\textbackslash{}pi\_\{\textbackslash{}circ\}(X')\textbackslash{}longrightarrow \textbackslash{}pi\_\{\textbackslash{}circ\}(Y')\par
\textbf{GPT-6.1 Sol.} nothing reported here\par
a fix from one model only\par
\par\noindent\hypertarget{DP0479}{}\textbf{DP0479 -- XII, p. 105 (one\_model\_fix)}\par
6.5 (ii) b''), display\par
\textbf{Claude Opus 5.5.} Place: 6.5 (ii) b''), display; Reported print: \ensuremath{\pi}\_i(Y',y') \ensuremath{\to} \ensuremath{\pi}\_i(X',x'); Source TeX: \textbackslash{}Pi\_i(Y',y')\textbackslash{}longrightarrow \textbackslash{}Pi\_i(X', x'); Note: The print has a small \ensuremath{\pi} (\ensuremath{\pi}\_0, \ensuremath{\pi}\_1), not a capital \ensuremath{\Pi}.; Proposed edit: \textbackslash{}Pi\_i(Y',y')\textbackslash{}longrightarrow \textbackslash{}Pi\_i(X', x') =\allowbreak > \textbackslash{}pi\_i(Y',y')\textbackslash{}longrightarrow \textbackslash{}pi\_i(X', x')\par
\textbf{GPT-6.1 Sol.} nothing reported here\par
a fix from one model only\par
\par\noindent\hypertarget{DP0480}{}\textbf{DP0480 -- XII, p. 109 (one\_model\_fix)}\par
Références après « grâce à IX 1.2 (v) »\par
\textbf{Claude Opus 5.5.} nothing reported here\par
\textbf{GPT-6.1 Sol.} Place: Références après « grâce à IX 1.2 (v) »; Reported print: IX 1.2 (v), 1.6 (iii); Source TeX: \textbackslash{}hyperrefext[IX][IX.prop1.2]\{1.2\} (v), \textbackslash{}hyperrefext[IX][IX.prop1.6]\{1.6\} (iii); Note: The second macro adds another exposé numeral IX; the print supplies IX only before 1.2.; Proposed edit: \textbackslash{}hyperrefext[IX][IX.prop1.6]\{1.6\} =\allowbreak > 1.6\par
a fix from one model only\par
\par\noindent\hypertarget{DP0481}{}\textbf{DP0481 -- XII, p. 110 (one\_model\_fix)}\par
« remplacer F par le produit des … »\par
\textbf{Claude Opus 5.5.} nothing reported here\par
\textbf{GPT-6.1 Sol.} Place: « remplacer F par le produit des … »; Reported print: p\_\{i\_\{\textbackslash{}ast\}\}(G\_\{i X'\_i\}); Source TeX: \$p\_\{i\textasciicircum{}\{\textbackslash{}ast\}\}(G\_\{i X'\_i\})\$; Note: The printed star denotes direct image and is a subscript, not a superscript.; Proposed edit: p\_\{i\textasciicircum{}\{\textbackslash{}ast\}\} =\allowbreak > p\_\{i\_\{\textbackslash{}ast\}\}\par
a fix from one model only\par
\par\noindent\hypertarget{DP0482}{}\textbf{DP0482 -- XII, p. 112 (one\_model\_fix)}\par
Remarque 6.7 3), «pour IL =\allowbreak  \ensuremath{\varnothing}»\par
\textbf{Claude Opus 5.5.} Place: Remarque 6.7 3), «pour IL =\allowbreak  \ensuremath{\varnothing}»; Reported print: \textbackslash{}LL =\allowbreak  \textbackslash{}emptyset (empty set); Source TeX: \$\textbackslash{}LL=\allowbreak \textbackslash{}phi\$; Note: The print has the empty-set sign, not the Greek letter phi.; Proposed edit: \$\textbackslash{}LL=\allowbreak \textbackslash{}phi\$ =\allowbreak > \$\textbackslash{}LL=\allowbreak \textbackslash{}emptyset\$\par
\textbf{GPT-6.1 Sol.} nothing reported here\par
a fix from one model only\par
\par\noindent\hypertarget{DP0483}{}\textbf{DP0483 -- XII, p. 113 (one\_model\_fix)}\par
Proposition 6.8, introduction, morphisme canonique vers Y\par
\textbf{Claude Opus 5.5.} nothing reported here\par
\textbf{GPT-6.1 Sol.} Place: Proposition 6.8, introduction, morphisme canonique vers Y; Reported print: p\_Y; Source TeX: p\_y; Note: L’indice imprimé est un Y majuscule.; Proposed edit: p\_y: =\allowbreak > p\_Y:\par
a fix from one model only\par
\par\noindent\hypertarget{DP0484}{}\textbf{DP0484 -- XII, p. 117 (one\_model\_fix)}\par
Corollaire 6.9, premier homomorphisme de changement de base\par
\textbf{Claude Opus 5.5.} nothing reported here\par
\textbf{GPT-6.1 Sol.} Place: Corollaire 6.9, premier homomorphisme de changement de base; Reported print: p\_\{Y\textbackslash{}ast\}(\textbackslash{}bar\{h\}\textasciicircum{}\{\textbackslash{}ast\}(\textbackslash{}bar\{F\})); Source TeX: p\_\{Y\textasciicircum{}\{\textbackslash{}ast\}\}(\textbackslash{}bar\{h\}\textasciicircum{}\{\textbackslash{}ast\}(\textbackslash{}bar\{F\})); Note: Dans le print, Y et l'étoile sont tous deux en indice de p ; le TeX fait de l'étoile un exposant de Y.; Proposed edit: p\_\{Y\textasciicircum{}\{\textbackslash{}ast\}\}(\textbackslash{}bar\{h\}\textasciicircum{}\{\textbackslash{}ast\}(\textbackslash{}bar\{F\}))\$ est injectif =\allowbreak > p\_\{Y\textbackslash{}ast\}(\textbackslash{}bar\{h\}\textasciicircum{}\{\textbackslash{}ast\}(\textbackslash{}bar\{F\}))\$ est injectif\par
a fix from one model only\par
\par\noindent\hypertarget{DP0485}{}\textbf{DP0485 -- XII, p. 117 (one\_model\_fix)}\par
Corollaire 6.9, homomorphisme portant sur R\textasciicircum{}i\par
\textbf{Claude Opus 5.5.} nothing reported here\par
\textbf{GPT-6.1 Sol.} Place: Corollaire 6.9, homomorphisme portant sur R\textasciicircum{}i; Reported print: R\textasciicircum{}\{i\}p\_\{Y\textbackslash{}ast\}(\textbackslash{}bar\{h\}\textasciicircum{}\{\textbackslash{}ast\}(\textbackslash{}bar\{F\})); Source TeX: R\textasciicircum{}\{i\}p\_\{Y\textasciicircum{}\{\textbackslash{}ast\}\}(\textbackslash{}bar\{h\}\textasciicircum{}\{\textbackslash{}ast\}(\textbackslash{}bar\{F\})); Note: L'étoile du foncteur d'image directe est en indice avec Y, non en exposant de Y.; Proposed edit: R\textasciicircum{}\{i\}p\_\{Y\textasciicircum{}\{\textbackslash{}ast\}\} =\allowbreak > R\textasciicircum{}\{i\}p\_\{Y\textbackslash{}ast\}\par
a fix from one model only\par
\par\noindent\hypertarget{DP0486}{}\textbf{DP0486 -- XII, p. 118 (one\_model\_fix)}\par
Corollaire 6.10 (i), condition 2°, premier homomorphisme\par
\textbf{Claude Opus 5.5.} nothing reported here\par
\textbf{GPT-6.1 Sol.} Place: Corollaire 6.10 (i), condition 2°, premier homomorphisme; Reported print: h\textasciicircum{}\{\textbackslash{}ast\}(p\_\{\textbackslash{}ast\}(\textbackslash{}bar\{F\}))\textbackslash{}longrightarrow p\_\{Y\textbackslash{}ast\}(\textbackslash{}bar\{h\}\textasciicircum{}\{\textbackslash{}ast\}(\textbackslash{}bar\{F\})); Source TeX: h\textasciicircum{}\{\textbackslash{}ast\}(p\_\{\textbackslash{}ast\}(\textbackslash{}bar\{F\}))\textbackslash{}rightarrow p\_\{Y\textasciicircum{}\{\textbackslash{}ast\}\}(\textbackslash{}bar\{h\}\textasciicircum{}\{\textbackslash{}ast\}(\textbackslash{}bar\{F\})); Note: L'étoile est imprimée en indice avec Y, non en exposant de Y.; Proposed edit: \textbackslash{}rightarrow p\_\{Y\textasciicircum{}\{\textbackslash{}ast\}\} =\allowbreak > \textbackslash{}rightarrow p\_\{Y\textbackslash{}ast\}\par
a fix from one model only\par
\par\noindent\hypertarget{DP0487}{}\textbf{DP0487 -- XII, p. 118 (one\_model\_fix)}\par
Corollaire 6.10 (i), condition 2°, homomorphisme portant sur R\textasciicircum{}i\par
\textbf{Claude Opus 5.5.} nothing reported here\par
\textbf{GPT-6.1 Sol.} Place: Corollaire 6.10 (i), condition 2°, homomorphisme portant sur R\textasciicircum{}i; Reported print: R\textasciicircum{}\{i\}p\_\{Y\textbackslash{}ast\}(\textbackslash{}bar\{h\}\textasciicircum{}\{\textbackslash{}ast\}(\textbackslash{}bar\{F\})); Source TeX: R\textasciicircum{}\{i\}p\_\{Y\textasciicircum{}\{\textbackslash{}ast\}\}(\textbackslash{}bar\{h\}\textasciicircum{}\{\textbackslash{}ast\}(\textbackslash{}bar\{F\})); Note: Le print met l'étoile en indice avec Y.; Proposed edit: R\textasciicircum{}\{i\}p\_\{Y\textasciicircum{}\{\textbackslash{}ast\}\} =\allowbreak > R\textasciicircum{}\{i\}p\_\{Y\textbackslash{}ast\}\par
a fix from one model only\par
\par\noindent\hypertarget{DP0488}{}\textbf{DP0488 -- XII, p. 118 (conflict)}\par
3°): H\textasciicircum{}i(Ȳ,\ensuremath{\bar{h}}*(p*(F)) | Corollaire 6.10 (i), condition 3°, premier groupe\par
\textbf{Claude Opus 5.5.} Place: 3°): H\textasciicircum{}i(Ȳ,\ensuremath{\bar{h}}*(p*(F)); Reported print: H\textasciicircum{}i(Ȳ,\ensuremath{\bar{h}}*(p*(F)) with one closing parenthesis missing; Source TeX: H\textasciicircum{}\{i\}(\textbackslash{}bar\{Y\}, \textbackslash{}bar\{h\}\textasciicircum{}\{\textbackslash{}ast\}(p\textasciicircum{}\{\textbackslash{}ast\}(F))); Proposed correction: H\textasciicircum{}\{i\}(\textbackslash{}bar\{Y\},\textbackslash{}bar\{h\}\textasciicircum{}\{\textbackslash{}ast\}(p\textasciicircum{}\{\textbackslash{}ast\}(F))); Note: The print's parentheses are unbalanced. The TeX silently corrects this.\par
\textbf{GPT-6.1 Sol.} Place: Corollaire 6.10 (i), condition 3°, premier groupe; Reported print: H\textasciicircum{}\{i\}(\textbackslash{}bar\{Y\},p\_Y\textasciicircum{}\{\textbackslash{}ast\}(h\textasciicircum{}\{\textbackslash{}ast\}(F))); Source TeX: H\textasciicircum{}\{i\}(\textbackslash{}bar\{Y\},p\_Y\textbackslash{}ast(h\textasciicircum{}\{\textbackslash{}ast\}(F))); Note: L'étoile imprimée est l'exposant du foncteur image inverse p\_Y\textasciicircum{}\{\textbackslash{}ast\} ; le TeX la laisse sur la ligne.; Proposed edit: p\_Y\textbackslash{}ast(h\textasciicircum{}\{\textbackslash{}ast\}(F)) =\allowbreak > p\_Y\textasciicircum{}\{\textbackslash{}ast\}(h\textasciicircum{}\{\textbackslash{}ast\}(F))\par
the two models differ at this place: one model fixes the TeX, the other calls it a misprint of the print\par
\par\noindent\hypertarget{DP0489}{}\textbf{DP0489 -- XII, p. 120 (cut)}\par
Fin de page, proposition 6.11\par
\textbf{Claude Opus 5.5.} nothing reported here\par
\textbf{GPT-6.1 Sol.} Place: Fin de page, proposition 6.11; Reported print: La page finit par « faisceau \ensuremath{\bar{F}} de ind-LL-groupes, l'homomorphisme de changement de base ».; Source TeX: \$h\textasciicircum{}\{\textbackslash{}ast\}(R\textasciicircum{}1p\_\{\textbackslash{}ast\}(\textbackslash{}bar\{F\}))\textbackslash{}rightarrow R\textasciicircum{}1p\_\{Y\textasciicircum{}\{\textbackslash{}ast\}\}(\textbackslash{}bar\{h\}\textasciicircum{}\{\textbackslash{}ast\}(\textbackslash{}bar\{F\}))\$\textbackslash{}pageoriginale est bijectif; Note: La formule entière des lignes source 1758–1759 appartient à la page 121. Le commentaire du fragment n'autorise du texte de la page suivante que dans sa dernière ligne, alors que la formule commence dès l'avant-dernière. Ne pas déplacer le texte ; rectifier la coupe ou son indication.\par
a model reports that the automatic cut gives lines of this page to its neighbour, or the neighbour's lines to this page: check the lines it names against the scan of the page that prints them (that page's checker did not have their TeX)\par
\par\noindent\hypertarget{DP0490}{}\textbf{DP0490 -- XII, p. 121 (cut)}\par
Début de la première ligne\par
\textbf{Claude Opus 5.5.} nothing reported here\par
\textbf{GPT-6.1 Sol.} Place: Début de la première ligne; Reported print: h\textasciicircum{}\{\textbackslash{}ast\}(R\textasciicircum{}1p\_\{\textbackslash{}ast\}(\textbackslash{}bar\{F\}))\textbackslash{}longrightarrow R\textasciicircum{}1p\_\{Y\textasciicircum{}\{\textbackslash{}ast\}\}(\textbackslash{}bar\{h\}\textasciicircum{}\{\textbackslash{}ast\}(\textbackslash{}bar\{F\})); Source TeX: R\textasciicircum{}1p\_\{Y\textasciicircum{}\{\textbackslash{}ast\}\}(\textbackslash{}bar\{h\}\textasciicircum{}\{\textbackslash{}ast\}(\textbackslash{}bar\{F\})); Note: La page commence par le membre gauche et la flèche de cet homomorphisme, avant le début de l'extrait fourni. La coupe doit reprendre la ligne précédente du source ; il ne s'agit pas d'une omission à réparer dans le texte de l'édition.\par
a model reports that the automatic cut gives lines of this page to its neighbour, or the neighbour's lines to this page: check the lines it names against the scan of the page that prints them (that page's checker did not have their TeX)\par
\par\noindent\hypertarget{DP0491}{}\textbf{DP0491 -- XII, p. 122 (conflict)}\par
second display, middle term of the composite | Homomorphisme composé, terme central\par
\textbf{Claude Opus 5.5.} Place: second display, middle term of the composite; Reported print: R\ensuremath{\Gamma}\_Y(h*(Rp\_*(\ensuremath{\bar{F}}))); Source TeX: \textbackslash{}RR\textbackslash{}Gamma\_Y(h\textasciicircum{}\{\textbackslash{}ast\}(\textbackslash{}bar\{F\})); Note: The print applies h* to Rp\_*(\ensuremath{\bar{F}}) in the middle term. The TeX drops Rp\_*.; Proposed edit: \textbackslash{}RR\textbackslash{}Gamma\_Y(h\textasciicircum{}\{\textbackslash{}ast\}(\textbackslash{}bar\{F\}))\textbackslash{}longrightarrow =\allowbreak > \textbackslash{}RR\textbackslash{}Gamma\_Y(h\textasciicircum{}\{\textbackslash{}ast\}(Rp\_\{\textbackslash{}ast\}(\textbackslash{}bar\{F\})))\textbackslash{}longrightarrow\par
\textbf{GPT-6.1 Sol.} Place: Homomorphisme composé, terme central; Reported print: \textbackslash{}RR\textbackslash{}Gamma\_Y(h\textasciicircum{}\{\textbackslash{}ast\}(Rp\_\{\textbackslash{}ast\}(\textbackslash{}bar\{F\})); Source TeX: \textbackslash{}RR\textbackslash{}Gamma\_Y(h\textasciicircum{}\{\textbackslash{}ast\}(\textbackslash{}bar\{F\})); Note: Le TeX omet Rp\_* à l'intérieur de h\textasciicircum{}*. Le fix restitue la lecture imprimée ; la parenthèse extérieure manquante dans cette lecture est une coquille du print signalée séparément.; Proposed edit: \textbackslash{}RR\textbackslash{}Gamma\_Y(h\textasciicircum{}\{\textbackslash{}ast\}(\textbackslash{}bar\{F\})) =\allowbreak > \textbackslash{}RR\textbackslash{}Gamma\_Y(h\textasciicircum{}\{\textbackslash{}ast\}(Rp\_\{\textbackslash{}ast\}(\textbackslash{}bar\{F\}))\par
the two models differ at this place\par
\par\noindent\hypertarget{DP0492}{}\textbf{DP0492 -- XII, p. 130 (conflict)}\par
tout nombre premier 1 | Démonstration de 8.5, phrase suivant la définition de \textbackslash{}varphi\textasciicircum{}q\par
\textbf{Claude Opus 5.5.} Place: tout nombre premier 1; Reported print: premier 1; Source TeX: premier 1; Proposed correction: premier \$\textbackslash{}ell\$; Note: The digit 1 stands for \ensuremath{\ell}. The TeX keeps the printed 1.\par
\textbf{GPT-6.1 Sol.} Place: Démonstration de 8.5, phrase suivant la définition de \textbackslash{}varphi\textasciicircum{}q; Reported print: tout nombre premier l; Source TeX: tout nombre premier 1; Note: Le l droit de machine désigne le nombre premier \textbackslash{}ell des formules. Son glyphe est très proche du chiffre 1 dans cette frappe ; le TeX donne le chiffre 1.; Proposed edit: tout nombre premier 1 =\allowbreak > tout nombre premier \$\textbackslash{}ell\$\par
the two models differ at this place: one model fixes the TeX, the other calls it a misprint of the print\par
\par\noindent\hypertarget{DP0493}{}\textbf{DP0493 -- XIII, p. 132 (conflict)}\par
Lemme 1.2, définition de \ensuremath{\varepsilon} | Lemme 1.2, définition du foncteur epsilon\par
\textbf{Claude Opus 5.5.} Place: Lemme 1.2, définition de \ensuremath{\varepsilon}; Reported print: (Sch)/S \ensuremath{\longrightarrow} (Ens); Source TeX: \$(\textbackslash{}Sch)/S(\textbackslash{}Ens)\$; Note: The arrow between (Sch)/S and (Ens) is missing in the TeX.; Proposed edit: \$(\textbackslash{}Sch)/S(\textbackslash{}Ens)\$ =\allowbreak > \$(\textbackslash{}Sch)/S\textbackslash{}rightarrow(\textbackslash{}Ens)\$\par
\textbf{GPT-6.1 Sol.} Place: Lemme 1.2, définition du foncteur epsilon; Reported print: \$(\textbackslash{}Sch)/S\textbackslash{}longrightarrow (\textbackslash{}Ens)\$; Source TeX: \$(\textbackslash{}Sch)/S(\textbackslash{}Ens)\$; Note: La flèche entre la catégorie source et la catégorie des ensembles manque.; Proposed edit: \$(\textbackslash{}Sch)/S(\textbackslash{}Ens)\$ =\allowbreak > \$(\textbackslash{}Sch)/S\textbackslash{}longrightarrow (\textbackslash{}Ens)\$\par
the two models differ at this place\par
\par\noindent\hypertarget{DP0494}{}\textbf{DP0494 -- XIII, p. 132 (fix\_unplaced)}\par
Signature sous le titre\par
\textbf{Claude Opus 5.5.} nothing reported here\par
\textbf{GPT-6.1 Sol.} Place: Signature sous le titre; Reported print: par M. ARTIN; Source TeX: M. Artin; Note: Le mot « par » manque.; Proposed edit: M. Artin =\allowbreak > par M. Artin\par
a fix whose old text is not found once on the page\par
\par\noindent\hypertarget{DP0495}{}\textbf{DP0495 -- XIII, p. 133 (one\_model\_fix)}\par
Lemme 1.3, formule définissant E'\par
\textbf{Claude Opus 5.5.} nothing reported here\par
\textbf{GPT-6.1 Sol.} Place: Lemme 1.3, formule définissant E'; Reported print: \$E'=\allowbreak E\textbackslash{}otimes\_\{\textbackslash{}Oo\_S\}\textbackslash{}Oo\_\{S'\}\$; Source TeX: \$E'=\allowbreak E\textbackslash{}otimes\_\{O\_S\}\textbackslash{}Oo\_\{S'\}\$; Note: Le O sous le produit tensoriel est souligné dans le print, comme les autres faisceaux structuraux ; le TeX donne ici un O ordinaire.; Proposed edit: E\textbackslash{}otimes\_\{O\_S\}\textbackslash{}Oo\_\{S'\} =\allowbreak > E\textbackslash{}otimes\_\{\textbackslash{}Oo\_S\}\textbackslash{}Oo\_\{S'\}\par
a fix from one model only\par
\par\noindent\hypertarget{DP0496}{}\textbf{DP0496 -- XIII, p. 139 (one\_model\_fix)}\par
Démonstration du lemme 2.2, définition de J₀, dernière variable\par
\textbf{Claude Opus 5.5.} nothing reported here\par
\textbf{GPT-6.1 Sol.} Place: Démonstration du lemme 2.2, définition de J₀, dernière variable; Reported print: k[X\_\{\textbackslash{}circ\},\textbackslash{}ldots,x\_N]; Source TeX: k[X\_\{\textbackslash{}circ\},\textbackslash{}ldots,X\_N]; Note: La dernière variable est imprimée avec un x minuscule.; Proposed edit: k[X\_\{\textbackslash{}circ\},\textbackslash{}ldots,X\_N] =\allowbreak > k[X\_\{\textbackslash{}circ\},\textbackslash{}ldots,x\_N]\par
a fix from one model only\par
\par\noindent\hypertarget{DP0497}{}\textbf{DP0497 -- XIII, p. 140 (conflict)}\par
Démonstration of 2.3, «de dimension 1» | Démonstration de 2.3, cas de dimension nulle\par
\textbf{Claude Opus 5.5.} Place: Démonstration of 2.3, «de dimension 1»; Reported print: de dimension 1,; Source TeX: \textbackslash{}hyperref[XIII.sec1]\{1\}, le; Note: «1» is a dimension, not a reference; the TeX turns it into a spurious link to §1. The printed output is the same.; Proposed edit: \textbackslash{}hyperref[XIII.sec1]\{1\}, le =\allowbreak > 1, le\par
\textbf{GPT-6.1 Sol.} Place: Démonstration de 2.3, cas de dimension nulle; Reported print: le cas où X est de dimension nulle étant triviale; Source TeX: cas où \$X\$ est de dimension nulle étant triviale; Proposed correction: le cas où \$X\$ est de dimension nulle étant trivial; Note: Le TeX conserve « triviale » du tirage ; l'adjectif doit s'accorder avec « cas ».\par
the two models differ at this place: one model fixes the TeX, the other calls it a misprint of the print\par
\par\noindent\hypertarget{DP0498}{}\textbf{DP0498 -- XIII, p. 141 (one\_model\_fix)}\par
Remarque 2.4, the two occurrences of K in the final parenthesis\par
\textbf{Claude Opus 5.5.} nothing reported here\par
\textbf{GPT-6.1 Sol.} Place: Remarque 2.4, the two occurrences of K in the final parenthesis; Reported print: \$Y=\allowbreak V(\textbackslash{}underline\{K\})\$, \$J\textbackslash{}cap \textbackslash{}underline\{K\}=\allowbreak 0\$; Source TeX: \$Y=\allowbreak V(\textbackslash{}bar\{K\})\$, \$J\textbackslash{}cap \textbackslash{}bar\{K\}=\allowbreak 0\$; Note: Both strokes are below K, in the typescript's sheaf notation, not overbars. This is mathematical notation rather than emphasis.; Proposed edit: \$Y=\allowbreak V(\textbackslash{}bar\{K\})\$, \$J\textbackslash{}cap \textbackslash{}bar\{K\}=\allowbreak 0\$ =\allowbreak > \$Y=\allowbreak V(\textbackslash{}underline\{K\})\$, \$J\textbackslash{}cap \textbackslash{}underline\{K\}=\allowbreak 0\$\par
a fix from one model only\par
\par\noindent\hypertarget{DP0499}{}\textbf{DP0499 -- XIII, p. 141 (fix\_unplaced)}\par
Remarque 2.4, last line\par
\textbf{Claude Opus 5.5.} Place: Remarque 2.4, last line; Reported print: Y =\allowbreak  V(\ensuremath{\underline{K}}), J\ensuremath{\cap}\ensuremath{\underline{K}} =\allowbreak  0 (K underlined: sheaf of ideals); Source TeX: Y=\allowbreak V(\textbackslash{}bar\{K\}), J\textbackslash{}cap \textbackslash{}bar\{K\}=\allowbreak 0; Note: The print underlines K, the typescript's mark for a script (sheaf) letter, as it does for \ensuremath{\underline{O}}\_X. The TeX's overline gives another sign. If the edition has \textbackslash{}Kk, use it instead of \textbackslash{}mathcal\{K\}. Check the exact delimiters of "old" against the source: the given substring assumes the TeX reads '\$Y=\allowbreak V(\textbackslash{}bar\{K\})\$, \$J\textbackslash{}cap \textbackslash{}bar\{K\}=\allowbreak 0\$'.; Proposed edit: Y=\allowbreak V(\textbackslash{}bar\{K\}), \$J\textbackslash{}cap \textbackslash{}bar\{K\}=\allowbreak 0 =\allowbreak > Y=\allowbreak V(\textbackslash{}mathcal\{K\})\$, \$J\textbackslash{}cap \textbackslash{}mathcal\{K\}=\allowbreak 0\par
\textbf{GPT-6.1 Sol.} nothing reported here\par
a fix whose old text is not found once on the page\par
\par\noindent\hypertarget{DP0500}{}\textbf{DP0500 -- XIII, p. 143 (one\_model\_fix)}\par
Introduction to Corollaire 3.2, reference after «la démonstration de»\par
\textbf{Claude Opus 5.5.} nothing reported here\par
\textbf{GPT-6.1 Sol.} Place: Introduction to Corollaire 3.2, reference after «la démonstration de»; Reported print: 5.1 (iii); Source TeX: \textbackslash{}hyperrefext[XII][XII.thm5.1]\{5.1\} (iii); Note: The print does not repeat XII after the preceding reference to XII 8.6. The external-reference macro adds that numeral.; Proposed edit: \textbackslash{}hyperrefext[XII][XII.thm5.1]\{5.1\} =\allowbreak > \textbackslash{}hyperref[XII.thm5.1]\{5.1\}\par
a fix from one model only\par
\par\noindent\hypertarget{DP0501}{}\textbf{DP0501 -- XIV, p. 145 (one\_model\_fix)}\par
Byline below the exposé title\par
\textbf{Claude Opus 5.5.} nothing reported here\par
\textbf{GPT-6.1 Sol.} Place: Byline below the exposé title; Reported print: par M. ARTIN; Source TeX: M. Artin; Note: The word « par » is omitted.; Proposed edit: \textbackslash{}begin\{center\} M. Artin =\allowbreak > \textbackslash{}begin\{center\} par M. Artin\par
a fix from one model only\par
\par\noindent\hypertarget{DP0502}{}\textbf{DP0502 -- XV, p. 168 (one\_model\_fix)}\par
Ligne d’auteur, sous le titre\par
\textbf{Claude Opus 5.5.} nothing reported here\par
\textbf{GPT-6.1 Sol.} Place: Ligne d’auteur, sous le titre; Reported print: par M. ARTIN; Source TeX: M. Artin; Note: Le mot « par » manque devant le nom de l’auteur.; Proposed edit: \textbackslash{}begin\{center\} M. Artin =\allowbreak > \textbackslash{}begin\{center\} par M. Artin\par
a fix from one model only\par
\par\noindent\hypertarget{DP0503}{}\textbf{DP0503 -- XV, p. 171 (conflict)}\par
D'après (i), … | Proof of 1.1, sentence beginning « D'après (i) »\par
\textbf{Claude Opus 5.5.} Place: D'après (i), …; Reported print: D'après (i), on a f\_*(F) \ensuremath{\to} g\_*g\textasciicircum{}*f\textasciicircum{}*F).; Source TeX: \textbackslash{}sisi\{le morphisme d'adjonction \$f\_\{\textbackslash{}ast\}F\textbackslash{}rightarrow g\_\{\textbackslash{}ast\}g\textasciicircum{}\{\textbackslash{}ast\}f\_\{\textbackslash{}ast\}F\$ est injectif (resp.\textasciitilde{}bijectif)\}\{on a \$f\_\{\textbackslash{}ast\}(F)\textbackslash{}rightarrow g\_\{\textbackslash{}ast\}g\textasciicircum{}\{\textbackslash{}ast\}f\textasciicircum{}\{\textbackslash{}ast\}F\$)\}; Note: The two arguments of \textbackslash{}sisi are swapped: the print's reading (a garbled phrase, with f\textasciicircum{}*F for f\_*F) stands in the second argument and the editors' correction in the first. The same swap recurs in this file (see coordinator).; Proposed edit: \textbackslash{}sisi\{le morphisme d'adjonction \$f\_\{\textbackslash{}ast\}F\textbackslash{}rightarrow g\_\{\textbackslash{}ast\}g\textasciicircum{}\{\textbackslash{}ast\}f\_\{\textbackslash{}ast\}F\$ est injectif (resp.\textasciitilde{}bijectif)\}\{on a \$f\_\{\textbackslash{}ast\}(F)\textbackslash{}rightarrow g\_\{\textbackslash{}ast\}g\textasciicircum{}\{\textbackslash{}ast\}f\textasciicircum{}\{\textbackslash{}ast\}F\$)\} =\allowbreak > \textbackslash{}sisi\{on a \$f\_\{\textbackslash{}ast\}(F)\textbackslash{}rightarrow g\_\{\textbackslash{}ast\}g\textasciicircum{}\{\textbackslash{}ast\}f\textasciicircum{}\{\textbackslash{}ast\}F\$)\}\{le morphisme d'adjonction \$f\_\{\textbackslash{}ast\}F\textbackslash{}rightarrow g\_\{\textbackslash{}ast\}g\textasciicircum{}\{\textbackslash{}ast\}f\_\{\textbackslash{}ast\}F\$ est injectif (resp.\textasciitilde{}bijectif)\}\par
\textbf{GPT-6.1 Sol.} Place: Proof of 1.1, sentence beginning « D'après (i) »; Reported print: on a \$f\_\{\textbackslash{}ast\}(F)\textbackslash{}rightarrow g\_\{\textbackslash{}ast\}g\textasciicircum{}\{\textbackslash{}ast\}f\textasciicircum{}\{\textbackslash{}ast\}F\$).; Source TeX: le morphisme d'adjonction \$f\_\{\textbackslash{}ast\}F\textbackslash{}rightarrow g\_\{\textbackslash{}ast\}g\textasciicircum{}\{\textbackslash{}ast\}f\_\{\textbackslash{}ast\}F\$ est injectif (resp.\textasciitilde{}bijectif); Note: The first argument of \textbackslash{}sisi rewrites the printed clause and adds an explicit injectivity/bijectivity assertion. This edit restores the printed wording while retaining the corrections to the print's typographical errors recorded separately below.; Proposed edit: le morphisme d'adjonction \$f\_\{\textbackslash{}ast\}F\textbackslash{}rightarrow g\_\{\textbackslash{}ast\}g\textasciicircum{}\{\textbackslash{}ast\}f\_\{\textbackslash{}ast\}F\$ est injectif (resp.\textasciitilde{}bijectif) =\allowbreak > on a \$f\_\{\textbackslash{}ast\}(F)\textbackslash{}rightarrow g\_\{\textbackslash{}ast\}g\textasciicircum{}\{\textbackslash{}ast\}f\_\{\textbackslash{}ast\}F\$\par
the two models differ at this place\par
\par\noindent\hypertarget{DP0504}{}\textbf{DP0504 -- XV, p. 171 (one\_model\_fix)}\par
est bijectif VIII 5.6.\par
\textbf{Claude Opus 5.5.} Place: est bijectif VIII 5.6.; Reported print: bijectif VIII 5.6. D'après; Source TeX: \textbackslash{}hyperrefext[VIII][VIII.cor5.6]\{5.6\} D'après; Note: The period that ends the sentence is missing.; Proposed edit: \textbackslash{}hyperrefext[VIII][VIII.cor5.6]\{5.6\} D'après =\allowbreak > \textbackslash{}hyperrefext[VIII][VIII.cor5.6]\{5.6\}. D'après\par
\textbf{GPT-6.1 Sol.} nothing reported here\par
a fix from one model only\par
\par\noindent\hypertarget{DP0505}{}\textbf{DP0505 -- XV, p. 174 (one\_model\_fix)}\par
Proof of 1.6, implication (ii) \ensuremath{\Rightarrow} (i), after « le faisceau associé est nul »\par
\textbf{Claude Opus 5.5.} nothing reported here\par
\textbf{GPT-6.1 Sol.} Place: Proof of 1.6, implication (ii) \ensuremath{\Rightarrow} (i), after « le faisceau associé est nul »; Reported print: donc le faisceau associé est nul, i.e. (ii) \ensuremath{\Rightarrow} (i); Source TeX: donc le faisceau associé est nul\textbackslash{}sisi\{ si \$1 \textbackslash{}leq q \textbackslash{}leq n\$\}\{\}, \textbackslash{}ie (ii) \$\textbackslash{}Rightarrow\$ (i); Note: The printed text passes directly from « nul » to the comma and « i.e. »; the restriction in the first argument of \textbackslash{}sisi is not printed.; Proposed edit: \textbackslash{}sisi\{ si \$1 \textbackslash{}leq q \textbackslash{}leq n\$\}\{\} =\allowbreak > \par
a fix from one model only\par
\par\noindent\hypertarget{DP0506}{}\textbf{DP0506 -- XV, p. 174 (one\_model\_fix)}\par
Proof of (i) \ensuremath{\Rightarrow} (iii), sentence introducing f′\par
\textbf{Claude Opus 5.5.} nothing reported here\par
\textbf{GPT-6.1 Sol.} Place: Proof of (i) \ensuremath{\Rightarrow} (iii), sentence introducing f′; Reported print: Soit \$f':X'\textbackslash{}rightarrow X\$ le morphisme.; Source TeX: Soit \$f':X'\textbackslash{}rightarrow X\$ le morphisme \textbackslash{}sisi\{déduit de\textasciitilde{}\$f\$ par le changement de base\textasciitilde{}\$g\$\}\{\}.; Note: The printed sentence ends after « le morphisme »; the descriptive clause in the first argument of \textbackslash{}sisi is absent.; Proposed edit: le morphisme \textbackslash{}sisi\{déduit de\textasciitilde{}\$f\$ par le changement de base\textasciitilde{}\$g\$\}\{\}. =\allowbreak > le morphisme.\par
a fix from one model only\par
\par\noindent\hypertarget{DP0507}{}\textbf{DP0507 -- XV, p. 176 (one\_model\_fix)}\par
Prop. 1.9 (i), « de ind L-groupes »\par
\textbf{Claude Opus 5.5.} Place: Prop. 1.9 (i), « de ind L-groupes »; Reported print: (resp. de ind \ensuremath{𝕃}-groupes, resp.; Source TeX: (\textbackslash{}resp de ind\textbackslash{}sisi\{-\}\{ \}\$\textbackslash{}LL\$-groupes, \textbackslash{}resp abélien de; Note: The first argument of \textbackslash{}sisi is meant to hold the print's reading, but it gives « ind-\ensuremath{𝕃} » with a hyphen. The print has a space, « ind \ensuremath{𝕃}-groupes ». So the \textbackslash{}sisi records a print reading that is not on the page.; Proposed edit: ind\textbackslash{}sisi\{-\}\{ \}\$\textbackslash{}LL\$-groupes =\allowbreak > ind \$\textbackslash{}LL\$-groupes\par
\textbf{GPT-6.1 Sol.} nothing reported here\par
a fix from one model only\par
\par\noindent\hypertarget{DP0508}{}\textbf{DP0508 -- XV, p. 178 (one\_model\_fix)}\par
Proposition 1.10 (ii), renvois précédant le premier morphisme de changement de base\par
\textbf{Claude Opus 5.5.} nothing reported here\par
\textbf{GPT-6.1 Sol.} Place: Proposition 1.10 (ii), renvois précédant le premier morphisme de changement de base; Reported print: XII 4.1, 4.2; Source TeX: \textbackslash{}hyperrefext[XII][XII.eq4.1]\{4.1\}, \textbackslash{}hyperrefext[XII][XII.eq4.2]\{4.2\}; Note: Le second renvoi ajoute un XII que l'imprimé ne répète pas. Convention signalée une seule fois ; trois exemples et les autres occurrences à traiter sont indiqués au coordinateur.; Proposed edit: \textbackslash{}hyperrefext[XII][XII.eq4.2]\{4.2\} =\allowbreak > \textbackslash{}ref\{XII.eq4.2\}\par
a fix from one model only\par
\par\noindent\hypertarget{DP0509}{}\textbf{DP0509 -- XV, p. 185 (conflict)}\par
footnote to Démonstration de 1.15 | Note de bas de page attachée à « Démonstration de 1.15 »\par
\textbf{Claude Opus 5.5.} Place: footnote to Démonstration de 1.15; Reported print: à se reporter à 1.18.; Source TeX: \textbackslash{}sisi\{la remarque \}\{\}\textbackslash{}hyperref[XV.rem1.18]\{1.18\}; Note: The \textbackslash{}sisi gives 'la remarque' as the printed reading, but the print has only 'à se reporter à 1.18.' with no 'la remarque'. The edition's output, B =\allowbreak  empty, is the same either way.; Proposed edit: \textbackslash{}sisi\{la remarque \}\{\} =\allowbreak > \par
\textbf{GPT-6.1 Sol.} Place: Note de bas de page attachée à « Démonstration de 1.15 »; Reported print: est invité à se reporter à 1.18.; Source TeX: est invité à se reporter à \textbackslash{}sisi\{la remarque \}\{\}\textbackslash{}hyperref[XV.rem1.18]\{1.18\}.; Note: Les mots « la remarque » sont absents de l'imprimé. Ils figurent à tort dans le premier argument de \textbackslash{}sisi, celui qui doit représenter la lecture imprimée ; le second argument est déjà vide.; Proposed edit: \textbackslash{}sisi\{la remarque \}\{\} =\allowbreak > \textbackslash{}sisi\{\}\{\}\par
the two models differ at this place\par
\par\noindent\hypertarget{DP0510}{}\textbf{DP0510 -- XV, p. 193 (one\_model\_fix)}\par
Démonstration de 2.1, reference to SGA 1\par
\textbf{Claude Opus 5.5.} Place: Démonstration de 2.1, reference to SGA 1; Reported print: supposer SGA1 II 1.1 que g; Source TeX: \textbackslash{}sisi\{(\}\{\}\textbackslash{}SGA 1 II 1.1\textbackslash{}sisi\{)\}\{\}; Note: The A argument of \textbackslash{}sisi records parentheses as the print's reading, but the print has none around « SGA1 II 1.1 ». The edition's printed output matches the print; only the recorded reading is wrong.; Proposed edit: \textbackslash{}sisi\{(\}\{\}\textbackslash{}SGA 1 II 1.1\textbackslash{}sisi\{)\}\{\} =\allowbreak > \textbackslash{}SGA 1 II 1.1\par
\textbf{GPT-6.1 Sol.} nothing reported here\par
a fix from one model only\par
\par\noindent\hypertarget{DP0511}{}\textbf{DP0511 -- XV, p. 194 (agreed\_fix\_held)}\par
« section a de O\_Ỹ » | Paragraphe initial, « En relevant a\textasciicircum{}0 en une section a de… »\par
\textbf{Claude Opus 5.5.} Place: « section a de O\_Ỹ »; Reported print: O\_\{Ỹ\} (tilde on Y); Source TeX: \textbackslash{}mathcal\{O\}\_\{\textbackslash{}sisi\{Y\}\{\textbackslash{}widetilde\{Y\}\}\}\$ et; Note: The A argument of \textbackslash{}sisi gives Y as the print's reading, but the print already has Ỹ with a tilde. The printed output of the edition agrees with the print.; Proposed edit: \$\textbackslash{}mathcal\{O\}\_\{\textbackslash{}sisi\{Y\}\{\textbackslash{}widetilde\{Y\}\}\}\$ et =\allowbreak > \$\textbackslash{}mathcal\{O\}\_\{\textbackslash{}widetilde\{Y\}\}\$ et\par
\textbf{GPT-6.1 Sol.} Place: Paragraphe initial, « En relevant a\textasciicircum{}0 en une section a de… »; Reported print: \textbackslash{}mathcal\{O\}\_\{\textbackslash{}widetilde\{Y\}\}; Source TeX: \textbackslash{}mathcal\{O\}\_\{\textbackslash{}sisi\{Y\}\{\textbackslash{}widetilde\{Y\}\}\}; Note: Le Y en indice porte déjà un tilde manuscrit dans l'imprimé ; le premier argument de \textbackslash{}sisi ne reproduit donc pas la lecture imprimée.; Proposed edit: \textbackslash{}mathcal\{O\}\_\{\textbackslash{}sisi\{Y\}\{\textbackslash{}widetilde\{Y\}\}\}\$ =\allowbreak > \textbackslash{}mathcal\{O\}\_\{\textbackslash{}widetilde\{Y\}\}\$\par
both models give the same fix, held: it changes the editors' apparatus (\textbackslash{}sisi): a \textbackslash{}sisi correction, note, label, cross-reference or citation is the editors' record and is not removed or retargeted on the models' word\par
\par\noindent\hypertarget{DP0512}{}\textbf{DP0512 -- XV, p. 194 (agreed\_fix\_held)}\par
Posons A =\allowbreak  \ensuremath{\Gamma}(Ỹ, O\_Ỹ) | Paragraphe initial, formule après « Posons »\par
\textbf{Claude Opus 5.5.} Place: Posons A =\allowbreak  \ensuremath{\Gamma}(Ỹ, O\_Ỹ); Reported print: \ensuremath{\Gamma}(Ỹ,O\_Ỹ) (tildes); Source TeX: (\textbackslash{}sisi\{Y\}\{\textbackslash{}widetilde\{Y\}\},\textbackslash{}mathcal\{O\}\_\{\textbackslash{}sisi\{Y\}\{\textbackslash{}widetilde\{Y\}\}\}); Note: The print has tildes, or tilde-like accents, on both Y. The \textbackslash{}sisi A argument records an untilded Y that the print does not show.; Proposed edit: (\textbackslash{}sisi\{Y\}\{\textbackslash{}widetilde\{Y\}\},\textbackslash{}mathcal\{O\}\_\{\textbackslash{}sisi\{Y\}\{\textbackslash{}widetilde\{Y\}\}\}) =\allowbreak > (\textbackslash{}widetilde\{Y\},\textbackslash{}mathcal\{O\}\_\{\textbackslash{}widetilde\{Y\}\})\par
\textbf{GPT-6.1 Sol.} Place: Paragraphe initial, formule après « Posons »; Reported print: A=\allowbreak \textbackslash{}Gamma(\textbackslash{}widetilde\{Y\},\textbackslash{}mathcal\{O\}\_\{\textbackslash{}widetilde\{Y\}\}); Source TeX: A=\allowbreak \textbackslash{}Gamma (\textbackslash{}sisi\{Y\}\{\textbackslash{}widetilde\{Y\}\},\textbackslash{}mathcal\{O\}\_\{\textbackslash{}sisi\{Y\}\{\textbackslash{}widetilde\{Y\}\}\}); Note: Les deux Y portent des tildes dans l'imprimé, absents des premières lectures données par les \textbackslash{}sisi.; Proposed edit: (\textbackslash{}sisi\{Y\}\{\textbackslash{}widetilde\{Y\}\},\textbackslash{}mathcal\{O\}\_\{\textbackslash{}sisi\{Y\}\{\textbackslash{}widetilde\{Y\}\}\}) =\allowbreak > (\textbackslash{}widetilde\{Y\},\textbackslash{}mathcal\{O\}\_\{\textbackslash{}widetilde\{Y\}\})\par
both models give the same fix, held: it changes the editors' apparatus (\textbackslash{}sisi): a \textbackslash{}sisi correction, note, label, cross-reference or citation is the editors' record and is not removed or retargeted on the models' word\par
\par\noindent\hypertarget{DP0513}{}\textbf{DP0513 -- XV, p. 195 (one\_model\_fix)}\par
reference to EGA IV after « excellent »\par
\textbf{Claude Opus 5.5.} Place: reference to EGA IV after « excellent »; Reported print: excellent EGA IV 7.8.3 (iii), 7.8.6 (i) , donc; Source TeX: \textbackslash{}sisi\{(\}\{\}\textbackslash{}EGA IV 7.8.3 (iii), 7.8.6 (i)\textbackslash{}sisi\{)\}\{\}; Note: The \textbackslash{}sisi A argument records enclosing parentheses that the print does not have. The edition's output matches the print.; Proposed edit: \textbackslash{}sisi\{(\}\{\}\textbackslash{}EGA IV 7.8.3 (iii), 7.8.6 (i)\textbackslash{}sisi\{)\}\{\} =\allowbreak > \textbackslash{}EGA IV 7.8.3 (iii), 7.8.6 (i)\par
\textbf{GPT-6.1 Sol.} nothing reported here\par
a fix from one model only\par
\par\noindent\hypertarget{DP0514}{}\textbf{DP0514 -- XV, p. 201 (one\_model\_fix)}\par
diagram, label of the top horizontal arrow\par
\textbf{Claude Opus 5.5.} Place: diagram, label of the top horizontal arrow; Reported print: i (no hat visible); Source TeX: \textbackslash{}ar[r]\textasciicircum{}\{\textbackslash{}sisi\{\textbackslash{}hat\{i\}\}\{i\}\}; Note: The \textbackslash{}sisi gives \textbackslash{}hat\{i\} as the print's reading, but the label over the arrow Û \ensuremath{\to} \ensuremath{\hat{X}} looks like a plain i in the scan. The mathematically right reading is î, which is the form used in the text below (î\_*(Ĉ)). The two arguments look swapped. The stroke is small and a hat cannot be ruled out.; Proposed edit: \textbackslash{}ar[r]\textasciicircum{}\{\textbackslash{}sisi\{\textbackslash{}hat\{i\}\}\{i\}\} =\allowbreak > \textbackslash{}ar[r]\textasciicircum{}\{\textbackslash{}sisi\{i\}\{\textbackslash{}hat\{i\}\}\}\par
\textbf{GPT-6.1 Sol.} nothing reported here\par
a fix from one model only\par
\par\noindent\hypertarget{DP0515}{}\textbf{DP0515 -- XVI, p. 207 (cut)}\par
Start of the supplied TeX cut, before §1\par
\textbf{Claude Opus 5.5.} nothing reported here\par
\textbf{GPT-6.1 Sol.} Place: Start of the supplied TeX cut, before §1; Reported print: 1. Le théorème de changement de base par un morphisme lisse.; Source TeX: \textbackslash{}chapter\{Théorème de changement de base par un morphisme lisse, et applications\}; Note: The printed page begins with §1. The leading exposé title and author block belong before this page, but the cut's comment does not identify them as neighbouring-page material. Recut at the section heading; do not delete this material from the edition.\par
a model reports that the automatic cut gives lines of this page to its neighbour, or the neighbour's lines to this page: check the lines it names against the scan of the page that prints them (that page's checker did not have their TeX)\par
\par\noindent\hypertarget{DP0516}{}\textbf{DP0516 -- XVI, p. 212 (conflict)}\par
où \$f:X\textbackslash{}rightarrow \textbackslash{}Spec(k)\$, ..., \$g:\textbackslash{}Spec(K)\textbackslash{}rightarrow \textbackslash{}Spec(k)\$ | Premier paragraphe, définition du morphisme g\par
\textbf{Claude Opus 5.5.} Place: où \$f:X\textbackslash{}rightarrow \textbackslash{}Spec(k)\$, ..., \$g:\textbackslash{}Spec(K)\textbackslash{}rightarrow \textbackslash{}Spec(k)\$; Reported print: g =\allowbreak  Speo(K) \ensuremath{\longrightarrow} Spec(k); Source TeX: \$g:\textbackslash{}Spec(K)\textbackslash{}rightarrow \textbackslash{}Spec(k)\$; Proposed correction: \$g:\textbackslash{}Spec(K)\textbackslash{}rightarrow \textbackslash{}Spec(k)\$; Note: The print seems to have '=\allowbreak ' where the parallel f, f', g' have ':'. The sign is doubtful in the scan. The TeX silently corrects to ':'.\par
\textbf{GPT-6.1 Sol.} Place: Premier paragraphe, définition du morphisme g; Reported print: \$g =\allowbreak  \textbackslash{}Spec(K)\textbackslash{}longrightarrow \textbackslash{}Spec(k)\$; Source TeX: \$g:\textbackslash{}Spec(K)\textbackslash{}rightarrow \textbackslash{}Spec(k)\$; Note: L'imprimé emploie un signe égal après g, contrairement aux deux-points du TeX.; Proposed edit: g:\textbackslash{}Spec(K) =\allowbreak > g=\allowbreak \textbackslash{}Spec(K)\par
the two models differ at this place: one model fixes the TeX, the other calls it a misprint of the print\par
\par\noindent\hypertarget{DP0517}{}\textbf{DP0517 -- XVI, p. 225 (one\_model\_fix)}\par
Troisième paragraphe, référence après « localement 1-asphérique »\par
\textbf{Claude Opus 5.5.} nothing reported here\par
\textbf{GPT-6.1 Sol.} Place: Troisième paragraphe, référence après « localement 1-asphérique »; Reported print: (XV 2.1, 1.11 (ii)); Source TeX: \textbackslash{}hyperrefext[XV][XV.thm2.1]\{2.1\}, \textbackslash{}hyperrefext[XV][XV.def1.11]\{1.11\}; Note: Le second renvoi n'est pas précédé d'un second « XV » dans l'imprimé ; \textbackslash{}hyperrefext l'ajoute.; Proposed edit: \textbackslash{}hyperrefext[XV][XV.def1.11]\{1.11\} =\allowbreak > \textbackslash{}hyperref[XV.def1.11]\{1.11\}\par
a fix from one model only\par
\par\noindent\hypertarget{DP0518}{}\textbf{DP0518 -- XVI, p. 225 (cut)}\par
Fin de page, dernière ligne\par
\textbf{Claude Opus 5.5.} nothing reported here\par
\textbf{GPT-6.1 Sol.} Place: Fin de page, dernière ligne; Reported print: « pour conclure que V induit un re- »; Source TeX: (\textbackslash{}SGA 1 X 3.6) pour conclure que \$V\$ induit un; Note: L'imprimé commence encore le mot suivant par « re- », tandis que la coupe TeX s'arrête à « un ». Vérifier la ligne source suivante avec la p. 226 et inclure la ligne partagée dans la coupe ; ne pas ajouter un second exemplaire du mot au texte.\par
a model reports that the automatic cut gives lines of this page to its neighbour, or the neighbour's lines to this page: check the lines it names against the scan of the page that prints them (that page's checker did not have their TeX)\par
\par\noindent\hypertarget{DP0519}{}\textbf{DP0519 -- XVI, p. 235 (cut)}\par
end of page\par
\textbf{Claude Opus 5.5.} Place: end of page; Reported print: et (R\textasciicircum{}q g\_\{cl*\}) j\_! F \ensuremath{\simeq} (R\textasciicircum{}q ḡ\_\{cl*\}) \ensuremath{\bar{j}}\_! \ensuremath{\bar{F}} ,; Source TeX: et  ; Note: The last display, (R\textasciicircum{}q g\_cl*) j\_!F \ensuremath{\simeq} (R\textasciicircum{}q ḡ\_cl*) \ensuremath{\bar{j}}\_!\ensuremath{\bar{F}} , is printed at the foot of p. 235. This page's TeX ends at "et", so the display stands in the next cut, which really begins after it.\par
\textbf{GPT-6.1 Sol.} nothing reported here\par
a model reports that the automatic cut gives lines of this page to its neighbour, or the neighbour's lines to this page: check the lines it names against the scan of the page that prints them (that page's checker did not have their TeX)\par
\par\noindent\hypertarget{DP0520}{}\textbf{DP0520 -- XVI, p. 235 (cut)}\par
Bottom of page, after the final « et » in the supplied TeX\par
\textbf{Claude Opus 5.5.} nothing reported here\par
\textbf{GPT-6.1 Sol.} Place: Bottom of page, after the final « et » in the supplied TeX; Reported print: (R\textasciicircum{}\{q\}g\_\{\textbackslash{}cl\textbackslash{}ast\})j\_\{!\}F\textbackslash{}simeq(R\textasciicircum{}\{q\}\textbackslash{}bar\{g\}\_\{\textbackslash{}cl\textbackslash{}ast\})\textbackslash{}bar\{j\}\_\{!\}\textbackslash{}bar\{F\},; Source TeX: The supplied cut ends at « et » after the étale comparison.; Note: The page also prints the classical comparison displayed here. The cut ends too early: p. 235 ends after this display and its comma. Include the following source lines in this page's cut rather than inserting duplicate text.\par
a model reports that the automatic cut gives lines of this page to its neighbour, or the neighbour's lines to this page: check the lines it names against the scan of the page that prints them (that page's checker did not have their TeX)\par
\par\noindent\hypertarget{DP0521}{}\textbf{DP0521 -- XVI, p. 236 (cut)}\par
start of the page's TeX, before 'd'où 4.2.'\par
\textbf{Claude Opus 5.5.} Place: start of the page's TeX, before 'd'où 4.2.'; Reported print: the page begins with 'd'où 4.2 .'; Source TeX: \$\$ (R\textasciicircum{}\{q\}g\_\{\textbackslash{}cl \textbackslash{}ast\})j\_\{!\}F\textbackslash{}simeq (R\textasciicircum{}\{q\}\textbackslash{}bar\{g\}\_\{\textbackslash{}cl \textbackslash{}ast\})\textbackslash{}bar\{j\}\_\{!\}\textbackslash{}bar\{F\}\{\textbackslash{}quoi\}, \$\$; Note: The opening display is printed on p. 235, not here. Page 236 really begins with 'd'où 4.2.' (source line 1083). The comment line does not flag this display as possibly belonging to the neighbouring page.\par
\textbf{GPT-6.1 Sol.} nothing reported here\par
a model reports that the automatic cut gives lines of this page to its neighbour, or the neighbour's lines to this page: check the lines it names against the scan of the page that prints them (that page's checker did not have their TeX)\par
\par\noindent\hypertarget{DP0522}{}\textbf{DP0522 -- XVI, p. 236 (cut)}\par
Début de l'extrait\par
\textbf{Claude Opus 5.5.} nothing reported here\par
\textbf{GPT-6.1 Sol.} Place: Début de l'extrait; Reported print: La page commence par « d'où 4.2. ».; Source TeX: \$\$ (R\textasciicircum{}\{q\}g\_\{\textbackslash{}cl \textbackslash{}ast\})j\_\{!\}F\textbackslash{}simeq (R\textasciicircum{}\{q\}\textbackslash{}bar\{g\}\_\{\textbackslash{}cl \textbackslash{}ast\})\textbackslash{}bar\{j\}\_\{!\}\textbackslash{}bar\{F\}\{\textbackslash{}quoi\}, \$\$; Note: L'affichage placé avant « d'où 4.2. » appartient à la page 235. Ce débordement initial n'est pas annoncé dans le commentaire de découpage.\par
a model reports that the automatic cut gives lines of this page to its neighbour, or the neighbour's lines to this page: check the lines it names against the scan of the page that prints them (that page's checker did not have their TeX)\par
\par\noindent\hypertarget{DP0523}{}\textbf{DP0523 -- XVI, p. 237 (conflict)}\par
Leray spectral sequence diagram, bottom left | Diagramme des suites spectrales de Leray, terme inférieur gauche\par
\textbf{Claude Opus 5.5.} Place: Leray spectral sequence diagram, bottom left; Reported print: (R\textasciicircum{}p b\_\{cl*\})(R\textasciicircum{}q a\_\{cl*\}) \ensuremath{\varepsilon}\textasciicircum{}*F; Source TeX: (R\textasciicircum{}\{p\}b\_\{\textbackslash{}text\{cl\}\textbackslash{}ast\}(R\textasciicircum{}\{q\}\{a\_\{\textbackslash{}text\{cl\}\textbackslash{}ast\}\})) \textbackslash{}varepsilon\textasciicircum{}\{\textbackslash{}ast\}F; Note: The print has two separate parenthesised factors. The TeX nests them instead.; Proposed edit: (R\textasciicircum{}\{p\}b\_\{\textbackslash{}text\{cl\}\textbackslash{}ast\}(R\textasciicircum{}\{q\}\{a\_\{\textbackslash{}text\{cl\}\textbackslash{}ast\}\})) =\allowbreak > (R\textasciicircum{}\{p\}b\_\{\textbackslash{}text\{cl\}\textbackslash{}ast\})(R\textasciicircum{}\{q\}a\_\{\textbackslash{}text\{cl\}\textbackslash{}ast\})\par
\textbf{GPT-6.1 Sol.} Place: Diagramme des suites spectrales de Leray, terme inférieur gauche; Reported print: (R\textasciicircum{}\{p\}b\_\{\textbackslash{}mathrm\{cl\}\textbackslash{}ast\})(R\textasciicircum{}\{q\}a\_\{\textbackslash{}mathrm\{cl\}\textbackslash{}ast\})\textbackslash{}varepsilon\textasciicircum{}\{\textbackslash{}ast\}F; Source TeX: (R\textasciicircum{}\{p\}b\_\{\textbackslash{}text\{cl\}\textbackslash{}ast\}(R\textasciicircum{}\{q\}\{a\_\{\textbackslash{}text\{cl\}\textbackslash{}ast\}\}))\textbackslash{}varepsilon\textasciicircum{}\{\textbackslash{}ast\}F; Note: Le print parenthèse séparément les deux facteurs ; le TeX emboîte le second dans le premier.; Proposed edit: (R\textasciicircum{}\{p\}b\_\{\textbackslash{}text\{cl\}\textbackslash{}ast\}(R\textasciicircum{}\{q\}\{a\_\{\textbackslash{}text\{cl\}\textbackslash{}ast\}\})) =\allowbreak > (R\textasciicircum{}\{p\}b\_\{\textbackslash{}text\{cl\}\textbackslash{}ast\})(R\textasciicircum{}\{q\}\{a\_\{\textbackslash{}text\{cl\}\textbackslash{}ast\}\})\par
the two models differ at this place\par
\par\noindent\hypertarget{DP0524}{}\textbf{DP0524 -- XVI, p. 239 (one\_model\_fix)}\par
Fin de la démonstration du lemme 4.4, renvois à VIII\par
\textbf{Claude Opus 5.5.} nothing reported here\par
\textbf{GPT-6.1 Sol.} Place: Fin de la démonstration du lemme 4.4, renvois à VIII; Reported print: (VIII 5.5 et 5.8); Source TeX: (\textbackslash{}hyperrefext[VIII][VIII.prop5.5]\{5.5\} et \textbackslash{}hyperrefext[VIII][VIII.rem5.8]\{5.8\}); Note: Le print ne répète pas « VIII » devant le second numéro ; le second \textbackslash{}hyperrefext l'ajoute.; Proposed edit: \textbackslash{}hyperrefext[VIII][VIII.rem5.8]\{5.8\} =\allowbreak > \textbackslash{}hyperref[VIII.rem5.8]\{5.8\}\par
a fix from one model only\par
\par\noindent\hypertarget{DP0525}{}\textbf{DP0525 -- XVI, p. 247 (one\_model\_fix)}\par
Proof of 5.1, paired references to VIII 5.5 and 5.8\par
\textbf{Claude Opus 5.5.} nothing reported here\par
\textbf{GPT-6.1 Sol.} Place: Proof of 5.1, paired references to VIII 5.5 and 5.8; Reported print: (VIII 5.5 et 5.8); Source TeX: (VIII 5.5 et VIII 5.8); Note: The second external-reference macro repeats « VIII », which is not repeated in the print.; Proposed edit: \textbackslash{}hyperrefext[VIII][VIII.rem5.8]\{5.8\} =\allowbreak > \textbackslash{}hyperref[VIII.rem5.8]\{5.8\}\par
a fix from one model only\par
\par\noindent\hypertarget{DP0526}{}\textbf{DP0526 -- XVII, p. 252 (one\_model\_fix)}\par
Ligne d’auteur sous le titre\par
\textbf{Claude Opus 5.5.} nothing reported here\par
\textbf{GPT-6.1 Sol.} Place: Ligne d’auteur sous le titre; Reported print: par P. Deligne (*); Source TeX: P. Deligne; Note: Il manque « par » et l’appel de la note, dont le texte est déjà présent.; Proposed edit: \textbackslash{}begin\{center\} P. Deligne =\allowbreak > \textbackslash{}begin\{center\} par P. Deligne\textbackslash{}footnotemark[1]\par
a fix from one model only\par
\par\noindent\hypertarget{DP0527}{}\textbf{DP0527 -- XVII, p. 266 (one\_model\_fix)}\par
1.1.8, first line, and throughout the page\par
\textbf{Claude Opus 5.5.} Place: 1.1.8, first line, and throughout the page; Reported print: \ensuremath{\psi} : I \ensuremath{\longrightarrow} J (handwritten small psi, with descender); Source TeX: \$\textbackslash{}Psi:I\textbackslash{}rightarrow J\$; Note: The print's handwritten letter looks like a small \ensuremath{\psi}, not a capital \ensuremath{\Psi}. The same holds for every \textbackslash{}Psi on the page (\textbackslash{}Psi soit, \textbackslash{}Psi(i)=\allowbreak j, the exponent in 1.1.8.1, the 'avec' line).; Proposed edit: \$\textbackslash{}Psi:I\textbackslash{}rightarrow J\$ =\allowbreak > \$\textbackslash{}psi:I\textbackslash{}rightarrow J\$\par
\textbf{GPT-6.1 Sol.} nothing reported here\par
a fix from one model only\par
\par\noindent\hypertarget{DP0528}{}\textbf{DP0528 -- XVII, p. 270 (agreed\_fix\_held)}\par
display after 'Soit' (before 1.1.11.2) | Premier affichage introduisant le complexe de foncteurs multiadditifs\par
\textbf{Claude Opus 5.5.} Place: display after 'Soit' (before 1.1.11.2); Reported print: F =\allowbreak  (F\textasciicircum{}n,d\textasciicircum{}n)\_\{n\ensuremath{\in}Z\} (no dot on F); Source TeX: F\textasciicircum{}\{\textbackslash{}bigdot\}=\allowbreak (F\textasciicircum{}\{n\},d\textasciicircum{}\{n\})\_\{n\textbackslash{}in \textbackslash{}ZZ\}; Note: In both scans the F of this first display has no dot exponent, unlike the F˙ of 1.1.12. The dot may be an omission by the typist.; Proposed edit: Soit \textbackslash{}[ F\textasciicircum{}\{\textbackslash{}bigdot\} =\allowbreak > Soit \textbackslash{}[ F\par
\textbf{GPT-6.1 Sol.} Place: Premier affichage introduisant le complexe de foncteurs multiadditifs; Reported print: F=\allowbreak (F\textasciicircum{}n,d\textasciicircum{}n)\_\{n\textbackslash{}in\textbackslash{}ZZ\}; Source TeX: F\textasciicircum{}\{\textbackslash{}bigdot\}=\allowbreak (F\textasciicircum{}\{n\},d\textasciicircum{}\{n\})\_\{n\textbackslash{}in \textbackslash{}ZZ\}; Note: Ce premier affichage porte F sans exposant dans les deux scans, bien que les occurrences imprimées ultérieures emploient F avec un point.; Proposed edit: Soit \textbackslash{}[ F\textasciicircum{}\{\textbackslash{}bigdot\} =\allowbreak > Soit \textbackslash{}[ F\par
both models give the same fix, held: a model marks it low confidence\par
\par\noindent\hypertarget{DP0529}{}\textbf{DP0529 -- XVII, p. 271 (conflict)}\par
formula (1.1.12.2) | Adjunction formula following diagram (1.1.12.1)\par
\textbf{Claude Opus 5.5.} Place: formula (1.1.12.2); Reported print: (1.1.12.2) Hom(F\textasciicircum{}n(A),B) \ensuremath{\leftrightarrow} Hom(A,F\_n(B)); Source TeX: equation* without a number; Note: The print numbers this formula (1.1.12.2); the TeX drops the number.; Proposed edit: \textbackslash{}Hom(F\textasciicircum{}\{n\}(A),B)\textbackslash{}overset =\allowbreak > \textbackslash{}tag\{1.1.12.2\}\textbackslash{}Hom(F\textasciicircum{}\{n\}(A),B)\textbackslash{}overset\par
\textbf{GPT-6.1 Sol.} Place: Adjunction formula following diagram (1.1.12.1); Reported print: (1.1.12.2); Source TeX: Unnumbered equation; Note: The printed formula number is missing.; Proposed edit: \textbackslash{}begin\{equation*\} \textbackslash{}Hom(F\textasciicircum{}\{n\}(A),B) =\allowbreak > \textbackslash{}begin\{equation*\} \textbackslash{}tag\{1.1.12.2\} \textbackslash{}Hom(F\textasciicircum{}\{n\}(A),B)\par
the two models differ at this place\par
\par\noindent\hypertarget{DP0530}{}\textbf{DP0530 -- XVII, p. 271 (conflict)}\par
formula (1.1.12.3) | Isomorphism obtained from conventions 1.1.3\par
\textbf{Claude Opus 5.5.} Place: formula (1.1.12.3); Reported print: (1.1.12.3) Hom\textasciicircum{}\ensuremath{\bullet}(F\textasciicircum{}\ensuremath{\bullet}(K),L) \ensuremath{\leftrightarrow} Hom(K,F\_\ensuremath{\bullet}(L)); Source TeX: equation* without a number; Note: The print numbers this formula (1.1.12.3); the TeX drops the number.; Proposed edit: \textbackslash{}Hom\textasciicircum{}\{\textbackslash{}bigdot\}(F\textasciicircum{}\{\textbackslash{}bigdot\}(K),L) =\allowbreak > \textbackslash{}tag\{1.1.12.3\}\textbackslash{}Hom\textasciicircum{}\{\textbackslash{}bigdot\}(F\textasciicircum{}\{\textbackslash{}bigdot\}(K),L)\par
\textbf{GPT-6.1 Sol.} Place: Isomorphism obtained from conventions 1.1.3; Reported print: (1.1.12.3); Source TeX: Unnumbered equation; Note: The printed formula number is missing.; Proposed edit: \textbackslash{}begin\{equation*\} \textbackslash{}Hom\textasciicircum{}\{\textbackslash{}bigdot\}(F\textasciicircum{}\{\textbackslash{}bigdot\}(K),L) =\allowbreak > \textbackslash{}begin\{equation*\} \textbackslash{}tag\{1.1.12.3\} \textbackslash{}Hom\textasciicircum{}\{\textbackslash{}bigdot\}(F\textasciicircum{}\{\textbackslash{}bigdot\}(K),L)\par
the two models differ at this place\par
\par\noindent\hypertarget{DP0531}{}\textbf{DP0531 -- XVII, p. 271 (one\_model\_fix)}\par
Introduction of the complex of functors\par
\textbf{Claude Opus 5.5.} nothing reported here\par
\textbf{GPT-6.1 Sol.} Place: Introduction of the complex of functors; Reported print: soit F\_\{\textbackslash{}bigdot\} le complexe; Source TeX: soit \$F\$. le complexe; Note: The point is the lower index denoting the complex, not sentence punctuation after F.; Proposed edit: \$F\$. le complexe =\allowbreak > \$F\_\{\textbackslash{}bigdot\}\$ le complexe\par
a fix from one model only\par
\par\noindent\hypertarget{DP0532}{}\textbf{DP0532 -- XVII, p. 271 (one\_model\_fix)}\par
The two equations defining the transposition functor\par
\textbf{Claude Opus 5.5.} nothing reported here\par
\textbf{GPT-6.1 Sol.} Place: The two equations defining the transposition functor; Reported print: et; Source TeX: \textbackslash{}quav; Note: The print joins the equations with « et », whereas \textbackslash{}quav supplies « avec ».; Proposed edit: \textbackslash{}quav =\allowbreak > \textbackslash{}quad\textbackslash{}text\{et\}\textbackslash{}quad\par
a fix from one model only\par
\par\noindent\hypertarget{DP0533}{}\textbf{DP0533 -- XVII, p. 273 (one\_model\_fix)}\par
Definition 1.2.1(i), label above the upper horizontal arrow\par
\textbf{Claude Opus 5.5.} nothing reported here\par
\textbf{GPT-6.1 Sol.} Place: Definition 1.2.1(i), label above the upper horizontal arrow; Reported print: (A\_i)\textbackslash{}mapsto F(A\_i\textasciicircum{}+); Source TeX: (A\_i)\textbackslash{}sisi\{\textbackslash{}simarrow\}\{\textbackslash{}mapsto\} F(A\_i\textasciicircum{}+); Note: Both scans already show an assignment arrow. The first argument of \textbackslash{}sisi incorrectly records an isomorphism arrow.; Proposed edit: \textbackslash{}sisi\{\textbackslash{}simarrow\}\{\textbackslash{}mapsto\} =\allowbreak > \textbackslash{}mapsto\par
a fix from one model only\par
\par\noindent\hypertarget{DP0534}{}\textbf{DP0534 -- XVII, p. 274 (conflict)}\par
Prop. 1.2.2 (ii), second argument | Proposition 1.2.2(ii), second family at which RF is defined\par
\textbf{Claude Opus 5.5.} Place: Prop. 1.2.2 (ii), second argument; Reported print: et en (A''\_1,(A\_i)); Source TeX: et en \$(A''\_1, (A\_1))\$; Note: The print has the index i in the family (A\_i); the TeX has A\_1.; Proposed edit: (A''\_1, (A\_1)) =\allowbreak > (A''\_1, (A\_i))\par
\textbf{GPT-6.1 Sol.} Place: Proposition 1.2.2(ii), second family at which RF is defined; Reported print: (A''\_1,(A\_1)); Source TeX: (A''\_1, (A\_1)); Proposed correction: (A''\_1,(A\_i)); Note: The inner index is printed as 1, also in the witness, but must be i for the remaining arguments. The TeX keeps the misprint.\par
the two models differ at this place: one model fixes the TeX, the other calls it a misprint of the print\par
\par\noindent\hypertarget{DP0535}{}\textbf{DP0535 -- XVII, p. 280 (one\_model\_fix)}\par
proof of 1.2.6, "\$\textbackslash{}Iim(d\_i)\textbackslash{}in \textbackslash{}PP\$"\par
\textbf{Claude Opus 5.5.} Place: proof of 1.2.6, "\$\textbackslash{}Iim(d\_i)\textbackslash{}in \textbackslash{}PP\$"; Reported print: Im(d\_i) \ensuremath{\in} P (underlined P, the same letter as elsewhere in 1.2.6); Source TeX: \$\textbackslash{}Iim(d\_i)\textbackslash{}in \textbackslash{}PP\$; Note: Everywhere else in 1.2.6 and its proof the TeX renders the print's underlined P as \ensuremath{𝒫}. Here it uses \textbackslash{}PP, which may print another alphabet.; Proposed edit: \$\textbackslash{}Iim(d\_i)\textbackslash{}in \textbackslash{}PP\$ =\allowbreak > \$\textbackslash{}Iim(d\_i)\textbackslash{}in \ensuremath{𝒫}\$\par
\textbf{GPT-6.1 Sol.} nothing reported here\par
a fix from one model only\par
\par\noindent\hypertarget{DP0536}{}\textbf{DP0536 -- XVII, p. 280 (cut)}\par
Début de page, suite exacte de la proposition 1.2.6 (b) ; coupure 279/280\par
\textbf{Claude Opus 5.5.} nothing reported here\par
\textbf{GPT-6.1 Sol.} Place: Début de page, suite exacte de la proposition 1.2.6 (b) ; coupure 279/280; Reported print: \$0\textbackslash{}rightarrow FP\textbackslash{}rightarrow FQ\textbackslash{}rightarrow FR\textbackslash{}rightarrow 0\$ est exacte.; Source TeX: 0\$\textbackslash{}pageoriginale est exacte.; Note: La page 280 commence au premier 0 de la suite exacte, non au dernier. Le début de cette formule est dans la dernière ligne de la coupe p. 279 (ligne source 1179), sans que son commentaire annonce ce partage. La page 279 finit après « 0-suite ». Refaire les coupes, sans déplacer ni ajouter de texte dans la source.\par
a model reports that the automatic cut gives lines of this page to its neighbour, or the neighbour's lines to this page: check the lines it names against the scan of the page that prints them (that page's checker did not have their TeX)\par
\par\noindent\hypertarget{DP0537}{}\textbf{DP0537 -- XVII, p. 301 (one\_model\_fix)}\par
factorisation of the functor, line 'il se factorise en \ensuremath{\psi} : A(S\textasciicircum{}\{-1\}) \ensuremath{\to} A\textasciicircum{}1(S\textasciicircum{}\{1-1\})'\par
\textbf{Claude Opus 5.5.} Place: factorisation of the functor, line 'il se factorise en \ensuremath{\psi} : A(S\textasciicircum{}\{-1\}) \ensuremath{\to} A\textasciicircum{}1(S\textasciicircum{}\{1-1\})'; Reported print: \ensuremath{\psi} (small psi); Source TeX: \textbackslash{}Psi; Note: The print has a small psi in all three places; the TeX has a capital Psi. Edit together with the next two findings.; Proposed edit: \$\textbackslash{}Psi:\textbackslash{}Aa(S\textasciicircum{}\{-1\}) =\allowbreak > \$\textbackslash{}psi:\textbackslash{}Aa(S\textasciicircum{}\{-1\})\par
\textbf{GPT-6.1 Sol.} nothing reported here\par
a fix from one model only\par
\par\noindent\hypertarget{DP0538}{}\textbf{DP0538 -- XVII, p. 301 (one\_model\_fix)}\par
'Le composé Φ \ensuremath{\psi} est isomorphe à l'identité'\par
\textbf{Claude Opus 5.5.} Place: 'Le composé Φ \ensuremath{\psi} est isomorphe à l'identité'; Reported print: Φ\ensuremath{\psi}; Source TeX: \textbackslash{}Phi\textbackslash{}Psi; Note: Small psi in the print.; Proposed edit: \$\textbackslash{}Phi\textbackslash{}Psi\$ =\allowbreak > \$\textbackslash{}Phi\textbackslash{}psi\$\par
\textbf{GPT-6.1 Sol.} nothing reported here\par
a fix from one model only\par
\par\noindent\hypertarget{DP0539}{}\textbf{DP0539 -- XVII, p. 301 (one\_model\_fix)}\par
'On vérifie de même que \ensuremath{\psi} Φ est isomorphe'\par
\textbf{Claude Opus 5.5.} Place: 'On vérifie de même que \ensuremath{\psi} Φ est isomorphe'; Reported print: \ensuremath{\psi}Φ; Source TeX: \textbackslash{}Psi\textbackslash{}Phi; Note: Small psi in the print.; Proposed edit: \$\textbackslash{}Psi\textbackslash{}Phi\$ =\allowbreak > \$\textbackslash{}psi\textbackslash{}Phi\$\par
\textbf{GPT-6.1 Sol.} nothing reported here\par
a fix from one model only\par
\par\noindent\hypertarget{DP0540}{}\textbf{DP0540 -- XVII, p. 311 (one\_model\_fix)}\par
condition (c), first word\par
\textbf{Claude Opus 5.5.} Place: condition (c), first word; Reported print: (c) Pour tout diagramme commutatif du type suivant; Source TeX: \textbackslash{}item[(c)] pour tout diagramme commutatif du type suivant; Note: The print capitalizes 'Pour', as in (c') on p. 312; the TeX has 'pour'.; Proposed edit: \textbackslash{}item[(c)] pour tout =\allowbreak > \textbackslash{}item[(c)] Pour tout\par
\textbf{GPT-6.1 Sol.} nothing reported here\par
a fix from one model only\par
\par\noindent\hypertarget{DP0541}{}\textbf{DP0541 -- XVII, p. 316 (one\_model\_fix)}\par
Hexagone, arête verticale extérieure droite\par
\textbf{Claude Opus 5.5.} nothing reported here\par
\textbf{GPT-6.1 Sol.} Place: Hexagone, arête verticale extérieure droite; Reported print: Trait vertical sans pointe entre \$\textbackslash{}bar\{\textbackslash{}bar\{g\}\}\_2pj'\_1\$ et \$p\textbackslash{}bar\{\textbackslash{}bar\{g\}\}\_1j'\_1\$; Source TeX: \textbackslash{}ar@\{-\}[ld]\textbackslash{}ar[dd]; Note: Les deux scans montrent un trait non orienté ; le TeX lui ajoute une pointe descendante.; Proposed edit: \textbackslash{}ar@\{-\}[ld]\textbackslash{}ar[dd] =\allowbreak > \textbackslash{}ar@\{-\}[ld]\textbackslash{}ar@\{-\}[dd]\par
a fix from one model only\par
\par\noindent\hypertarget{DP0542}{}\textbf{DP0542 -- XVII, p. 345 (cut)}\par
end of page, Théorème 5.1.8 (ii)\par
\textbf{Claude Opus 5.5.} Place: end of page, Théorème 5.1.8 (ii); Reported print: on retrouve la théorie de variance 5.1.4 (resp. 5.1.3 ).; Source TeX: on retrouve la théorie de variance; Note: The page prints '5.1.4 (resp. 5.1.3).' as the end of its last line, but this page's TeX stops at 'variance'. The rest is presumably in source line 3929, in the next page's chunk. The page really ends after '5.1.4 (resp. 5.1.3).'\par
\textbf{GPT-6.1 Sol.} nothing reported here\par
a model reports that the automatic cut gives lines of this page to its neighbour, or the neighbour's lines to this page: check the lines it names against the scan of the page that prints them (that page's checker did not have their TeX)\par
\par\noindent\hypertarget{DP0543}{}\textbf{DP0543 -- XVII, p. 345 (cut)}\par
End of Theorem 5.1.8(a)(ii), page cut\par
\textbf{Claude Opus 5.5.} nothing reported here\par
\textbf{GPT-6.1 Sol.} Place: End of Theorem 5.1.8(a)(ii), page cut; Reported print: on retrouve la théorie de variance 5.1.4 (resp. 5.1.3).; Source TeX: immersions ouvertes) on retrouve la théorie de variance; Note: The supplied cut stops before the final references, although they finish this printed page. Extend the extraction through “5.1.4 (resp. 5.1.3).”; do not insert duplicate text into the source.\par
a model reports that the automatic cut gives lines of this page to its neighbour, or the neighbour's lines to this page: check the lines it names against the scan of the page that prints them (that page's checker did not have their TeX)\par
\par\noindent\hypertarget{DP0544}{}\textbf{DP0544 -- XVII, p. 346 (cut)}\par
first line of the TeX, before item (iii)\par
\textbf{Claude Opus 5.5.} Place: first line of the TeX, before item (iii); Reported print: the page opens with '(iii) Pour tout diagramme (5.1.5.1)…'; the words '5.1.4 (resp. 5.1.3).' are not on p. 346; Source TeX: \textbackslash{}hyperref[XVII.subsubsec5.1.4]\{5.1.4\} (\textbackslash{}resp\textasciitilde{}\textbackslash{}hyperref[XVII.subsubsec5.1.3]\{5.1.3\}).; Note: Source line 3929 is the end of item (ii), which is printed on p. 345. Page 346 really begins with item (iii), 'Pour tout diagramme', where its \textbackslash{}pageoriginale stands. The comment line announces only an overlap at the end of the page, not at its start.\par
\textbf{GPT-6.1 Sol.} nothing reported here\par
a model reports that the automatic cut gives lines of this page to its neighbour, or the neighbour's lines to this page: check the lines it names against the scan of the page that prints them (that page's checker did not have their TeX)\par
\par\noindent\hypertarget{DP0545}{}\textbf{DP0545 -- XVII, p. 346 (cut)}\par
Début du fragment TeX, avant l’alinéa (iii)\par
\textbf{Claude Opus 5.5.} nothing reported here\par
\textbf{GPT-6.1 Sol.} Place: Début du fragment TeX, avant l’alinéa (iii); Reported print: La page commence par « (iii) Pour tout diagramme (5.1.5.1)… ».; Source TeX: \textbackslash{}hyperref[XVII.subsubsec5.1.4]\{5.1.4\} (\textbackslash{}resp\textasciitilde{}\textbackslash{}hyperref[XVII.subsubsec5.1.3]\{5.1.3\}).; Note: Cette ligne initiale appartient à la page 345, contrairement à la délimitation annoncée dans le commentaire. La page 346 commence à « Pour tout diagramme ».\par
a model reports that the automatic cut gives lines of this page to its neighbour, or the neighbour's lines to this page: check the lines it names against the scan of the page that prints them (that page's checker did not have their TeX)\par
\par\noindent\hypertarget{DP0546}{}\textbf{DP0546 -- XVII, p. 348 (one\_model\_fix)}\par
diagram (5.1.12.1), punctuation after the diagram\par
\textbf{Claude Opus 5.5.} Place: diagram (5.1.12.1), punctuation after the diagram; Reported print: comma after the diagram (the sentence continues 'où \textbackslash{}bar f\_1 =\allowbreak  …'); Source TeX: Y   \& \& .; Note: The print ends the diagram with a comma, not a full stop.; Proposed edit: Y   \& \& . =\allowbreak > Y   \& \& ,\par
\textbf{GPT-6.1 Sol.} nothing reported here\par
a fix from one model only\par
\par\noindent\hypertarget{DP0547}{}\textbf{DP0547 -- XVII, p. 355 (cut)}\par
end of page: cartesian diagram and '\$L'=\allowbreak ...\$ et'\par
\textbf{Claude Opus 5.5.} Place: end of page: cartesian diagram and '\$L'=\allowbreak ...\$ et'; Reported print: page ends with 'Soit le diagramme cartésien'; rest of page blank; Source TeX: \textbackslash{}[ \textbackslash{}xymatrix\{   \textbackslash{}bar\{X\}' \textbackslash{}ar[ddd]\textasciicircum{}\{p'\_1\} ... \$L'=\allowbreak \{p'\_2\}\textasciicircum{}\{\textbackslash{}ast\}L\$...\textbackslash{}pageoriginale et; Note: Printed p. 355 ends after 'Soit le diagramme cartésien'. The following cartesian diagram and '\$L'=\allowbreak ...\$' belong to p. 356, but they sit in this page's cut ahead of p. 356's \textbackslash{}pageoriginale, and the header comment does not flag them.\par
\textbf{GPT-6.1 Sol.} nothing reported here\par
a model reports that the automatic cut gives lines of this page to its neighbour, or the neighbour's lines to this page: check the lines it names against the scan of the page that prints them (that page's checker did not have their TeX)\par
\par\noindent\hypertarget{DP0548}{}\textbf{DP0548 -- XVII, p. 355 (cut)}\par
Fin de page, après « Soit le diagramme cartésien »\par
\textbf{Claude Opus 5.5.} nothing reported here\par
\textbf{GPT-6.1 Sol.} Place: Fin de page, après « Soit le diagramme cartésien »; Reported print: La page se termine après « Soit le diagramme cartésien », sans le diagramme annoncé.; Source TeX: La coupe contient encore le grand diagramme suivant, lignes source 4325–4336, puis la ligne 4337.; Note: Le grand diagramme appartient à la p. 356. Le chevauchement annoncé pour la seule dernière ligne ne couvre pas ce diagramme. La p. 355 finit après la phrase introductive ; corriger la coupe, sans déplacer ni modifier le texte.\par
a model reports that the automatic cut gives lines of this page to its neighbour, or the neighbour's lines to this page: check the lines it names against the scan of the page that prints them (that page's checker did not have their TeX)\par
\par\noindent\hypertarget{DP0549}{}\textbf{DP0549 -- XVII, p. 357 (conflict)}\par
sentence before Théorème 5.2.6 | Sentence introducing Theorem 5.2.6\par
\textbf{Claude Opus 5.5.} Place: sentence before Théorème 5.2.6; Reported print: On résume les énoncés qui précédent en le; Source TeX: On résume les énoncés qui précédent en le; Proposed correction: On résume les énoncés qui précèdent en le; Note: The verb should be 'précèdent' (grave accent). The print has 'précédent' and the TeX keeps it.\par
\textbf{GPT-6.1 Sol.} Place: Sentence introducing Theorem 5.2.6; Reported print: qui précèdent en le; Source TeX: qui précédent en le; Note: The printed verb has a grave accent on its second accented e.; Proposed edit: qui précédent en le =\allowbreak > qui précèdent en le\par
the two models differ at this place: one model fixes the TeX, the other calls it a misprint of the print\par
\par\noindent\hypertarget{DP0550}{}\textbf{DP0550 -- XVII, p. 370 (one\_model\_fix)}\par
Lemme 5.4.3.5, lower-left vertex\par
\textbf{Claude Opus 5.5.} nothing reported here\par
\textbf{GPT-6.1 Sol.} Place: Lemme 5.4.3.5, lower-left vertex; Reported print: Rf\_\{!\}(K\textbackslash{}overset\{L\}\{\textbackslash{}otimes\}\_\{\textbackslash{}Aa\}f\textasciicircum{}\{\textbackslash{}ast\}Rg\_\{!\}L); Source TeX: \textbackslash{}sisi\{\textbackslash{}boxtimes\}\{\textbackslash{}otimes\}; Note: The print already has \textbackslash{}otimes. The first argument of \textbackslash{}sisi incorrectly records it as \textbackslash{}boxtimes.; Proposed edit: \textbackslash{}sisi\{\textbackslash{}boxtimes\}\{\textbackslash{}otimes\} =\allowbreak > \textbackslash{}otimes\par
a fix from one model only\par
\par\noindent\hypertarget{DP0551}{}\textbf{DP0551 -- XVII, p. 370 (fix\_unplaced)}\par
Lemme 5.4.3.5 diagram, bottom left\par
\textbf{Claude Opus 5.5.} Place: Lemme 5.4.3.5 diagram, bottom left; Reported print: Rf\_!(K \ensuremath{\otimes}\textasciicircum{}L\_A f*Rg\_!L); Source TeX: \textbackslash{}sisi\{\textbackslash{}boxtimes\}\{\textbackslash{}otimes\}\_\{\textbackslash{}Aa\}; Note: The print has \ensuremath{\otimes} here. The \textbackslash{}sisi gives \ensuremath{\boxtimes} as the printed reading, which is wrong; the \ensuremath{\boxtimes} of the print is in the top right corner.; Proposed edit: \textbackslash{}sisi\{\textbackslash{}boxtimes\}\{\textbackslash{}otimes\}\_\{\textbackslash{}Aa\} =\allowbreak > \textbackslash{}otimes\_\{\textbackslash{}Aa\}\par
\textbf{GPT-6.1 Sol.} nothing reported here\par
a fix whose old text is not found once on the page\par
\par\noindent\hypertarget{DP0552}{}\textbf{DP0552 -- XVII, p. 388 (agreed\_fix\_held)}\par
Lemme 5.5.13.2, after 'Tor-dimension \ensuremath{\leq} 0' | Lemme 5.5.13.2, after « Tor-dimension \ensuremath{\leq} 0 »\par
\textbf{Claude Opus 5.5.} Place: Lemme 5.5.13.2, after 'Tor-dimension \ensuremath{\leq} 0'; Reported print: des complexes de Tor-dimension \ensuremath{\leq} 0 est équivalente; Source TeX: de \textbackslash{}tTor-dimension \$\textbackslash{}leq 0\$ \textbackslash{}emph\{\textbackslash{}hyperref[XVII.def4.1.9]\{4.1.9\}\}; Note: The print has no reference 4.1.9 here. It was added to the TeX without an editor's note.; Proposed edit:  \$\textbackslash{}leq 0\$ \textbackslash{}emph\{\textbackslash{}hyperref[XVII.def4.1.9]\{4.1.9\}\} =\allowbreak >  \$\textbackslash{}leq 0\$\par
\textbf{GPT-6.1 Sol.} Place: Lemme 5.5.13.2, after « Tor-dimension \ensuremath{\leq} 0 »; Reported print: Tor-dimension \$\textbackslash{}leq 0\$ est équivalente; Source TeX: \$\textbackslash{}leq 0\$ \textbackslash{}emph\{\textbackslash{}hyperref[XVII.def4.1.9]\{4.1.9\}\}; Note: The reference 4.1.9 is not printed; the sentence proceeds directly from the inequality to « est équivalente ».; Proposed edit:  \textbackslash{}emph\{\textbackslash{}hyperref[XVII.def4.1.9]\{4.1.9\}\} =\allowbreak > \par
both models give the same fix, held: a deletion in the first or last three lines of the page (the page cut may be off); it changes the editors' apparatus (\textbackslash{}hyperref, the target of a cross-reference): a \textbackslash{}sisi correction, note, label, cross-reference or citation is the editors' record and is not removed or retargeted on the models' word\par
\par\noindent\hypertarget{DP0553}{}\textbf{DP0553 -- XVII, p. 390 (one\_model\_fix)}\par
reference to Dold-Puppe\par
\textbf{Claude Opus 5.5.} Place: reference to Dold-Puppe; Reported print: DOLD PUPPE[1.2.9 pg.213]; Source TeX: \textbackslash{}cite[2.9 \textbackslash{}sisi\{pg.213\}\{p.\textasciitilde{}213\}]\{SGA4-17:DP\}; Note: The print cites 1.2.9; the TeX has 2.9.; Proposed edit: [2.9 \textbackslash{}sisi\{pg.213\} =\allowbreak > [1.2.9 \textbackslash{}sisi\{pg.213\}\par
\textbf{GPT-6.1 Sol.} nothing reported here\par
a fix from one model only\par
\par\noindent\hypertarget{DP0554}{}\textbf{DP0554 -- XVII, p. 391 (one\_model\_fix)}\par
item (i)\par
\textbf{Claude Opus 5.5.} Place: item (i); Reported print: (i) Condition (5.5.13.1).; Source TeX: \textbackslash{}item condition\textbackslash{}pageoriginale (\textbackslash{}hyperref[XVII.par5.5.13.1]\{5.5.13.1\}).; Note: The print capitalizes 'Condition' at the start of item (i); the TeX has lower case.; Proposed edit: \textbackslash{}item condition\textbackslash{}pageoriginale =\allowbreak > \textbackslash{}item Condition\textbackslash{}pageoriginale\par
\textbf{GPT-6.1 Sol.} nothing reported here\par
a fix from one model only\par
\par\noindent\hypertarget{DP0555}{}\textbf{DP0555 -- XVII, p. 394 (conflict)}\par
composite after (5.5.19.2), second arrow target | Composé définissant (5.5.19.2), arrivée de l'isomorphisme (5.5.19.1)\par
\textbf{Claude Opus 5.5.} Place: composite after (5.5.19.2), second arrow target; Reported print: R Sym\textasciicircum{}n\_S(ū)\_* Sym\textasciicircum{}n\_S(j)\_! L\ensuremath{\Gamma}\textasciicircum{}n\_ext; Source TeX: R\textbackslash{}Sym\textasciicircum{}n\_s(\textbackslash{}bar\{u\})\_*\textbackslash{}Sym\textasciicircum{}\{n\}\_\{S\}(j)\_\{!\}L\textbackslash{}Gamma\textasciicircum{}\{n\}\_\{\textbackslash{}ext\}; Note: The print's subscript is a capital S (the base scheme); the TeX has a lower-case s.; Proposed edit: R\textbackslash{}Sym\textasciicircum{}n\_s(\textbackslash{}bar\{u\})\_* =\allowbreak > R\textbackslash{}Sym\textasciicircum{}\{n\}\_\{S\}(\textbackslash{}bar\{u\})\_\{\textbackslash{}ast\}\par
\textbf{GPT-6.1 Sol.} Place: Composé définissant (5.5.19.2), arrivée de l'isomorphisme (5.5.19.1); Reported print: R\textbackslash{}Sym\textasciicircum{}\{n\}\_\{S\}(\textbackslash{}bar\{u\})\_\{\textbackslash{}ast\}; Source TeX: R\textbackslash{}Sym\textasciicircum{}n\_s(\textbackslash{}bar\{u\})\_*; Note: Le sous-indice imprimé est le S majuscule du schéma de base, non s.; Proposed edit: R\textbackslash{}Sym\textasciicircum{}n\_s(\textbackslash{}bar\{u\})\_* =\allowbreak > R\textbackslash{}Sym\textasciicircum{}n\_S(\textbackslash{}bar\{u\})\_*\par
the two models differ at this place\par
\par\noindent\hypertarget{DP0556}{}\textbf{DP0556 -- XVII, p. 394 (one\_model\_fix)}\par
Composé définissant (5.5.19.2), arrivée de la flèche (5.5.17.2)\par
\textbf{Claude Opus 5.5.} nothing reported here\par
\textbf{GPT-6.1 Sol.} Place: Composé définissant (5.5.19.2), arrivée de la flèche (5.5.17.2); Reported print: L\textbackslash{}Gamma\textasciicircum{}\{n\}\_\{\textbackslash{}ext\}j\_\{!\}; Source TeX: L\textbackslash{}Gamma\textasciicircum{}\{n\}\_\{\textbackslash{}ext j!\}; Note: Le j\_! est un facteur composé après L\ensuremath{\Gamma}\_ext, non une partie de l'indice ext. Sa position matériellement abaissée dans le tapuscrit ne doit pas être reproduite comme un indice supplémentaire.; Proposed edit: L\textbackslash{}Gamma\textasciicircum{}\{n\}\_\{\textbackslash{}ext j!\} =\allowbreak > L\textbackslash{}Gamma\textasciicircum{}\{n\}\_\{\textbackslash{}ext\}j\_\{!\}\par
a fix from one model only\par
\par\noindent\hypertarget{DP0557}{}\textbf{DP0557 -- XVII, p. 394 (one\_model\_fix)}\par
Étiquette de la dernière flèche du composé, lecture originale A de \textbackslash{}sisi\par
\textbf{Claude Opus 5.5.} nothing reported here\par
\textbf{GPT-6.1 Sol.} Place: Étiquette de la dernière flèche du composé, lecture originale A de \textbackslash{}sisi; Reported print: (5.5.13.5 (vii'')); Source TeX: \textbackslash{}text\{(5.5.13.5 (vii'))\}; Note: Le tirage et le témoin portent deux primes. Cet edit restitue uniquement la lecture originale A, sans modifier B. Le caractère erroné du renvoi imprimé est signalé séparément.; Proposed edit: \textbackslash{}text\{(5.5.13.5 (vii'))\} =\allowbreak > \textbackslash{}text\{(5.5.13.5 (vii''))\}\par
a fix from one model only\par
\par\noindent\hypertarget{DP0558}{}\textbf{DP0558 -- XVII, p. 396 (cut)}\par
end of page: second commutative diagram of (vi'') (source lines 6019-6027)\par
\textbf{Claude Opus 5.5.} Place: end of page: second commutative diagram of (vi'') (source lines 6019-6027); Reported print: Page 396 ends with « le diagramme de foncteurs suivant est commutatif : ». The diagram itself, Sym\textasciicircum{}n\_S(g)\textasciicircum{}*L\ensuremath{\Gamma}\textasciicircum{}n\_ext Ru\_! … R Sym\textasciicircum{}n\_S(u')\_!L\ensuremath{\Gamma}\textasciicircum{}n\_ext g'\textasciicircum{}*, is printed at the top of p. 397 (typescript XVII-146).; Source TeX: \textbackslash{}[ \textbackslash{}xymatrix\{ \textbackslash{}Sym\textasciicircum{}\{n\}\_S(g)\textasciicircum{}\{\textbackslash{}ast\}L \textbackslash{}Gamma\textasciicircum{}\{n\}\_\{\textbackslash{}ext\} Ru\_\{!\} ... \textbackslash{}]; Note: Page 396's cut holds the display at source lines 6019-6027, which belongs to p. 397. Its first line declares only lines 41-43 as possibly belonging to p. 397. Page 397's TeX already starts at 6019 and contains this display, so nothing is missing. Page 396 really ends after « suivant est commutatif: » (source line 6018).\par
\textbf{GPT-6.1 Sol.} nothing reported here\par
a model reports that the automatic cut gives lines of this page to its neighbour, or the neighbour's lines to this page: check the lines it names against the scan of the page that prints them (that page's checker did not have their TeX)\par
\par\noindent\hypertarget{DP0559}{}\textbf{DP0559 -- XVII, p. 396 (cut)}\par
Fin de (vi''), après « le diagramme de foncteurs suivant est commutatif: »\par
\textbf{Claude Opus 5.5.} nothing reported here\par
\textbf{GPT-6.1 Sol.} Place: Fin de (vi''), après « le diagramme de foncteurs suivant est commutatif: »; Reported print: La page finit après cette annonce. Le diagramme annoncé ouvre la page 397.; Source TeX: le diagramme de foncteurs suivant est commutatif: \textbackslash{}[ \textbackslash{}xymatrix\{; Note: La coupe de 396 comprend indûment le diagramme des lignes source 6019–6027. Ces lignes ne sont pas dans la plage 6028–6030 signalée comme incertaine par le commentaire. La page 396 finit à la ligne 6018; 397 commence au diagramme de la ligne 6019. Aucun déplacement de texte proposé.\par
a model reports that the automatic cut gives lines of this page to its neighbour, or the neighbour's lines to this page: check the lines it names against the scan of the page that prints them (that page's checker did not have their TeX)\par
\par\noindent\hypertarget{DP0560}{}\textbf{DP0560 -- XVII, p. 405 (one\_model\_fix)}\par
final diagram, after Spec(k)\par
\textbf{Claude Opus 5.5.} Place: final diagram, after Spec(k); Reported print: Spec(k)   ,; Source TeX: \& \textbackslash{}Spec(k), \textbackslash{}ar[l],; Note: The second comma after \textbackslash{}ar[l] typesets a second comma in the cell. The print has only one.; Proposed edit: \textbackslash{}Spec(k), \textbackslash{}ar[l], =\allowbreak > \textbackslash{}Spec(k), \textbackslash{}ar[l]\par
\textbf{GPT-6.1 Sol.} nothing reported here\par
a fix from one model only\par
\par\noindent\hypertarget{DP0561}{}\textbf{DP0561 -- XVII, p. 415 (one\_model\_fix)}\par
6.1.1 (iii)\par
\textbf{Claude Opus 5.5.} Place: 6.1.1 (iii); Reported print: \ensuremath{\varphi}\_X(U), \ensuremath{\varphi}\_X, \ensuremath{\varphi} (typed small phi); Source TeX: \textbackslash{}Phi\_\{X\}(U), \textbackslash{}Phi\_\{X\}, \textbackslash{}Phi; Note: The typed glyph looks like a small phi with a descender, not a capital Phi. It occurs three times in (iii); it may be only the typewriter's form of Phi. Only the first occurrence is fixed here.; Proposed edit: désignons par \$\textbackslash{}Phi\_\{X\}(U)\$ =\allowbreak > désignons par \$\textbackslash{}phi\_\{X\}(U)\$\par
\textbf{GPT-6.1 Sol.} nothing reported here\par
a fix from one model only\par
\par\noindent\hypertarget{DP0562}{}\textbf{DP0562 -- XVII, p. 446 (agreed\_fix\_held)}\par
6.3.16, first line | 6.3.16, première phrase\par
\textbf{Claude Opus 5.5.} Place: 6.3.16, first line; Reported print: Soit maintenant f : T \ensuremath{\longrightarrow} S un morphisme quasi-fini plat de présentation finie; Source TeX: \textbackslash{}sisi\{séparé \}\{\}quasi-fini; Note: The \textbackslash{}sisi gives « séparé » as the print's reading, but the print has no « séparé » here. The edition's output (empty B) agrees with the print; only the recorded reading A is wrong.; Proposed edit: \textbackslash{}sisi\{séparé \}\{\}quasi-fini =\allowbreak > quasi-fini\par
\textbf{GPT-6.1 Sol.} Place: 6.3.16, première phrase; Reported print: un morphisme quasi-fini plat de présentation finie; Source TeX: un morphisme \textbackslash{}sisi\{séparé \}\{\}quasi-fini plat de présentation finie; Note: Le mot « séparé » ne figure pas dans l'imprimé ; le premier argument de \textbackslash{}sisi lui attribue une lecture inexistante.; Proposed edit: \textbackslash{}sisi\{séparé \}\{\} =\allowbreak > \par
both models give the same fix, held: it changes the editors' apparatus (\textbackslash{}sisi): a \textbackslash{}sisi correction, note, label, cross-reference or citation is the editors' record and is not removed or retargeted on the models' word\par
\par\noindent\hypertarget{DP0563}{}\textbf{DP0563 -- XVII, p. 448 (conflict)}\par
diagram (6.3.17.1), label of the top left arrow | Diagramme (6.3.17.1), première flèche horizontale\par
\textbf{Claude Opus 5.5.} Place: diagram (6.3.17.1), label of the top left arrow; Reported print: f\_! ch; Source TeX: \textbackslash{}ar[r]\textasciicircum{}\{f\_\{!\}c\}; Note: The print labels the arrow f\_!ch, not f\_!c.; Proposed edit: \textbackslash{}ar[r]\textasciicircum{}\{f\_\{!\}c\} =\allowbreak > \textbackslash{}ar[r]\textasciicircum{}\{f\_\{!\}ch\}\par
\textbf{GPT-6.1 Sol.} Place: Diagramme (6.3.17.1), première flèche horizontale; Reported print: f\_\{!\}ch; Source TeX: f\_\{!\}c; Note: Les deux scans portent les deux lettres « ch » dans cette étiquette, bien que le morphisme soit nommé c dans le texte précédent.; Proposed edit: \textasciicircum{}\{f\_\{!\}c\} =\allowbreak > \textasciicircum{}\{f\_\{!\}\textbackslash{}mathrm\{ch\}\}\par
the two models differ at this place\par
\par\noindent\hypertarget{DP0564}{}\textbf{DP0564 -- XVII, p. 450 (one\_model\_fix)}\par
Preuve (i), parenthèse après \ensuremath{\sigma}(D) dans le même terme\par
\textbf{Claude Opus 5.5.} nothing reported here\par
\textbf{GPT-6.1 Sol.} Place: Preuve (i), parenthèse après \ensuremath{\sigma}(D) dans le même terme; Reported print: \textbackslash{}sigma(D)\textasciicircum{}\{\textbackslash{}ast\}; Source TeX: \textbackslash{}sigma(D))\textasciicircum{}\textbackslash{}ast; Note: La seconde parenthèse fermante n'est pas imprimée. Elle appartient à la correction B et doit être conditionnelle pour ne pas apparaître dans la lecture A.; Proposed edit: \textbackslash{}sigma(D))\textasciicircum{}\textbackslash{}ast =\allowbreak > \textbackslash{}sigma(D)\textbackslash{}sisi\{\}\{)\}\textasciicircum{}\textbackslash{}ast\par
a fix from one model only\par
\par\noindent\hypertarget{DP0565}{}\textbf{DP0565 -- XVII, p. 465 (conflict)}\par
reference after «une amélioration de» | Renvoi précédant le lemme 7.0.6\par
\textbf{Claude Opus 5.5.} Place: reference after «une amélioration de»; Reported print: (XII 7.1); Source TeX: \textbackslash{}sisi\{\textbackslash{}hyperref[XVII.subsec7.1]\{7.1\}\}\{\textbackslash{}hyperrefext[XII][XII.lem7.1]\{7.1\}\}; Note: The \textbackslash{}sisi gives as the print's reading a reference to 7.1 of this exposé. The print already reads «XII 7.1», so the claimed original reading is wrong.; Proposed edit: \textbackslash{}sisi\{\textbackslash{}hyperref[XVII.subsec7.1]\{7.1\}\}\{\textbackslash{}hyperrefext[XII][XII.lem7.1]\{7.1\}\} =\allowbreak > \textbackslash{}hyperrefext[XII][XII.lem7.1]\{7.1\}\par
\textbf{GPT-6.1 Sol.} Place: Renvoi précédant le lemme 7.0.6; Reported print: XII 7.1; Source TeX: \textbackslash{}hyperref[XVII.subsec7.1]\{7.1\}; Note: Le premier argument de \textbackslash{}sisi omet XII et renvoie à XVII 7.1. L'imprimé porte bien XII 7.1 ; seul ce premier argument est à modifier.; Proposed edit: \textbackslash{}hyperref[XVII.subsec7.1]\{7.1\} =\allowbreak > \textbackslash{}hyperrefext[XII][XII.lem7.1]\{7.1\}\par
the two models differ at this place\par
\par\noindent\hypertarget{DP0566}{}\textbf{DP0566 -- XVII, p. 469 (conflict)}\par
7.1.4.2, On dispose d'un diagramme commutatif | 7.1.4.2, phrase introduisant (7.1.4.2.1)\par
\textbf{Claude Opus 5.5.} Place: 7.1.4.2, On dispose d'un diagramme commutatif; Reported print: On dispose d'un diagramme commutatif, où toutes les flèches se déduisent naturellement de; Source TeX: commutatif, où \$Z\_2''=\allowbreak \textbackslash{}fprod\{T''\_2\}\{Z'\_2\}\{T'\_2\}\$ et toutes les flèches; Note: Editorial addition, not marked by \textbackslash{}sisi, that the print does not have.; Proposed edit: commutatif, où \$Z\_2''=\allowbreak \textbackslash{}fprod\{T''\_2\}\{Z'\_2\}\{T'\_2\}\$ et toutes =\allowbreak > commutatif, où toutes\par
\textbf{GPT-6.1 Sol.} Place: 7.1.4.2, phrase introduisant (7.1.4.2.1); Reported print: où toutes les flèches se déduisent naturellement; Source TeX: où \$Z\_2''=\allowbreak \textbackslash{}fprod\{T''\_2\}\{Z'\_2\}\{T'\_2\}\$ et toutes les flèches se; Note: La définition de Z\_2'' et le mot « et » sont absents de l'imprimé. Cet ajout lié à la modification éditoriale doit être marqué comme tel; le remplacement conserve l'ajout dans la lecture corrigée seulement.; Proposed edit: où \$Z\_2''=\allowbreak \textbackslash{}fprod\{T''\_2\}\{Z'\_2\}\{T'\_2\}\$ et toutes =\allowbreak > où \textbackslash{}sisi\{\}\{\$Z\_2''=\allowbreak \textbackslash{}fprod\{T''\_2\}\{Z'\_2\}\{T'\_2\}\$ et \}toutes\par
the two models differ at this place\par
\par\noindent\hypertarget{DP0567}{}\textbf{DP0567 -- XVII, p. 475 (agreed\_fix\_held)}\par
'On définit (XVII 6.2.9)' | Démonstration de 7.2.1, référence après «On définit»\par
\textbf{Claude Opus 5.5.} Place: 'On définit (XVII 6.2.9)'; Reported print: (XVII 6.2.9); Source TeX: (\textbackslash{}hyperref[XVII.prop6.2.9]\{6.2.9\}); Note: The print gives the exposé numeral XVII; the TeX drops it.; Proposed edit: On définit (\textbackslash{}hyperref[XVII.prop6.2.9]\{6.2.9\}) =\allowbreak > On définit (XVII \textbackslash{}hyperref[XVII.prop6.2.9]\{6.2.9\})\par
\textbf{GPT-6.1 Sol.} Place: Démonstration de 7.2.1, référence après «On définit»; Reported print: (XVII 6.2.9); Source TeX: (\textbackslash{}hyperref[XVII.prop6.2.9]\{6.2.9\}); Note: Le numéro d'exposé XVII, expressément imprimé, manque dans le texte du renvoi.; Proposed edit: (\textbackslash{}hyperref[XVII.prop6.2.9]\{6.2.9\}) =\allowbreak > (XVII \textbackslash{}hyperref[XVII.prop6.2.9]\{6.2.9\})\par
both models give the same fix, held: a model marks it low confidence\par
\par\noindent\hypertarget{DP0568}{}\textbf{DP0568 -- XVII, p. 475 (agreed\_fix\_held)}\par
'point par point (XVII (4.1.9))' | Démonstration de 7.2.2, référence après «point par point»\par
\textbf{Claude Opus 5.5.} Place: 'point par point (XVII (4.1.9))'; Reported print: (XVII (4.1.9)); Source TeX: (\textbackslash{}hyperref[XVII.def4.1.9]\{4.1.9\}); Note: The print gives XVII and inner parentheses; the TeX drops both.; Proposed edit: (\textbackslash{}hyperref[XVII.def4.1.9]\{4.1.9\}) =\allowbreak > (XVII (\textbackslash{}hyperref[XVII.def4.1.9]\{4.1.9\}))\par
\textbf{GPT-6.1 Sol.} Place: Démonstration de 7.2.2, référence après «point par point»; Reported print: (XVII (4.1.9)); Source TeX: (\textbackslash{}hyperref[XVII.def4.1.9]\{4.1.9\}); Note: Le numéro d'exposé XVII manque dans le renvoi transcrit.; Proposed edit: (\textbackslash{}hyperref[XVII.def4.1.9]\{4.1.9\}) =\allowbreak > (XVII (\textbackslash{}hyperref[XVII.def4.1.9]\{4.1.9\}))\par
both models give the same fix, held: a model marks it low confidence\par
\par\noindent\hypertarget{DP0569}{}\textbf{DP0569 -- XVII, p. 475 (cut)}\par
end of page: exact sequence display\par
\textbf{Claude Opus 5.5.} Place: end of page: exact sequence display; Reported print: page ends with '... est prise dans C\textasciicircum{}+(C\textasciicircum{}+(Mod(X,A\_X))) :'; Source TeX: \textbackslash{}[ \textbackslash{}sisi\{0\textbackslash{}longrightarrow (\textbackslash{}sigma\_\{\textbackslash{}leq ...; Note: The exact-sequence display at the end of p. 475's TeX is not printed on p. 475; it belongs to p. 476.\par
\textbf{GPT-6.1 Sol.} nothing reported here\par
a model reports that the automatic cut gives lines of this page to its neighbour, or the neighbour's lines to this page: check the lines it names against the scan of the page that prints them (that page's checker did not have their TeX)\par
\par\noindent\hypertarget{DP0570}{}\textbf{DP0570 -- XVII, p. 475 (cut)}\par
Après le deux-points final\par
\textbf{Claude Opus 5.5.} nothing reported here\par
\textbf{GPT-6.1 Sol.} Place: Après le deux-points final; Reported print: La page finit au deux-points après C\textasciicircum{}\{+\}(C\textasciicircum{}\{+\}(\textbackslash{}Mod(X,\textbackslash{}mathcal\{A\}\_\{X\}))).; Source TeX: Le display contenant la suite exacte 0\textbackslash{}longrightarrow\textbackslash{}cdots\textbackslash{}longrightarrow0 figure encore dans cette coupe.; Note: La suite exacte n'est pas imprimée sur la page 475 : elle appartient à la page suivante. Le commentaire de coupe n'annonce pas ce débordement. Ne pas contrôler ici les deux arguments de ce \textbackslash{}sisi ; recaler la fin de la coupe au deux-points.\par
a model reports that the automatic cut gives lines of this page to its neighbour, or the neighbour's lines to this page: check the lines it names against the scan of the page that prints them (that page's checker did not have their TeX)\par
\par\noindent\hypertarget{DP0571}{}\textbf{DP0571 -- XVII, p. 476 (agreed\_fix\_held)}\par
distinguished triangle, arrow from K\textasciicircum{}\{n+1,\ensuremath{\cdot}\}\_\ensuremath{𝔘} to the lower left | Triangle distingué, origine de la flèche oblique gauche\par
\textbf{Claude Opus 5.5.} Place: distinguished triangle, arrow from K\textasciicircum{}\{n+1,\ensuremath{\cdot}\}\_\ensuremath{𝔘} to the lower left; Reported print: arrow with a small circle at its tail (degree +1 arrow); Source TeX: \textbackslash{}sisi\{\textbackslash{}ar[ld]\}; Note: The print marks the degree-one arrow of the triangle with a circle. The printed reading A does not show it.; Proposed edit: \textbackslash{}sisi\{\textbackslash{}ar[ld]\} =\allowbreak > \textbackslash{}sisi\{\textbackslash{}ar@\{o->\}[ld]\}\par
\textbf{GPT-6.1 Sol.} Place: Triangle distingué, origine de la flèche oblique gauche; Reported print: La flèche issue de K\textasciicircum{}\{n+1,\ensuremath{\cdot}\} vers le sommet inférieur gauche porte un petit cercle à son origine.; Source TeX: \textbackslash{}sisi\{\textbackslash{}ar[ld]\}\{\textbackslash{}ar[rd]\}; Note: Le cercle fait partie du signe de la flèche dans les deux scans. Il manque dans la lecture originale enregistrée par le premier argument de \textbackslash{}sisi.; Proposed edit: \textbackslash{}sisi\{\textbackslash{}ar[ld]\}\{\textbackslash{}ar[rd]\} =\allowbreak > \textbackslash{}sisi\{\textbackslash{}ar@\{o->\}[ld]\}\{\textbackslash{}ar[rd]\}\par
both models give the same fix, held: a model marks it low confidence\par
\par\noindent\hypertarget{DP0572}{}\textbf{DP0572 -- XVII, p. 476 (one\_model\_fix)}\par
Paragraphe « Utilisant le fait », fin de la première phrase\par
\textbf{Claude Opus 5.5.} nothing reported here\par
\textbf{GPT-6.1 Sol.} Place: Paragraphe « Utilisant le fait », fin de la première phrase; Reported print: pour tout \$n\$; Source TeX: pour tout \$\textbackslash{}sim\$; Note: Le signe manuscrit est la lettre n, indice du complexe considéré, et non une relation tilde.; Proposed edit: pour tout \$\textbackslash{}sim\$ =\allowbreak > pour tout \$n\$\par
a fix from one model only\par
\par\noindent\hypertarget{DP0573}{}\textbf{DP0573 -- XVII, p. 476 (one\_model\_fix)}\par
Renvoi à 5.2.10 après « pour tout n »\par
\textbf{Claude Opus 5.5.} nothing reported here\par
\textbf{GPT-6.1 Sol.} Place: Renvoi à 5.2.10 après « pour tout n »; Reported print: (XVII (5.2.10)); Source TeX: (\textbackslash{}hyperref[XVII.thm5.2.10]\{5.2.10\}); Note: Suppression récurrente du préfixe d'exposé dans les renvois internes, signalée une seule fois comme convention ; trois exemples sont indiqués au coordinator.; Proposed edit: \textbackslash{}hyperref[XVII.thm5.2.10]\{5.2.10\} =\allowbreak > XVII \textbackslash{}hyperref[XVII.thm5.2.10]\{5.2.10\}\par
a fix from one model only\par
\par\noindent\hypertarget{DP0574}{}\textbf{DP0574 -- XVII, p. 476 (cut)}\par
top of page, before « d'où un triangle distingué »\par
\textbf{Claude Opus 5.5.} Place: top of page, before « d'où un triangle distingué »; Reported print: 0 \ensuremath{\to} (\ensuremath{\sigma}\_\{\ensuremath{\leq}n\}(K\ensuremath{\cdot}\ensuremath{\cdot}\_\ensuremath{𝔘}))\_s \ensuremath{\to} (\ensuremath{\sigma}\_\{\ensuremath{\leq}n+1\}(K\ensuremath{\cdot}\ensuremath{\cdot}\_\ensuremath{𝔘}))\_s \ensuremath{\to} (-1)\textasciicircum{}\{n+1\}K\textasciicircum{}\{n+1,\ensuremath{\cdot}\}\_\ensuremath{𝔘} \ensuremath{\to} 0; Source TeX: (absent: the TeX of p. 476 begins with « d'où\textbackslash{}pageoriginale »); Note: The page opens with the display of the short exact sequence. It is not in this page's cut (source lines 9291ff.) and is presumably at the end of p. 475's TeX. Page 476 really begins with that display, just before « d'où ».\par
\textbf{GPT-6.1 Sol.} nothing reported here\par
a model reports that the automatic cut gives lines of this page to its neighbour, or the neighbour's lines to this page: check the lines it names against the scan of the page that prints them (that page's checker did not have their TeX)\par
\par\noindent\hypertarget{DP0575}{}\textbf{DP0575 -- XVII, p. 476 (cut)}\par
Début de page, avant « d'où un triangle distingué »\par
\textbf{Claude Opus 5.5.} nothing reported here\par
\textbf{GPT-6.1 Sol.} Place: Début de page, avant « d'où un triangle distingué »; Reported print: La page commence par la suite exacte 0 \ensuremath{\to} (\ensuremath{\sigma}\ensuremath{\leq}n(K))s \ensuremath{\to} (\ensuremath{\sigma}\ensuremath{\leq}n+1(K))s \ensuremath{\to} (-1)\textasciicircum{}\{n+1\}K\textasciicircum{}\{n+1,\ensuremath{\cdot}\} \ensuremath{\to} 0.; Source TeX: L'extrait commence à « d'où\textbackslash{}pageoriginale un triangle distingué ».; Note: La suite exacte initiale manque dans la coupe fournie. Vérifier la fin de la coupe précédente et faire commencer celle de p. 476 avant ce display, sans ajouter une seconde copie du texte.\par
a model reports that the automatic cut gives lines of this page to its neighbour, or the neighbour's lines to this page: check the lines it names against the scan of the page that prints them (that page's checker did not have their TeX)\par
\par\noindent\hypertarget{DP0576}{}\textbf{DP0576 -- XVII, p. 477 (one\_model\_fix)}\par
line after the cube diagram\par
\textbf{Claude Opus 5.5.} Place: line after the cube diagram; Reported print: la flèche ch\_\{K\ensuremath{\cdot}\}; Source TeX: la flèche \$\textbackslash{}mathrm\{ch\}\_\{K\}\$.; Note: The print has the dot of K\ensuremath{\cdot}, as in the rest of the page.; Proposed edit: la flèche \$\textbackslash{}mathrm\{ch\}\_\{K\}\$. =\allowbreak > la flèche \$\textbackslash{}mathrm\{ch\}\_\{K\textasciicircum{}\{\textbackslash{}bigdot\}\}\$.\par
\textbf{GPT-6.1 Sol.} nothing reported here\par
a fix from one model only\par
\par\noindent\hypertarget{DP0577}{}\textbf{DP0577 -- XVII, p. 478 (one\_model\_fix)}\par
diagram, second row\par
\textbf{Claude Opus 5.5.} Place: diagram, second row; Reported print: Ker S\textasciicircum{}\{2d\}(F); Source TeX: \textbackslash{}ker S\textasciicircum{}\{2d\}(F); Note: \textbackslash{}ker prints a lower-case « ker »; the print has Ker, as in the first row.; Proposed edit: \textbackslash{}ker S\textasciicircum{}\{2d\}(F) =\allowbreak > \textbackslash{}Ker S\textasciicircum{}\{2d\}(F)\par
\textbf{GPT-6.1 Sol.} nothing reported here\par
a fix from one model only\par
\par\noindent\hypertarget{DP0578}{}\textbf{DP0578 -- XVII, p. 478 (one\_model\_fix)}\par
diagram, third row\par
\textbf{Claude Opus 5.5.} Place: diagram, third row; Reported print: Ker S\textasciicircum{}\{2d\}(F''); Source TeX: \textbackslash{}ker S\textasciicircum{}\{2d\}(F''); Note: Same as the second row.; Proposed edit: \textbackslash{}ker S\textasciicircum{}\{2d\}(F'') =\allowbreak > \textbackslash{}Ker S\textasciicircum{}\{2d\}(F'')\par
\textbf{GPT-6.1 Sol.} nothing reported here\par
a fix from one model only\par
\par\noindent\hypertarget{DP0579}{}\textbf{DP0579 -- XVIII, p. 481 (one\_model\_fix)}\par
Attribution sous le titre\par
\textbf{Claude Opus 5.5.} nothing reported here\par
\textbf{GPT-6.1 Sol.} Place: Attribution sous le titre; Reported print: par P. DELIGNE; Source TeX: P. Deligne; Note: Le mot « par » manque devant le nom de l'auteur.; Proposed edit: \textbackslash{}begin\{center\} P. Deligne =\allowbreak > \textbackslash{}begin\{center\} par P. Deligne\par
a fix from one model only\par
\par\noindent\hypertarget{DP0580}{}\textbf{DP0580 -- XVIII, p. 490 (conflict)}\par
est exacte, d'après (XVII ...) | Reference immediately after the opening exact sequence\par
\textbf{Claude Opus 5.5.} Place: est exacte, d'après (XVII ...); Reported print: (XVII 6.2.8); Source TeX: 6.2.9; Note: The print cites XVII 6.2.8. The TeX prints 6.2.9 without a \textbackslash{}sisi. If this is an intended correction, it should be marked as one.; Proposed edit: \textbackslash{}hyperrefext[XVII][XVII.prop6.2.9]\{6.2.9\} =\allowbreak > \textbackslash{}hyperrefext[XVII][XVII.prop6.2.9]\{6.2.8\}\par
\textbf{GPT-6.1 Sol.} Place: Reference immediately after the opening exact sequence; Reported print: XVII 6.2.8; Source TeX: \textbackslash{}hyperrefext[XVII][XVII.prop6.2.9]\{6.2.9\}; Note: The printed reference is 6.2.8, also clearly so in the witness. The replacement avoids inventing an unavailable target label.; Proposed edit: \textbackslash{}hyperrefext[XVII][XVII.prop6.2.9]\{6.2.9\} =\allowbreak > XVII 6.2.8\par
the two models differ at this place\par
\par\noindent\hypertarget{DP0581}{}\textbf{DP0581 -- XVIII, p. 493 (one\_model\_fix)}\par
Démonstration de 1.2.3, décomposition du faisceau inversible\par
\textbf{Claude Opus 5.5.} nothing reported here\par
\textbf{GPT-6.1 Sol.} Place: Démonstration de 1.2.3, décomposition du faisceau inversible; Reported print: \textbackslash{}Ll\textbackslash{}sim\textbackslash{}pr\_2\textasciicircum{}\{\textbackslash{}ast\}(x\_\{G\textasciicircum{}\{\textbackslash{}ast\}\})\textasciicircum{}\{\textbackslash{}ast\}\textbackslash{}Ll\textbackslash{}otimes\textbackslash{}Mm; Source TeX: \textbackslash{}Ll\textbackslash{}sim  \textbackslash{}pr\textasciicircum{}\{\textbackslash{}ast\}\_\{2\}(x\textbackslash{}times G\textasciicircum{}\textbackslash{}ast)\textasciicircum{}\{\textbackslash{}ast\}\textbackslash{}Ll\textbackslash{}otimes \textbackslash{}Mm; Note: L'imprimé porte x indicé par G*, notation du changement de base, et non le produit x \ensuremath{\times} G*.; Proposed edit: (x\textbackslash{}times G\textasciicircum{}\textbackslash{}ast) =\allowbreak > (x\_\{G\textasciicircum{}\textbackslash{}ast\})\par
a fix from one model only\par
\par\noindent\hypertarget{DP0582}{}\textbf{DP0582 -- XVIII, p. 495 (conflict)}\par
proof of 1.2.6, before the definition of \ensuremath{\Psi}\_1 | Démonstration de 1.2.6, après la translation à droite\par
\textbf{Claude Opus 5.5.} Place: proof of 1.2.6, before the definition of \ensuremath{\Psi}\_1; Reported print: un homomorphisme de W\_n dans G\_n qui transforme o en e; Source TeX: un morphisme de \$W\_\{n\}\$ dans \$G\_\{n\}\$; Note: The print reads 'homomorphisme'; the TeX has 'morphisme'.; Proposed edit: un morphisme de \$W\_\{n\}\$ =\allowbreak > un homomorphisme de \$W\_\{n\}\$\par
\textbf{GPT-6.1 Sol.} Place: Démonstration de 1.2.6, après la translation à droite; Reported print: un homomorphisme de W\_n dans G\_n qui transforme 0 en e; Source TeX: un morphisme de \$W\_\{n\}\$ dans \$G\_\{n\}\$ qui transforme \$0\$ en \$e\$; Proposed correction: un morphisme de \$W\_\{n\}\$ dans \$G\_\{n\}\$ qui transforme \$0\$ en \$e\$; Note: Le TeX corrige silencieusement « homomorphisme » en « morphisme ». La construction fournit un morphisme de schémas pointés, non un homomorphisme de groupes ; ces voisinages infinitésimaux ne sont pas, en général, des sous-groupes.\par
the two models differ at this place: one model fixes the TeX, the other calls it a misprint of the print\par
\par\noindent\hypertarget{DP0583}{}\textbf{DP0583 -- XVIII, p. 496 (one\_model\_fix)}\par
Fin de 1.2.7\par
\textbf{Claude Opus 5.5.} nothing reported here\par
\textbf{GPT-6.1 Sol.} Place: Fin de 1.2.7; Reported print: Le paragraphe se termine après « sur P\textasciicircum{}r\_L. », sans note correspondante.; Source TeX: \textbackslash{}acorriger\{Le lecteur pourra trouver une introduction au langage de la géométrie; Note: Tout le passage inclus dans \textbackslash{}acorriger est absent de la page imprimée.; Proposed edit: \textbackslash{}acorriger\{Le lecteur pourra trouver une introduction au langage de la géométrie analytique rigide dans \textbackslash{}cite\{SGA4-18:5bis\}. L'inexistence de sections non-triviales pour les faisceaux \$\textbackslash{}Oo(n)\$ pour \$n<0\$ dans ce contexte résulte par exemple d'une version rigide de GAGA, voir \textbackslash{}cite\{SGA4-18:3ter\}.\} =\allowbreak > \par
a fix from one model only\par
\par\noindent\hypertarget{DP0584}{}\textbf{DP0584 -- XVIII, p. 496 (one\_model\_fix)}\par
Preuve de 1.2.8, fin de la réduction aux groupes élémentaires\par
\textbf{Claude Opus 5.5.} nothing reported here\par
\textbf{GPT-6.1 Sol.} Place: Preuve de 1.2.8, fin de la réduction aux groupes élémentaires; Reported print: \ensuremath{𝛍}\_\{\textbackslash{}ell\_\{i\}\}; Source TeX: \ensuremath{𝛍}\_\{\textbackslash{}ell i\}; Note: L'indice est \ensuremath{\ell}\_i : i doit être en indice de \ensuremath{\ell}.; Proposed edit: \ensuremath{𝛍}\_\{\textbackslash{}ell i\} =\allowbreak > \ensuremath{𝛍}\_\{\textbackslash{}ell\_\{i\}\}\par
a fix from one model only\par
\par\noindent\hypertarget{DP0585}{}\textbf{DP0585 -- XVIII, p. 498 (conflict)}\par
Pour G plat sur S, cotangent complex | Paragraphe sur le complexe de Lie, première désignation du complexe cotangent\par
\textbf{Claude Opus 5.5.} Place: Pour G plat sur S, cotangent complex; Reported print: Si L\ensuremath{\Omega}(G\_s/k(s)) est le complexe cotangent; Source TeX: Si \$L\textbackslash{}cap (G\_\{s\}/k(s))\$; Note: The handwritten sign after L is a capital Omega (with feet), not an intersection sign. Same at the second occurrence.; Proposed edit: Si \$L\textbackslash{}cap (G\_\{s\}/k(s))\$ =\allowbreak > Si \$L\textbackslash{}Omega (G\_\{s\}/k(s))\$\par
\textbf{GPT-6.1 Sol.} Place: Paragraphe sur le complexe de Lie, première désignation du complexe cotangent; Reported print: L\textbackslash{}Omega(G\_s/k(s)); Source TeX: \$L\textbackslash{}cap (G\_\{s\}/k(s))\$; Note: Les deux scans montrent un oméga majuscule, non un signe d'intersection.; Proposed edit: \$L\textbackslash{}cap (G\_\{s\}/k(s))\$ =\allowbreak > \$L\textbackslash{}Omega(G\_\{s\}/k(s))\$\par
the two models differ at this place\par
\par\noindent\hypertarget{DP0586}{}\textbf{DP0586 -- XVIII, p. 498 (conflict)}\par
est le dual de Le\textasciicircum{}*(...) | Même paragraphe, expression dont L est le dual\par
\textbf{Claude Opus 5.5.} Place: est le dual de Le\textasciicircum{}*(...); Reported print: Le\textasciicircum{}*(L\ensuremath{\Omega}(G\_s/k(s))); Source TeX: \$Le\textasciicircum{}\{\textbackslash{}ast\}(L\textbackslash{}cap (G\_\{s\}/k(s)))\$; Note: Same handwritten Omega as at the first occurrence.; Proposed edit: \$Le\textasciicircum{}\{\textbackslash{}ast\}(L\textbackslash{}cap (G\_\{s\}/k(s)))\$ =\allowbreak > \$Le\textasciicircum{}\{\textbackslash{}ast\}(L\textbackslash{}Omega (G\_\{s\}/k(s)))\$\par
\textbf{GPT-6.1 Sol.} Place: Même paragraphe, expression dont L est le dual; Reported print: Le\textasciicircum{}\{\textbackslash{}ast\}(L\textbackslash{}Omega(G\_s/k(s))); Source TeX: Le\textasciicircum{}\{\textbackslash{}ast\}(L\textbackslash{}cap (G\_\{s\}/k(s))); Note: La seconde occurrence porte également \ensuremath{\Omega} dans les deux scans.; Proposed edit: Le\textasciicircum{}\{\textbackslash{}ast\}(L\textbackslash{}cap (G\_\{s\}/k(s))) =\allowbreak > Le\textasciicircum{}\{\textbackslash{}ast\}(L\textbackslash{}Omega(G\_\{s\}/k(s)))\par
the two models differ at this place\par
\par\noindent\hypertarget{DP0587}{}\textbf{DP0587 -- XVIII, p. 498 (agreed\_fix\_held)}\par
Les déformations ... (Illusie [1]) | Référence à Illusie après la classification des déformations\par
\textbf{Claude Opus 5.5.} Place: Les déformations ... (Illusie [1]); Reported print: (Illusie [1]); Source TeX: (\textbackslash{}textsc\{Illusie\} \textbackslash{}cite[VII\textasciitilde{}2.4.4\textasciitilde{}(ii)]\{SGA4-18:3\}); Note: The print cites only [1], with no location. 'VII 2.4.4 (ii)' is an editorial addition not marked by \textbackslash{}sisi.; Proposed edit: \textbackslash{}cite[VII\textasciitilde{}2.4.4\textasciitilde{}(ii)]\{SGA4-18:3\} =\allowbreak > \textbackslash{}cite\{SGA4-18:3\}\par
\textbf{GPT-6.1 Sol.} Place: Référence à Illusie après la classification des déformations; Reported print: Illusie [1]; Source TeX: \textbackslash{}cite[VII\textasciitilde{}2.4.4\textasciitilde{}(ii)]\{SGA4-18:3\}; Note: L'imprimé ne donne pas le complément « VII 2.4.4 (ii) ».; Proposed edit: \textbackslash{}cite[VII\textasciitilde{}2.4.4\textasciitilde{}(ii)]\{SGA4-18:3\} =\allowbreak > \textbackslash{}cite\{SGA4-18:3\}\par
both models give the same fix, held: it changes the editors' apparatus (the target of a \textbackslash{}cite): a \textbackslash{}sisi correction, note, label, cross-reference or citation is the editors' record and is not removed or retargeted on the models' word\par
\par\noindent\hypertarget{DP0588}{}\textbf{DP0588 -- XVIII, p. 499 (agreed\_fix\_held)}\par
first line, formal isomorphism \ensuremath{\Psi}\textasciicircum{} | Première formule, source de l'isomorphisme formel\par
\textbf{Claude Opus 5.5.} Place: first line, formal isomorphism \ensuremath{\Psi}\textasciicircum{}; Reported print: \ensuremath{\Psi}\textasciicircum{} : p\textasciicircum{}*s\textasciicircum{}*K + \ensuremath{\phi}\textasciicircum{}O(1) \ensuremath{⥲} K\textasciicircum{}; Source TeX: \textbackslash{}widehat\{p\textasciicircum{}\{\textbackslash{}ast\}s\textasciicircum{}\{\textbackslash{}ast\}K\}+ \textbackslash{}widehat\{\textbackslash{}varphi\}\textbackslash{}widehat\{\textbackslash{}Oo(1)\}; Note: In the print only \ensuremath{\phi} (and K on the right) carry the hat. p\textasciicircum{}*s\textasciicircum{}*K and O(1) are unhatted; elsewhere the print writes the postfix \textasciicircum{} when O(1) is completed.; Proposed edit: \textbackslash{}widehat\{p\textasciicircum{}\{\textbackslash{}ast\}s\textasciicircum{}\{\textbackslash{}ast\}K\}+ \textbackslash{}widehat\{\textbackslash{}varphi\}\textbackslash{}widehat\{\textbackslash{}Oo(1)\} =\allowbreak > p\textasciicircum{}\{\textbackslash{}ast\}s\textasciicircum{}\{\textbackslash{}ast\}K+ \textbackslash{}widehat\{\textbackslash{}varphi\}\textbackslash{}Oo(1)\par
\textbf{GPT-6.1 Sol.} Place: Première formule, source de l'isomorphisme formel; Reported print: p\textasciicircum{}\{\textbackslash{}ast\}s\textasciicircum{}\{\textbackslash{}ast\}K+\textbackslash{}widehat\{\textbackslash{}varphi\}\textbackslash{}Oo(1); Source TeX: \textbackslash{}widehat\{p\textasciicircum{}\{\textbackslash{}ast\}s\textasciicircum{}\{\textbackslash{}ast\}K\}+ \textbackslash{}widehat\{\textbackslash{}varphi\}\textbackslash{}widehat\{\textbackslash{}Oo(1)\}; Note: L'imprimé ne porte de chapeau ni sur p*s*K ni sur O(1) dans cette formule.; Proposed edit: \textbackslash{}widehat\{p\textasciicircum{}\{\textbackslash{}ast\}s\textasciicircum{}\{\textbackslash{}ast\}K\}+ \textbackslash{}widehat\{\textbackslash{}varphi\}\textbackslash{}widehat\{\textbackslash{}Oo(1)\} =\allowbreak > p\textasciicircum{}\{\textbackslash{}ast\}s\textasciicircum{}\{\textbackslash{}ast\}K+ \textbackslash{}widehat\{\textbackslash{}varphi\}\textbackslash{}Oo(1)\par
both models give the same fix, held: a deletion in the first or last three lines of the page (the page cut may be off)\par
\par\noindent\hypertarget{DP0589}{}\textbf{DP0589 -- XVIII, p. 499 (agreed\_fix\_held)}\par
heading 1.2.12, word | Titre de 1.2.12\par
\textbf{Claude Opus 5.5.} Place: heading 1.2.12, word; Reported print: Algébraïcité; Source TeX: Algébricité; Note: The print spells it Algébraïcité. The edition normalises the spelling.; Proposed edit: Algébricité de \$\textbackslash{}widehat\{\textbackslash{}varphi\}\$: vérification =\allowbreak > Algébraïcité de \$\textbackslash{}widehat\{\textbackslash{}varphi\}\$: vérification\par
\textbf{GPT-6.1 Sol.} Place: Titre de 1.2.12; Reported print: Algébraïcité; Source TeX: Algébricité; Note: Le titre imprimé emploie la graphie « Algébraïcité ».; Proposed edit: Algébricité de \$\textbackslash{}widehat\{\textbackslash{}varphi\}\$: =\allowbreak > Algébraïcité de \$\textbackslash{}widehat\{\textbackslash{}varphi\}\$:\par
both models give the same fix, held: a model marks it low confidence\par
\par\noindent\hypertarget{DP0590}{}\textbf{DP0590 -- XVIII, p. 499 (agreed\_fix\_held)}\par
heading 1.2.13, word | Titre de 1.2.13\par
\textbf{Claude Opus 5.5.} Place: heading 1.2.13, word; Reported print: Algébraïcité; Source TeX: Algébricité; Note: Same spelling as in the heading of 1.2.12.; Proposed edit: Algébricité de \$\textbackslash{}widehat\{\textbackslash{}varphi\}\$\} =\allowbreak > Algébraïcité de \$\textbackslash{}widehat\{\textbackslash{}varphi\}\$\}\par
\textbf{GPT-6.1 Sol.} Place: Titre de 1.2.13; Reported print: Algébraïcité; Source TeX: Algébricité; Note: Ce second titre porte lui aussi « Algébraïcité ».; Proposed edit: Algébricité de \$\textbackslash{}widehat\{\textbackslash{}varphi\}\$\} =\allowbreak > Algébraïcité de \$\textbackslash{}widehat\{\textbackslash{}varphi\}\$\}\par
both models give the same fix, held: a model marks it low confidence\par
\par\noindent\hypertarget{DP0591}{}\textbf{DP0591 -- XVIII, p. 499 (one\_model\_fix)}\par
1.2.12 b), « Le morphisme formel… de 1.2.11 »\par
\textbf{Claude Opus 5.5.} nothing reported here\par
\textbf{GPT-6.1 Sol.} Place: 1.2.12 b), « Le morphisme formel… de 1.2.11 »; Reported print: \textbackslash{}widehat\{\textbackslash{}varphi\}; Source TeX: \textbackslash{}widehat\{\textbackslash{}Psi\}; Note: L'imprimé désigne l'homomorphisme \ensuremath{\hat{\phi}}, non l'isomorphisme \ensuremath{\hat{\Psi}}.; Proposed edit: formel \$\textbackslash{}widehat\{\textbackslash{}Psi\}\$ de =\allowbreak > formel \$\textbackslash{}widehat\{\textbackslash{}varphi\}\$ de\par
a fix from one model only\par
\par\noindent\hypertarget{DP0592}{}\textbf{DP0592 -- XVIII, p. 501 (agreed\_fix\_held)}\par
1.3.1.1 (iii) | 1.3.1.1 (iii), première référence après « soit les conditions » | 1.3.1.1 (iii), dernière référence\par
\textbf{Claude Opus 5.5.} Place: 1.3.1.1 (iii); Reported print: (XVII 6.3.24.1) et (XVII 6.3.24.2); Source TeX: (\textbackslash{}ref\{XVII.par6.3.24.1\}) \textbackslash{}emph\{et\} (\textbackslash{}ref\{XVII.par6.3.24.2\}); Note: \textbackslash{}ref prints only 6.3.24.1 and 6.3.24.2. The print gives the exposé numeral XVII with both references.; Proposed edit: (\textbackslash{}ref\{XVII.par6.3.24.1\}) \textbackslash{}emph\{et\} (\textbackslash{}ref\{XVII.par6.3.24.2\}) =\allowbreak > (\textbackslash{}hyperrefext[XVII][XVII.par6.3.24.1]\{6.3.24.1\}) \textbackslash{}emph\{et\} (\textbackslash{}hyperrefext[XVII][XVII.par6.3.24.2]\{6.3.24.2\})\par
\textbf{GPT-6.1 Sol.} Place: 1.3.1.1 (iii), première référence après « soit les conditions »; Reported print: (XVII 6.3.24.1); Source TeX: (\textbackslash{}ref\{XVII.par6.3.24.1\}); Note: Le simple \textbackslash{}ref omet le numéro d'exposé XVII imprimé devant la référence.; Proposed edit: \textbackslash{}ref\{XVII.par6.3.24.1\} =\allowbreak > \textbackslash{}hyperrefext[XVII][XVII.par6.3.24.1]\{6.3.24.1\} || Place: 1.3.1.1 (iii), dernière référence; Reported print: (XVII 6.3.24.2); Source TeX: (\textbackslash{}ref\{XVII.par6.3.24.2\}); Note: Le simple \textbackslash{}ref omet ici aussi le numéro d'exposé XVII.; Proposed edit: \textbackslash{}ref\{XVII.par6.3.24.2\} =\allowbreak > \textbackslash{}hyperrefext[XVII][XVII.par6.3.24.2]\{6.3.24.2\}\par
both models give the same fix, held: it changes the editors' apparatus (\textbackslash{}hyperrefext, \textbackslash{}ref, the target of a cross-reference): a \textbackslash{}sisi correction, note, label, cross-reference or citation is the editors' record and is not removed or retargeted on the models' word\par
\par\noindent\hypertarget{DP0593}{}\textbf{DP0593 -- XVIII, p. 501 (one\_model\_fix)}\par
1.3.1.1 (ii'), exposant du second Hom\par
\textbf{Claude Opus 5.5.} nothing reported here\par
\textbf{GPT-6.1 Sol.} Place: 1.3.1.1 (ii'), exposant du second Hom; Reported print: \textbackslash{}mathfrak\{S\}\_n; Source TeX: \textbackslash{}sigma\_n; Note: Le signe manuscrit du groupe symétrique est un S gothique, non un sigma minuscule.; Proposed edit: \textbackslash{}sigma\_n =\allowbreak > \textbackslash{}mathfrak\{S\}\_n\par
a fix from one model only\par
\par\noindent\hypertarget{DP0594}{}\textbf{DP0594 -- XVIII, p. 502 (one\_model\_fix)}\par
1.3.2.1, second display\par
\textbf{Claude Opus 5.5.} Place: 1.3.2.1, second display; Reported print: \ensuremath{\geqslant} 1 si g =\allowbreak  0, et deg\_\{X\_s\}(\ensuremath{\mathcal{L}}) \ensuremath{\geqslant} 2; Source TeX: g=\allowbreak 0,\textbackslash{}quaet; Note: \textbackslash{}quaet looks like a typo for \textbackslash{}quad et. Unless macros.tex defines it, the word « et » is lost or the line does not compile.; Proposed edit: g=\allowbreak 0,\textbackslash{}quaet =\allowbreak > g=\allowbreak 0,\textbackslash{}quad\textbackslash{}mbox\{et\}\textbackslash{}quad\par
\textbf{GPT-6.1 Sol.} nothing reported here\par
a fix from one model only\par
\par\noindent\hypertarget{DP0595}{}\textbf{DP0595 -- XVIII, p. 504 (one\_model\_fix)}\par
display before 1.3.3.8\par
\textbf{Claude Opus 5.5.} Place: display before 1.3.3.8; Reported print: u\_s\textasciicircum{}*O(1) + U\_s\textasciicircum{}\ensuremath{\vee}; Source TeX: u\textasciicircum{}\{\textbackslash{}ast\}\textbackslash{}Oo(1)\textbackslash{}simeq \textbackslash{}mathcal\{U\}\textasciicircum{}\{\textbackslash{}vee\}; Note: The print puts the index s on both u\textasciicircum{}* and U\textasciicircum{}\ensuremath{\vee}. The TeX drops it.; Proposed edit: u\textasciicircum{}\{\textbackslash{}ast\}\textbackslash{}Oo(1)\textbackslash{}simeq \textbackslash{}mathcal\{U\}\textasciicircum{}\{\textbackslash{}vee\} =\allowbreak > u\textasciicircum{}\{\textbackslash{}ast\}\_\{s\}\textbackslash{}Oo(1)\textbackslash{}simeq \textbackslash{}mathcal\{U\}\textasciicircum{}\{\textbackslash{}vee\}\_\{s\}\par
\textbf{GPT-6.1 Sol.} nothing reported here\par
a fix from one model only\par
\par\noindent\hypertarget{DP0596}{}\textbf{DP0596 -- XVIII, p. 505 (one\_model\_fix)}\par
Diagrammes (1.3.5.2) et (1.3.5.3), quatre isomorphismes horizontaux\par
\textbf{Claude Opus 5.5.} nothing reported here\par
\textbf{GPT-6.1 Sol.} Place: Diagrammes (1.3.5.2) et (1.3.5.3), quatre isomorphismes horizontaux; Reported print: \textbackslash{}simeq, sans pointe de flèche; Source TeX: \textbackslash{}ar[r]\textasciicircum{}\{\textbackslash{}sim\}; Note: Convention locale répétée : les quatre signes d'isomorphisme horizontaux sont transformés en flèches. Le correctif donné supprime la pointe du premier; la même correction est à appliquer aux trois autres.; Proposed edit: \textbackslash{}langle \textbackslash{}Oo(D) , K\_1 + K\_2] \textbackslash{}ar[r]\textasciicircum{}\{\textbackslash{}sim\} =\allowbreak > \textbackslash{}langle \textbackslash{}Oo(D) , K\_1 + K\_2] \textbackslash{}ar@\{-\}[r]\textasciicircum{}\{\textbackslash{}sim\}\par
a fix from one model only\par
\par\noindent\hypertarget{DP0597}{}\textbf{DP0597 -- XVIII, p. 510 (conflict)}\par
D'après XVII 6.3., on dispose ainsi… | Référence précédant la formule (1.3.10.2)\par
\textbf{Claude Opus 5.5.} Place: D'après XVII 6.3., on dispose ainsi…; Reported print: D'après XVII 6.3.,; Source TeX: \textbackslash{}hyperrefext[XVII][XVII.prop6.3.9]\{6.3.9\}; Note: The print cites only XVII 6.3. The TeX makes it XVII 6.3.9 without a \textbackslash{}sisi.; Proposed edit: [XVII.prop6.3.9]\{6.3.9\} =\allowbreak > [XVII.prop6.3.9]\{6.3\}\par
\textbf{GPT-6.1 Sol.} Place: Référence précédant la formule (1.3.10.2); Reported print: D'après XVII 6.3.; Source TeX: D'après \textbackslash{}hyperrefext[XVII][XVII.prop6.3.9]\{6.3.9\}; Note: L'imprimé renvoie à la section XVII 6.3, sans numéro de proposition ; le TeX ajoute « .9 ».; Proposed edit: \textbackslash{}hyperrefext[XVII][XVII.prop6.3.9]\{6.3.9\} =\allowbreak > XVII 6.3\par
the two models differ at this place\par
\par\noindent\hypertarget{DP0598}{}\textbf{DP0598 -- XVIII, p. 511 (conflict)}\par
second paragraph | Premier paragraphe, « un isomorphisme de foncteur »\par
\textbf{Claude Opus 5.5.} Place: second paragraph; Reported print: un isomorphisme de foncteur; Source TeX: un isomorphisme de foncteurs; Proposed correction: de foncteurs; Note: The print has the singular 'foncteur'. The TeX corrects it to 'foncteurs'; on p. 512 it keeps the printed 'isomorphisme de foncteur'.\par
\textbf{GPT-6.1 Sol.} Place: Premier paragraphe, « un isomorphisme de foncteur »; Reported print: de foncteur entre les restrictions; Source TeX: de foncteurs entre les restrictions; Note: Le tirage emploie le singulier « foncteur », confirmé par le scan témoin.; Proposed edit: de foncteurs entre =\allowbreak > de foncteur entre\par
the two models differ at this place: one model fixes the TeX, the other calls it a misprint of the print\par
\par\noindent\hypertarget{DP0599}{}\textbf{DP0599 -- XVIII, p. 518 (agreed\_fix\_held)}\par
«D'après Serre [1]» | Citation de Serre avant (1.3.18.1)\par
\textbf{Claude Opus 5.5.} Place: «D'après Serre [1]»; Reported print: Serre [1]; Source TeX: \textbackslash{}cite[Ch.\textasciitilde{}III, prop.\textasciitilde{}6, p.\textasciitilde{}44]\{SGA4-18:6\}; Note: The location «Ch. III, prop. 6, p. 44» is not in the print.; Proposed edit: \textbackslash{}cite[Ch.\textasciitilde{}III, prop.\textasciitilde{}6, p.\textasciitilde{}44]\{SGA4-18:6\} =\allowbreak > \textbackslash{}cite\{SGA4-18:6\}\par
\textbf{GPT-6.1 Sol.} Place: Citation de Serre avant (1.3.18.1); Reported print: D'après Serre [1], on a; Source TeX: \textbackslash{}cite[Ch.\textasciitilde{}III, prop.\textasciitilde{}6, p.\textasciitilde{}44]\{SGA4-18:6\}; Note: L'imprimé ne donne ni chapitre, ni proposition, ni page dans cette citation.; Proposed edit: \textbackslash{}cite[Ch.\textasciitilde{}III, prop.\textasciitilde{}6, p.\textasciitilde{}44]\{SGA4-18:6\} =\allowbreak > \textbackslash{}cite\{SGA4-18:6\}\par
both models give the same fix, held: it changes the editors' apparatus (the target of a \textbackslash{}cite): a \textbackslash{}sisi correction, note, label, cross-reference or citation is the editors' record and is not removed or retargeted on the models' word\par
\par\noindent\hypertarget{DP0600}{}\textbf{DP0600 -- XVIII, p. 518 (conflict)}\par
«On en tire que», second line of the computation | Calcul après « On en tire que », deuxième ligne\par
\textbf{Claude Opus 5.5.} Place: «On en tire que», second line of the computation; Reported print: (f\_1\textasciicircum{}\{v\_x(g\_2)\} . f\_2\textasciicircum{}\{-v\_x(f\_1)\})(x); Source TeX: \textbackslash{}cdot g\_\{2\}\textasciicircum{}\{-v\_\{x\}(f\_\{1\})\})(x); Note: The print has f\_2 where the TeX has g\_2.; Proposed edit: \textbackslash{}cdot g\_\{2\}\textasciicircum{}\{-v\_\{x\}(f\_\{1\})\})(x) =\allowbreak > \textbackslash{}cdot f\_\{2\}\textasciicircum{}\{-v\_\{x\}(f\_\{1\})\})(x)\par
\textbf{GPT-6.1 Sol.} Place: Calcul après « On en tire que », deuxième ligne; Reported print: f\_\{2\}\textasciicircum{}\{-v\_x(f\_1)\}; Source TeX: g\_\{2\}\textasciicircum{}\{-v\_\{x\}(f\_\{1\})\}; Proposed correction: g\_\{2\}\textasciicircum{}\{-v\_\{x\}(f\_\{1\})\}; Note: L'imprimé porte bien f₂, également dans le témoin. Il faut g₂, comme dans (1.3.18.1) ; le TeX corrige silencieusement cette lettre.\par
the two models differ at this place: one model fixes the TeX, the other calls it a misprint of the print\par
\par\noindent\hypertarget{DP0601}{}\textbf{DP0601 -- XVIII, p. 522 (one\_model\_fix)}\par
end of the paragraph after Définition 1.4.5\par
\textbf{Claude Opus 5.5.} Place: end of the paragraph after Définition 1.4.5; Reported print: et Aut(e) est un faisceau abélien; Source TeX: et \$\textbackslash{}sheafAut(e)\$ est un faisceau abélien; Note: The print writes plain Aut(e), not underlined, so it does not mark the sheaf Aut here. The TeX uses the sheaf macro instead.; Proposed edit: \$\textbackslash{}sheafAut(e)\$ =\allowbreak > \$\textbackslash{}Aut(e)\$\par
\textbf{GPT-6.1 Sol.} nothing reported here\par
a fix from one model only\par
\par\noindent\hypertarget{DP0602}{}\textbf{DP0602 -- XVIII, p. 530 (conflict)}\par
(b) | Construction 1.4.20 (b), before the final formula\par
\textbf{Claude Opus 5.5.} Place: (b); Reported print: et f : K'\_1\ensuremath{\otimes}K'\_2 \ensuremath{\to} L tel que; Source TeX: tels que; Proposed correction: tels que; Note: The print has the singular 'tel que'. The TeX silently corrects it to 'tels que'.\par
\textbf{GPT-6.1 Sol.} Place: Construction 1.4.20 (b), before the final formula; Reported print: tel que; Source TeX: tels que; Note: The print has singular « tel que » after the morphism f; TeX pluralizes it.; Proposed edit: tels que =\allowbreak > tel que\par
the two models differ at this place: one model fixes the TeX, the other calls it a misprint of the print\par
\par\noindent\hypertarget{DP0603}{}\textbf{DP0603 -- XVIII, p. 537 (agreed\_fix\_held)}\par
Lemme 1.5.7, reference to Grothendieck | Lemme 1.5.7, référence à Grothendieck\par
\textbf{Claude Opus 5.5.} Place: Lemme 1.5.7, reference to Grothendieck; Reported print: (Grothendieck[1]); Source TeX: \textbackslash{}cite[Corollaire\textasciitilde{}11.9]\{SGA4-18:2\}; Note: The print gives only the bibliography reference [1]. The precision « Corollaire 11.9 » is added by the TeX.; Proposed edit: \textbackslash{}cite[Corollaire\textasciitilde{}11.9]\{SGA4-18:2\} =\allowbreak > \textbackslash{}cite\{SGA4-18:2\}\par
\textbf{GPT-6.1 Sol.} Place: Lemme 1.5.7, référence à Grothendieck; Reported print: Grothendieck suivi de la seule référence bibliographique, sans indication de corollaire; Source TeX: \textbackslash{}cite[Corollaire\textasciitilde{}11.9]\{SGA4-18:2\}; Note: L'indication « Corollaire 11.9 » est ajoutée au texte imprimé.; Proposed edit: \textbackslash{}cite[Corollaire\textasciitilde{}11.9]\{SGA4-18:2\} =\allowbreak > \textbackslash{}cite\{SGA4-18:2\}\par
both models give the same fix, held: it changes the editors' apparatus (the target of a \textbackslash{}cite): a \textbackslash{}sisi correction, note, label, cross-reference or citation is the editors' record and is not removed or retargeted on the models' word\par
\par\noindent\hypertarget{DP0604}{}\textbf{DP0604 -- XVIII, p. 538 (conflict)}\par
(1.5.10.1) | Fin de la condition (1.5.10.1)\par
\textbf{Claude Opus 5.5.} Place: (1.5.10.1); Reported print: valant 1 sur Y\_1U; Source TeX: \textbackslash{}fprod\{Y\}\{U\}\{S\}; Proposed correction: \textbackslash{}fprod\{Y\}\{U\}\{S\}; Note: The print has the garbled « Y\_1U »; the TeX silently corrects it to Y \ensuremath{\times}\_S U.\par
\textbf{GPT-6.1 Sol.} Place: Fin de la condition (1.5.10.1); Reported print: Y\textbackslash{}cap U; Source TeX: \textbackslash{}fprod\{Y\}\{U\}\{S\}; Note: L'imprimé donne l'intersection avec U, lequel vient d'être défini comme l'image réciproque d'un voisinage de s, et non un produit fibré sur S.; Proposed edit: \textbackslash{}fprod\{Y\}\{U\}\{S\} =\allowbreak > Y\textbackslash{}cap U\par
the two models differ at this place: one model fixes the TeX, the other calls it a misprint of the print\par
\par\noindent\hypertarget{DP0605}{}\textbf{DP0605 -- XVIII, p. 538 (one\_model\_fix)}\par
Condition (1.5.10.4), droite affine\par
\textbf{Claude Opus 5.5.} nothing reported here\par
\textbf{GPT-6.1 Sol.} Place: Condition (1.5.10.4), droite affine; Reported print: g:\textbackslash{}EE\textasciicircum{}\{1\}\_\{S\}\textbackslash{}rightarrow S; Source TeX: g:\textbackslash{}EE\textasciicircum{}\{1\}\_\{s\}\textbackslash{}rightarrow S; Note: L'indice imprimé est un S majuscule, confirmé par le témoin.; Proposed edit: g:\textbackslash{}EE\textasciicircum{}\{1\}\_\{s\}\textbackslash{}rightarrow S =\allowbreak > g:\textbackslash{}EE\textasciicircum{}\{1\}\_\{S\}\textbackslash{}rightarrow S\par
a fix from one model only\par
\par\noindent\hypertarget{DP0606}{}\textbf{DP0606 -- XVIII, p. 541 (conflict)}\par
1.5.15, line before the display \ensuremath{\langle}L, tr(K)] | 1.5.15, hypothesis on the line bundle\par
\textbf{Claude Opus 5.5.} Place: 1.5.15, line before the display \ensuremath{\langle}L, tr(K)]; Reported print: pour \ensuremath{𝓛} faisceau inversible sur X\_2 et K torseur sur X\_1,; Source TeX: pour \$\textbackslash{}Ll\$ faisceau inversible sur \$\textbackslash{}bar\{X\}\_\{2\}\$ rigidifié le long de \$Y\_2\$ et \$K\$ torseur sur \$X\_\{1\}\$,; Note: The TeX adds the clause 'rigidifié le long de Y\_2' and changes X\_2 to \textbackslash{}bar X\_2, without \textbackslash{}sisi.; Proposed edit: pour \$\textbackslash{}Ll\$ faisceau inversible sur \$\textbackslash{}bar\{X\}\_\{2\}\$ rigidifié le long de \$Y\_2\$ =\allowbreak > pour \$\textbackslash{}Ll\$ faisceau inversible sur \$X\_\{2\}\$\par
\textbf{GPT-6.1 Sol.} Place: 1.5.15, hypothesis on the line bundle; Reported print: pour \ensuremath{\mathcal{L}} faisceau inversible sur X\_2 et K torseur sur X\_1; Source TeX: rigidifié le long de \$Y\_2\$; Note: The printed sentence contains no rigidification clause.; Proposed edit:  rigidifié le long de \$Y\_2\$ =\allowbreak > \par
the two models differ at this place\par
\par\noindent\hypertarget{DP0607}{}\textbf{DP0607 -- XVIII, p. 541 (one\_model\_fix)}\par
Direct images throughout this page and the diagrams on p. 542\par
\textbf{Claude Opus 5.5.} nothing reported here\par
\textbf{GPT-6.1 Sol.} Place: Direct images throughout this page and the diagrams on p. 542; Reported print: f\_\{\textbackslash{}ast\} and R\textasciicircum{}\{1\}f\_\{\textbackslash{}ast\}, without a bar on f; Source TeX: \textbackslash{}bar\{f\}\_\{\textbackslash{}ast\} and R\textasciicircum{}\{1\}\textbackslash{}bar\{f\}\_\{\textbackslash{}ast\}; Note: Systematic notational difference, reported once. This edit fixes a sample occurrence; the coordinator gives three examples and the remaining scope.; Proposed edit: \$\textbackslash{}bar\{f\}\_\{\textbackslash{}ast\}\textbackslash{}GG\_\{m\}(\textbackslash{}mathfrak\{m\})\$ dans =\allowbreak > \$f\_\{\textbackslash{}ast\}\textbackslash{}GG\_\{m\}(\textbackslash{}mathfrak\{m\})\$ dans\par
a fix from one model only\par
\par\noindent\hypertarget{DP0608}{}\textbf{DP0608 -- XVIII, p. 542 (one\_model\_fix)}\par
The four left-hand vertices of diagrams (1.5.15.1) and (1.5.15.2)\par
\textbf{Claude Opus 5.5.} nothing reported here\par
\textbf{GPT-6.1 Sol.} Place: The four left-hand vertices of diagrams (1.5.15.1) and (1.5.15.2); Reported print: \textbackslash{}GG\_m(\textbackslash{}mathfrak\{m\}), with no index on the modulus; Source TeX: \textbackslash{}GG\_m(\textbackslash{}mathfrak\{m\}\_1) or \textbackslash{}GG\_m(\textbackslash{}mathfrak\{m\}\_2); Note: The edition systematically adds modulus indices in these diagrams. This sample edit is to be extended to all four vertices, as specified in the coordinator.; Proposed edit: \textbackslash{}Ext\textasciicircum{}\{1\}(R\textasciicircum{}1\textbackslash{}bar\{f\}\_\{1\textbackslash{}ast\} \textbackslash{}GG\_m (\textbackslash{}mathfrak\{m\}\_1), G) =\allowbreak > \textbackslash{}Ext\textasciicircum{}\{1\}(R\textasciicircum{}1\textbackslash{}bar\{f\}\_\{1\textbackslash{}ast\} \textbackslash{}GG\_m (\textbackslash{}mathfrak\{m\}), G)\par
a fix from one model only\par
\par\noindent\hypertarget{DP0609}{}\textbf{DP0609 -- XVIII, p. 542 (one\_model\_fix)}\par
Diagram (1.5.15.1), label of the left vertical arrow\par
\textbf{Claude Opus 5.5.} nothing reported here\par
\textbf{GPT-6.1 Sol.} Place: Diagram (1.5.15.1), label of the left vertical arrow; Reported print: \textbackslash{}Ext\textasciicircum{}\{1\}(u\textasciicircum{}\{\textbackslash{}ast\},\textbackslash{}Id); Source TeX: \textbackslash{}Ext\textasciicircum{}\{1\}(\textbackslash{}bar\{u\}\textasciicircum{}\{\textbackslash{}ast\},\textbackslash{}Id); Note: The printed arrow label has no bar on u.; Proposed edit: \textbackslash{}Ext\textasciicircum{}\{1\}(\textbackslash{}bar\{u\}\textasciicircum{}\{\textbackslash{}ast\},\textbackslash{}Id) =\allowbreak > \textbackslash{}Ext\textasciicircum{}\{1\}(u\textasciicircum{}\{\textbackslash{}ast\},\textbackslash{}Id)\par
a fix from one model only\par
\par\noindent\hypertarget{DP0610}{}\textbf{DP0610 -- XVIII, p. 542 (one\_model\_fix)}\par
Diagram (1.5.15.2), label of the left vertical arrow\par
\textbf{Claude Opus 5.5.} nothing reported here\par
\textbf{GPT-6.1 Sol.} Place: Diagram (1.5.15.2), label of the left vertical arrow; Reported print: \textbackslash{}Ext\textasciicircum{}\{1\}(N\_\{X\_1/X\_2\},\textbackslash{}Id); Source TeX: \textbackslash{}Ext\textasciicircum{}\{1\}(N\_\{\textbackslash{}bar\{X\}\_1 /\textbackslash{}bar\{X\}\_2\}, \textbackslash{}Id); Note: Unlike the norm in the preceding displayed isomorphism, the printed diagram label has no bars on either X.; Proposed edit: \textbackslash{}Ext\textasciicircum{}\{1\}(N\_\{\textbackslash{}bar\{X\}\_1 /\textbackslash{}bar\{X\}\_2\}, \textbackslash{}Id) =\allowbreak > \textbackslash{}Ext\textasciicircum{}\{1\}(N\_\{X\_1 /X\_2\}, \textbackslash{}Id)\par
a fix from one model only\par
\par\noindent\hypertarget{DP0611}{}\textbf{DP0611 -- XVIII, p. 543 (one\_model\_fix)}\par
1.6.2, after the first representability assertion ending « localement de type fini »\par
\textbf{Claude Opus 5.5.} nothing reported here\par
\textbf{GPT-6.1 Sol.} Place: 1.6.2, after the first representability assertion ending « localement de type fini »; Reported print: The assertion ends with a period and is followed directly by « Si P est un groupe algébrique… »; there is no footnote.; Source TeX: \textbackslash{}acorriger\{Dans le cas où ... un schéma en groupe lisse.\}; Note: The entire additional discussion in this \textbackslash{}acorriger block is absent from the printed page.; Proposed edit: \textbackslash{}acorriger\{Dans le cas où \$G'\_1=\allowbreak G'\_2=\allowbreak 0\$, c'est évident. Dans le cas où \$G''\_1=\allowbreak G''\_2=\allowbreak 0\$, en utilisant \textbackslash{}SGA 3 V 10.1.1 (nouvelle édition), on peut se ramener au cas où \$G'\_1\textbackslash{}to G'\_2\$ est un monomorphisme, qui est traité dans (\textbackslash{}SGA 3 VI\$\{\}\_\{\textbackslash{}text\{A\}\}\$ 3.2). On est ainsi ramené au cas où \$G'\_1=\allowbreak G''\_1=\allowbreak 0\$. On peut alors considérer le faisceau \$G\$ comme un torseur sous \$G'\_2\$ au-dessus du schéma \$G''\_2\$. Si on suppose que \$G'\_2\$ est affine (situation dans laquelle on se retrouve d'après (\textbackslash{}SGA 3 VI\$\{\}\_\{\textbackslash{}text\{B\}\}\$ 11.17) si avant de faire des réductions on a supposé \$G'\_1\$ et \$G'\_2\$ affines), alors la représentabilité de \$G\$ résulte aussitôt de la descente fidèlement plate. Ce cas particulier est suffisant pour les utilisations faites dans ce paragraphe. Si \$G'\_2\$ est quasi-projectif, on peut obtenir la représentabilité de \$G\$ en utilisant (\textbackslash{}SGA 1 VIII 7.7) et des arguments semblables à ceux de \textbackslash{}cite[Lemme\textasciitilde{}3.3.1]\{SGA4-18:4bis\}. En fait, tout schéma en groupes \$G\$ de type fini sur un corps est quasi-projectif : on peut supposer que \$G\$ est connexe, c'est alors \textbackslash{}cite[Corollary\textasciitilde{}1.2]\{SGA4-18:0bis\} si \$G\$ est lisse et (M. Raynaud) le cas général s'y ramène parce que d'après (\textbackslash{}SGA 3 VII\$\{\}\_\{\textbackslash{}text\{A\}\}\$ 8.3) il existe un morphisme fini de \$G\$ vers un schéma en groupe lisse.\} =\allowbreak > \par
a fix from one model only\par
\par\noindent\hypertarget{DP0612}{}\textbf{DP0612 -- XVIII, p. 544 (conflict)}\par
1.6.3, paragraph after (1.6.3.1) | Paragraph immediately after (1.6.3.1), second direct image\par
\textbf{Claude Opus 5.5.} Place: 1.6.3, paragraph after (1.6.3.1); Reported print: Les faisceaux p\_* G\_\{m\_X\} et p\_* G\_\{m\_Y\}; Source TeX: \$p\_\{\textbackslash{}ast\}\textbackslash{}GG\_\{m\_\{X\}\}\$ et \$(pi)\_\{\textbackslash{}ast\}\textbackslash{}GG\_\{m\_\{Y\}\}\$; Proposed correction: (pi)\_\{\textbackslash{}ast\}\textbackslash{}GG\_\{m\_\{Y\}\}; Note: The print has p\_* for the second sheaf; the TeX silently corrects it to (pi)\_*, as in (1.6.3.1).\par
\textbf{GPT-6.1 Sol.} Place: Paragraph immediately after (1.6.3.1), second direct image; Reported print: p\_\{\textbackslash{}ast\}\textbackslash{}GG\_\{m\_Y\}; Source TeX: (pi)\_\{\textbackslash{}ast\}\textbackslash{}GG\_\{m\_Y\}; Note: The printed prose uses p\_\{\textbackslash{}ast\} for this second direct image, whereas the TeX writes the composite explicitly.; Proposed edit: \$(pi)\_\{\textbackslash{}ast\}\textbackslash{}GG\_\{m\_\{Y\}\}\$ sont =\allowbreak > \$p\_\{\textbackslash{}ast\}\textbackslash{}GG\_\{m\_\{Y\}\}\$ sont\par
the two models differ at this place: one model fixes the TeX, the other calls it a misprint of the print\par
\par\noindent\hypertarget{DP0613}{}\textbf{DP0613 -- XVIII, p. 546 (conflict)}\par
Rappelons (SGA 1 …) que le revêtement X' de X se relève | Reference after « Rappelons »\par
\textbf{Claude Opus 5.5.} Place: Rappelons (SGA 1 …) que le revêtement X' de X se relève; Reported print: Rappelons (SGA 1 XII) que; Source TeX: Rappelons (\textbackslash{}SGA 1 XIII 2.4); Note: The print cites SGA 1 XII with no number. The TeX gives XIII 2.4 without a \textbackslash{}sisi.; Proposed edit: Rappelons (\textbackslash{}SGA 1 XIII 2.4) =\allowbreak > Rappelons (\textbackslash{}SGA 1 XII)\par
\textbf{GPT-6.1 Sol.} Place: Reference after « Rappelons »; Reported print: (SGA 1 XII); Source TeX: (\textbackslash{}SGA 1 XIII 2.4); Proposed correction: (\textbackslash{}SGA 1 XIII 2.4); Note: The lifting result is in SGA 1 XIII, not XII. The TeX silently corrects and specifies the printed reference.\par
the two models differ at this place: one model fixes the TeX, the other calls it a misprint of the print\par
\par\noindent\hypertarget{DP0614}{}\textbf{DP0614 -- XVIII, p. 548 (one\_model\_fix)}\par
Diagram (1.6.9.1), vertical arrow from f\textasciicircum{}\{-1\}(V) to V\par
\textbf{Claude Opus 5.5.} nothing reported here\par
\textbf{GPT-6.1 Sol.} Place: Diagram (1.6.9.1), vertical arrow from f\textasciicircum{}\{-1\}(V) to V; Reported print: \$f\_V\$; Source TeX: \$f\_v\$; Note: The printed subscript is capital V, as confirmed by the witness scan.; Proposed edit: \textbackslash{}ar[d]\textasciicircum{}\{f\_v\} =\allowbreak > \textbackslash{}ar[d]\textasciicircum{}\{f\_V\}\par
a fix from one model only\par
\par\noindent\hypertarget{DP0615}{}\textbf{DP0615 -- XVIII, p. 548 (one\_model\_fix)}\par
Target of the displayed trace Tr\_u in Lemma 1.6.9\par
\textbf{Claude Opus 5.5.} nothing reported here\par
\textbf{GPT-6.1 Sol.} Place: Target of the displayed trace Tr\_u in Lemma 1.6.9; Reported print: \$R\textasciicircum{}\{1\}f\_\{V!\}\textbackslash{}sheaf\{\textbackslash{}ZZ\}/n\$; Source TeX: \$R\textasciicircum{}\{1\}f\_\{v!\}\textbackslash{}sheaf\{\textbackslash{}ZZ\}/n\$; Note: Here too the printed subscript is capital V.; Proposed edit: R\textasciicircum{}\{1\}f\_\{v!\}\textbackslash{}sheaf\{\textbackslash{}ZZ\}/n =\allowbreak > R\textasciicircum{}\{1\}f\_\{V!\}\textbackslash{}sheaf\{\textbackslash{}ZZ\}/n\par
a fix from one model only\par
\par\noindent\hypertarget{DP0616}{}\textbf{DP0616 -- XVIII, p. 553 (one\_model\_fix)}\par
Renvois aux cup-produits, après (2.8.2)\par
\textbf{Claude Opus 5.5.} nothing reported here\par
\textbf{GPT-6.1 Sol.} Place: Renvois aux cup-produits, après (2.8.2); Reported print: XVII 5.4.2.2 et 5.4.3.5; Source TeX: XVII 5.4.2.2 et \textbackslash{}hyperrefext[XVII][XVII.lem5.4.3.5]\{5.4.3.5\}; Note: Le macro ajoute un second « XVII », absent de l'imprimé.; Proposed edit: \textbackslash{}hyperrefext[XVII][XVII.lem5.4.3.5]\{5.4.3.5\} =\allowbreak > \textbackslash{}hyperref[XVII.lem5.4.3.5]\{5.4.3.5\}\par
a fix from one model only\par
\par\noindent\hypertarget{DP0617}{}\textbf{DP0617 -- XVIII, p. 555 (one\_model\_fix)}\par
Début de b), étiquette de la condition\par
\textbf{Claude Opus 5.5.} nothing reported here\par
\textbf{GPT-6.1 Sol.} Place: Début de b), étiquette de la condition; Reported print: (\textbackslash{}ast\textbackslash{}ast)\_\{d\}; Source TeX: condition: \$(\textbackslash{}ast)\_\{d\}\$; Note: L'imprimé porte deux astérisques pour cette condition plus forte, distincte de la condition (\textbackslash{}ast)\_\{d\} du théorème.; Proposed edit: condition: \$(\textbackslash{}ast)\_\{d\}\$ =\allowbreak > condition: \$(\textbackslash{}ast\textbackslash{}ast)\_\{d\}\$\par
a fix from one model only\par
\par\noindent\hypertarget{DP0618}{}\textbf{DP0618 -- XVIII, p. 563 (conflict)}\par
« le cas d =\allowbreak  1 n'est autre que 1.4.7 » | Reference for the case d=\allowbreak 1\par
\textbf{Claude Opus 5.5.} Place: « le cas d =\allowbreak  1 n'est autre que 1.4.7 »; Reported print: 1.4.7; Source TeX: \textbackslash{}hyperref[XVIII.lemfond1.6.9]\{1.6.9\}; Note: The print refers to 1.4.7; the TeX prints 1.6.9.; Proposed edit: \{1.6.9\} =\allowbreak > \{1.4.7\}\par
\textbf{GPT-6.1 Sol.} Place: Reference for the case d=\allowbreak 1; Reported print: 1.4.7; Source TeX: \textbackslash{}hyperref[XVIII.lemfond1.6.9]\{1.6.9\}; Proposed correction: \textbackslash{}hyperref[XVIII.lemfond1.6.9]\{1.6.9\}; Note: The printed reference is wrong. The TeX silently corrects it to the fundamental lemma 1.6.9.\par
the two models differ at this place: one model fixes the TeX, the other calls it a misprint of the print\par
\par\noindent\hypertarget{DP0619}{}\textbf{DP0619 -- XVIII, p. 565 (conflict)}\par
« composé des isomorphismes Tr\_\{g'\_0\} et Rg'\_0 Tr\_\{h'\_1\} » | Last trace-composition formula\par
\textbf{Claude Opus 5.5.} Place: « composé des isomorphismes Tr\_\{g'\_0\} et Rg'\_0 Tr\_\{h'\_1\} »; Reported print: Rg'\_0 Tr\_\{h'\_1\}; Source TeX: R\textasciicircum{}\{2d'\}g'\_\{0!\}\textbackslash{}Tr\_\{h'\_\{1\}\}; Note: The print has Rg'\_0 with no exponent 2d' and no !; the TeX rewrites the formula.; Proposed edit: R\textasciicircum{}\{2d'\}g'\_\{0!\}\textbackslash{}Tr\_\{h'\_\{1\}\} =\allowbreak > Rg'\_\{0\}\textbackslash{}Tr\_\{h'\_\{1\}\}\par
\textbf{GPT-6.1 Sol.} Place: Last trace-composition formula; Reported print: Rg'\_0\textbackslash{}Tr\_\{h'\_1\}; Source TeX: R\textasciicircum{}\{2d'\}g'\_\{0!\}\textbackslash{}Tr\_\{h'\_\{1\}\}; Proposed correction: R\textasciicircum{}\{2d'\}g'\_\{0!\}\textbackslash{}Tr\_\{h'\_\{1\}\}; Note: The printed expression omits both the cohomological degree 2d' and the proper-support sign !. The TeX silently supplies both for the composition of the top-degree trace maps.\par
the two models differ at this place: one model fixes the TeX, the other calls it a misprint of the print\par
\par\noindent\hypertarget{DP0620}{}\textbf{DP0620 -- XVIII, p. 572 (conflict)}\par
«et (iii) est XVII 6.2» | Proof of Proposition 3.1.8, reference for assertion (iii)\par
\textbf{Claude Opus 5.5.} Place: «et (iii) est XVII 6.2»; Reported print: et (iii) est XVII 6.2.; Source TeX: et (iii) est \textbackslash{}hyperrefext[XVII][XVII.prop6.2.11]\{6.2.11\}.; Note: The print refers to XVII 6.2. The TeX prints XVII 6.2.11, with no \textbackslash{}sisi.; Proposed edit: \textbackslash{}hyperrefext[XVII][XVII.prop6.2.11]\{6.2.11\} =\allowbreak > \textbackslash{}hyperrefext[XVII][XVII.prop6.2.11]\{6.2\}\par
\textbf{GPT-6.1 Sol.} Place: Proof of Proposition 3.1.8, reference for assertion (iii); Reported print: XVII 6.2; Source TeX: \textbackslash{}hyperrefext[XVII][XVII.prop6.2.11]\{6.2.11\}; Note: The print refers to subsection 6.2, not specifically to Proposition 6.2.11. The broader printed reference is not a misprint.; Proposed edit: \textbackslash{}hyperrefext[XVII][XVII.prop6.2.11]\{6.2.11\} =\allowbreak > XVII\textasciitilde{}6.2\par
the two models differ at this place\par
\par\noindent\hypertarget{DP0621}{}\textbf{DP0621 -- XVIII, p. 575 (conflict)}\par
proof of 3.1.12.1, «M =\allowbreak  A» | Proof of Corollary 3.1.12.1, final choice of M\par
\textbf{Claude Opus 5.5.} Place: proof of 3.1.12.1, «M =\allowbreak  A»; Reported print: M =\allowbreak  \ensuremath{𝒜} .; Source TeX: \$M=\allowbreak \textbackslash{}mathcal\{B\}\$; Note: The print reads M =\allowbreak  \ensuremath{𝒜}. The TeX changes it to \ensuremath{𝓑} without \textbackslash{}sisi.; Proposed edit: \$M=\allowbreak \textbackslash{}mathcal\{B\}\$ =\allowbreak > \$M=\allowbreak \textbackslash{}mathcal\{A\}\$\par
\textbf{GPT-6.1 Sol.} Place: Proof of Corollary 3.1.12.1, final choice of M; Reported print: M=\allowbreak \textbackslash{}mathcal\{A\}; Source TeX: M=\allowbreak \textbackslash{}mathcal\{B\}; Proposed correction: M=\allowbreak \textbackslash{}mathcal\{B\}; Note: To obtain restriction of scalars, M must be the regular B-module. The printed A is a misprint silently corrected by TeX.\par
the two models differ at this place: one model fixes the TeX, the other calls it a misprint of the print\par
\par\noindent\hypertarget{DP0622}{}\textbf{DP0622 -- XVIII, p. 576 (one\_model\_fix)}\par
(3.1.13.2), right-hand side\par
\textbf{Claude Opus 5.5.} nothing reported here\par
\textbf{GPT-6.1 Sol.} Place: (3.1.13.2), right-hand side; Reported print: R\{g'\}\textasciicircum{}\{!\}Rf\_\{!\}; Source TeX: Rg\textasciicircum{}\{!\} Rf\_\{!\}; Note: The printed right-hand g has a prime; the TeX omits it.; Proposed edit: \textbackslash{}longrightarrow Rg\textasciicircum{}\{!\} Rf\_\{!\} =\allowbreak > \textbackslash{}longrightarrow R\{g'\}\textasciicircum{}\{!\} Rf\_\{!\}\par
a fix from one model only\par
\par\noindent\hypertarget{DP0623}{}\textbf{DP0623 -- XVIII, p. 578 (conflict)}\par
3.1.15.2 | 3.1.15.2, identity involving the site morphisms\par
\textbf{Claude Opus 5.5.} Place: 3.1.15.2; Reported print: Il vérifie f\_\ensuremath{\Gamma}\ensuremath{\gamma} =\allowbreak  f; Source TeX: f\_\{\textbackslash{}Gamma\}\textbackslash{}gamma\textbackslash{}simeq f; Note: The print has an equals sign here. Elsewhere on these pages the typescript writes \ensuremath{\simeq} with a tilde over the bar, and that is not what this sign shows.; Proposed edit: \textbackslash{}gamma\textbackslash{}simeq f =\allowbreak > \textbackslash{}gamma=\allowbreak  f\par
\textbf{GPT-6.1 Sol.} Place: 3.1.15.2, identity involving the site morphisms; Reported print: f\_\{\textbackslash{}Gamma\}\textbackslash{}gamma=\allowbreak f; Source TeX: f\_\{\textbackslash{}Gamma\}\textbackslash{}gamma\textbackslash{}simeq f; Note: The print has an equality sign, not an isomorphism sign.; Proposed edit: \$f\_\{\textbackslash{}Gamma\}\textbackslash{}gamma\textbackslash{}simeq f\$ =\allowbreak > \$f\_\{\textbackslash{}Gamma\}\textbackslash{}gamma=\allowbreak f\$\par
the two models differ at this place\par
\par\noindent\hypertarget{DP0624}{}\textbf{DP0624 -- XVIII, p. 582 (one\_model\_fix)}\par
(3.2.1.1), terme j\_\{!\}R\textasciicircum{}\{2d\}\textbackslash{}varphi\_\{!\}\textbackslash{}ZZ/n\par
\textbf{Claude Opus 5.5.} nothing reported here\par
\textbf{GPT-6.1 Sol.} Place: (3.2.1.1), terme j\_\{!\}R\textasciicircum{}\{2d\}\textbackslash{}varphi\_\{!\}\textbackslash{}ZZ/n; Reported print: \textbackslash{}varphi\_\{!\}\textbackslash{}ZZ/n; Source TeX: \textbackslash{}varphi ! \textbackslash{}ZZ/n; Note: Le point d'exclamation est un indice de \ensuremath{\phi} dans le print, et non un symbole placé sur la ligne.; Proposed edit: \textbackslash{}varphi ! \textbackslash{}ZZ/n =\allowbreak > \textbackslash{}varphi\_\{!\}\textbackslash{}ZZ/n\par
a fix from one model only\par
\par\noindent\hypertarget{DP0625}{}\textbf{DP0625 -- XVIII, p. 583 (conflict)}\par
Remarque 3.2.2 | Remarque 3.2.2, « alors t est encore défini »\par
\textbf{Claude Opus 5.5.} Place: Remarque 3.2.2; Reported print: alors t est encore défini; Source TeX: \$t\_f\$ est encore défini; Proposed correction: \$t\_f\$; Note: The print has t without its index f; the TeX adds the index (corrects).\par
\textbf{GPT-6.1 Sol.} Place: Remarque 3.2.2, « alors t est encore défini »; Reported print: alors t est encore défini; Source TeX: \$t\_f\$ est encore défini; Note: Le print emploie ici t sans indice ; le TeX ajoute l'indice f.; Proposed edit: \$t\_f\$ est encore défini =\allowbreak > \$t\$ est encore défini\par
the two models differ at this place: one model fixes the TeX, the other calls it a misprint of the print\par
\par\noindent\hypertarget{DP0626}{}\textbf{DP0626 -- XVIII, p. 583 (one\_model\_fix)}\par
Après le lemme 3.2.3, première phrase de la démonstration\par
\textbf{Claude Opus 5.5.} nothing reported here\par
\textbf{GPT-6.1 Sol.} Place: Après le lemme 3.2.3, première phrase de la démonstration; Reported print: une variante de la définition de t, en termes; Source TeX: définition de \$t\_f\$,; Note: Le print emploie de nouveau t sans indice ; le TeX ajoute f.; Proposed edit: définition de \$t\_f\$, =\allowbreak > définition de \$t\$,\par
a fix from one model only\par
\par\noindent\hypertarget{DP0627}{}\textbf{DP0627 -- XVIII, p. 586 (one\_model\_fix)}\par
3.2.6, first line\par
\textbf{Claude Opus 5.5.} Place: 3.2.6, first line; Reported print: Soient S le spectre d'un corps algébriquement clos k, X un schéma; Source TeX: Soient \$S\$ le spectre d'un corps algébriquement clos, \$k\$, \$X\$ un; Note: The TeX adds a comma after « clos »; in the print k is the name of the field, not a separate item of the list.; Proposed edit: algébriquement clos, \$k\$, \$X\$ =\allowbreak > algébriquement clos \$k\$, \$X\$\par
\textbf{GPT-6.1 Sol.} nothing reported here\par
a fix from one model only\par
\par\noindent\hypertarget{DP0628}{}\textbf{DP0628 -- XIX, p. 588 (one\_model\_fix)}\par
after the title and author, before « Dans cet exposé »\par
\textbf{Claude Opus 5.5.} Place: after the title and author, before « Dans cet exposé »; Reported print: Sommaire. 1. Pureté pour l'anneau k[[x\_1, ..., x\_n]]. 2. Le cas d'un anneau strictement local. 3. Pureté. 4. Acyclicité locale d'un morphisme régulier. 5. Théorème de finitude. 6. Dimension cohomologique des morphismes affines. 7. Morphismes affines - fin de la démonstration.; Source TeX: (absent); Note: The printed page has a table of contents (Sommaire) that the TeX leaves out. The edition may have replaced it with its own table of contents on purpose; coordinator to decide.; Proposed edit: \textbackslash{}newcommand\{\textbackslash{}UU\}\{\textbackslash{}mathbf\{U\}\} =\allowbreak > \textbackslash{}newcommand\{\textbackslash{}UU\}\{\textbackslash{}mathbf\{U\}\}  \textbackslash{}emph\{Sommaire.\} \textbackslash{}begin\{enumerate\} \textbackslash{}item Pureté pour l'anneau \$k[[x\_\{1\},\textbackslash{}ldots,x\_\{n\}]]\$. \textbackslash{}item Le cas d'un anneau strictement local. \textbackslash{}item Pureté. \textbackslash{}item Acyclicité locale d'un morphisme régulier. \textbackslash{}item Théorème de finitude. \textbackslash{}item Dimension cohomologique des morphismes affines. \textbackslash{}item Morphismes affines - fin de la démonstration. \textbackslash{}end\{enumerate\}\par
\textbf{GPT-6.1 Sol.} nothing reported here\par
a fix from one model only\par
\par\noindent\hypertarget{DP0629}{}\textbf{DP0629 -- XIX, p. 588 (one\_model\_fix)}\par
Author line\par
\textbf{Claude Opus 5.5.} nothing reported here\par
\textbf{GPT-6.1 Sol.} Place: Author line; Reported print: par M. ARTIN; Source TeX: M. Artin; Note: The word « par » is omitted; the capitalization of the author's name is not at issue.; Proposed edit: \textbackslash{}begin\{center\} M. Artin =\allowbreak > \textbackslash{}begin\{center\} par M. Artin\par
a fix from one model only\par
\par\noindent\hypertarget{DP0630}{}\textbf{DP0630 -- XIX, p. 588 (one\_model\_fix)}\par
Between the author line and the introductory paragraph\par
\textbf{Claude Opus 5.5.} nothing reported here\par
\textbf{GPT-6.1 Sol.} Place: Between the author line and the introductory paragraph; Reported print: Sommaire. 1. Pureté pour l'anneau k[[x₁, …, xₙ]]. 2. Le cas d'un anneau strictement local. 3. Pureté. 4. Acyclicité locale d'un morphisme régulier. 5. Théorème de finitude. 6. Dimension cohomologique des morphismes affines. 7. Morphismes affines - fin de la démonstration.; Source TeX: The Sommaire and its seven entries are absent.; Note: The entire seven-item contents list is missing.; Proposed edit: Dans\textbackslash{}pageoriginale cet exposé =\allowbreak > \textbackslash{}noindent Sommaire. \textbackslash{}begin\{enumerate\} \textbackslash{}item Pureté pour l'anneau \$k[[x\_\{1\},\textbackslash{}ldots,x\_\{n\}]]\$. \textbackslash{}item Le cas d'un anneau strictement local. \textbackslash{}item Pureté. \textbackslash{}item Acyclicité locale d'un morphisme régulier. \textbackslash{}item Théorème de finitude. \textbackslash{}item Dimension cohomologique des morphismes affines. \textbackslash{}item Morphismes affines -- fin de la démonstration. \textbackslash{}end\{enumerate\}  Dans\textbackslash{}pageoriginale cet exposé\par
a fix from one model only\par
\par\noindent\hypertarget{DP0631}{}\textbf{DP0631 -- XIX, p. 589 (one\_model\_fix)}\par
« partout de codimension 1 »\par
\textbf{Claude Opus 5.5.} Place: « partout de codimension 1 »; Reported print: codimension 1 (plain number); Source TeX: codimension \textbackslash{}hyperref[XIX.sec1]\{1\}; Note: This is a dimension, not a reference, but the TeX links it to section 1. The printed text is the same; only the link is wrong.; Proposed edit: codimension \textbackslash{}hyperref[XIX.sec1]\{1\} =\allowbreak > codimension \$1\$\par
\textbf{GPT-6.1 Sol.} nothing reported here\par
a fix from one model only\par
\par\noindent\hypertarget{DP0632}{}\textbf{DP0632 -- XIX, p. 589 (one\_model\_fix)}\par
« même en dimension 2 »\par
\textbf{Claude Opus 5.5.} Place: « même en dimension 2 »; Reported print: dimension 2 (plain number); Source TeX: \textbackslash{}hyperref[XIX.sec2]\{2\}, où on dispose; Note: This is a dimension, not a reference, but the TeX links it to section 2. The printed text is the same; only the link is wrong.; Proposed edit: \textbackslash{}hyperref[XIX.sec2]\{2\}, où on dispose =\allowbreak > \$2\$, où on dispose\par
\textbf{GPT-6.1 Sol.} nothing reported here\par
a fix from one model only\par
\par\noindent\hypertarget{DP0633}{}\textbf{DP0633 -- XIX, p. 595 (one\_model\_fix)}\par
first « d'après ([3] 1.4) »\par
\textbf{Claude Opus 5.5.} Place: first « d'après ([3] 1.4) »; Reported print: ([3] 1.4); Source TeX: \textbackslash{}cite\{SGA4-19:3\}\textasciitilde{}\textbackslash{}hyperref[XIX.subsec1.4]\{1.4\}; Note: The printed text is the same, but 1.4 here is a number in reference [3]. The TeX links it to 1.4 of this exposé, which is a wrong reference target.; Proposed edit: d'après (\textbackslash{}cite\{SGA4-19:3\}\textasciitilde{}\textbackslash{}hyperref[XIX.subsec1.4]\{1.4\}) =\allowbreak > d'après (\textbackslash{}cite\{SGA4-19:3\}\textasciitilde{}1.4)\par
\textbf{GPT-6.1 Sol.} nothing reported here\par
a fix from one model only\par
\par\noindent\hypertarget{DP0634}{}\textbf{DP0634 -- XIX, p. 595 (one\_model\_fix)}\par
second « donc ([3] 1.4) »\par
\textbf{Claude Opus 5.5.} Place: second « donc ([3] 1.4) »; Reported print: ([3] 1.4); Source TeX: \textbackslash{}cite\{SGA4-19:3\}\textasciitilde{}\textbackslash{}hyperref[XIX.subsec1.4]\{1.4\}; Note: Same as the first occurrence: the link points to XIX 1.4 instead of 1.4 of [3].; Proposed edit: donc (\textbackslash{}cite\{SGA4-19:3\}\textasciitilde{}\textbackslash{}hyperref[XIX.subsec1.4]\{1.4\}) =\allowbreak > donc (\textbackslash{}cite\{SGA4-19:3\}\textasciitilde{}1.4)\par
\textbf{GPT-6.1 Sol.} nothing reported here\par
a fix from one model only\par
\par\noindent\hypertarget{DP0635}{}\textbf{DP0635 -- XIX, p. 605 (conflict)}\par
« il s'agit donc de démontrer que » | Paragraph beginning « Mais on sait déjà », coefficient in the vanishing assertion\par
\textbf{Claude Opus 5.5.} Place: « il s'agit donc de démontrer que »; Reported print: H\textasciicircum{}q(U, /u\_n) =\allowbreak  0; Source TeX: \$H\textasciicircum{}\{q\}(U,/u\_\{n\})=\allowbreak 0\$; Proposed correction: \$H\textasciicircum{}\{q\}(U,\ensuremath{𝛍}\_\{n\})=\allowbreak 0\$; Note: The handwritten \ensuremath{\mu}\_n is malformed in the print. The TeX keeps « /u\_n ».\par
\textbf{GPT-6.1 Sol.} Place: Paragraph beginning « Mais on sait déjà », coefficient in the vanishing assertion; Reported print: H\textasciicircum{}\{q\}(U,\textbackslash{}mu\_\{n\})=\allowbreak 0; Source TeX: H\textasciicircum{}\{q\}(U,/u\_\{n\})=\allowbreak 0; Note: The handwritten \ensuremath{\mu} has a strongly oblique descender. The TeX instead gives a slash followed by Latin u.; Proposed edit: /u\_\{n\} =\allowbreak > \ensuremath{𝛍}\_\{n\}\par
the two models differ at this place: one model fixes the TeX, the other calls it a misprint of the print\par
\par\noindent\hypertarget{DP0636}{}\textbf{DP0636 -- XIX, p. 606 (conflict)}\par
l. 7, définition de \textbackslash{}hat B | Premier paragraphe, définition de B complété\par
\textbf{Claude Opus 5.5.} Place: l. 7, définition de \textbackslash{}hat B; Reported print: ou \ensuremath{\hat{B}} =\allowbreak  B \ensuremath{\otimes}\_A Â; Source TeX: ou \$\textbackslash{}hat\{B\}=\allowbreak B\textbackslash{}otimes\_\{A\}\textbackslash{}hat\{A\}\$; Proposed correction: où \$\textbackslash{}hat\{B\}=\allowbreak B\textbackslash{}otimes\_\{A\}\textbackslash{}hat\{A\}\$; Note: The print has 'ou' without its accent where 'où' is meant; the TeX keeps the printed reading.\par
\textbf{GPT-6.1 Sol.} Place: Premier paragraphe, définition de B complété; Reported print: où \$\textbackslash{}hat\{B\}=\allowbreak B\textbackslash{}otimes\_\{A\}\textbackslash{}hat\{A\}\$; Source TeX: ou \$\textbackslash{}hat\{B\}=\allowbreak B\textbackslash{}otimes\_\{A\}\textbackslash{}hat\{A\}\$; Note: L'accent grave du relatif « où » est léger mais visible dans le scan.; Proposed edit: , ou \$\textbackslash{}hat\{B\}=\allowbreak  =\allowbreak > , où \$\textbackslash{}hat\{B\}=\allowbreak \par
the two models differ at this place: one model fixes the TeX, the other calls it a misprint of the print\par
\par\noindent\hypertarget{DP0637}{}\textbf{DP0637 -- XIX, p. 610 (one\_model\_fix)}\par
Condition (ii), sommet supérieur droit du carré\par
\textbf{Claude Opus 5.5.} nothing reported here\par
\textbf{GPT-6.1 Sol.} Place: Condition (ii), sommet supérieur droit du carré; Reported print: (R\textasciicircum{}\{2c\}i'\{\}\textasciicircum{}\{!\}); Source TeX: (R\textasciicircum{}\{2c\}i\textasciicircum{}\{!\}); Note: Le tirage porte i prime pour l'immersion du couple (Y',X').; Proposed edit: \& \& (R\textasciicircum{}\{2c\}i\textasciicircum{}\{!\}) =\allowbreak > \& \& (R\textasciicircum{}\{2c\}i'\{\}\textasciicircum{}\{!\})\par
a fix from one model only\par
\par\noindent\hypertarget{DP0638}{}\textbf{DP0638 -- XIX, p. 611 (conflict)}\par
diagram (3.4.1), bottom right object and bottom arrow | Diagramme (3.4.1), objet inférieur droit | Diagramme (3.4.1), flèche horizontale inférieure\par
\textbf{Claude Opus 5.5.} Place: diagram (3.4.1), bottom right object and bottom arrow; Reported print: (R\textasciicircum{}\{2c\}w\textasciicircum{}!)\ensuremath{\mu}\_n\textasciicircum{}\{\ensuremath{\otimes}c\} \ensuremath{\leftarrow}ξ— (R\textasciicircum{}\{2b\}v\textasciicircum{}!)(R\textasciicircum{}\{2a\}u\textasciicircum{}!)\ensuremath{\mu}\_n\textasciicircum{}\{\ensuremath{\otimes}c\}; Source TeX: (R\textasciicircum{}\{2b\}v\textasciicircum{}\{!\}(R\textasciicircum{}\{2a\}u\textasciicircum{}\{!\})\ensuremath{𝛍}\_\{n\}\textasciicircum{}\{\textbackslash{}boxtimes c\})\textbackslash{}ar[l]; Note: Three differences, fixed in one edit: (1) the parenthesis after v\textasciicircum{}! is misplaced, closing at the end instead; (2) the tensor sign is \textbackslash{}boxtimes instead of \textbackslash{}otimes; (3) the label ξ of the bottom arrow is missing.; Proposed edit: (R\textasciicircum{}\{2b\}v\textasciicircum{}\{!\}(R\textasciicircum{}\{2a\}u\textasciicircum{}\{!\})\ensuremath{𝛍}\_\{n\}\textasciicircum{}\{\textbackslash{}boxtimes c\})\textbackslash{}ar[l] =\allowbreak > (R\textasciicircum{}\{2b\}v\textasciicircum{}\{!\})(R\textasciicircum{}\{2a\}u\textasciicircum{}\{!\})\ensuremath{𝛍}\_\{n\}\textasciicircum{}\{\textbackslash{}otimes c\}\textbackslash{}ar[l]\_\{\textbackslash{}xi\}\par
\textbf{GPT-6.1 Sol.} Place: Diagramme (3.4.1), objet inférieur droit; Reported print: (R\textasciicircum{}\{2b\}v\textasciicircum{}\{!\})(R\textasciicircum{}\{2a\}u\textasciicircum{}\{!\})\ensuremath{𝛍}\_\{n\}\textasciicircum{}\{\textbackslash{}otimes c\}; Source TeX: (R\textasciicircum{}\{2b\}v\textasciicircum{}\{!\}(R\textasciicircum{}\{2a\}u\textasciicircum{}\{!\})\ensuremath{𝛍}\_\{n\}\textasciicircum{}\{\textbackslash{}boxtimes c\}); Note: L'imprimé ferme la première parenthèse après v\textasciicircum{}\{!\}, non à la fin de l'expression, et donne une puissance tensorielle.; Proposed edit: (R\textasciicircum{}\{2b\}v\textasciicircum{}\{!\}(R\textasciicircum{}\{2a\}u\textasciicircum{}\{!\})\ensuremath{𝛍}\_\{n\}\textasciicircum{}\{\textbackslash{}boxtimes c\}) =\allowbreak > (R\textasciicircum{}\{2b\}v\textasciicircum{}\{!\})(R\textasciicircum{}\{2a\}u\textasciicircum{}\{!\})\ensuremath{𝛍}\_\{n\}\textasciicircum{}\{\textbackslash{}otimes c\} || Place: Diagramme (3.4.1), flèche horizontale inférieure; Reported print: \textbackslash{}xleftarrow\{\textbackslash{}xi\}; Source TeX: \textbackslash{}ar[l]; Note: L'étiquette xi de la flèche manque.; Proposed edit: \textbackslash{}ar[l] =\allowbreak > \textbackslash{}ar[l]\textasciicircum{}\{\textbackslash{}xi\}\par
the two models differ at this place\par
\par\noindent\hypertarget{DP0639}{}\textbf{DP0639 -- XIX, p. 618 (one\_model\_fix)}\par
Théorème 6.1 and 6.1 bis, footnote (*)\par
\textbf{Claude Opus 5.5.} Place: Théorème 6.1 and 6.1 bis, footnote (*); Reported print: (*) Dépend de la résolution des singularités (cf. Introduction).; Source TeX: \textbackslash{}footnotemark. Soit \$f:X\textbackslash{}rightarrow Y\$; Note: Page 618 prints the footnote (*) again, for both 6.1 and 6.1 bis. The TeX has only \textbackslash{}footnotemark and \textbackslash{}footnotemark[1], so the page prints no footnote text.; Proposed edit: \textbackslash{}footnotemark. Soit \$f:X\textbackslash{}rightarrow Y\$ =\allowbreak > \textbackslash{}footnote\{Dépend de la résolution des singularités (\textbackslash{}cf Introduction).\}. Soit \$f:X\textbackslash{}rightarrow Y\$\par
\textbf{GPT-6.1 Sol.} nothing reported here\par
a fix from one model only\par
\par\noindent\hypertarget{DP0640}{}\textbf{DP0640 -- XIX, p. 622 (conflict)}\par
Cela veut dire que si z \ensuremath{\in} E\textasciicircum{}\{N-1\} est un point tel que \ensuremath{\delta}(z) … d-q | Premier paragraphe, condition sur le point z\par
\textbf{Claude Opus 5.5.} Place: Cela veut dire que si z \ensuremath{\in} E\textasciicircum{}\{N-1\} est un point tel que \ensuremath{\delta}(z) … d-q; Reported print: \ensuremath{\delta}(z) > d-q; Source TeX: \textbackslash{}delta(z)\textbackslash{}geq d-q; Note: The print has a strict >, the same glyph as in 'q > d', which agrees with \ensuremath{\delta}(R\textasciicircum{}q g\_*F) \ensuremath{\leq} d-q. The TeX has \ensuremath{\geq}.; Proposed edit: \$\textbackslash{}delta(z)\textbackslash{}geq d-q\$ =\allowbreak > \$\textbackslash{}delta(z)> d-q\$\par
\textbf{GPT-6.1 Sol.} Place: Premier paragraphe, condition sur le point z; Reported print: \$\textbackslash{}delta(z)>d-q\$; Source TeX: \$\textbackslash{}delta(z)\textbackslash{}geq d-q\$; Note: L'inégalité imprimée est stricte.; Proposed edit: \textbackslash{}delta(z)\textbackslash{}geq d-q =\allowbreak > \textbackslash{}delta(z)>d-q\par
the two models differ at this place\par
\par\noindent\hypertarget{DP0641}{}\textbf{DP0641 -- XIX, p. 623 (agreed\_fix\_held)}\par
Notons que U est déduit … x \ensuremath{\in} \ensuremath{\Gamma}(X,O\_X) | Même phrase, symbole du faisceau structural dans \ensuremath{\Gamma}\par
\textbf{Claude Opus 5.5.} Place: Notons que U est déduit … x \ensuremath{\in} \ensuremath{\Gamma}(X,O\_X); Reported print: \ensuremath{\Gamma}(X,O\_X); Source TeX: \textbackslash{}Gamma(X,0\_\{X\}); Note: The structure sheaf. The TeX types the digit zero; in the typewriter face O and 0 are hard to tell apart. Elsewhere (O\_Y in Lemme 6.5) the TeX uses the letter O.; Proposed edit: \textbackslash{}Gamma(X,0\_\{X\}) =\allowbreak > \textbackslash{}Gamma(X,O\_\{X\})\par
\textbf{GPT-6.1 Sol.} Place: Même phrase, symbole du faisceau structural dans \ensuremath{\Gamma}; Reported print: \$\textbackslash{}Gamma(X,O\_\{X\})\$; Source TeX: \$\textbackslash{}Gamma(X,0\_\{X\})\$; Note: Le symbole imprimé est le O du faisceau structural, non le chiffre 0.; Proposed edit: \textbackslash{}Gamma(X,0\_\{X\}) =\allowbreak > \textbackslash{}Gamma(X,O\_\{X\})\par
both models give the same fix, held: a model marks it low confidence\par
\par\noindent\hypertarget{DP0642}{}\textbf{DP0642 -- XIX, p. 633 (one\_model\_fix)}\par
fin de la suite spectrale de Leray\par
\textbf{Claude Opus 5.5.} Place: fin de la suite spectrale de Leray; Reported print: \ensuremath{\Longrightarrow} H\textasciicircum{}\{p+q\}(U,Z / \ensuremath{\ell}) ,  on se ramène; Source TeX: H\textasciicircum{}\{p+q\}(U,\textbackslash{}ZZ / \textbackslash{}ell)\{\textbackslash{}quoi\}.; Note: The print has a comma here: the sentence continues with « on se ramène ». The TeX puts a full stop.; Proposed edit: H\textasciicircum{}\{p+q\}(U,\textbackslash{}ZZ / \textbackslash{}ell)\{\textbackslash{}quoi\}. =\allowbreak > H\textasciicircum{}\{p+q\}(U,\textbackslash{}ZZ / \textbackslash{}ell)\{\textbackslash{}quoi\},\par
\textbf{GPT-6.1 Sol.} nothing reported here\par
a fix from one model only\par
\par\noindent\hypertarget{DP0643}{}\textbf{DP0643 -- XIX, p. 638 (one\_model\_fix)}\par
bibliography [3]\par
\textbf{Claude Opus 5.5.} Place: bibliography [3]; Reported print: Amer. Journal Math. vol 88, n°4; Source TeX: Amer, Journal Math. vol 88 \$\{\textbackslash{}rm n\textasciicircum{}\{\textbackslash{}circ\}\}\$ 4; Note: The print abbreviates 'Amer.' with a period and puts a comma after 'vol 88'. The TeX has a comma after 'Amer' and no comma after 88.; Proposed edit: Amer, Journal Math. vol 88 \$\{\textbackslash{}rm =\allowbreak > Amer. Journal Math. vol 88, \$\{\textbackslash{}rm\par
\textbf{GPT-6.1 Sol.} nothing reported here\par
a fix from one model only\par
\endgroup
